\documentclass[11pt,a4paper]{article}

% ============================================================
% PHÉNOMÈNES PHYSIQUES DU QUOTIDIEN
% De l'observation à la mise en équation
% Niveau Terminale -> CPGE -> Bac+3
% ============================================================

\usepackage[utf8]{inputenc}
\usepackage[T1]{fontenc}
\usepackage[french]{babel}
\usepackage{lmodern}
\usepackage{geometry}
\geometry{margin=1.7cm}
\usepackage{amsmath,amssymb}
\usepackage{siunitx}
\usepackage{array}
\usepackage{tabularx}
\usepackage{longtable}
\usepackage{booktabs}
\usepackage{multicol}
\usepackage{xcolor}
\usepackage{enumitem}
\usepackage{fancyhdr}
\usepackage{hyperref}
\usepackage[most]{tcolorbox}
\usepackage{tikz}
\usepackage{pgfplots}
\pgfplotsset{compat=1.18}

\usetikzlibrary{arrows.meta,calc,angles,quotes,patterns,decorations.pathmorphing}

\hypersetup{
  colorlinks=true,
  linkcolor=blue!60!black,
  urlcolor=blue!60!black
}

\pagestyle{fancy}
\fancyhf{}
\lhead{Phénomènes physiques du quotidien}
\rhead{Terminale -- CPGE -- Bac+3}
\cfoot{\thepage}

\setlist[itemize]{leftmargin=1.2em}
\setlist[enumerate]{leftmargin=1.5em}

\definecolor{bleuprepa}{RGB}{20,55,110}
\definecolor{vertprepa}{RGB}{20,105,70}
\definecolor{rougeprepa}{RGB}{160,35,35}
\definecolor{orangeprepa}{RGB}{180,95,20}
\definecolor{violetprepa}{RGB}{90,45,135}
\definecolor{grisclair}{RGB}{245,247,250}

\tcbset{
  enhanced,
  breakable,
  sharp corners=downhill,
  boxrule=0.7pt,
  colback=white,
  fonttitle=\bfseries
}

\newtcolorbox{definitionbox}[1][]{
  colframe=bleuprepa,
  colback=blue!2!white,
  title={Définition #1}
}

\newtcolorbox{theorembox}[1][]{
  colframe=vertprepa,
  colback=green!2!white,
  title={Loi / Théorème #1}
}

\newtcolorbox{methodebox}[1][]{
  colframe=vertprepa,
  colback=green!3!white,
  title={Méthode #1}
}

\newtcolorbox{exemplebox}[1][]{
  colframe=bleuprepa,
  colback=cyan!2!white,
  title={Exemple #1}
}

\newtcolorbox{attentionbox}[1][]{
  colframe=rougeprepa,
  colback=red!2!white,
  title={Attention #1}
}

\newtcolorbox{approfondissementbox}[1][]{
  colframe=violetprepa,
  colback=violet!3!white,
  title={Approfondissement #1}
}

\newtcolorbox{exercicebox}[1][]{
  colframe=orangeprepa,
  colback=orange!2!white,
  title={Exercice #1}
}

\newtcolorbox{corrigebox}[1][]{
  colframe=vertprepa,
  colback=green!1!white,
  title={Correction / Indication #1}
}

\newtcolorbox{odgbox}[1][]{
  colframe=black,
  colback=gray!4!white,
  title={Ordre de grandeur #1}
}

\newcommand{\dd}{\,\mathrm d}
\newcommand{\vect}[1]{\overrightarrow{#1}}
\newcommand{\e}{\mathrm e}

\title{\Huge\textbf{Phénomènes physiques du quotidien}\\[2mm]
\Large De l'observation à la mise en équation\\[1mm]
\large Niveau Terminale, CPGE et Bac+3}
\author{Polycopié pédagogique de physique}
\date{}

\begin{document}
\maketitle

\begin{attentionbox}[Ambition du document]
Ce document est volontairement dense. Il part d'observations du quotidien puis construit des modèles physiques progressifs, depuis le niveau Terminale jusqu'à des encadrés de niveau CPGE/Bac+3. Les modèles avancés sont donnés pour comprendre la structure des phénomènes, sans exiger une maîtrise complète de tous les outils mathématiques dès la première lecture.
\end{attentionbox}

\tableofcontents
\newpage


% ============================================================
\section{Partie I — Outils fondamentaux de physique}
% ============================================================

\subsection{Grandeurs, unités et ordres de grandeur}

\begin{definitionbox}[Grandeur physique]
Une grandeur physique est une quantité mesurable associée à une dimension : longueur, durée, masse, intensité électrique, température, quantité de matière, intensité lumineuse. Une équation physique n'a de sens que si elle est homogène, c'est-à-dire si tous ses termes ont la même dimension.
\end{definitionbox}

\begin{methodebox}[Tester une formule]
Pour tester une relation :
\begin{enumerate}
  \item identifier la dimension du membre de gauche ;
  \item identifier la dimension de chaque terme du membre de droite ;
  \item vérifier que les fonctions mathématiques portent sur des quantités sans dimension ;
  \item vérifier les ordres de grandeur numériques.
\end{enumerate}
\end{methodebox}

\begin{odgbox}[Ordres de grandeur utiles]
\begin{center}
\begin{tabularx}{\textwidth}{lllX}
\toprule
Grandeur & Symbole & Valeur typique & Commentaire\\
\midrule
Vitesse de la lumière & $c$ & $\SI{3.00e8}{m.s^{-1}}$ & propagation dans le vide\\
Vitesse du son dans l'air & $c_s$ & $\SI{340}{m.s^{-1}}$ & dépend de la température\\
Longueur d'onde visible & $\lambda$ & $\SI{400}{nm}$ à $\SI{700}{nm}$ & violet vers rouge\\
Indice de l'air & $n$ & $1{,}0003$ & proche de 1 mais variable\\
Indice de l'eau & $n$ & $1{,}33$ & responsable de la réfraction\\
Indice du verre & $n$ & $1{,}5$ environ & dépend du verre\\
Taille d'une goutte de pluie & $R$ & $\SI{0.5}{mm}$ à $\SI{2}{mm}$ & arc-en-ciel\\
Fréquence audible & $f$ & $\SI{20}{Hz}$ à $\SI{20}{kHz}$ & oreille humaine\\
\bottomrule
\end{tabularx}
\end{center}
\end{odgbox}

\begin{approfondissementbox}[Niveau prépa / Bac+3]
L'analyse dimensionnelle peut s'étendre au théorème de Buckingham. Si un phénomène dépend de grandeurs dimensionnées, on peut souvent réduire le problème à quelques nombres sans dimension. Par exemple, en mécanique des fluides, le nombre de Reynolds
\[
\mathrm{Re}=\frac{\rho vL}{\eta}
\]
compare les effets inertiels aux effets visqueux. Il permet de prévoir si un écoulement est plutôt laminaire ou turbulent.
\end{approfondissementbox}

\subsection{Ondes : vocabulaire commun}

\begin{definitionbox}[Onde progressive]
Une onde progressive transporte une perturbation sans transport global de matière. Elle est caractérisée par une célérité $v$, une longueur d'onde $\lambda$, une fréquence $f$ et une période $T$.
\[
v=\lambda f,\qquad f=\frac1T.
\]
\end{definitionbox}

\begin{exemplebox}[Son de concert]
Pour un La de fréquence $\SI{440}{Hz}$ dans l'air, avec $v=\SI{340}{m.s^{-1}}$ :
\[
\lambda=\frac{v}{f}\simeq \frac{340}{440}\simeq \SI{0.77}{m}.
\]
La longueur d'onde sonore est donc de l'ordre du mètre, ce qui explique que le son contourne facilement des obstacles de taille humaine.
\end{exemplebox}

\begin{approfondissementbox}[Niveau prépa / Bac+3]
Une onde unidimensionnelle idéale peut être décrite par l'équation d'onde :
\[
\frac{\partial^2 s}{\partial t^2}=c^2\frac{\partial^2 s}{\partial x^2}.
\]
Les solutions de la forme $s(x,t)=f(x-ct)+g(x+ct)$ représentent deux ondes progressives, l'une vers la droite, l'autre vers la gauche. Les phénomènes d'interférences apparaissent par superposition linéaire.
\end{approfondissementbox}

\subsection{Lumière : rayon, onde, photon}

\begin{methodebox}[Choisir le bon modèle]
\begin{center}
\begin{tabularx}{\textwidth}{p{4cm}p{4cm}X}
\toprule
Phénomène & Modèle adapté & Loi ou outil\\
\midrule
Lentille, miroir, fibre & rayon lumineux & Snell-Descartes, réflexion\\
Diffraction, interférences & onde & différence de marche, phase\\
Effet photoélectrique, laser & photon & énergie $E=h\nu$\\
Arc-en-ciel & rayon + dispersion & réfraction, réflexion interne\\
Ciel bleu & onde + diffusion & loi en $1/\lambda^4$\\
\bottomrule
\end{tabularx}
\end{center}
\end{methodebox}

\begin{approfondissementbox}[Niveau prépa / Bac+3]
Le modèle du rayon lumineux correspond à une limite de très courte longueur d'onde devant les dimensions du problème. On parle d'approximation eikonale. À l'inverse, lorsque la taille des obstacles est comparable à $\lambda$, la diffraction devient incontournable et le modèle géométrique échoue.
\end{approfondissementbox}


% ============================================================
\section{Partie — Optique géométrique}
% ============================================================


% ------------------------------------------------------------
\subsection{1. Mirage inférieur sur une route chaude}
% ------------------------------------------------------------

\begin{exemplebox}[Observation concrète]
Sur une route chauffée par le Soleil, on croit voir une flaque brillante au loin.
\end{exemplebox}

\begin{methodebox}[Description qualitative]
L'air proche du sol est plus chaud, moins dense et possède un indice légèrement plus faible. Les rayons se courbent vers les zones d'indice plus élevé.
La stratégie de compréhension est toujours la même : isoler la grandeur qui varie, identifier la loi physique, puis vérifier le sens du phénomène par un ordre de grandeur.
\end{methodebox}

\begin{theorembox}[Mise en équation niveau Terminale]
La relation centrale utilisée ici est :
\[
n=n(z),\qquad n_k\sin i_k=n_{k+1}\sin i_{k+1}
\]
Cette formule n'est pas seulement un outil de calcul : elle permet de prévoir comment le phénomène change lorsqu'on modifie une grandeur expérimentale.
\end{theorembox}

\begin{odgbox}[Ordre de grandeur]
La variation de l'indice de l'air est très faible, de l'ordre de $10^{-4}$, mais suffisante sur de grandes distances.
\end{odgbox}

\begin{attentionbox}[Limites du modèle]
Le modèle présenté repose sur des hypothèses simplificatrices : milieu homogène ou localement homogène, grandeurs moyennes, géométrie idéalisée, absence de couplages secondaires. En situation réelle, les frottements, gradients, turbulences, imperfections géométriques ou effets non linéaires peuvent modifier le résultat.
\end{attentionbox}

\begin{approfondissementbox}[Niveau prépa / Bac+3]
Dans un milieu stratifié, la quantité $n(z)\sin i(z)$ se conserve localement. Le rayon se courbe continûment au lieu de se briser en segments.
On peut souvent améliorer le modèle en remplaçant une relation algébrique simple par une équation différentielle, une loi locale ou un bilan intégral. Cela permet de passer d'une explication qualitative à une prédiction quantitative.
\end{approfondissementbox}

\begin{exercicebox}[Exercice guidé]
Modéliser trois couches d'air avec $n_1>n_2>n_3$ et déterminer qualitativement le sens de courbure.
\begin{enumerate}
  \item Identifier la grandeur principale.
  \item Donner son unité.
  \item Écrire la loi physique pertinente.
  \item Tester l'homogénéité de la relation.
  \item Interpréter le résultat dans le contexte.
\end{enumerate}
\end{exercicebox}

\begin{corrigebox}[Indication]
Le rayon se rapproche de la normale quand il entre dans un indice plus élevé et s'en éloigne dans un indice plus faible ; par couches successives il se courbe.
\end{corrigebox}


% ------------------------------------------------------------
\subsection{2. Mirage supérieur}
% ------------------------------------------------------------

\begin{exemplebox}[Observation concrète]
Un bateau ou un paysage semble flotter au-dessus de l'horizon.
\end{exemplebox}

\begin{methodebox}[Description qualitative]
Une inversion de température crée une couche d'air plus froide près du sol ou de la mer, donc plus réfringente.
La stratégie de compréhension est toujours la même : isoler la grandeur qui varie, identifier la loi physique, puis vérifier le sens du phénomène par un ordre de grandeur.
\end{methodebox}

\begin{theorembox}[Mise en équation niveau Terminale]
La relation centrale utilisée ici est :
\[
n_{\mathrm{bas}}>n_{\mathrm{haut}}\quad\Rightarrow\quad \text{rayon courbé vers le bas}
\]
Cette formule n'est pas seulement un outil de calcul : elle permet de prévoir comment le phénomène change lorsqu'on modifie une grandeur expérimentale.
\end{theorembox}

\begin{odgbox}[Ordre de grandeur]
Le phénomène est fréquent au-dessus de mers froides ou de surfaces glacées.
\end{odgbox}

\begin{attentionbox}[Limites du modèle]
Le modèle présenté repose sur des hypothèses simplificatrices : milieu homogène ou localement homogène, grandeurs moyennes, géométrie idéalisée, absence de couplages secondaires. En situation réelle, les frottements, gradients, turbulences, imperfections géométriques ou effets non linéaires peuvent modifier le résultat.
\end{attentionbox}

\begin{approfondissementbox}[Niveau prépa / Bac+3]
Le rayon lumineux peut suivre approximativement la courbure terrestre si le gradient d'indice compense la géométrie.
On peut souvent améliorer le modèle en remplaçant une relation algébrique simple par une équation différentielle, une loi locale ou un bilan intégral. Cela permet de passer d'une explication qualitative à une prédiction quantitative.
\end{approfondissementbox}

\begin{exercicebox}[Exercice guidé]
Expliquer pourquoi un objet caché par la courbure terrestre peut devenir visible.
\begin{enumerate}
  \item Identifier la grandeur principale.
  \item Donner son unité.
  \item Écrire la loi physique pertinente.
  \item Tester l'homogénéité de la relation.
  \item Interpréter le résultat dans le contexte.
\end{enumerate}
\end{exercicebox}

\begin{corrigebox}[Indication]
Les rayons issus de l'objet sont ramenés vers l'observateur par réfraction atmosphérique.
\end{corrigebox}


% ------------------------------------------------------------
\subsection{3. Fata Morgana}
% ------------------------------------------------------------

\begin{exemplebox}[Observation concrète]
Des objets lointains sont étirés, inversés ou empilés verticalement.
\end{exemplebox}

\begin{methodebox}[Description qualitative]
Des couches d'air multiples d'indices différents produisent plusieurs trajets lumineux.
La stratégie de compréhension est toujours la même : isoler la grandeur qui varie, identifier la loi physique, puis vérifier le sens du phénomène par un ordre de grandeur.
\end{methodebox}

\begin{theorembox}[Mise en équation niveau Terminale]
La relation centrale utilisée ici est :
\[
\text{plusieurs chemins optiques}\quad\Rightarrow\quad \text{images multiples}
\]
Cette formule n'est pas seulement un outil de calcul : elle permet de prévoir comment le phénomène change lorsqu'on modifie une grandeur expérimentale.
\end{theorembox}

\begin{odgbox}[Ordre de grandeur]
Les gradients de température près de la mer peuvent être très structurés.
\end{odgbox}

\begin{attentionbox}[Limites du modèle]
Le modèle présenté repose sur des hypothèses simplificatrices : milieu homogène ou localement homogène, grandeurs moyennes, géométrie idéalisée, absence de couplages secondaires. En situation réelle, les frottements, gradients, turbulences, imperfections géométriques ou effets non linéaires peuvent modifier le résultat.
\end{attentionbox}

\begin{approfondissementbox}[Niveau prépa / Bac+3]
La description complète demande une équation de rayon dans un indice $n(z)$ non monotone.
On peut souvent améliorer le modèle en remplaçant une relation algébrique simple par une équation différentielle, une loi locale ou un bilan intégral. Cela permet de passer d'une explication qualitative à une prédiction quantitative.
\end{approfondissementbox}

\begin{exercicebox}[Exercice guidé]
Dessiner deux rayons issus d'un même point atteignant l'œil par deux chemins.
\begin{enumerate}
  \item Identifier la grandeur principale.
  \item Donner son unité.
  \item Écrire la loi physique pertinente.
  \item Tester l'homogénéité de la relation.
  \item Interpréter le résultat dans le contexte.
\end{enumerate}
\end{exercicebox}

\begin{corrigebox}[Indication]
Deux rayons différents peuvent former deux images apparentes d'un même objet.
\end{corrigebox}


% ============================================================
\section{Partie — Optique atmosphérique}
% ============================================================


% ------------------------------------------------------------
\subsection{4. Arc-en-ciel primaire}
% ------------------------------------------------------------

\begin{exemplebox}[Observation concrète]
Un arc coloré apparaît face au Soleil après la pluie.
\end{exemplebox}

\begin{methodebox}[Description qualitative]
Chaque goutte réfracte, réfléchit puis réfracte la lumière. La dispersion sépare les couleurs.
La stratégie de compréhension est toujours la même : isoler la grandeur qui varie, identifier la loi physique, puis vérifier le sens du phénomène par un ordre de grandeur.
\end{methodebox}

\begin{theorembox}[Mise en équation niveau Terminale]
La relation centrale utilisée ici est :
\[
D(i)=\pi+2i-4r,\qquad \sin i=n\sin r
\]
Cette formule n'est pas seulement un outil de calcul : elle permet de prévoir comment le phénomène change lorsqu'on modifie une grandeur expérimentale.
\end{theorembox}

\begin{odgbox}[Ordre de grandeur]
L'angle d'observation du rouge est proche de $42^\circ$.
\end{odgbox}

\begin{attentionbox}[Limites du modèle]
Le modèle présenté repose sur des hypothèses simplificatrices : milieu homogène ou localement homogène, grandeurs moyennes, géométrie idéalisée, absence de couplages secondaires. En situation réelle, les frottements, gradients, turbulences, imperfections géométriques ou effets non linéaires peuvent modifier le résultat.
\end{attentionbox}

\begin{approfondissementbox}[Niveau prépa / Bac+3]
L'arc correspond à un extremum de déviation : beaucoup de rayons émergent avec des directions voisines.
On peut souvent améliorer le modèle en remplaçant une relation algébrique simple par une équation différentielle, une loi locale ou un bilan intégral. Cela permet de passer d'une explication qualitative à une prédiction quantitative.
\end{approfondissementbox}

\begin{exercicebox}[Exercice guidé]
À partir de $D(i)$, expliquer pourquoi on cherche un extremum.
\begin{enumerate}
  \item Identifier la grandeur principale.
  \item Donner son unité.
  \item Écrire la loi physique pertinente.
  \item Tester l'homogénéité de la relation.
  \item Interpréter le résultat dans le contexte.
\end{enumerate}
\end{exercicebox}

\begin{corrigebox}[Indication]
Un extremum concentre les rayons et augmente l'intensité observée.
\end{corrigebox}


% ------------------------------------------------------------
\subsection{5. Arc-en-ciel secondaire}
% ------------------------------------------------------------

\begin{exemplebox}[Observation concrète]
Un second arc plus pâle apparaît parfois, avec les couleurs inversées.
\end{exemplebox}

\begin{methodebox}[Description qualitative]
Il résulte de deux réflexions internes dans la goutte.
La stratégie de compréhension est toujours la même : isoler la grandeur qui varie, identifier la loi physique, puis vérifier le sens du phénomène par un ordre de grandeur.
\end{methodebox}

\begin{theorembox}[Mise en équation niveau Terminale]
La relation centrale utilisée ici est :
\[
D_2(i)=2\pi+2i-6r
\]
Cette formule n'est pas seulement un outil de calcul : elle permet de prévoir comment le phénomène change lorsqu'on modifie une grandeur expérimentale.
\end{theorembox}

\begin{odgbox}[Ordre de grandeur]
Il est plus faible car chaque réflexion interne perd une partie de l'énergie.
\end{odgbox}

\begin{attentionbox}[Limites du modèle]
Le modèle présenté repose sur des hypothèses simplificatrices : milieu homogène ou localement homogène, grandeurs moyennes, géométrie idéalisée, absence de couplages secondaires. En situation réelle, les frottements, gradients, turbulences, imperfections géométriques ou effets non linéaires peuvent modifier le résultat.
\end{attentionbox}

\begin{approfondissementbox}[Niveau prépa / Bac+3]
La bande sombre d'Alexandre sépare les angles où aucun rayon primaire ou secondaire intense n'arrive.
On peut souvent améliorer le modèle en remplaçant une relation algébrique simple par une équation différentielle, une loi locale ou un bilan intégral. Cela permet de passer d'une explication qualitative à une prédiction quantitative.
\end{approfondissementbox}

\begin{exercicebox}[Exercice guidé]
Comparer qualitativement une et deux réflexions internes.
\begin{enumerate}
  \item Identifier la grandeur principale.
  \item Donner son unité.
  \item Écrire la loi physique pertinente.
  \item Tester l'homogénéité de la relation.
  \item Interpréter le résultat dans le contexte.
\end{enumerate}
\end{exercicebox}

\begin{corrigebox}[Indication]
Deux réflexions inversent l'ordre apparent des couleurs et diminuent l'intensité.
\end{corrigebox}


% ------------------------------------------------------------
\subsection{6. Bande sombre d'Alexandre}
% ------------------------------------------------------------

\begin{exemplebox}[Observation concrète]
La zone entre l'arc primaire et secondaire paraît plus sombre.
\end{exemplebox}

\begin{methodebox}[Description qualitative]
Les directions de sortie des rayons ne remplissent pas uniformément toutes les zones angulaires.
La stratégie de compréhension est toujours la même : isoler la grandeur qui varie, identifier la loi physique, puis vérifier le sens du phénomène par un ordre de grandeur.
\end{methodebox}

\begin{theorembox}[Mise en équation niveau Terminale]
La relation centrale utilisée ici est :
\[
\theta_{\mathrm{primaire}}\simeq42^\circ,\qquad \theta_{\mathrm{secondaire}}\simeq51^\circ
\]
Cette formule n'est pas seulement un outil de calcul : elle permet de prévoir comment le phénomène change lorsqu'on modifie une grandeur expérimentale.
\end{theorembox}

\begin{odgbox}[Ordre de grandeur]
La bande sombre se situe entre ces deux régions.
\end{odgbox}

\begin{attentionbox}[Limites du modèle]
Le modèle présenté repose sur des hypothèses simplificatrices : milieu homogène ou localement homogène, grandeurs moyennes, géométrie idéalisée, absence de couplages secondaires. En situation réelle, les frottements, gradients, turbulences, imperfections géométriques ou effets non linéaires peuvent modifier le résultat.
\end{attentionbox}

\begin{approfondissementbox}[Niveau prépa / Bac+3]
Une étude complète revient à analyser la densité angulaire des rayons émergents.
On peut souvent améliorer le modèle en remplaçant une relation algébrique simple par une équation différentielle, une loi locale ou un bilan intégral. Cela permet de passer d'une explication qualitative à une prédiction quantitative.
\end{approfondissementbox}

\begin{exercicebox}[Exercice guidé]
Expliquer pourquoi l'absence de rayons concentrés rend une zone plus sombre.
\begin{enumerate}
  \item Identifier la grandeur principale.
  \item Donner son unité.
  \item Écrire la loi physique pertinente.
  \item Tester l'homogénéité de la relation.
  \item Interpréter le résultat dans le contexte.
\end{enumerate}
\end{exercicebox}

\begin{corrigebox}[Indication]
L'intensité perçue dépend du nombre de rayons arrivant dans une même direction.
\end{corrigebox}


% ------------------------------------------------------------
\subsection{7. Halo solaire}
% ------------------------------------------------------------

\begin{exemplebox}[Observation concrète]
Un cercle lumineux apparaît autour du Soleil, souvent à $22^\circ$.
\end{exemplebox}

\begin{methodebox}[Description qualitative]
La lumière est réfractée par des cristaux de glace hexagonaux.
La stratégie de compréhension est toujours la même : isoler la grandeur qui varie, identifier la loi physique, puis vérifier le sens du phénomène par un ordre de grandeur.
\end{methodebox}

\begin{theorembox}[Mise en équation niveau Terminale]
La relation centrale utilisée ici est :
\[
\text{déviation minimale dans un prisme de glace}
\]
Cette formule n'est pas seulement un outil de calcul : elle permet de prévoir comment le phénomène change lorsqu'on modifie une grandeur expérimentale.
\end{theorembox}

\begin{odgbox}[Ordre de grandeur]
Le halo de $22^\circ$ est lié à la géométrie hexagonale.
\end{odgbox}

\begin{attentionbox}[Limites du modèle]
Le modèle présenté repose sur des hypothèses simplificatrices : milieu homogène ou localement homogène, grandeurs moyennes, géométrie idéalisée, absence de couplages secondaires. En situation réelle, les frottements, gradients, turbulences, imperfections géométriques ou effets non linéaires peuvent modifier le résultat.
\end{attentionbox}

\begin{approfondissementbox}[Niveau prépa / Bac+3]
On peut modéliser un cristal comme un prisme et chercher une déviation minimale.
On peut souvent améliorer le modèle en remplaçant une relation algébrique simple par une équation différentielle, une loi locale ou un bilan intégral. Cela permet de passer d'une explication qualitative à une prédiction quantitative.
\end{approfondissementbox}

\begin{exercicebox}[Exercice guidé]
Comparer le halo à l'arc-en-ciel : gouttes sphériques ou cristaux ?
\begin{enumerate}
  \item Identifier la grandeur principale.
  \item Donner son unité.
  \item Écrire la loi physique pertinente.
  \item Tester l'homogénéité de la relation.
  \item Interpréter le résultat dans le contexte.
\end{enumerate}
\end{exercicebox}

\begin{corrigebox}[Indication]
L'arc-en-ciel vient de gouttes liquides ; le halo vient de cristaux solides.
\end{corrigebox}


% ------------------------------------------------------------
\subsection{8. Parhélies}
% ------------------------------------------------------------

\begin{exemplebox}[Observation concrète]
Deux taches brillantes apparaissent parfois à gauche et à droite du Soleil.
\end{exemplebox}

\begin{methodebox}[Description qualitative]
Des cristaux de glace orientés horizontalement réfractent préférentiellement la lumière.
La stratégie de compréhension est toujours la même : isoler la grandeur qui varie, identifier la loi physique, puis vérifier le sens du phénomène par un ordre de grandeur.
\end{methodebox}

\begin{theorembox}[Mise en équation niveau Terminale]
La relation centrale utilisée ici est :
\[
\text{orientation statistique des cristaux}\Rightarrow\text{directions privilégiées}
\]
Cette formule n'est pas seulement un outil de calcul : elle permet de prévoir comment le phénomène change lorsqu'on modifie une grandeur expérimentale.
\end{theorembox}

\begin{odgbox}[Ordre de grandeur]
Les parhélies sont fréquents par temps froid avec nuages élevés.
\end{odgbox}

\begin{attentionbox}[Limites du modèle]
Le modèle présenté repose sur des hypothèses simplificatrices : milieu homogène ou localement homogène, grandeurs moyennes, géométrie idéalisée, absence de couplages secondaires. En situation réelle, les frottements, gradients, turbulences, imperfections géométriques ou effets non linéaires peuvent modifier le résultat.
\end{attentionbox}

\begin{approfondissementbox}[Niveau prépa / Bac+3]
Le phénomène demande une moyenne statistique sur les orientations cristallines.
On peut souvent améliorer le modèle en remplaçant une relation algébrique simple par une équation différentielle, une loi locale ou un bilan intégral. Cela permet de passer d'une explication qualitative à une prédiction quantitative.
\end{approfondissementbox}

\begin{exercicebox}[Exercice guidé]
Pourquoi les parhélies ne forment-ils pas toujours un cercle complet ?
\begin{enumerate}
  \item Identifier la grandeur principale.
  \item Donner son unité.
  \item Écrire la loi physique pertinente.
  \item Tester l'homogénéité de la relation.
  \item Interpréter le résultat dans le contexte.
\end{enumerate}
\end{exercicebox}

\begin{corrigebox}[Indication]
Parce que certains cristaux ont des orientations dominantes.
\end{corrigebox}


% ============================================================
\section{Partie — Atmosphère}
% ============================================================


% ------------------------------------------------------------
\subsection{9. Ciel bleu}
% ------------------------------------------------------------

\begin{exemplebox}[Observation concrète]
Le ciel est bleu en journée loin du Soleil.
\end{exemplebox}

\begin{methodebox}[Description qualitative]
Les molécules diffusent davantage les courtes longueurs d'onde.
La stratégie de compréhension est toujours la même : isoler la grandeur qui varie, identifier la loi physique, puis vérifier le sens du phénomène par un ordre de grandeur.
\end{methodebox}

\begin{theorembox}[Mise en équation niveau Terminale]
La relation centrale utilisée ici est :
\[
I_{\mathrm{diff}}\propto \frac{1}{\lambda^4}
\]
Cette formule n'est pas seulement un outil de calcul : elle permet de prévoir comment le phénomène change lorsqu'on modifie une grandeur expérimentale.
\end{theorembox}

\begin{odgbox}[Ordre de grandeur]
Le bleu est environ $(650/450)^4\simeq4{,}4$ fois plus diffusé que le rouge.
\end{odgbox}

\begin{attentionbox}[Limites du modèle]
Le modèle présenté repose sur des hypothèses simplificatrices : milieu homogène ou localement homogène, grandeurs moyennes, géométrie idéalisée, absence de couplages secondaires. En situation réelle, les frottements, gradients, turbulences, imperfections géométriques ou effets non linéaires peuvent modifier le résultat.
\end{attentionbox}

\begin{approfondissementbox}[Niveau prépa / Bac+3]
La diffusion Rayleigh provient de dipôles induits oscillants dans les molécules.
On peut souvent améliorer le modèle en remplaçant une relation algébrique simple par une équation différentielle, une loi locale ou un bilan intégral. Cela permet de passer d'une explication qualitative à une prédiction quantitative.
\end{approfondissementbox}

\begin{exercicebox}[Exercice guidé]
Comparer la diffusion du bleu et du rouge.
\begin{enumerate}
  \item Identifier la grandeur principale.
  \item Donner son unité.
  \item Écrire la loi physique pertinente.
  \item Tester l'homogénéité de la relation.
  \item Interpréter le résultat dans le contexte.
\end{enumerate}
\end{exercicebox}

\begin{corrigebox}[Indication]
Le bleu est plus diffusé, donc arrive à l'œil depuis toutes les directions.
\end{corrigebox}


% ------------------------------------------------------------
\subsection{10. Coucher de soleil rouge}
% ------------------------------------------------------------

\begin{exemplebox}[Observation concrète]
Le Soleil devient rouge près de l'horizon.
\end{exemplebox}

\begin{methodebox}[Description qualitative]
La lumière traverse une plus grande épaisseur d'atmosphère ; le bleu est diffusé hors du faisceau direct.
La stratégie de compréhension est toujours la même : isoler la grandeur qui varie, identifier la loi physique, puis vérifier le sens du phénomène par un ordre de grandeur.
\end{methodebox}

\begin{theorembox}[Mise en équation niveau Terminale]
La relation centrale utilisée ici est :
\[
I=I_0\e^{-\alpha(\lambda)L}
\]
Cette formule n'est pas seulement un outil de calcul : elle permet de prévoir comment le phénomène change lorsqu'on modifie une grandeur expérimentale.
\end{theorembox}

\begin{odgbox}[Ordre de grandeur]
Le trajet atmosphérique peut être plusieurs dizaines de fois plus long qu'au zénith.
\end{odgbox}

\begin{attentionbox}[Limites du modèle]
Le modèle présenté repose sur des hypothèses simplificatrices : milieu homogène ou localement homogène, grandeurs moyennes, géométrie idéalisée, absence de couplages secondaires. En situation réelle, les frottements, gradients, turbulences, imperfections géométriques ou effets non linéaires peuvent modifier le résultat.
\end{attentionbox}

\begin{approfondissementbox}[Niveau prépa / Bac+3]
Le coefficient $lpha$ dépend fortement de la longueur d'onde.
On peut souvent améliorer le modèle en remplaçant une relation algébrique simple par une équation différentielle, une loi locale ou un bilan intégral. Cela permet de passer d'une explication qualitative à une prédiction quantitative.
\end{approfondissementbox}

\begin{exercicebox}[Exercice guidé]
Pourquoi le Soleil n'est-il pas bleu au coucher ?
\begin{enumerate}
  \item Identifier la grandeur principale.
  \item Donner son unité.
  \item Écrire la loi physique pertinente.
  \item Tester l'homogénéité de la relation.
  \item Interpréter le résultat dans le contexte.
\end{enumerate}
\end{exercicebox}

\begin{corrigebox}[Indication]
Parce que le bleu est retiré du faisceau transmis par diffusion.
\end{corrigebox}


% ------------------------------------------------------------
\subsection{11. Nuages blancs}
% ------------------------------------------------------------

\begin{exemplebox}[Observation concrète]
Les nuages apparaissent blancs ou gris.
\end{exemplebox}

\begin{methodebox}[Description qualitative]
Les gouttelettes diffusent toutes les longueurs d'onde visibles de manière comparable.
La stratégie de compréhension est toujours la même : isoler la grandeur qui varie, identifier la loi physique, puis vérifier le sens du phénomène par un ordre de grandeur.
\end{methodebox}

\begin{theorembox}[Mise en équation niveau Terminale]
La relation centrale utilisée ici est :
\[
a\gg\lambda\quad\Rightarrow\quad \text{diffusion de Mie}
\]
Cette formule n'est pas seulement un outil de calcul : elle permet de prévoir comment le phénomène change lorsqu'on modifie une grandeur expérimentale.
\end{theorembox}

\begin{odgbox}[Ordre de grandeur]
Les gouttelettes ont des tailles micrométriques, plus grandes que les longueurs d'onde visibles.
\end{odgbox}

\begin{attentionbox}[Limites du modèle]
Le modèle présenté repose sur des hypothèses simplificatrices : milieu homogène ou localement homogène, grandeurs moyennes, géométrie idéalisée, absence de couplages secondaires. En situation réelle, les frottements, gradients, turbulences, imperfections géométriques ou effets non linéaires peuvent modifier le résultat.
\end{attentionbox}

\begin{approfondissementbox}[Niveau prépa / Bac+3]
La diffusion de Mie ne suit pas simplement une loi en $1/\lambda^4$.
On peut souvent améliorer le modèle en remplaçant une relation algébrique simple par une équation différentielle, une loi locale ou un bilan intégral. Cela permet de passer d'une explication qualitative à une prédiction quantitative.
\end{approfondissementbox}

\begin{exercicebox}[Exercice guidé]
Pourquoi un nuage épais devient-il gris ?
\begin{enumerate}
  \item Identifier la grandeur principale.
  \item Donner son unité.
  \item Écrire la loi physique pertinente.
  \item Tester l'homogénéité de la relation.
  \item Interpréter le résultat dans le contexte.
\end{enumerate}
\end{exercicebox}

\begin{corrigebox}[Indication]
Il diffuse beaucoup mais transmet peu de lumière : l'intensité arrivant dessous diminue.
\end{corrigebox}


% ------------------------------------------------------------
\subsection{12. Brouillard}
% ------------------------------------------------------------

\begin{exemplebox}[Observation concrète]
La visibilité chute lorsque l'air contient de fines gouttelettes.
\end{exemplebox}

\begin{methodebox}[Description qualitative]
Les gouttelettes diffusent la lumière dans toutes les directions, ce qui diminue le contraste.
La stratégie de compréhension est toujours la même : isoler la grandeur qui varie, identifier la loi physique, puis vérifier le sens du phénomène par un ordre de grandeur.
\end{methodebox}

\begin{theorembox}[Mise en équation niveau Terminale]
La relation centrale utilisée ici est :
\[
I(L)=I_0\e^{-\beta L}
\]
Cette formule n'est pas seulement un outil de calcul : elle permet de prévoir comment le phénomène change lorsqu'on modifie une grandeur expérimentale.
\end{theorembox}

\begin{odgbox}[Ordre de grandeur]
Une visibilité de quelques dizaines de mètres correspond à une forte atténuation.
\end{odgbox}

\begin{attentionbox}[Limites du modèle]
Le modèle présenté repose sur des hypothèses simplificatrices : milieu homogène ou localement homogène, grandeurs moyennes, géométrie idéalisée, absence de couplages secondaires. En situation réelle, les frottements, gradients, turbulences, imperfections géométriques ou effets non linéaires peuvent modifier le résultat.
\end{attentionbox}

\begin{approfondissementbox}[Niveau prépa / Bac+3]
Le transport radiatif décrit l'atténuation et la diffusion multiple.
On peut souvent améliorer le modèle en remplaçant une relation algébrique simple par une équation différentielle, une loi locale ou un bilan intégral. Cela permet de passer d'une explication qualitative à une prédiction quantitative.
\end{approfondissementbox}

\begin{exercicebox}[Exercice guidé]
Relier brouillard et loi exponentielle d'atténuation.
\begin{enumerate}
  \item Identifier la grandeur principale.
  \item Donner son unité.
  \item Écrire la loi physique pertinente.
  \item Tester l'homogénéité de la relation.
  \item Interpréter le résultat dans le contexte.
\end{enumerate}
\end{exercicebox}

\begin{corrigebox}[Indication]
La lumière directe diminue exponentiellement avec la distance parcourue.
\end{corrigebox}


% ------------------------------------------------------------
\subsection{13. Diffusion de la lumière dans l'atmosphère}
% ------------------------------------------------------------

\begin{exemplebox}[Observation concrète]
La lumière solaire est redistribuée dans le ciel.
\end{exemplebox}

\begin{methodebox}[Description qualitative]
Molécules, aérosols et gouttelettes diffusent la lumière.
La stratégie de compréhension est toujours la même : isoler la grandeur qui varie, identifier la loi physique, puis vérifier le sens du phénomène par un ordre de grandeur.
\end{methodebox}

\begin{theorembox}[Mise en équation niveau Terminale]
La relation centrale utilisée ici est :
\[
I_{\mathrm{diff}}=I_{\mathrm{Rayleigh}}+I_{\mathrm{Mie}}
\]
Cette formule n'est pas seulement un outil de calcul : elle permet de prévoir comment le phénomène change lorsqu'on modifie une grandeur expérimentale.
\end{theorembox}

\begin{odgbox}[Ordre de grandeur]
La couleur dépend de la taille des diffuseurs.
\end{odgbox}

\begin{attentionbox}[Limites du modèle]
Le modèle présenté repose sur des hypothèses simplificatrices : milieu homogène ou localement homogène, grandeurs moyennes, géométrie idéalisée, absence de couplages secondaires. En situation réelle, les frottements, gradients, turbulences, imperfections géométriques ou effets non linéaires peuvent modifier le résultat.
\end{attentionbox}

\begin{approfondissementbox}[Niveau prépa / Bac+3]
Une description complète nécessite une équation de transfert radiatif.
On peut souvent améliorer le modèle en remplaçant une relation algébrique simple par une équation différentielle, une loi locale ou un bilan intégral. Cela permet de passer d'une explication qualitative à une prédiction quantitative.
\end{approfondissementbox}

\begin{exercicebox}[Exercice guidé]
Classer molécules, poussières et gouttes selon leur taille relative à $\lambda$.
\begin{enumerate}
  \item Identifier la grandeur principale.
  \item Donner son unité.
  \item Écrire la loi physique pertinente.
  \item Tester l'homogénéité de la relation.
  \item Interpréter le résultat dans le contexte.
\end{enumerate}
\end{exercicebox}

\begin{corrigebox}[Indication]
Molécules : Rayleigh ; gouttes : Mie ; grosses particules : diffusion moins sélective.
\end{corrigebox}


% ============================================================
\section{Partie — Électricité atmosphérique}
% ============================================================


% ------------------------------------------------------------
\subsection{14. Éclair}
% ------------------------------------------------------------

\begin{exemplebox}[Observation concrète]
Une décharge lumineuse intense traverse l'air pendant un orage.
\end{exemplebox}

\begin{methodebox}[Description qualitative]
Un champ électrique très fort ionise l'air et crée un plasma conducteur.
La stratégie de compréhension est toujours la même : isoler la grandeur qui varie, identifier la loi physique, puis vérifier le sens du phénomène par un ordre de grandeur.
\end{methodebox}

\begin{theorembox}[Mise en équation niveau Terminale]
La relation centrale utilisée ici est :
\[
E_{\mathrm{rupture}}\sim \SI{3e6}{V.m^{-1}}
\]
Cette formule n'est pas seulement un outil de calcul : elle permet de prévoir comment le phénomène change lorsqu'on modifie une grandeur expérimentale.
\end{theorembox}

\begin{odgbox}[Ordre de grandeur]
Un éclair peut transporter des dizaines de kiloampères.
\end{odgbox}

\begin{attentionbox}[Limites du modèle]
Le modèle présenté repose sur des hypothèses simplificatrices : milieu homogène ou localement homogène, grandeurs moyennes, géométrie idéalisée, absence de couplages secondaires. En situation réelle, les frottements, gradients, turbulences, imperfections géométriques ou effets non linéaires peuvent modifier le résultat.
\end{attentionbox}

\begin{approfondissementbox}[Niveau prépa / Bac+3]
La décharge implique plasma, ionisation, échauffement et propagation non linéaire.
On peut souvent améliorer le modèle en remplaçant une relation algébrique simple par une équation différentielle, une loi locale ou un bilan intégral. Cela permet de passer d'une explication qualitative à une prédiction quantitative.
\end{approfondissementbox}

\begin{exercicebox}[Exercice guidé]
Estimer la tension nécessaire pour claquer un mètre d'air sec.
\begin{enumerate}
  \item Identifier la grandeur principale.
  \item Donner son unité.
  \item Écrire la loi physique pertinente.
  \item Tester l'homogénéité de la relation.
  \item Interpréter le résultat dans le contexte.
\end{enumerate}
\end{exercicebox}

\begin{corrigebox}[Indication]
Avec $\SI{3e6}{V.m^{-1}}$, il faut environ $\SI{3}{MV}$ par mètre.
\end{corrigebox}


% ------------------------------------------------------------
\subsection{15. Tonnerre}
% ------------------------------------------------------------

\begin{exemplebox}[Observation concrète]
Un bruit violent suit l'éclair.
\end{exemplebox}

\begin{methodebox}[Description qualitative]
L'air chauffé brutalement se dilate et crée une onde de pression.
La stratégie de compréhension est toujours la même : isoler la grandeur qui varie, identifier la loi physique, puis vérifier le sens du phénomène par un ordre de grandeur.
\end{methodebox}

\begin{theorembox}[Mise en équation niveau Terminale]
La relation centrale utilisée ici est :
\[
d\simeq c_s\Delta t
\]
Cette formule n'est pas seulement un outil de calcul : elle permet de prévoir comment le phénomène change lorsqu'on modifie une grandeur expérimentale.
\end{theorembox}

\begin{odgbox}[Ordre de grandeur]
Un délai de 3 s correspond à environ 1 km.
\end{odgbox}

\begin{attentionbox}[Limites du modèle]
Le modèle présenté repose sur des hypothèses simplificatrices : milieu homogène ou localement homogène, grandeurs moyennes, géométrie idéalisée, absence de couplages secondaires. En situation réelle, les frottements, gradients, turbulences, imperfections géométriques ou effets non linéaires peuvent modifier le résultat.
\end{attentionbox}

\begin{approfondissementbox}[Niveau prépa / Bac+3]
Le tonnerre est une onde acoustique issue d'un canal étendu, pas d'un point unique.
On peut souvent améliorer le modèle en remplaçant une relation algébrique simple par une équation différentielle, une loi locale ou un bilan intégral. Cela permet de passer d'une explication qualitative à une prédiction quantitative.
\end{approfondissementbox}

\begin{exercicebox}[Exercice guidé]
Si le tonnerre arrive 6 s après l'éclair, estimer la distance.
\begin{enumerate}
  \item Identifier la grandeur principale.
  \item Donner son unité.
  \item Écrire la loi physique pertinente.
  \item Tester l'homogénéité de la relation.
  \item Interpréter le résultat dans le contexte.
\end{enumerate}
\end{exercicebox}

\begin{corrigebox}[Indication]
$d\simeq340	imes6\simeq\SI{2040}{m}$.
\end{corrigebox}


% ------------------------------------------------------------
\subsection{16. Estimation de la distance d'un orage}
% ------------------------------------------------------------

\begin{exemplebox}[Observation concrète]
On compte les secondes entre éclair et tonnerre.
\end{exemplebox}

\begin{methodebox}[Description qualitative]
La lumière arrive quasi instantanément à notre échelle ; le son met plus longtemps.
La stratégie de compréhension est toujours la même : isoler la grandeur qui varie, identifier la loi physique, puis vérifier le sens du phénomène par un ordre de grandeur.
\end{methodebox}

\begin{theorembox}[Mise en équation niveau Terminale]
La relation centrale utilisée ici est :
\[
d=c_s\Delta t
\]
Cette formule n'est pas seulement un outil de calcul : elle permet de prévoir comment le phénomène change lorsqu'on modifie une grandeur expérimentale.
\end{theorembox}

\begin{odgbox}[Ordre de grandeur]
Une règle pratique : 3 s par km.
\end{odgbox}

\begin{attentionbox}[Limites du modèle]
Le modèle présenté repose sur des hypothèses simplificatrices : milieu homogène ou localement homogène, grandeurs moyennes, géométrie idéalisée, absence de couplages secondaires. En situation réelle, les frottements, gradients, turbulences, imperfections géométriques ou effets non linéaires peuvent modifier le résultat.
\end{attentionbox}

\begin{approfondissementbox}[Niveau prépa / Bac+3]
La vitesse du son dépend de la température : $c_s\simeq331+0{,}6T$ en m/s.
On peut souvent améliorer le modèle en remplaçant une relation algébrique simple par une équation différentielle, une loi locale ou un bilan intégral. Cela permet de passer d'une explication qualitative à une prédiction quantitative.
\end{approfondissementbox}

\begin{exercicebox}[Exercice guidé]
Donner la distance pour $\Delta t=\SI{9}{s}$.
\begin{enumerate}
  \item Identifier la grandeur principale.
  \item Donner son unité.
  \item Écrire la loi physique pertinente.
  \item Tester l'homogénéité de la relation.
  \item Interpréter le résultat dans le contexte.
\end{enumerate}
\end{exercicebox}

\begin{corrigebox}[Indication]
Environ $\SI{3}{km}$.
\end{corrigebox}


% ============================================================
\section{Partie — Thermique}
% ============================================================


% ------------------------------------------------------------
\subsection{17. Givre}
% ------------------------------------------------------------

\begin{exemplebox}[Observation concrète]
Une couche blanche se forme sur des surfaces froides.
\end{exemplebox}

\begin{methodebox}[Description qualitative]
La vapeur d'eau passe directement à l'état solide lorsque la surface est sous $0^\circ$C.
La stratégie de compréhension est toujours la même : isoler la grandeur qui varie, identifier la loi physique, puis vérifier le sens du phénomène par un ordre de grandeur.
\end{methodebox}

\begin{theorembox}[Mise en équation niveau Terminale]
La relation centrale utilisée ici est :
\[
\text{vapeur}\rightarrow\text{solide}
\]
Cette formule n'est pas seulement un outil de calcul : elle permet de prévoir comment le phénomène change lorsqu'on modifie une grandeur expérimentale.
\end{theorembox}

\begin{odgbox}[Ordre de grandeur]
Le pare-brise peut être sous la température de l'air par rayonnement nocturne.
\end{odgbox}

\begin{attentionbox}[Limites du modèle]
Le modèle présenté repose sur des hypothèses simplificatrices : milieu homogène ou localement homogène, grandeurs moyennes, géométrie idéalisée, absence de couplages secondaires. En situation réelle, les frottements, gradients, turbulences, imperfections géométriques ou effets non linéaires peuvent modifier le résultat.
\end{attentionbox}

\begin{approfondissementbox}[Niveau prépa / Bac+3]
Le bilan thermique de surface inclut rayonnement, convection et chaleur latente.
On peut souvent améliorer le modèle en remplaçant une relation algébrique simple par une équation différentielle, une loi locale ou un bilan intégral. Cela permet de passer d'une explication qualitative à une prédiction quantitative.
\end{approfondissementbox}

\begin{exercicebox}[Exercice guidé]
Pourquoi le givre se forme-t-il surtout la nuit ?
\begin{enumerate}
  \item Identifier la grandeur principale.
  \item Donner son unité.
  \item Écrire la loi physique pertinente.
  \item Tester l'homogénéité de la relation.
  \item Interpréter le résultat dans le contexte.
\end{enumerate}
\end{exercicebox}

\begin{corrigebox}[Indication]
La surface se refroidit par rayonnement vers le ciel.
\end{corrigebox}


% ------------------------------------------------------------
\subsection{18. Rosée}
% ------------------------------------------------------------

\begin{exemplebox}[Observation concrète]
Des gouttes se forment le matin sur l'herbe.
\end{exemplebox}

\begin{methodebox}[Description qualitative]
L'air au contact d'une surface froide atteint la saturation et l'eau condense.
La stratégie de compréhension est toujours la même : isoler la grandeur qui varie, identifier la loi physique, puis vérifier le sens du phénomène par un ordre de grandeur.
\end{methodebox}

\begin{theorembox}[Mise en équation niveau Terminale]
La relation centrale utilisée ici est :
\[
T_{\mathrm{surface}}<T_{\mathrm{rosée}}
\]
Cette formule n'est pas seulement un outil de calcul : elle permet de prévoir comment le phénomène change lorsqu'on modifie une grandeur expérimentale.
\end{theorembox}

\begin{odgbox}[Ordre de grandeur]
Le point de rosée dépend de l'humidité relative.
\end{odgbox}

\begin{attentionbox}[Limites du modèle]
Le modèle présenté repose sur des hypothèses simplificatrices : milieu homogène ou localement homogène, grandeurs moyennes, géométrie idéalisée, absence de couplages secondaires. En situation réelle, les frottements, gradients, turbulences, imperfections géométriques ou effets non linéaires peuvent modifier le résultat.
\end{attentionbox}

\begin{approfondissementbox}[Niveau prépa / Bac+3]
La thermodynamique de l'air humide utilise la pression de vapeur saturante.
On peut souvent améliorer le modèle en remplaçant une relation algébrique simple par une équation différentielle, une loi locale ou un bilan intégral. Cela permet de passer d'une explication qualitative à une prédiction quantitative.
\end{approfondissementbox}

\begin{exercicebox}[Exercice guidé]
Pourquoi la rosée est-elle plus fréquente par ciel clair ?
\begin{enumerate}
  \item Identifier la grandeur principale.
  \item Donner son unité.
  \item Écrire la loi physique pertinente.
  \item Tester l'homogénéité de la relation.
  \item Interpréter le résultat dans le contexte.
\end{enumerate}
\end{exercicebox}

\begin{corrigebox}[Indication]
Le refroidissement radiatif nocturne est plus efficace.
\end{corrigebox}


% ------------------------------------------------------------
\subsection{19. Buée sur une vitre}
% ------------------------------------------------------------

\begin{exemplebox}[Observation concrète]
Une vitre froide se couvre de gouttelettes.
\end{exemplebox}

\begin{methodebox}[Description qualitative]
L'air humide chaud au contact se refroidit et dépasse la saturation.
La stratégie de compréhension est toujours la même : isoler la grandeur qui varie, identifier la loi physique, puis vérifier le sens du phénomène par un ordre de grandeur.
\end{methodebox}

\begin{theorembox}[Mise en équation niveau Terminale]
La relation centrale utilisée ici est :
\[
\mathrm{HR}=100\%\quad\Rightarrow\quad \text{condensation}
\]
Cette formule n'est pas seulement un outil de calcul : elle permet de prévoir comment le phénomène change lorsqu'on modifie une grandeur expérimentale.
\end{theorembox}

\begin{odgbox}[Ordre de grandeur]
Une vitre simple est plus froide qu'un double vitrage.
\end{odgbox}

\begin{attentionbox}[Limites du modèle]
Le modèle présenté repose sur des hypothèses simplificatrices : milieu homogène ou localement homogène, grandeurs moyennes, géométrie idéalisée, absence de couplages secondaires. En situation réelle, les frottements, gradients, turbulences, imperfections géométriques ou effets non linéaires peuvent modifier le résultat.
\end{attentionbox}

\begin{approfondissementbox}[Niveau prépa / Bac+3]
La diffusion de vapeur et le transfert thermique couplent deux équations.
On peut souvent améliorer le modèle en remplaçant une relation algébrique simple par une équation différentielle, une loi locale ou un bilan intégral. Cela permet de passer d'une explication qualitative à une prédiction quantitative.
\end{approfondissementbox}

\begin{exercicebox}[Exercice guidé]
Pourquoi souffler sur une vitre peut faire de la buée ?
\begin{enumerate}
  \item Identifier la grandeur principale.
  \item Donner son unité.
  \item Écrire la loi physique pertinente.
  \item Tester l'homogénéité de la relation.
  \item Interpréter le résultat dans le contexte.
\end{enumerate}
\end{exercicebox}

\begin{corrigebox}[Indication]
L'air expiré est chaud et humide ; en refroidissant il condense.
\end{corrigebox}


% ------------------------------------------------------------
\subsection{20. Vapeur visible au-dessus d'une casserole}
% ------------------------------------------------------------

\begin{exemplebox}[Observation concrète]
Un panache blanc apparaît au-dessus de l'eau chaude.
\end{exemplebox}

\begin{methodebox}[Description qualitative]
La vapeur d'eau invisible condense en microgouttelettes dans l'air plus froid.
La stratégie de compréhension est toujours la même : isoler la grandeur qui varie, identifier la loi physique, puis vérifier le sens du phénomène par un ordre de grandeur.
\end{methodebox}

\begin{theorembox}[Mise en équation niveau Terminale]
La relation centrale utilisée ici est :
\[
\text{vapeur invisible}\rightarrow\text{aérosol liquide visible}
\]
Cette formule n'est pas seulement un outil de calcul : elle permet de prévoir comment le phénomène change lorsqu'on modifie une grandeur expérimentale.
\end{theorembox}

\begin{odgbox}[Ordre de grandeur]
La vapeur d'eau pure est invisible.
\end{odgbox}

\begin{attentionbox}[Limites du modèle]
Le modèle présenté repose sur des hypothèses simplificatrices : milieu homogène ou localement homogène, grandeurs moyennes, géométrie idéalisée, absence de couplages secondaires. En situation réelle, les frottements, gradients, turbulences, imperfections géométriques ou effets non linéaires peuvent modifier le résultat.
\end{attentionbox}

\begin{approfondissementbox}[Niveau prépa / Bac+3]
La nucléation dépend de la sursaturation et des particules présentes.
On peut souvent améliorer le modèle en remplaçant une relation algébrique simple par une équation différentielle, une loi locale ou un bilan intégral. Cela permet de passer d'une explication qualitative à une prédiction quantitative.
\end{approfondissementbox}

\begin{exercicebox}[Exercice guidé]
Pourquoi le nuage blanc n'est-il pas exactement à la surface ?
\begin{enumerate}
  \item Identifier la grandeur principale.
  \item Donner son unité.
  \item Écrire la loi physique pertinente.
  \item Tester l'homogénéité de la relation.
  \item Interpréter le résultat dans le contexte.
\end{enumerate}
\end{exercicebox}

\begin{corrigebox}[Indication]
Il faut que la vapeur se refroidisse suffisamment pour condenser.
\end{corrigebox}


% ------------------------------------------------------------
\subsection{21. Ébullition de l'eau}
% ------------------------------------------------------------

\begin{exemplebox}[Observation concrète]
Des bulles se forment dans tout le liquide.
\end{exemplebox}

\begin{methodebox}[Description qualitative]
La pression de vapeur saturante atteint la pression extérieure.
La stratégie de compréhension est toujours la même : isoler la grandeur qui varie, identifier la loi physique, puis vérifier le sens du phénomène par un ordre de grandeur.
\end{methodebox}

\begin{theorembox}[Mise en équation niveau Terminale]
La relation centrale utilisée ici est :
\[
P_{\mathrm{sat}}(T_{\mathrm{éb}})=P_{\mathrm{ext}}
\]
Cette formule n'est pas seulement un outil de calcul : elle permet de prévoir comment le phénomène change lorsqu'on modifie une grandeur expérimentale.
\end{theorembox}

\begin{odgbox}[Ordre de grandeur]
À pression atmosphérique, l'eau bout vers $100^\circ$C.
\end{odgbox}

\begin{attentionbox}[Limites du modèle]
Le modèle présenté repose sur des hypothèses simplificatrices : milieu homogène ou localement homogène, grandeurs moyennes, géométrie idéalisée, absence de couplages secondaires. En situation réelle, les frottements, gradients, turbulences, imperfections géométriques ou effets non linéaires peuvent modifier le résultat.
\end{attentionbox}

\begin{approfondissementbox}[Niveau prépa / Bac+3]
La relation de Clausius-Clapeyron relie $P_{\mathrm{sat}}$ à $T$.
On peut souvent améliorer le modèle en remplaçant une relation algébrique simple par une équation différentielle, une loi locale ou un bilan intégral. Cela permet de passer d'une explication qualitative à une prédiction quantitative.
\end{approfondissementbox}

\begin{exercicebox}[Exercice guidé]
Pourquoi l'ébullition n'est-elle pas seulement une évaporation de surface ?
\begin{enumerate}
  \item Identifier la grandeur principale.
  \item Donner son unité.
  \item Écrire la loi physique pertinente.
  \item Tester l'homogénéité de la relation.
  \item Interpréter le résultat dans le contexte.
\end{enumerate}
\end{exercicebox}

\begin{corrigebox}[Indication]
Parce que des bulles de vapeur se forment dans le volume.
\end{corrigebox}


% ------------------------------------------------------------
\subsection{22. Température d'ébullition en altitude}
% ------------------------------------------------------------

\begin{exemplebox}[Observation concrète]
L'eau bout à plus basse température en montagne.
\end{exemplebox}

\begin{methodebox}[Description qualitative]
La pression atmosphérique diminue avec l'altitude.
La stratégie de compréhension est toujours la même : isoler la grandeur qui varie, identifier la loi physique, puis vérifier le sens du phénomène par un ordre de grandeur.
\end{methodebox}

\begin{theorembox}[Mise en équation niveau Terminale]
La relation centrale utilisée ici est :
\[
P_{\mathrm{sat}}(T)=P_{\mathrm{atm}}(z)
\]
Cette formule n'est pas seulement un outil de calcul : elle permet de prévoir comment le phénomène change lorsqu'on modifie une grandeur expérimentale.
\end{theorembox}

\begin{odgbox}[Ordre de grandeur]
Au sommet de l'Everest, l'eau bout bien sous $100^\circ$C.
\end{odgbox}

\begin{attentionbox}[Limites du modèle]
Le modèle présenté repose sur des hypothèses simplificatrices : milieu homogène ou localement homogène, grandeurs moyennes, géométrie idéalisée, absence de couplages secondaires. En situation réelle, les frottements, gradients, turbulences, imperfections géométriques ou effets non linéaires peuvent modifier le résultat.
\end{attentionbox}

\begin{approfondissementbox}[Niveau prépa / Bac+3]
Un modèle barométrique donne $P(z)=P_0\e^{-z/H}$.
On peut souvent améliorer le modèle en remplaçant une relation algébrique simple par une équation différentielle, une loi locale ou un bilan intégral. Cela permet de passer d'une explication qualitative à une prédiction quantitative.
\end{approfondissementbox}

\begin{exercicebox}[Exercice guidé]
Pourquoi les pâtes cuisent-elles moins vite en altitude ?
\begin{enumerate}
  \item Identifier la grandeur principale.
  \item Donner son unité.
  \item Écrire la loi physique pertinente.
  \item Tester l'homogénéité de la relation.
  \item Interpréter le résultat dans le contexte.
\end{enumerate}
\end{exercicebox}

\begin{corrigebox}[Indication]
L'eau bouillante y est moins chaude.
\end{corrigebox}


% ------------------------------------------------------------
\subsection{23. Cocotte-minute}
% ------------------------------------------------------------

\begin{exemplebox}[Observation concrète]
Les aliments cuisent plus vite dans une cocotte fermée.
\end{exemplebox}

\begin{methodebox}[Description qualitative]
La pression augmente, donc la température d'ébullition augmente.
La stratégie de compréhension est toujours la même : isoler la grandeur qui varie, identifier la loi physique, puis vérifier le sens du phénomène par un ordre de grandeur.
\end{methodebox}

\begin{theorembox}[Mise en équation niveau Terminale]
La relation centrale utilisée ici est :
\[
P_{\mathrm{int}}>P_{\mathrm{atm}}\Rightarrow T_{\mathrm{éb}}>100^\circ C
\]
Cette formule n'est pas seulement un outil de calcul : elle permet de prévoir comment le phénomène change lorsqu'on modifie une grandeur expérimentale.
\end{theorembox}

\begin{odgbox}[Ordre de grandeur]
Une cocotte peut atteindre environ $120^\circ$C.
\end{odgbox}

\begin{attentionbox}[Limites du modèle]
Le modèle présenté repose sur des hypothèses simplificatrices : milieu homogène ou localement homogène, grandeurs moyennes, géométrie idéalisée, absence de couplages secondaires. En situation réelle, les frottements, gradients, turbulences, imperfections géométriques ou effets non linéaires peuvent modifier le résultat.
\end{attentionbox}

\begin{approfondissementbox}[Niveau prépa / Bac+3]
Le bilan énergétique inclut chaleur sensible et chaleur latente.
On peut souvent améliorer le modèle en remplaçant une relation algébrique simple par une équation différentielle, une loi locale ou un bilan intégral. Cela permet de passer d'une explication qualitative à une prédiction quantitative.
\end{approfondissementbox}

\begin{exercicebox}[Exercice guidé]
Pourquoi la soupape est-elle indispensable ?
\begin{enumerate}
  \item Identifier la grandeur principale.
  \item Donner son unité.
  \item Écrire la loi physique pertinente.
  \item Tester l'homogénéité de la relation.
  \item Interpréter le résultat dans le contexte.
\end{enumerate}
\end{exercicebox}

\begin{corrigebox}[Indication]
Elle limite la pression et évite un danger mécanique.
\end{corrigebox}


% ============================================================
\section{Partie — Climat}
% ============================================================


% ------------------------------------------------------------
\subsection{24. Effet de serre}
% ------------------------------------------------------------

\begin{exemplebox}[Observation concrète]
La Terre est plus chaude que sans atmosphère.
\end{exemplebox}

\begin{methodebox}[Description qualitative]
L'atmosphère laisse passer une partie du visible et absorbe l'infrarouge émis par la surface.
La stratégie de compréhension est toujours la même : isoler la grandeur qui varie, identifier la loi physique, puis vérifier le sens du phénomène par un ordre de grandeur.
\end{methodebox}

\begin{theorembox}[Mise en équation niveau Terminale]
La relation centrale utilisée ici est :
\[
P=\sigma S T^4
\]
Cette formule n'est pas seulement un outil de calcul : elle permet de prévoir comment le phénomène change lorsqu'on modifie une grandeur expérimentale.
\end{theorembox}

\begin{odgbox}[Ordre de grandeur]
Sans effet de serre naturel, la température moyenne serait beaucoup plus basse.
\end{odgbox}

\begin{attentionbox}[Limites du modèle]
Le modèle présenté repose sur des hypothèses simplificatrices : milieu homogène ou localement homogène, grandeurs moyennes, géométrie idéalisée, absence de couplages secondaires. En situation réelle, les frottements, gradients, turbulences, imperfections géométriques ou effets non linéaires peuvent modifier le résultat.
\end{attentionbox}

\begin{approfondissementbox}[Niveau prépa / Bac+3]
Un modèle à une couche donne un équilibre radiatif couplé surface-atmosphère.
On peut souvent améliorer le modèle en remplaçant une relation algébrique simple par une équation différentielle, une loi locale ou un bilan intégral. Cela permet de passer d'une explication qualitative à une prédiction quantitative.
\end{approfondissementbox}

\begin{exercicebox}[Exercice guidé]
Pourquoi le CO2 agit-il surtout sur l'infrarouge ?
\begin{enumerate}
  \item Identifier la grandeur principale.
  \item Donner son unité.
  \item Écrire la loi physique pertinente.
  \item Tester l'homogénéité de la relation.
  \item Interpréter le résultat dans le contexte.
\end{enumerate}
\end{exercicebox}

\begin{corrigebox}[Indication]
Ses modes vibrationnels absorbent certaines longueurs d'onde IR.
\end{corrigebox}


% ============================================================
\section{Partie — Thermique}
% ============================================================


% ------------------------------------------------------------
\subsection{25. Route qui chauffe au soleil}
% ------------------------------------------------------------

\begin{exemplebox}[Observation concrète]
L'asphalte devient brûlant en été.
\end{exemplebox}

\begin{methodebox}[Description qualitative]
Il absorbe fortement le rayonnement solaire et a une faible réflectivité.
La stratégie de compréhension est toujours la même : isoler la grandeur qui varie, identifier la loi physique, puis vérifier le sens du phénomène par un ordre de grandeur.
\end{methodebox}

\begin{theorembox}[Mise en équation niveau Terminale]
La relation centrale utilisée ici est :
\[
Q=mc\Delta T
\]
Cette formule n'est pas seulement un outil de calcul : elle permet de prévoir comment le phénomène change lorsqu'on modifie une grandeur expérimentale.
\end{theorembox}

\begin{odgbox}[Ordre de grandeur]
Une surface sombre peut dépasser largement la température de l'air.
\end{odgbox}

\begin{attentionbox}[Limites du modèle]
Le modèle présenté repose sur des hypothèses simplificatrices : milieu homogène ou localement homogène, grandeurs moyennes, géométrie idéalisée, absence de couplages secondaires. En situation réelle, les frottements, gradients, turbulences, imperfections géométriques ou effets non linéaires peuvent modifier le résultat.
\end{attentionbox}

\begin{approfondissementbox}[Niveau prépa / Bac+3]
Le bilan de surface inclut absorption, émission IR, convection et conduction.
On peut souvent améliorer le modèle en remplaçant une relation algébrique simple par une équation différentielle, une loi locale ou un bilan intégral. Cela permet de passer d'une explication qualitative à une prédiction quantitative.
\end{approfondissementbox}

\begin{exercicebox}[Exercice guidé]
Pourquoi une route noire chauffe plus qu'un mur clair ?
\begin{enumerate}
  \item Identifier la grandeur principale.
  \item Donner son unité.
  \item Écrire la loi physique pertinente.
  \item Tester l'homogénéité de la relation.
  \item Interpréter le résultat dans le contexte.
\end{enumerate}
\end{exercicebox}

\begin{corrigebox}[Indication]
Son albédo est plus faible.
\end{corrigebox}


% ------------------------------------------------------------
\subsection{26. Sable brûlant}
% ------------------------------------------------------------

\begin{exemplebox}[Observation concrète]
Le sable peut brûler les pieds à la plage.
\end{exemplebox}

\begin{methodebox}[Description qualitative]
Il absorbe le rayonnement et conduit mal la chaleur vers la profondeur.
La stratégie de compréhension est toujours la même : isoler la grandeur qui varie, identifier la loi physique, puis vérifier le sens du phénomène par un ordre de grandeur.
\end{methodebox}

\begin{theorembox}[Mise en équation niveau Terminale]
La relation centrale utilisée ici est :
\[
\Phi=-\lambda_{\mathrm{th}}\frac{\Delta T}{e}
\]
Cette formule n'est pas seulement un outil de calcul : elle permet de prévoir comment le phénomène change lorsqu'on modifie une grandeur expérimentale.
\end{theorembox}

\begin{odgbox}[Ordre de grandeur]
La surface chauffe vite car peu d'énergie est évacuée.
\end{odgbox}

\begin{attentionbox}[Limites du modèle]
Le modèle présenté repose sur des hypothèses simplificatrices : milieu homogène ou localement homogène, grandeurs moyennes, géométrie idéalisée, absence de couplages secondaires. En situation réelle, les frottements, gradients, turbulences, imperfections géométriques ou effets non linéaires peuvent modifier le résultat.
\end{attentionbox}

\begin{approfondissementbox}[Niveau prépa / Bac+3]
L'équation de la chaleur décrit la diffusion thermique dans le sol.
On peut souvent améliorer le modèle en remplaçant une relation algébrique simple par une équation différentielle, une loi locale ou un bilan intégral. Cela permet de passer d'une explication qualitative à une prédiction quantitative.
\end{approfondissementbox}

\begin{exercicebox}[Exercice guidé]
Pourquoi quelques centimètres sous le sable il fait plus frais ?
\begin{enumerate}
  \item Identifier la grandeur principale.
  \item Donner son unité.
  \item Écrire la loi physique pertinente.
  \item Tester l'homogénéité de la relation.
  \item Interpréter le résultat dans le contexte.
\end{enumerate}
\end{exercicebox}

\begin{corrigebox}[Indication]
La diffusion thermique est lente et amortit les variations rapides.
\end{corrigebox}


% ------------------------------------------------------------
\subsection{27. Métal plus froid que bois}
% ------------------------------------------------------------

\begin{exemplebox}[Observation concrète]
À même température, le métal semble plus froid au toucher.
\end{exemplebox}

\begin{methodebox}[Description qualitative]
Le métal conduit plus vite la chaleur hors de la peau.
La stratégie de compréhension est toujours la même : isoler la grandeur qui varie, identifier la loi physique, puis vérifier le sens du phénomène par un ordre de grandeur.
\end{methodebox}

\begin{theorembox}[Mise en équation niveau Terminale]
La relation centrale utilisée ici est :
\[
\Phi \propto \lambda_{\mathrm{th}}\Delta T
\]
Cette formule n'est pas seulement un outil de calcul : elle permet de prévoir comment le phénomène change lorsqu'on modifie une grandeur expérimentale.
\end{theorembox}

\begin{odgbox}[Ordre de grandeur]
La conductivité thermique du métal est beaucoup plus grande.
\end{odgbox}

\begin{attentionbox}[Limites du modèle]
Le modèle présenté repose sur des hypothèses simplificatrices : milieu homogène ou localement homogène, grandeurs moyennes, géométrie idéalisée, absence de couplages secondaires. En situation réelle, les frottements, gradients, turbulences, imperfections géométriques ou effets non linéaires peuvent modifier le résultat.
\end{attentionbox}

\begin{approfondissementbox}[Niveau prépa / Bac+3]
La sensation dépend d'un problème transitoire de diffusion thermique.
On peut souvent améliorer le modèle en remplaçant une relation algébrique simple par une équation différentielle, une loi locale ou un bilan intégral. Cela permet de passer d'une explication qualitative à une prédiction quantitative.
\end{approfondissementbox}

\begin{exercicebox}[Exercice guidé]
Pourquoi le métal n'est-il pas réellement plus froid ?
\begin{enumerate}
  \item Identifier la grandeur principale.
  \item Donner son unité.
  \item Écrire la loi physique pertinente.
  \item Tester l'homogénéité de la relation.
  \item Interpréter le résultat dans le contexte.
\end{enumerate}
\end{exercicebox}

\begin{corrigebox}[Indication]
Il est à la même température, mais extrait la chaleur plus vite.
\end{corrigebox}


% ------------------------------------------------------------
\subsection{28. Refroidissement par évaporation}
% ------------------------------------------------------------

\begin{exemplebox}[Observation concrète]
Une peau mouillée ou un linge humide se refroidit.
\end{exemplebox}

\begin{methodebox}[Description qualitative]
L'évaporation consomme de l'énergie sous forme de chaleur latente.
La stratégie de compréhension est toujours la même : isoler la grandeur qui varie, identifier la loi physique, puis vérifier le sens du phénomène par un ordre de grandeur.
\end{methodebox}

\begin{theorembox}[Mise en équation niveau Terminale]
La relation centrale utilisée ici est :
\[
Q=mL_v
\]
Cette formule n'est pas seulement un outil de calcul : elle permet de prévoir comment le phénomène change lorsqu'on modifie une grandeur expérimentale.
\end{theorembox}

\begin{odgbox}[Ordre de grandeur]
La chaleur latente de vaporisation de l'eau est très grande.
\end{odgbox}

\begin{attentionbox}[Limites du modèle]
Le modèle présenté repose sur des hypothèses simplificatrices : milieu homogène ou localement homogène, grandeurs moyennes, géométrie idéalisée, absence de couplages secondaires. En situation réelle, les frottements, gradients, turbulences, imperfections géométriques ou effets non linéaires peuvent modifier le résultat.
\end{attentionbox}

\begin{approfondissementbox}[Niveau prépa / Bac+3]
Le flux évaporatif dépend de l'humidité, du vent et de la température.
On peut souvent améliorer le modèle en remplaçant une relation algébrique simple par une équation différentielle, une loi locale ou un bilan intégral. Cela permet de passer d'une explication qualitative à une prédiction quantitative.
\end{approfondissementbox}

\begin{exercicebox}[Exercice guidé]
Pourquoi le vent augmente-t-il le refroidissement ?
\begin{enumerate}
  \item Identifier la grandeur principale.
  \item Donner son unité.
  \item Écrire la loi physique pertinente.
  \item Tester l'homogénéité de la relation.
  \item Interpréter le résultat dans le contexte.
\end{enumerate}
\end{exercicebox}

\begin{corrigebox}[Indication]
Il renouvelle l'air humide près de la surface.
\end{corrigebox}


% ============================================================
\section{Partie — Biophysique}
% ============================================================


% ------------------------------------------------------------
\subsection{29. Transpiration}
% ------------------------------------------------------------

\begin{exemplebox}[Observation concrète]
Le corps transpire pour réguler sa température.
\end{exemplebox}

\begin{methodebox}[Description qualitative]
L'évaporation de la sueur emporte de l'énergie.
La stratégie de compréhension est toujours la même : isoler la grandeur qui varie, identifier la loi physique, puis vérifier le sens du phénomène par un ordre de grandeur.
\end{methodebox}

\begin{theorembox}[Mise en équation niveau Terminale]
La relation centrale utilisée ici est :
\[
P_{\mathrm{évap}}=\dot m L_v
\]
Cette formule n'est pas seulement un outil de calcul : elle permet de prévoir comment le phénomène change lorsqu'on modifie une grandeur expérimentale.
\end{theorembox}

\begin{odgbox}[Ordre de grandeur]
Quelques grammes d'eau évaporée retirent des kilojoules.
\end{odgbox}

\begin{attentionbox}[Limites du modèle]
Le modèle présenté repose sur des hypothèses simplificatrices : milieu homogène ou localement homogène, grandeurs moyennes, géométrie idéalisée, absence de couplages secondaires. En situation réelle, les frottements, gradients, turbulences, imperfections géométriques ou effets non linéaires peuvent modifier le résultat.
\end{attentionbox}

\begin{approfondissementbox}[Niveau prépa / Bac+3]
La régulation thermique peut se modéliser par une équation différentielle de bilan d'énergie.
On peut souvent améliorer le modèle en remplaçant une relation algébrique simple par une équation différentielle, une loi locale ou un bilan intégral. Cela permet de passer d'une explication qualitative à une prédiction quantitative.
\end{approfondissementbox}

\begin{exercicebox}[Exercice guidé]
Pourquoi transpirer dans un air humide refroidit moins ?
\begin{enumerate}
  \item Identifier la grandeur principale.
  \item Donner son unité.
  \item Écrire la loi physique pertinente.
  \item Tester l'homogénéité de la relation.
  \item Interpréter le résultat dans le contexte.
\end{enumerate}
\end{exercicebox}

\begin{corrigebox}[Indication]
L'évaporation est ralentie par l'humidité élevée.
\end{corrigebox}


% ============================================================
\section{Partie — Habitat}
% ============================================================


% ------------------------------------------------------------
\subsection{30. Isolation thermique d'une maison}
% ------------------------------------------------------------

\begin{exemplebox}[Observation concrète]
Une maison isolée perd moins de chaleur.
\end{exemplebox}

\begin{methodebox}[Description qualitative]
L'isolant augmente la résistance thermique.
La stratégie de compréhension est toujours la même : isoler la grandeur qui varie, identifier la loi physique, puis vérifier le sens du phénomène par un ordre de grandeur.
\end{methodebox}

\begin{theorembox}[Mise en équation niveau Terminale]
La relation centrale utilisée ici est :
\[
\Phi=\frac{\Delta T}{R_{\mathrm{th}}}
\]
Cette formule n'est pas seulement un outil de calcul : elle permet de prévoir comment le phénomène change lorsqu'on modifie une grandeur expérimentale.
\end{theorembox}

\begin{odgbox}[Ordre de grandeur]
Doubler la résistance thermique divise approximativement le flux par deux.
\end{odgbox}

\begin{attentionbox}[Limites du modèle]
Le modèle présenté repose sur des hypothèses simplificatrices : milieu homogène ou localement homogène, grandeurs moyennes, géométrie idéalisée, absence de couplages secondaires. En situation réelle, les frottements, gradients, turbulences, imperfections géométriques ou effets non linéaires peuvent modifier le résultat.
\end{attentionbox}

\begin{approfondissementbox}[Niveau prépa / Bac+3]
Les ponts thermiques imposent une modélisation multidimensionnelle.
On peut souvent améliorer le modèle en remplaçant une relation algébrique simple par une équation différentielle, une loi locale ou un bilan intégral. Cela permet de passer d'une explication qualitative à une prédiction quantitative.
\end{approfondissementbox}

\begin{exercicebox}[Exercice guidé]
Pourquoi l'isolation réduit-elle la puissance de chauffage ?
\begin{enumerate}
  \item Identifier la grandeur principale.
  \item Donner son unité.
  \item Écrire la loi physique pertinente.
  \item Tester l'homogénéité de la relation.
  \item Interpréter le résultat dans le contexte.
\end{enumerate}
\end{exercicebox}

\begin{corrigebox}[Indication]
Elle diminue le flux sortant pour un même écart de température.
\end{corrigebox}


% ------------------------------------------------------------
\subsection{31. Double vitrage}
% ------------------------------------------------------------

\begin{exemplebox}[Observation concrète]
Une fenêtre double vitrage limite les pertes et la buée.
\end{exemplebox}

\begin{methodebox}[Description qualitative]
Une lame de gaz peu conductrice sépare deux vitres.
La stratégie de compréhension est toujours la même : isoler la grandeur qui varie, identifier la loi physique, puis vérifier le sens du phénomène par un ordre de grandeur.
\end{methodebox}

\begin{theorembox}[Mise en équation niveau Terminale]
La relation centrale utilisée ici est :
\[
R_{\mathrm{tot}}=R_1+R_{\mathrm{gaz}}+R_2
\]
Cette formule n'est pas seulement un outil de calcul : elle permet de prévoir comment le phénomène change lorsqu'on modifie une grandeur expérimentale.
\end{theorembox}

\begin{odgbox}[Ordre de grandeur]
L'air immobile est un mauvais conducteur thermique.
\end{odgbox}

\begin{attentionbox}[Limites du modèle]
Le modèle présenté repose sur des hypothèses simplificatrices : milieu homogène ou localement homogène, grandeurs moyennes, géométrie idéalisée, absence de couplages secondaires. En situation réelle, les frottements, gradients, turbulences, imperfections géométriques ou effets non linéaires peuvent modifier le résultat.
\end{attentionbox}

\begin{approfondissementbox}[Niveau prépa / Bac+3]
Le rayonnement et la convection interne compliquent le modèle.
On peut souvent améliorer le modèle en remplaçant une relation algébrique simple par une équation différentielle, une loi locale ou un bilan intégral. Cela permet de passer d'une explication qualitative à une prédiction quantitative.
\end{approfondissementbox}

\begin{exercicebox}[Exercice guidé]
Pourquoi mettre un gaz rare peut améliorer le vitrage ?
\begin{enumerate}
  \item Identifier la grandeur principale.
  \item Donner son unité.
  \item Écrire la loi physique pertinente.
  \item Tester l'homogénéité de la relation.
  \item Interpréter le résultat dans le contexte.
\end{enumerate}
\end{exercicebox}

\begin{corrigebox}[Indication]
Certains gaz réduisent conduction et convection.
\end{corrigebox}


% ============================================================
\section{Partie — Technologie optique}
% ============================================================


% ------------------------------------------------------------
\subsection{32. Fibre optique}
% ------------------------------------------------------------

\begin{exemplebox}[Observation concrète]
La lumière est guidée sur des kilomètres dans un fil de verre.
\end{exemplebox}

\begin{methodebox}[Description qualitative]
Le cœur possède un indice supérieur à celui de la gaine ; la réflexion totale piège la lumière.
La stratégie de compréhension est toujours la même : isoler la grandeur qui varie, identifier la loi physique, puis vérifier le sens du phénomène par un ordre de grandeur.
\end{methodebox}

\begin{theorembox}[Mise en équation niveau Terminale]
La relation centrale utilisée ici est :
\[
\sin i_c=\frac{n_g}{n_c}
\]
Cette formule n'est pas seulement un outil de calcul : elle permet de prévoir comment le phénomène change lorsqu'on modifie une grandeur expérimentale.
\end{theorembox}

\begin{odgbox}[Ordre de grandeur]
Le temps dans 10 km vaut environ $t\simeq nL/cpprox\SI{50}{\micro s}$.
\end{odgbox}

\begin{attentionbox}[Limites du modèle]
Le modèle présenté repose sur des hypothèses simplificatrices : milieu homogène ou localement homogène, grandeurs moyennes, géométrie idéalisée, absence de couplages secondaires. En situation réelle, les frottements, gradients, turbulences, imperfections géométriques ou effets non linéaires peuvent modifier le résultat.
\end{attentionbox}

\begin{approfondissementbox}[Niveau prépa / Bac+3]
Les modes guidés se décrivent par résolution de l'équation de Helmholtz avec conditions aux limites.
On peut souvent améliorer le modèle en remplaçant une relation algébrique simple par une équation différentielle, une loi locale ou un bilan intégral. Cela permet de passer d'une explication qualitative à une prédiction quantitative.
\end{approfondissementbox}

\begin{exercicebox}[Exercice guidé]
Calculer l'angle limite pour $n_c=1{,}48$ et $n_g=1{,}46$.
\begin{enumerate}
  \item Identifier la grandeur principale.
  \item Donner son unité.
  \item Écrire la loi physique pertinente.
  \item Tester l'homogénéité de la relation.
  \item Interpréter le résultat dans le contexte.
\end{enumerate}
\end{exercicebox}

\begin{corrigebox}[Indication]
$i_c=rcsin(1{,}46/1{,}48)$.
\end{corrigebox}


% ------------------------------------------------------------
\subsection{33. Réflexion totale interne}
% ------------------------------------------------------------

\begin{exemplebox}[Observation concrète]
Un rayon dans l'eau ou le verre peut être totalement réfléchi.
\end{exemplebox}

\begin{methodebox}[Description qualitative]
Lorsque l'angle d'incidence dépasse l'angle limite, aucun rayon réfracté propagatif n'existe.
La stratégie de compréhension est toujours la même : isoler la grandeur qui varie, identifier la loi physique, puis vérifier le sens du phénomène par un ordre de grandeur.
\end{methodebox}

\begin{theorembox}[Mise en équation niveau Terminale]
La relation centrale utilisée ici est :
\[
\sin i_c=\frac{n_2}{n_1},\quad n_1>n_2
\]
Cette formule n'est pas seulement un outil de calcul : elle permet de prévoir comment le phénomène change lorsqu'on modifie une grandeur expérimentale.
\end{theorembox}

\begin{odgbox}[Ordre de grandeur]
Dans l'eau vers l'air, $i_c\simeq48{,}8^\circ$.
\end{odgbox}

\begin{attentionbox}[Limites du modèle]
Le modèle présenté repose sur des hypothèses simplificatrices : milieu homogène ou localement homogène, grandeurs moyennes, géométrie idéalisée, absence de couplages secondaires. En situation réelle, les frottements, gradients, turbulences, imperfections géométriques ou effets non linéaires peuvent modifier le résultat.
\end{attentionbox}

\begin{approfondissementbox}[Niveau prépa / Bac+3]
Une onde évanescente existe tout de même dans le second milieu.
On peut souvent améliorer le modèle en remplaçant une relation algébrique simple par une équation différentielle, une loi locale ou un bilan intégral. Cela permet de passer d'une explication qualitative à une prédiction quantitative.
\end{approfondissementbox}

\begin{exercicebox}[Exercice guidé]
Pourquoi la réflexion totale exige-t-elle $n_1>n_2$ ?
\begin{enumerate}
  \item Identifier la grandeur principale.
  \item Donner son unité.
  \item Écrire la loi physique pertinente.
  \item Tester l'homogénéité de la relation.
  \item Interpréter le résultat dans le contexte.
\end{enumerate}
\end{exercicebox}

\begin{corrigebox}[Indication]
Il faut que $\sin r$ dépasserait 1 pour certains angles.
\end{corrigebox}


% ============================================================
\section{Partie — Optique}
% ============================================================


% ------------------------------------------------------------
\subsection{34. Lentilles et formation des images}
% ------------------------------------------------------------

\begin{exemplebox}[Observation concrète]
Une lentille peut former une image nette sur un écran.
\end{exemplebox}

\begin{methodebox}[Description qualitative]
La réfraction aux surfaces courbes fait converger ou diverger les rayons.
La stratégie de compréhension est toujours la même : isoler la grandeur qui varie, identifier la loi physique, puis vérifier le sens du phénomène par un ordre de grandeur.
\end{methodebox}

\begin{theorembox}[Mise en équation niveau Terminale]
La relation centrale utilisée ici est :
\[
\frac{1}{OA'}-\frac{1}{OA}=\frac{1}{f'}
\]
Cette formule n'est pas seulement un outil de calcul : elle permet de prévoir comment le phénomène change lorsqu'on modifie une grandeur expérimentale.
\end{theorembox}

\begin{odgbox}[Ordre de grandeur]
Une lentille de focale 10 cm a une vergence 10 dioptries.
\end{odgbox}

\begin{attentionbox}[Limites du modèle]
Le modèle présenté repose sur des hypothèses simplificatrices : milieu homogène ou localement homogène, grandeurs moyennes, géométrie idéalisée, absence de couplages secondaires. En situation réelle, les frottements, gradients, turbulences, imperfections géométriques ou effets non linéaires peuvent modifier le résultat.
\end{attentionbox}

\begin{approfondissementbox}[Niveau prépa / Bac+3]
L'équation des lentilles vient de l'approximation de Gauss.
On peut souvent améliorer le modèle en remplaçant une relation algébrique simple par une équation différentielle, une loi locale ou un bilan intégral. Cela permet de passer d'une explication qualitative à une prédiction quantitative.
\end{approfondissementbox}

\begin{exercicebox}[Exercice guidé]
Trouver l'image d'un objet à 30 cm d'une lentille de focale 10 cm.
\begin{enumerate}
  \item Identifier la grandeur principale.
  \item Donner son unité.
  \item Écrire la loi physique pertinente.
  \item Tester l'homogénéité de la relation.
  \item Interpréter le résultat dans le contexte.
\end{enumerate}
\end{exercicebox}

\begin{corrigebox}[Indication]
On applique la relation de conjugaison avec signes algébriques.
\end{corrigebox}


% ============================================================
\section{Partie — Optique physiologique}
% ============================================================


% ------------------------------------------------------------
\subsection{35. Œil humain}
% ------------------------------------------------------------

\begin{exemplebox}[Observation concrète]
L'œil forme une image sur la rétine.
\end{exemplebox}

\begin{methodebox}[Description qualitative]
Cornée et cristallin jouent le rôle de système optique convergent.
La stratégie de compréhension est toujours la même : isoler la grandeur qui varie, identifier la loi physique, puis vérifier le sens du phénomène par un ordre de grandeur.
\end{methodebox}

\begin{theorembox}[Mise en équation niveau Terminale]
La relation centrale utilisée ici est :
\[
\text{image nette sur la rétine}
\]
Cette formule n'est pas seulement un outil de calcul : elle permet de prévoir comment le phénomène change lorsqu'on modifie une grandeur expérimentale.
\end{theorembox}

\begin{odgbox}[Ordre de grandeur]
La distance focale effective est de l'ordre de quelques cm.
\end{odgbox}

\begin{attentionbox}[Limites du modèle]
Le modèle présenté repose sur des hypothèses simplificatrices : milieu homogène ou localement homogène, grandeurs moyennes, géométrie idéalisée, absence de couplages secondaires. En situation réelle, les frottements, gradients, turbulences, imperfections géométriques ou effets non linéaires peuvent modifier le résultat.
\end{attentionbox}

\begin{approfondissementbox}[Niveau prépa / Bac+3]
Le modèle réduit de l'œil remplace les dioptres par une lentille équivalente.
On peut souvent améliorer le modèle en remplaçant une relation algébrique simple par une équation différentielle, une loi locale ou un bilan intégral. Cela permet de passer d'une explication qualitative à une prédiction quantitative.
\end{approfondissementbox}

\begin{exercicebox}[Exercice guidé]
Pourquoi l'accommodation est-elle nécessaire ?
\begin{enumerate}
  \item Identifier la grandeur principale.
  \item Donner son unité.
  \item Écrire la loi physique pertinente.
  \item Tester l'homogénéité de la relation.
  \item Interpréter le résultat dans le contexte.
\end{enumerate}
\end{exercicebox}

\begin{corrigebox}[Indication]
Pour ajuster la puissance optique aux objets proches.
\end{corrigebox}


% ------------------------------------------------------------
\subsection{36. Myopie}
% ------------------------------------------------------------

\begin{exemplebox}[Observation concrète]
Un myope voit flou de loin.
\end{exemplebox}

\begin{methodebox}[Description qualitative]
L'image d'un objet lointain se forme avant la rétine.
La stratégie de compréhension est toujours la même : isoler la grandeur qui varie, identifier la loi physique, puis vérifier le sens du phénomène par un ordre de grandeur.
\end{methodebox}

\begin{theorembox}[Mise en équation niveau Terminale]
La relation centrale utilisée ici est :
\[
\text{correction par lentille divergente}
\]
Cette formule n'est pas seulement un outil de calcul : elle permet de prévoir comment le phénomène change lorsqu'on modifie une grandeur expérimentale.
\end{theorembox}

\begin{odgbox}[Ordre de grandeur]
Une correction de -2 dioptries correspond à une focale -0,5 m.
\end{odgbox}

\begin{attentionbox}[Limites du modèle]
Le modèle présenté repose sur des hypothèses simplificatrices : milieu homogène ou localement homogène, grandeurs moyennes, géométrie idéalisée, absence de couplages secondaires. En situation réelle, les frottements, gradients, turbulences, imperfections géométriques ou effets non linéaires peuvent modifier le résultat.
\end{attentionbox}

\begin{approfondissementbox}[Niveau prépa / Bac+3]
La puissance optique s'ajoute en approximation de lentilles minces accolées.
On peut souvent améliorer le modèle en remplaçant une relation algébrique simple par une équation différentielle, une loi locale ou un bilan intégral. Cela permet de passer d'une explication qualitative à une prédiction quantitative.
\end{approfondissementbox}

\begin{exercicebox}[Exercice guidé]
Pourquoi une lentille divergente corrige-t-elle la myopie ?
\begin{enumerate}
  \item Identifier la grandeur principale.
  \item Donner son unité.
  \item Écrire la loi physique pertinente.
  \item Tester l'homogénéité de la relation.
  \item Interpréter le résultat dans le contexte.
\end{enumerate}
\end{exercicebox}

\begin{corrigebox}[Indication]
Elle fait diverger les rayons avant l'œil.
\end{corrigebox}


% ------------------------------------------------------------
\subsection{37. Hypermétropie}
% ------------------------------------------------------------

\begin{exemplebox}[Observation concrète]
Un hypermétrope fatigue en vision proche.
\end{exemplebox}

\begin{methodebox}[Description qualitative]
L'image tend à se former derrière la rétine.
La stratégie de compréhension est toujours la même : isoler la grandeur qui varie, identifier la loi physique, puis vérifier le sens du phénomène par un ordre de grandeur.
\end{methodebox}

\begin{theorembox}[Mise en équation niveau Terminale]
La relation centrale utilisée ici est :
\[
\text{correction par lentille convergente}
\]
Cette formule n'est pas seulement un outil de calcul : elle permet de prévoir comment le phénomène change lorsqu'on modifie une grandeur expérimentale.
\end{theorembox}

\begin{odgbox}[Ordre de grandeur]
La correction est en dioptries positives.
\end{odgbox}

\begin{attentionbox}[Limites du modèle]
Le modèle présenté repose sur des hypothèses simplificatrices : milieu homogène ou localement homogène, grandeurs moyennes, géométrie idéalisée, absence de couplages secondaires. En situation réelle, les frottements, gradients, turbulences, imperfections géométriques ou effets non linéaires peuvent modifier le résultat.
\end{attentionbox}

\begin{approfondissementbox}[Niveau prépa / Bac+3]
Le punctum proximum est repoussé et dépend de l'accommodation.
On peut souvent améliorer le modèle en remplaçant une relation algébrique simple par une équation différentielle, une loi locale ou un bilan intégral. Cela permet de passer d'une explication qualitative à une prédiction quantitative.
\end{approfondissementbox}

\begin{exercicebox}[Exercice guidé]
Pourquoi une lentille convergente aide-t-elle ?
\begin{enumerate}
  \item Identifier la grandeur principale.
  \item Donner son unité.
  \item Écrire la loi physique pertinente.
  \item Tester l'homogénéité de la relation.
  \item Interpréter le résultat dans le contexte.
\end{enumerate}
\end{exercicebox}

\begin{corrigebox}[Indication]
Elle augmente la convergence avant l'œil.
\end{corrigebox}


% ============================================================
\section{Partie — Photographie}
% ============================================================


% ------------------------------------------------------------
\subsection{38. Appareil photo}
% ------------------------------------------------------------

\begin{exemplebox}[Observation concrète]
Un appareil forme une image sur un capteur.
\end{exemplebox}

\begin{methodebox}[Description qualitative]
Une lentille focalise les rayons, le diaphragme contrôle la lumière et la profondeur de champ.
La stratégie de compréhension est toujours la même : isoler la grandeur qui varie, identifier la loi physique, puis vérifier le sens du phénomène par un ordre de grandeur.
\end{methodebox}

\begin{theorembox}[Mise en équation niveau Terminale]
La relation centrale utilisée ici est :
\[
\frac1{p'}-\frac1p=\frac1f
\]
Cette formule n'est pas seulement un outil de calcul : elle permet de prévoir comment le phénomène change lorsqu'on modifie une grandeur expérimentale.
\end{theorembox}

\begin{odgbox}[Ordre de grandeur]
Un capteur plein format mesure environ 36 mm par 24 mm.
\end{odgbox}

\begin{attentionbox}[Limites du modèle]
Le modèle présenté repose sur des hypothèses simplificatrices : milieu homogène ou localement homogène, grandeurs moyennes, géométrie idéalisée, absence de couplages secondaires. En situation réelle, les frottements, gradients, turbulences, imperfections géométriques ou effets non linéaires peuvent modifier le résultat.
\end{attentionbox}

\begin{approfondissementbox}[Niveau prépa / Bac+3]
La profondeur de champ dépend du cercle de confusion et de l'ouverture numérique.
On peut souvent améliorer le modèle en remplaçant une relation algébrique simple par une équation différentielle, une loi locale ou un bilan intégral. Cela permet de passer d'une explication qualitative à une prédiction quantitative.
\end{approfondissementbox}

\begin{exercicebox}[Exercice guidé]
Pourquoi fermer le diaphragme augmente-t-il la profondeur de champ ?
\begin{enumerate}
  \item Identifier la grandeur principale.
  \item Donner son unité.
  \item Écrire la loi physique pertinente.
  \item Tester l'homogénéité de la relation.
  \item Interpréter le résultat dans le contexte.
\end{enumerate}
\end{exercicebox}

\begin{corrigebox}[Indication]
Les faisceaux sont moins ouverts, donc le flou géométrique diminue.
\end{corrigebox}


% ============================================================
\section{Partie — Optique}
% ============================================================


% ------------------------------------------------------------
\subsection{39. Loupe}
% ------------------------------------------------------------

\begin{exemplebox}[Observation concrète]
Une loupe agrandit un objet proche.
\end{exemplebox}

\begin{methodebox}[Description qualitative]
Elle permet d'observer sous un angle plus grand.
La stratégie de compréhension est toujours la même : isoler la grandeur qui varie, identifier la loi physique, puis vérifier le sens du phénomène par un ordre de grandeur.
\end{methodebox}

\begin{theorembox}[Mise en équation niveau Terminale]
La relation centrale utilisée ici est :
\[
G\simeq \frac{25\,\mathrm{cm}}{f'}
\]
Cette formule n'est pas seulement un outil de calcul : elle permet de prévoir comment le phénomène change lorsqu'on modifie une grandeur expérimentale.
\end{theorembox}

\begin{odgbox}[Ordre de grandeur]
Une loupe de focale 5 cm grossit environ 5 fois.
\end{odgbox}

\begin{attentionbox}[Limites du modèle]
Le modèle présenté repose sur des hypothèses simplificatrices : milieu homogène ou localement homogène, grandeurs moyennes, géométrie idéalisée, absence de couplages secondaires. En situation réelle, les frottements, gradients, turbulences, imperfections géométriques ou effets non linéaires peuvent modifier le résultat.
\end{attentionbox}

\begin{approfondissementbox}[Niveau prépa / Bac+3]
Le grossissement angulaire dépend du punctum proximum.
On peut souvent améliorer le modèle en remplaçant une relation algébrique simple par une équation différentielle, une loi locale ou un bilan intégral. Cela permet de passer d'une explication qualitative à une prédiction quantitative.
\end{approfondissementbox}

\begin{exercicebox}[Exercice guidé]
Pourquoi l'objet doit-il être proche du foyer ?
\begin{enumerate}
  \item Identifier la grandeur principale.
  \item Donner son unité.
  \item Écrire la loi physique pertinente.
  \item Tester l'homogénéité de la relation.
  \item Interpréter le résultat dans le contexte.
\end{enumerate}
\end{exercicebox}

\begin{corrigebox}[Indication]
Pour obtenir une image virtuelle agrandie confortable.
\end{corrigebox}


% ------------------------------------------------------------
\subsection{40. Microscope}
% ------------------------------------------------------------

\begin{exemplebox}[Observation concrète]
Un microscope observe des objets très petits.
\end{exemplebox}

\begin{methodebox}[Description qualitative]
Objectif et oculaire combinent deux grossissements.
La stratégie de compréhension est toujours la même : isoler la grandeur qui varie, identifier la loi physique, puis vérifier le sens du phénomène par un ordre de grandeur.
\end{methodebox}

\begin{theorembox}[Mise en équation niveau Terminale]
La relation centrale utilisée ici est :
\[
G_{\mathrm{total}}=G_{\mathrm{objectif}}G_{\mathrm{oculaire}}
\]
Cette formule n'est pas seulement un outil de calcul : elle permet de prévoir comment le phénomène change lorsqu'on modifie une grandeur expérimentale.
\end{theorembox}

\begin{odgbox}[Ordre de grandeur]
La résolution est limitée par la diffraction.
\end{odgbox}

\begin{attentionbox}[Limites du modèle]
Le modèle présenté repose sur des hypothèses simplificatrices : milieu homogène ou localement homogène, grandeurs moyennes, géométrie idéalisée, absence de couplages secondaires. En situation réelle, les frottements, gradients, turbulences, imperfections géométriques ou effets non linéaires peuvent modifier le résultat.
\end{attentionbox}

\begin{approfondissementbox}[Niveau prépa / Bac+3]
Le critère d'Abbe donne $d\simeq \lambda/(2\mathrm{ON})$.
On peut souvent améliorer le modèle en remplaçant une relation algébrique simple par une équation différentielle, une loi locale ou un bilan intégral. Cela permet de passer d'une explication qualitative à une prédiction quantitative.
\end{approfondissementbox}

\begin{exercicebox}[Exercice guidé]
Pourquoi augmenter sans cesse le grossissement ne suffit pas ?
\begin{enumerate}
  \item Identifier la grandeur principale.
  \item Donner son unité.
  \item Écrire la loi physique pertinente.
  \item Tester l'homogénéité de la relation.
  \item Interpréter le résultat dans le contexte.
\end{enumerate}
\end{exercicebox}

\begin{corrigebox}[Indication]
La résolution optique impose une limite.
\end{corrigebox}


% ------------------------------------------------------------
\subsection{41. Télescope}
% ------------------------------------------------------------

\begin{exemplebox}[Observation concrète]
Un télescope collecte la lumière d'objets lointains.
\end{exemplebox}

\begin{methodebox}[Description qualitative]
Un grand objectif augmente la lumière collectée et la résolution.
La stratégie de compréhension est toujours la même : isoler la grandeur qui varie, identifier la loi physique, puis vérifier le sens du phénomène par un ordre de grandeur.
\end{methodebox}

\begin{theorembox}[Mise en équation niveau Terminale]
La relation centrale utilisée ici est :
\[
\theta_{\min}\simeq1{,}22\frac{\lambda}{D}
\]
Cette formule n'est pas seulement un outil de calcul : elle permet de prévoir comment le phénomène change lorsqu'on modifie une grandeur expérimentale.
\end{theorembox}

\begin{odgbox}[Ordre de grandeur]
Un diamètre plus grand améliore la résolution angulaire.
\end{odgbox}

\begin{attentionbox}[Limites du modèle]
Le modèle présenté repose sur des hypothèses simplificatrices : milieu homogène ou localement homogène, grandeurs moyennes, géométrie idéalisée, absence de couplages secondaires. En situation réelle, les frottements, gradients, turbulences, imperfections géométriques ou effets non linéaires peuvent modifier le résultat.
\end{attentionbox}

\begin{approfondissementbox}[Niveau prépa / Bac+3]
La turbulence atmosphérique limite souvent plus que la diffraction.
On peut souvent améliorer le modèle en remplaçant une relation algébrique simple par une équation différentielle, une loi locale ou un bilan intégral. Cela permet de passer d'une explication qualitative à une prédiction quantitative.
\end{approfondissementbox}

\begin{exercicebox}[Exercice guidé]
Pourquoi les grands télescopes ont-ils de grands miroirs ?
\begin{enumerate}
  \item Identifier la grandeur principale.
  \item Donner son unité.
  \item Écrire la loi physique pertinente.
  \item Tester l'homogénéité de la relation.
  \item Interpréter le résultat dans le contexte.
\end{enumerate}
\end{exercicebox}

\begin{corrigebox}[Indication]
Pour collecter plus de lumière et réduire la limite de diffraction.
\end{corrigebox}


% ============================================================
\section{Partie — Optique ondulatoire}
% ============================================================


% ------------------------------------------------------------
\subsection{42. Diffraction par une fente}
% ------------------------------------------------------------

\begin{exemplebox}[Observation concrète]
Un faisceau s'élargit après une petite ouverture.
\end{exemplebox}

\begin{methodebox}[Description qualitative]
Lorsque la fente est comparable à la longueur d'onde, les rayons ne restent pas rectilignes.
La stratégie de compréhension est toujours la même : isoler la grandeur qui varie, identifier la loi physique, puis vérifier le sens du phénomène par un ordre de grandeur.
\end{methodebox}

\begin{theorembox}[Mise en équation niveau Terminale]
La relation centrale utilisée ici est :
\[
\theta\simeq\frac{\lambda}{a}
\]
Cette formule n'est pas seulement un outil de calcul : elle permet de prévoir comment le phénomène change lorsqu'on modifie une grandeur expérimentale.
\end{theorembox}

\begin{odgbox}[Ordre de grandeur]
Pour un cheveu de 80 micromètres et un laser rouge, l'angle est de l'ordre de $10^{-2}$ rad.
\end{odgbox}

\begin{attentionbox}[Limites du modèle]
Le modèle présenté repose sur des hypothèses simplificatrices : milieu homogène ou localement homogène, grandeurs moyennes, géométrie idéalisée, absence de couplages secondaires. En situation réelle, les frottements, gradients, turbulences, imperfections géométriques ou effets non linéaires peuvent modifier le résultat.
\end{attentionbox}

\begin{approfondissementbox}[Niveau prépa / Bac+3]
La figure résulte d'une transformée de Fourier de l'ouverture.
On peut souvent améliorer le modèle en remplaçant une relation algébrique simple par une équation différentielle, une loi locale ou un bilan intégral. Cela permet de passer d'une explication qualitative à une prédiction quantitative.
\end{approfondissementbox}

\begin{exercicebox}[Exercice guidé]
Pourquoi une fente plus petite donne-t-elle une tache plus large ?
\begin{enumerate}
  \item Identifier la grandeur principale.
  \item Donner son unité.
  \item Écrire la loi physique pertinente.
  \item Tester l'homogénéité de la relation.
  \item Interpréter le résultat dans le contexte.
\end{enumerate}
\end{exercicebox}

\begin{corrigebox}[Indication]
$	heta$ est inversement proportionnel à $a$.
\end{corrigebox}


% ------------------------------------------------------------
\subsection{43. Mesure d'un cheveu par diffraction}
% ------------------------------------------------------------

\begin{exemplebox}[Observation concrète]
Un cheveu placé devant un laser donne une tache centrale.
\end{exemplebox}

\begin{methodebox}[Description qualitative]
Par le principe de Babinet, un obstacle fin diffracte comme une fente de même largeur.
La stratégie de compréhension est toujours la même : isoler la grandeur qui varie, identifier la loi physique, puis vérifier le sens du phénomène par un ordre de grandeur.
\end{methodebox}

\begin{theorembox}[Mise en équation niveau Terminale]
La relation centrale utilisée ici est :
\[
a\simeq\frac{2\lambda D}{L}
\]
Cette formule n'est pas seulement un outil de calcul : elle permet de prévoir comment le phénomène change lorsqu'on modifie une grandeur expérimentale.
\end{theorembox}

\begin{odgbox}[Ordre de grandeur]
On peut mesurer des dizaines de micromètres avec une règle.
\end{odgbox}

\begin{attentionbox}[Limites du modèle]
Le modèle présenté repose sur des hypothèses simplificatrices : milieu homogène ou localement homogène, grandeurs moyennes, géométrie idéalisée, absence de couplages secondaires. En situation réelle, les frottements, gradients, turbulences, imperfections géométriques ou effets non linéaires peuvent modifier le résultat.
\end{attentionbox}

\begin{approfondissementbox}[Niveau prépa / Bac+3]
La précision dépend surtout de la mesure de $L$ et $D$.
On peut souvent améliorer le modèle en remplaçant une relation algébrique simple par une équation différentielle, une loi locale ou un bilan intégral. Cela permet de passer d'une explication qualitative à une prédiction quantitative.
\end{approfondissementbox}

\begin{exercicebox}[Exercice guidé]
Calculer $a$ pour $\lambda=650$ nm, $D=2$ m, $L=3$ cm.
\begin{enumerate}
  \item Identifier la grandeur principale.
  \item Donner son unité.
  \item Écrire la loi physique pertinente.
  \item Tester l'homogénéité de la relation.
  \item Interpréter le résultat dans le contexte.
\end{enumerate}
\end{exercicebox}

\begin{corrigebox}[Indication]
$a\simeq 2\lambda D/L$.
\end{corrigebox}


% ------------------------------------------------------------
\subsection{44. Interférences de Young}
% ------------------------------------------------------------

\begin{exemplebox}[Observation concrète]
Des franges claires et sombres apparaissent après deux fentes.
\end{exemplebox}

\begin{methodebox}[Description qualitative]
Les ondes issues des fentes se superposent avec différence de marche.
La stratégie de compréhension est toujours la même : isoler la grandeur qui varie, identifier la loi physique, puis vérifier le sens du phénomène par un ordre de grandeur.
\end{methodebox}

\begin{theorembox}[Mise en équation niveau Terminale]
La relation centrale utilisée ici est :
\[
i=\frac{\lambda D}{a}
\]
Cette formule n'est pas seulement un outil de calcul : elle permet de prévoir comment le phénomène change lorsqu'on modifie une grandeur expérimentale.
\end{theorembox}

\begin{odgbox}[Ordre de grandeur]
L'interfrange millimétrique est mesurable sur écran.
\end{odgbox}

\begin{attentionbox}[Limites du modèle]
Le modèle présenté repose sur des hypothèses simplificatrices : milieu homogène ou localement homogène, grandeurs moyennes, géométrie idéalisée, absence de couplages secondaires. En situation réelle, les frottements, gradients, turbulences, imperfections géométriques ou effets non linéaires peuvent modifier le résultat.
\end{attentionbox}

\begin{approfondissementbox}[Niveau prépa / Bac+3]
La cohérence spatiale et temporelle conditionne la visibilité.
On peut souvent améliorer le modèle en remplaçant une relation algébrique simple par une équation différentielle, une loi locale ou un bilan intégral. Cela permet de passer d'une explication qualitative à une prédiction quantitative.
\end{approfondissementbox}

\begin{exercicebox}[Exercice guidé]
Que devient $i$ si on double $D$ ?
\begin{enumerate}
  \item Identifier la grandeur principale.
  \item Donner son unité.
  \item Écrire la loi physique pertinente.
  \item Tester l'homogénéité de la relation.
  \item Interpréter le résultat dans le contexte.
\end{enumerate}
\end{exercicebox}

\begin{corrigebox}[Indication]
Il double.
\end{corrigebox}


% ------------------------------------------------------------
\subsection{45. Couleurs des bulles de savon}
% ------------------------------------------------------------

\begin{exemplebox}[Observation concrète]
Une bulle présente des couleurs irisées.
\end{exemplebox}

\begin{methodebox}[Description qualitative]
Les réflexions sur les deux faces du film interfèrent.
La stratégie de compréhension est toujours la même : isoler la grandeur qui varie, identifier la loi physique, puis vérifier le sens du phénomène par un ordre de grandeur.
\end{methodebox}

\begin{theorembox}[Mise en équation niveau Terminale]
La relation centrale utilisée ici est :
\[
2ne=\left(k+\frac12\right)\lambda\quad\text{selon les phases}
\]
Cette formule n'est pas seulement un outil de calcul : elle permet de prévoir comment le phénomène change lorsqu'on modifie une grandeur expérimentale.
\end{theorembox}

\begin{odgbox}[Ordre de grandeur]
L'épaisseur du film est de l'ordre de la longueur d'onde.
\end{odgbox}

\begin{attentionbox}[Limites du modèle]
Le modèle présenté repose sur des hypothèses simplificatrices : milieu homogène ou localement homogène, grandeurs moyennes, géométrie idéalisée, absence de couplages secondaires. En situation réelle, les frottements, gradients, turbulences, imperfections géométriques ou effets non linéaires peuvent modifier le résultat.
\end{attentionbox}

\begin{approfondissementbox}[Niveau prépa / Bac+3]
Les conditions dépendent des déphasages à la réflexion.
On peut souvent améliorer le modèle en remplaçant une relation algébrique simple par une équation différentielle, une loi locale ou un bilan intégral. Cela permet de passer d'une explication qualitative à une prédiction quantitative.
\end{approfondissementbox}

\begin{exercicebox}[Exercice guidé]
Pourquoi les couleurs changent-elles avec l'épaisseur ?
\begin{enumerate}
  \item Identifier la grandeur principale.
  \item Donner son unité.
  \item Écrire la loi physique pertinente.
  \item Tester l'homogénéité de la relation.
  \item Interpréter le résultat dans le contexte.
\end{enumerate}
\end{exercicebox}

\begin{corrigebox}[Indication]
La condition d'interférence sélectionne une longueur d'onde.
\end{corrigebox}


% ------------------------------------------------------------
\subsection{46. Couleurs des taches d'huile}
% ------------------------------------------------------------

\begin{exemplebox}[Observation concrète]
Une tache d'huile sur l'eau produit des couleurs.
\end{exemplebox}

\begin{methodebox}[Description qualitative]
Une couche mince crée des interférences entre rayons réfléchis.
La stratégie de compréhension est toujours la même : isoler la grandeur qui varie, identifier la loi physique, puis vérifier le sens du phénomène par un ordre de grandeur.
\end{methodebox}

\begin{theorembox}[Mise en équation niveau Terminale]
La relation centrale utilisée ici est :
\[
\delta\simeq2ne\cos r
\]
Cette formule n'est pas seulement un outil de calcul : elle permet de prévoir comment le phénomène change lorsqu'on modifie une grandeur expérimentale.
\end{theorembox}

\begin{odgbox}[Ordre de grandeur]
Quelques centaines de nm suffisent pour sélectionner le visible.
\end{odgbox}

\begin{attentionbox}[Limites du modèle]
Le modèle présenté repose sur des hypothèses simplificatrices : milieu homogène ou localement homogène, grandeurs moyennes, géométrie idéalisée, absence de couplages secondaires. En situation réelle, les frottements, gradients, turbulences, imperfections géométriques ou effets non linéaires peuvent modifier le résultat.
\end{attentionbox}

\begin{approfondissementbox}[Niveau prépa / Bac+3]
Une modélisation complète utilise matrices de transfert optique.
On peut souvent améliorer le modèle en remplaçant une relation algébrique simple par une équation différentielle, une loi locale ou un bilan intégral. Cela permet de passer d'une explication qualitative à une prédiction quantitative.
\end{approfondissementbox}

\begin{exercicebox}[Exercice guidé]
Pourquoi la couleur dépend-elle de l'angle de vue ?
\begin{enumerate}
  \item Identifier la grandeur principale.
  \item Donner son unité.
  \item Écrire la loi physique pertinente.
  \item Tester l'homogénéité de la relation.
  \item Interpréter le résultat dans le contexte.
\end{enumerate}
\end{exercicebox}

\begin{corrigebox}[Indication]
La différence de marche dépend de $\cos r$.
\end{corrigebox}


% ------------------------------------------------------------
\subsection{47. Polarisation}
% ------------------------------------------------------------

\begin{exemplebox}[Observation concrète]
Certaines lunettes filtrent les reflets.
\end{exemplebox}

\begin{methodebox}[Description qualitative]
La lumière est une onde transverse et peut avoir une direction de polarisation.
La stratégie de compréhension est toujours la même : isoler la grandeur qui varie, identifier la loi physique, puis vérifier le sens du phénomène par un ordre de grandeur.
\end{methodebox}

\begin{theorembox}[Mise en équation niveau Terminale]
La relation centrale utilisée ici est :
\[
I=I_0\cos^2\theta
\]
Cette formule n'est pas seulement un outil de calcul : elle permet de prévoir comment le phénomène change lorsqu'on modifie une grandeur expérimentale.
\end{theorembox}

\begin{odgbox}[Ordre de grandeur]
Deux polariseurs croisés bloquent presque toute la lumière.
\end{odgbox}

\begin{attentionbox}[Limites du modèle]
Le modèle présenté repose sur des hypothèses simplificatrices : milieu homogène ou localement homogène, grandeurs moyennes, géométrie idéalisée, absence de couplages secondaires. En situation réelle, les frottements, gradients, turbulences, imperfections géométriques ou effets non linéaires peuvent modifier le résultat.
\end{attentionbox}

\begin{approfondissementbox}[Niveau prépa / Bac+3]
La polarisation se décrit par vecteurs de Jones ou paramètres de Stokes.
On peut souvent améliorer le modèle en remplaçant une relation algébrique simple par une équation différentielle, une loi locale ou un bilan intégral. Cela permet de passer d'une explication qualitative à une prédiction quantitative.
\end{approfondissementbox}

\begin{exercicebox}[Exercice guidé]
Calculer l'intensité transmise pour $	heta=60^\circ$.
\begin{enumerate}
  \item Identifier la grandeur principale.
  \item Donner son unité.
  \item Écrire la loi physique pertinente.
  \item Tester l'homogénéité de la relation.
  \item Interpréter le résultat dans le contexte.
\end{enumerate}
\end{exercicebox}

\begin{corrigebox}[Indication]
$I=I_0/4$.
\end{corrigebox}


% ============================================================
\section{Partie — Technologie optique}
% ============================================================


% ------------------------------------------------------------
\subsection{48. Lunettes polarisantes}
% ------------------------------------------------------------

\begin{exemplebox}[Observation concrète]
Les lunettes réduisent les reflets horizontaux sur l'eau ou la route.
\end{exemplebox}

\begin{methodebox}[Description qualitative]
Les reflets sont partiellement polarisés horizontalement.
La stratégie de compréhension est toujours la même : isoler la grandeur qui varie, identifier la loi physique, puis vérifier le sens du phénomène par un ordre de grandeur.
\end{methodebox}

\begin{theorembox}[Mise en équation niveau Terminale]
La relation centrale utilisée ici est :
\[
I=I_0\cos^2\theta
\]
Cette formule n'est pas seulement un outil de calcul : elle permet de prévoir comment le phénomène change lorsqu'on modifie une grandeur expérimentale.
\end{theorembox}

\begin{odgbox}[Ordre de grandeur]
L'effet est maximal près de l'angle de Brewster.
\end{odgbox}

\begin{attentionbox}[Limites du modèle]
Le modèle présenté repose sur des hypothèses simplificatrices : milieu homogène ou localement homogène, grandeurs moyennes, géométrie idéalisée, absence de couplages secondaires. En situation réelle, les frottements, gradients, turbulences, imperfections géométriques ou effets non linéaires peuvent modifier le résultat.
\end{attentionbox}

\begin{approfondissementbox}[Niveau prépa / Bac+3]
À l'angle de Brewster, la réflexion polarisée s'analyse par les coefficients de Fresnel.
On peut souvent améliorer le modèle en remplaçant une relation algébrique simple par une équation différentielle, une loi locale ou un bilan intégral. Cela permet de passer d'une explication qualitative à une prédiction quantitative.
\end{approfondissementbox}

\begin{exercicebox}[Exercice guidé]
Pourquoi tourner la tête change-t-il parfois l'intensité ?
\begin{enumerate}
  \item Identifier la grandeur principale.
  \item Donner son unité.
  \item Écrire la loi physique pertinente.
  \item Tester l'homogénéité de la relation.
  \item Interpréter le résultat dans le contexte.
\end{enumerate}
\end{exercicebox}

\begin{corrigebox}[Indication]
On change l'angle entre polarisation et axe du polariseur.
\end{corrigebox}


% ------------------------------------------------------------
\subsection{49. Écrans LCD}
% ------------------------------------------------------------

\begin{exemplebox}[Observation concrète]
Un écran LCD module la lumière avec des cristaux liquides.
\end{exemplebox}

\begin{methodebox}[Description qualitative]
Polariseurs et orientation moléculaire contrôlent la transmission.
La stratégie de compréhension est toujours la même : isoler la grandeur qui varie, identifier la loi physique, puis vérifier le sens du phénomène par un ordre de grandeur.
\end{methodebox}

\begin{theorembox}[Mise en équation niveau Terminale]
La relation centrale utilisée ici est :
\[
I=I_0\cos^2\theta
\]
Cette formule n'est pas seulement un outil de calcul : elle permet de prévoir comment le phénomène change lorsqu'on modifie une grandeur expérimentale.
\end{theorembox}

\begin{odgbox}[Ordre de grandeur]
Chaque pixel possède des sous-pixels rouge, vert, bleu.
\end{odgbox}

\begin{attentionbox}[Limites du modèle]
Le modèle présenté repose sur des hypothèses simplificatrices : milieu homogène ou localement homogène, grandeurs moyennes, géométrie idéalisée, absence de couplages secondaires. En situation réelle, les frottements, gradients, turbulences, imperfections géométriques ou effets non linéaires peuvent modifier le résultat.
\end{attentionbox}

\begin{approfondissementbox}[Niveau prépa / Bac+3]
Le modèle complet couple électrostatique et anisotropie optique.
On peut souvent améliorer le modèle en remplaçant une relation algébrique simple par une équation différentielle, une loi locale ou un bilan intégral. Cela permet de passer d'une explication qualitative à une prédiction quantitative.
\end{approfondissementbox}

\begin{exercicebox}[Exercice guidé]
Pourquoi un écran LCD a-t-il besoin d'un rétroéclairage ?
\begin{enumerate}
  \item Identifier la grandeur principale.
  \item Donner son unité.
  \item Écrire la loi physique pertinente.
  \item Tester l'homogénéité de la relation.
  \item Interpréter le résultat dans le contexte.
\end{enumerate}
\end{exercicebox}

\begin{corrigebox}[Indication]
Les cristaux liquides modulent mais n'émettent pas forcément la lumière.
\end{corrigebox}


% ------------------------------------------------------------
\subsection{50. Laser}
% ------------------------------------------------------------

\begin{exemplebox}[Observation concrète]
Un laser produit un faisceau intense, directionnel et monochromatique.
\end{exemplebox}

\begin{methodebox}[Description qualitative]
L'émission stimulée amplifie une onde cohérente.
La stratégie de compréhension est toujours la même : isoler la grandeur qui varie, identifier la loi physique, puis vérifier le sens du phénomène par un ordre de grandeur.
\end{methodebox}

\begin{theorembox}[Mise en équation niveau Terminale]
La relation centrale utilisée ici est :
\[
E=h\nu
\]
Cette formule n'est pas seulement un outil de calcul : elle permet de prévoir comment le phénomène change lorsqu'on modifie une grandeur expérimentale.
\end{theorembox}

\begin{odgbox}[Ordre de grandeur]
Un pointeur laser peut avoir quelques mW.
\end{odgbox}

\begin{attentionbox}[Limites du modèle]
Le modèle présenté repose sur des hypothèses simplificatrices : milieu homogène ou localement homogène, grandeurs moyennes, géométrie idéalisée, absence de couplages secondaires. En situation réelle, les frottements, gradients, turbulences, imperfections géométriques ou effets non linéaires peuvent modifier le résultat.
\end{attentionbox}

\begin{approfondissementbox}[Niveau prépa / Bac+3]
Les équations de taux décrivent population excitée et photons dans la cavité.
On peut souvent améliorer le modèle en remplaçant une relation algébrique simple par une équation différentielle, une loi locale ou un bilan intégral. Cela permet de passer d'une explication qualitative à une prédiction quantitative.
\end{approfondissementbox}

\begin{exercicebox}[Exercice guidé]
Pourquoi un laser diffracte-t-il peu mais pas zéro ?
\begin{enumerate}
  \item Identifier la grandeur principale.
  \item Donner son unité.
  \item Écrire la loi physique pertinente.
  \item Tester l'homogénéité de la relation.
  \item Interpréter le résultat dans le contexte.
\end{enumerate}
\end{exercicebox}

\begin{corrigebox}[Indication]
Toute ouverture finie impose une divergence de diffraction.
\end{corrigebox}


% ------------------------------------------------------------
\subsection{51. Lidar}
% ------------------------------------------------------------

\begin{exemplebox}[Observation concrète]
Un lidar mesure une distance par temps de vol.
\end{exemplebox}

\begin{methodebox}[Description qualitative]
Une impulsion laser est envoyée puis réfléchie.
La stratégie de compréhension est toujours la même : isoler la grandeur qui varie, identifier la loi physique, puis vérifier le sens du phénomène par un ordre de grandeur.
\end{methodebox}

\begin{theorembox}[Mise en équation niveau Terminale]
La relation centrale utilisée ici est :
\[
d=\frac{c\Delta t}{2}
\]
Cette formule n'est pas seulement un outil de calcul : elle permet de prévoir comment le phénomène change lorsqu'on modifie une grandeur expérimentale.
\end{theorembox}

\begin{odgbox}[Ordre de grandeur]
Un délai de 100 ns correspond à 15 m.
\end{odgbox}

\begin{attentionbox}[Limites du modèle]
Le modèle présenté repose sur des hypothèses simplificatrices : milieu homogène ou localement homogène, grandeurs moyennes, géométrie idéalisée, absence de couplages secondaires. En situation réelle, les frottements, gradients, turbulences, imperfections géométriques ou effets non linéaires peuvent modifier le résultat.
\end{attentionbox}

\begin{approfondissementbox}[Niveau prépa / Bac+3]
La résolution dépend de la durée d'impulsion et du rapport signal/bruit.
On peut souvent améliorer le modèle en remplaçant une relation algébrique simple par une équation différentielle, une loi locale ou un bilan intégral. Cela permet de passer d'une explication qualitative à une prédiction quantitative.
\end{approfondissementbox}

\begin{exercicebox}[Exercice guidé]
Pourquoi divise-t-on par 2 ?
\begin{enumerate}
  \item Identifier la grandeur principale.
  \item Donner son unité.
  \item Écrire la loi physique pertinente.
  \item Tester l'homogénéité de la relation.
  \item Interpréter le résultat dans le contexte.
\end{enumerate}
\end{exercicebox}

\begin{corrigebox}[Indication]
Le signal fait l'aller-retour.
\end{corrigebox}


% ------------------------------------------------------------
\subsection{52. Spectroscopie}
% ------------------------------------------------------------

\begin{exemplebox}[Observation concrète]
Un spectre révèle la composition d'une source ou solution.
\end{exemplebox}

\begin{methodebox}[Description qualitative]
Les atomes et molécules absorbent ou émettent à certaines longueurs d'onde.
La stratégie de compréhension est toujours la même : isoler la grandeur qui varie, identifier la loi physique, puis vérifier le sens du phénomène par un ordre de grandeur.
\end{methodebox}

\begin{theorembox}[Mise en équation niveau Terminale]
La relation centrale utilisée ici est :
\[
A=\varepsilon \ell C
\]
Cette formule n'est pas seulement un outil de calcul : elle permet de prévoir comment le phénomène change lorsqu'on modifie une grandeur expérimentale.
\end{theorembox}

\begin{odgbox}[Ordre de grandeur]
La loi de Beer-Lambert permet un dosage.
\end{odgbox}

\begin{attentionbox}[Limites du modèle]
Le modèle présenté repose sur des hypothèses simplificatrices : milieu homogène ou localement homogène, grandeurs moyennes, géométrie idéalisée, absence de couplages secondaires. En situation réelle, les frottements, gradients, turbulences, imperfections géométriques ou effets non linéaires peuvent modifier le résultat.
\end{attentionbox}

\begin{approfondissementbox}[Niveau prépa / Bac+3]
Les transitions quantiques expliquent les raies.
On peut souvent améliorer le modèle en remplaçant une relation algébrique simple par une équation différentielle, une loi locale ou un bilan intégral. Cela permet de passer d'une explication qualitative à une prédiction quantitative.
\end{approfondissementbox}

\begin{exercicebox}[Exercice guidé]
Pourquoi une solution colorée absorbe-t-elle certaines couleurs ?
\begin{enumerate}
  \item Identifier la grandeur principale.
  \item Donner son unité.
  \item Écrire la loi physique pertinente.
  \item Tester l'homogénéité de la relation.
  \item Interpréter le résultat dans le contexte.
\end{enumerate}
\end{exercicebox}

\begin{corrigebox}[Indication]
Elle absorbe les longueurs d'onde complémentaires de la couleur perçue.
\end{corrigebox}


% ============================================================
\section{Partie — Acoustique}
% ============================================================


% ------------------------------------------------------------
\subsection{53. Effet Doppler sonore}
% ------------------------------------------------------------

\begin{exemplebox}[Observation concrète]
Une ambulance paraît plus aiguë en approche puis plus grave en s'éloignant.
\end{exemplebox}

\begin{methodebox}[Description qualitative]
Le mouvement change la fréquence d'arrivée des fronts d'onde.
La stratégie de compréhension est toujours la même : isoler la grandeur qui varie, identifier la loi physique, puis vérifier le sens du phénomène par un ordre de grandeur.
\end{methodebox}

\begin{theorembox}[Mise en équation niveau Terminale]
La relation centrale utilisée ici est :
\[
f'=f\frac{c}{c-v_s}\quad\text{source approchante}
\]
Cette formule n'est pas seulement un outil de calcul : elle permet de prévoir comment le phénomène change lorsqu'on modifie une grandeur expérimentale.
\end{theorembox}

\begin{odgbox}[Ordre de grandeur]
À vitesse urbaine, l'effet est audible.
\end{odgbox}

\begin{attentionbox}[Limites du modèle]
Le modèle présenté repose sur des hypothèses simplificatrices : milieu homogène ou localement homogène, grandeurs moyennes, géométrie idéalisée, absence de couplages secondaires. En situation réelle, les frottements, gradients, turbulences, imperfections géométriques ou effets non linéaires peuvent modifier le résultat.
\end{attentionbox}

\begin{approfondissementbox}[Niveau prépa / Bac+3]
La formule générale dépend du mouvement source-observateur projeté sur la ligne de visée.
On peut souvent améliorer le modèle en remplaçant une relation algébrique simple par une équation différentielle, une loi locale ou un bilan intégral. Cela permet de passer d'une explication qualitative à une prédiction quantitative.
\end{approfondissementbox}

\begin{exercicebox}[Exercice guidé]
Pourquoi la fréquence baisse-t-elle après le passage ?
\begin{enumerate}
  \item Identifier la grandeur principale.
  \item Donner son unité.
  \item Écrire la loi physique pertinente.
  \item Tester l'homogénéité de la relation.
  \item Interpréter le résultat dans le contexte.
\end{enumerate}
\end{exercicebox}

\begin{corrigebox}[Indication]
La distance entre fronts d'onde reçus augmente.
\end{corrigebox}


% ============================================================
\section{Partie — Relativité/astrophysique}
% ============================================================


% ------------------------------------------------------------
\subsection{54. Effet Doppler lumineux}
% ------------------------------------------------------------

\begin{exemplebox}[Observation concrète]
La lumière d'une source qui s'éloigne est décalée vers le rouge.
\end{exemplebox}

\begin{methodebox}[Description qualitative]
La fréquence reçue dépend du mouvement relatif.
La stratégie de compréhension est toujours la même : isoler la grandeur qui varie, identifier la loi physique, puis vérifier le sens du phénomène par un ordre de grandeur.
\end{methodebox}

\begin{theorembox}[Mise en équation niveau Terminale]
La relation centrale utilisée ici est :
\[
\frac{\Delta\lambda}{\lambda}\simeq\frac{v}{c}\quad(v\ll c)
\]
Cette formule n'est pas seulement un outil de calcul : elle permet de prévoir comment le phénomène change lorsqu'on modifie une grandeur expérimentale.
\end{theorembox}

\begin{odgbox}[Ordre de grandeur]
Les vitesses stellaires se mesurent ainsi.
\end{odgbox}

\begin{attentionbox}[Limites du modèle]
Le modèle présenté repose sur des hypothèses simplificatrices : milieu homogène ou localement homogène, grandeurs moyennes, géométrie idéalisée, absence de couplages secondaires. En situation réelle, les frottements, gradients, turbulences, imperfections géométriques ou effets non linéaires peuvent modifier le résultat.
\end{attentionbox}

\begin{approfondissementbox}[Niveau prépa / Bac+3]
À grande vitesse, il faut la formule relativiste.
On peut souvent améliorer le modèle en remplaçant une relation algébrique simple par une équation différentielle, une loi locale ou un bilan intégral. Cela permet de passer d'une explication qualitative à une prédiction quantitative.
\end{approfondissementbox}

\begin{exercicebox}[Exercice guidé]
Que signifie un décalage vers le bleu ?
\begin{enumerate}
  \item Identifier la grandeur principale.
  \item Donner son unité.
  \item Écrire la loi physique pertinente.
  \item Tester l'homogénéité de la relation.
  \item Interpréter le résultat dans le contexte.
\end{enumerate}
\end{exercicebox}

\begin{corrigebox}[Indication]
La source se rapproche.
\end{corrigebox}


% ============================================================
\section{Partie — Acoustique}
% ============================================================


% ------------------------------------------------------------
\subsection{55. Mur du son}
% ------------------------------------------------------------

\begin{exemplebox}[Observation concrète]
Un avion supersonique crée un bang.
\end{exemplebox}

\begin{methodebox}[Description qualitative]
À vitesse supérieure au son, les ondes s'accumulent en cône de Mach.
La stratégie de compréhension est toujours la même : isoler la grandeur qui varie, identifier la loi physique, puis vérifier le sens du phénomène par un ordre de grandeur.
\end{methodebox}

\begin{theorembox}[Mise en équation niveau Terminale]
La relation centrale utilisée ici est :
\[
M=\frac{v}{c_s}
\]
Cette formule n'est pas seulement un outil de calcul : elle permet de prévoir comment le phénomène change lorsqu'on modifie une grandeur expérimentale.
\end{theorembox}

\begin{odgbox}[Ordre de grandeur]
Mach 1 correspond à environ 340 m/s dans l'air.
\end{odgbox}

\begin{attentionbox}[Limites du modèle]
Le modèle présenté repose sur des hypothèses simplificatrices : milieu homogène ou localement homogène, grandeurs moyennes, géométrie idéalisée, absence de couplages secondaires. En situation réelle, les frottements, gradients, turbulences, imperfections géométriques ou effets non linéaires peuvent modifier le résultat.
\end{attentionbox}

\begin{approfondissementbox}[Niveau prépa / Bac+3]
Le cône de Mach vérifie $\sin	heta=1/M$.
On peut souvent améliorer le modèle en remplaçant une relation algébrique simple par une équation différentielle, une loi locale ou un bilan intégral. Cela permet de passer d'une explication qualitative à une prédiction quantitative.
\end{approfondissementbox}

\begin{exercicebox}[Exercice guidé]
Calculer $M$ pour $v=680$ m/s.
\begin{enumerate}
  \item Identifier la grandeur principale.
  \item Donner son unité.
  \item Écrire la loi physique pertinente.
  \item Tester l'homogénéité de la relation.
  \item Interpréter le résultat dans le contexte.
\end{enumerate}
\end{exercicebox}

\begin{corrigebox}[Indication]
$M=2$.
\end{corrigebox}


% ------------------------------------------------------------
\subsection{56. Bang supersonique}
% ------------------------------------------------------------

\begin{exemplebox}[Observation concrète]
Un bruit sec est entendu au passage d'un avion supersonique.
\end{exemplebox}

\begin{methodebox}[Description qualitative]
L'onde de choc atteint l'observateur.
La stratégie de compréhension est toujours la même : isoler la grandeur qui varie, identifier la loi physique, puis vérifier le sens du phénomène par un ordre de grandeur.
\end{methodebox}

\begin{theorembox}[Mise en équation niveau Terminale]
La relation centrale utilisée ici est :
\[
\sin\theta=\frac{c_s}{v}
\]
Cette formule n'est pas seulement un outil de calcul : elle permet de prévoir comment le phénomène change lorsqu'on modifie une grandeur expérimentale.
\end{theorembox}

\begin{odgbox}[Ordre de grandeur]
Le bang n'est pas seulement au moment du franchissement du mur du son.
\end{odgbox}

\begin{attentionbox}[Limites du modèle]
Le modèle présenté repose sur des hypothèses simplificatrices : milieu homogène ou localement homogène, grandeurs moyennes, géométrie idéalisée, absence de couplages secondaires. En situation réelle, les frottements, gradients, turbulences, imperfections géométriques ou effets non linéaires peuvent modifier le résultat.
\end{attentionbox}

\begin{approfondissementbox}[Niveau prépa / Bac+3]
L'onde de choc est une discontinuité compressible.
On peut souvent améliorer le modèle en remplaçant une relation algébrique simple par une équation différentielle, une loi locale ou un bilan intégral. Cela permet de passer d'une explication qualitative à une prédiction quantitative.
\end{approfondissementbox}

\begin{exercicebox}[Exercice guidé]
Pourquoi entend-on un bang même si l'avion est déjà supersonique ?
\begin{enumerate}
  \item Identifier la grandeur principale.
  \item Donner son unité.
  \item Écrire la loi physique pertinente.
  \item Tester l'homogénéité de la relation.
  \item Interpréter le résultat dans le contexte.
\end{enumerate}
\end{exercicebox}

\begin{corrigebox}[Indication]
Le cône de choc balaie le sol.
\end{corrigebox}


% ============================================================
\section{Partie — Vibrations}
% ============================================================


% ------------------------------------------------------------
\subsection{57. Résonance d'un verre}
% ------------------------------------------------------------

\begin{exemplebox}[Observation concrète]
Un verre peut vibrer fortement à une fréquence précise.
\end{exemplebox}

\begin{methodebox}[Description qualitative]
Une excitation proche de la fréquence propre augmente l'amplitude.
La stratégie de compréhension est toujours la même : isoler la grandeur qui varie, identifier la loi physique, puis vérifier le sens du phénomène par un ordre de grandeur.
\end{methodebox}

\begin{theorembox}[Mise en équation niveau Terminale]
La relation centrale utilisée ici est :
\[
\omega_0=\sqrt{\frac{k}{m}}
\]
Cette formule n'est pas seulement un outil de calcul : elle permet de prévoir comment le phénomène change lorsqu'on modifie une grandeur expérimentale.
\end{theorembox}

\begin{odgbox}[Ordre de grandeur]
La résonance dépend de la forme et du matériau.
\end{odgbox}

\begin{attentionbox}[Limites du modèle]
Le modèle présenté repose sur des hypothèses simplificatrices : milieu homogène ou localement homogène, grandeurs moyennes, géométrie idéalisée, absence de couplages secondaires. En situation réelle, les frottements, gradients, turbulences, imperfections géométriques ou effets non linéaires peuvent modifier le résultat.
\end{attentionbox}

\begin{approfondissementbox}[Niveau prépa / Bac+3]
Un oscillateur amorti forcé a une amplitude maximale près de $\omega_0$.
On peut souvent améliorer le modèle en remplaçant une relation algébrique simple par une équation différentielle, une loi locale ou un bilan intégral. Cela permet de passer d'une explication qualitative à une prédiction quantitative.
\end{approfondissementbox}

\begin{exercicebox}[Exercice guidé]
Pourquoi l'amortissement limite-t-il l'amplitude ?
\begin{enumerate}
  \item Identifier la grandeur principale.
  \item Donner son unité.
  \item Écrire la loi physique pertinente.
  \item Tester l'homogénéité de la relation.
  \item Interpréter le résultat dans le contexte.
\end{enumerate}
\end{exercicebox}

\begin{corrigebox}[Indication]
Il dissipe l'énergie apportée.
\end{corrigebox}


% ------------------------------------------------------------
\subsection{58. Résonance d'un pont}
% ------------------------------------------------------------

\begin{exemplebox}[Observation concrète]
Un pont peut osciller sous le vent ou des pas synchronisés.
\end{exemplebox}

\begin{methodebox}[Description qualitative]
Une excitation périodique peut transférer efficacement de l'énergie.
La stratégie de compréhension est toujours la même : isoler la grandeur qui varie, identifier la loi physique, puis vérifier le sens du phénomène par un ordre de grandeur.
\end{methodebox}

\begin{theorembox}[Mise en équation niveau Terminale]
La relation centrale utilisée ici est :
\[
\omega\simeq\omega_0
\]
Cette formule n'est pas seulement un outil de calcul : elle permet de prévoir comment le phénomène change lorsqu'on modifie une grandeur expérimentale.
\end{theorembox}

\begin{odgbox}[Ordre de grandeur]
Les ingénieurs évitent les fréquences dangereuses.
\end{odgbox}

\begin{attentionbox}[Limites du modèle]
Le modèle présenté repose sur des hypothèses simplificatrices : milieu homogène ou localement homogène, grandeurs moyennes, géométrie idéalisée, absence de couplages secondaires. En situation réelle, les frottements, gradients, turbulences, imperfections géométriques ou effets non linéaires peuvent modifier le résultat.
\end{attentionbox}

\begin{approfondissementbox}[Niveau prépa / Bac+3]
Les modes propres d'une structure sont solutions d'un problème aux valeurs propres.
On peut souvent améliorer le modèle en remplaçant une relation algébrique simple par une équation différentielle, une loi locale ou un bilan intégral. Cela permet de passer d'une explication qualitative à une prédiction quantitative.
\end{approfondissementbox}

\begin{exercicebox}[Exercice guidé]
Pourquoi rompre le pas sur un pont suspendu ?
\begin{enumerate}
  \item Identifier la grandeur principale.
  \item Donner son unité.
  \item Écrire la loi physique pertinente.
  \item Tester l'homogénéité de la relation.
  \item Interpréter le résultat dans le contexte.
\end{enumerate}
\end{exercicebox}

\begin{corrigebox}[Indication]
Pour éviter une excitation cohérente.
\end{corrigebox}


% ============================================================
\section{Partie — Musique}
% ============================================================


% ------------------------------------------------------------
\subsection{59. Instrument à corde}
% ------------------------------------------------------------

\begin{exemplebox}[Observation concrète]
Une corde de guitare vibre à des fréquences précises.
\end{exemplebox}

\begin{methodebox}[Description qualitative]
Les extrémités imposent des nœuds et sélectionnent des modes.
La stratégie de compréhension est toujours la même : isoler la grandeur qui varie, identifier la loi physique, puis vérifier le sens du phénomène par un ordre de grandeur.
\end{methodebox}

\begin{theorembox}[Mise en équation niveau Terminale]
La relation centrale utilisée ici est :
\[
f_n=n\frac{v}{2L}
\]
Cette formule n'est pas seulement un outil de calcul : elle permet de prévoir comment le phénomène change lorsqu'on modifie une grandeur expérimentale.
\end{theorembox}

\begin{odgbox}[Ordre de grandeur]
Raccourcir la corde augmente la fréquence.
\end{odgbox}

\begin{attentionbox}[Limites du modèle]
Le modèle présenté repose sur des hypothèses simplificatrices : milieu homogène ou localement homogène, grandeurs moyennes, géométrie idéalisée, absence de couplages secondaires. En situation réelle, les frottements, gradients, turbulences, imperfections géométriques ou effets non linéaires peuvent modifier le résultat.
\end{attentionbox}

\begin{approfondissementbox}[Niveau prépa / Bac+3]
La vitesse dépend de la tension et de la masse linéique : $v=\sqrt{T/\mu}$.
On peut souvent améliorer le modèle en remplaçant une relation algébrique simple par une équation différentielle, une loi locale ou un bilan intégral. Cela permet de passer d'une explication qualitative à une prédiction quantitative.
\end{approfondissementbox}

\begin{exercicebox}[Exercice guidé]
Pourquoi tendre une corde rend le son plus aigu ?
\begin{enumerate}
  \item Identifier la grandeur principale.
  \item Donner son unité.
  \item Écrire la loi physique pertinente.
  \item Tester l'homogénéité de la relation.
  \item Interpréter le résultat dans le contexte.
\end{enumerate}
\end{exercicebox}

\begin{corrigebox}[Indication]
La vitesse d'onde augmente.
\end{corrigebox}


% ------------------------------------------------------------
\subsection{60. Instrument à vent}
% ------------------------------------------------------------

\begin{exemplebox}[Observation concrète]
Une colonne d'air résonne dans un tube.
\end{exemplebox}

\begin{methodebox}[Description qualitative]
Les conditions aux extrémités sélectionnent les modes.
La stratégie de compréhension est toujours la même : isoler la grandeur qui varie, identifier la loi physique, puis vérifier le sens du phénomène par un ordre de grandeur.
\end{methodebox}

\begin{theorembox}[Mise en équation niveau Terminale]
La relation centrale utilisée ici est :
\[
f_n=n\frac{c}{2L}\quad\text{tube ouvert-ouvert}
\]
Cette formule n'est pas seulement un outil de calcul : elle permet de prévoir comment le phénomène change lorsqu'on modifie une grandeur expérimentale.
\end{theorembox}

\begin{odgbox}[Ordre de grandeur]
Un tube plus long donne un son plus grave.
\end{odgbox}

\begin{attentionbox}[Limites du modèle]
Le modèle présenté repose sur des hypothèses simplificatrices : milieu homogène ou localement homogène, grandeurs moyennes, géométrie idéalisée, absence de couplages secondaires. En situation réelle, les frottements, gradients, turbulences, imperfections géométriques ou effets non linéaires peuvent modifier le résultat.
\end{attentionbox}

\begin{approfondissementbox}[Niveau prépa / Bac+3]
Les corrections d'extrémité modifient la longueur effective.
On peut souvent améliorer le modèle en remplaçant une relation algébrique simple par une équation différentielle, une loi locale ou un bilan intégral. Cela permet de passer d'une explication qualitative à une prédiction quantitative.
\end{approfondissementbox}

\begin{exercicebox}[Exercice guidé]
Pourquoi ouvrir des trous sur une flûte augmente la hauteur ?
\begin{enumerate}
  \item Identifier la grandeur principale.
  \item Donner son unité.
  \item Écrire la loi physique pertinente.
  \item Tester l'homogénéité de la relation.
  \item Interpréter le résultat dans le contexte.
\end{enumerate}
\end{exercicebox}

\begin{corrigebox}[Indication]
La longueur résonante effective diminue.
\end{corrigebox}


% ============================================================
\section{Partie — Acoustique}
% ============================================================


% ------------------------------------------------------------
\subsection{61. Écho}
% ------------------------------------------------------------

\begin{exemplebox}[Observation concrète]
Un son revient après réflexion sur une paroi.
\end{exemplebox}

\begin{methodebox}[Description qualitative]
L'onde sonore fait un aller-retour.
La stratégie de compréhension est toujours la même : isoler la grandeur qui varie, identifier la loi physique, puis vérifier le sens du phénomène par un ordre de grandeur.
\end{methodebox}

\begin{theorembox}[Mise en équation niveau Terminale]
La relation centrale utilisée ici est :
\[
d=\frac{c_s\Delta t}{2}
\]
Cette formule n'est pas seulement un outil de calcul : elle permet de prévoir comment le phénomène change lorsqu'on modifie une grandeur expérimentale.
\end{theorembox}

\begin{odgbox}[Ordre de grandeur]
Un délai de 0,1 s correspond à environ 17 m.
\end{odgbox}

\begin{attentionbox}[Limites du modèle]
Le modèle présenté repose sur des hypothèses simplificatrices : milieu homogène ou localement homogène, grandeurs moyennes, géométrie idéalisée, absence de couplages secondaires. En situation réelle, les frottements, gradients, turbulences, imperfections géométriques ou effets non linéaires peuvent modifier le résultat.
\end{attentionbox}

\begin{approfondissementbox}[Niveau prépa / Bac+3]
L'échographie utilise une idée similaire avec des ultrasons.
On peut souvent améliorer le modèle en remplaçant une relation algébrique simple par une équation différentielle, une loi locale ou un bilan intégral. Cela permet de passer d'une explication qualitative à une prédiction quantitative.
\end{approfondissementbox}

\begin{exercicebox}[Exercice guidé]
Pourquoi diviser par deux ?
\begin{enumerate}
  \item Identifier la grandeur principale.
  \item Donner son unité.
  \item Écrire la loi physique pertinente.
  \item Tester l'homogénéité de la relation.
  \item Interpréter le résultat dans le contexte.
\end{enumerate}
\end{exercicebox}

\begin{corrigebox}[Indication]
Le son parcourt l'aller et le retour.
\end{corrigebox}


% ------------------------------------------------------------
\subsection{62. Acoustique d'une salle}
% ------------------------------------------------------------

\begin{exemplebox}[Observation concrète]
Une salle peut être claire, sourde ou réverbérante.
\end{exemplebox}

\begin{methodebox}[Description qualitative]
Les réflexions multiples prolongent le son.
La stratégie de compréhension est toujours la même : isoler la grandeur qui varie, identifier la loi physique, puis vérifier le sens du phénomène par un ordre de grandeur.
\end{methodebox}

\begin{theorembox}[Mise en équation niveau Terminale]
La relation centrale utilisée ici est :
\[
T_R\simeq0{,}16\frac{V}{A}
\]
Cette formule n'est pas seulement un outil de calcul : elle permet de prévoir comment le phénomène change lorsqu'on modifie une grandeur expérimentale.
\end{theorembox}

\begin{odgbox}[Ordre de grandeur]
La formule de Sabine relie volume, absorption et temps de réverbération.
\end{odgbox}

\begin{attentionbox}[Limites du modèle]
Le modèle présenté repose sur des hypothèses simplificatrices : milieu homogène ou localement homogène, grandeurs moyennes, géométrie idéalisée, absence de couplages secondaires. En situation réelle, les frottements, gradients, turbulences, imperfections géométriques ou effets non linéaires peuvent modifier le résultat.
\end{attentionbox}

\begin{approfondissementbox}[Niveau prépa / Bac+3]
Une modélisation avancée utilise l'équation d'onde et la réponse impulsionnelle.
On peut souvent améliorer le modèle en remplaçant une relation algébrique simple par une équation différentielle, une loi locale ou un bilan intégral. Cela permet de passer d'une explication qualitative à une prédiction quantitative.
\end{approfondissementbox}

\begin{exercicebox}[Exercice guidé]
Pourquoi mettre des rideaux diminue-t-il la réverbération ?
\begin{enumerate}
  \item Identifier la grandeur principale.
  \item Donner son unité.
  \item Écrire la loi physique pertinente.
  \item Tester l'homogénéité de la relation.
  \item Interpréter le résultat dans le contexte.
\end{enumerate}
\end{exercicebox}

\begin{corrigebox}[Indication]
Ils augmentent l'absorption acoustique.
\end{corrigebox}


% ============================================================
\section{Partie — Sport}
% ============================================================


% ------------------------------------------------------------
\subsection{63. Ambulance qui change de son}
% ------------------------------------------------------------

\begin{exemplebox}[Observation concrète]
Le son d'une ambulance change brutalement lors du passage.
\end{exemplebox}

\begin{methodebox}[Description qualitative]
C'est une application directe de l'effet Doppler sonore.
La stratégie de compréhension est toujours la même : isoler la grandeur qui varie, identifier la loi physique, puis vérifier le sens du phénomène par un ordre de grandeur.
\end{methodebox}

\begin{theorembox}[Mise en équation niveau Terminale]
La relation centrale utilisée ici est :
\[
f'=f\frac{c}{c\mp v_s}
\]
Cette formule n'est pas seulement un outil de calcul : elle permet de prévoir comment le phénomène change lorsqu'on modifie une grandeur expérimentale.
\end{theorembox}

\begin{odgbox}[Ordre de grandeur]
À 50 km/h, le décalage est de quelques pourcents.
\end{odgbox}

\begin{attentionbox}[Limites du modèle]
Le modèle présenté repose sur des hypothèses simplificatrices : milieu homogène ou localement homogène, grandeurs moyennes, géométrie idéalisée, absence de couplages secondaires. En situation réelle, les frottements, gradients, turbulences, imperfections géométriques ou effets non linéaires peuvent modifier le résultat.
\end{attentionbox}

\begin{approfondissementbox}[Niveau prépa / Bac+3]
La projection de la vitesse sur la ligne de visée change continûment.
On peut souvent améliorer le modèle en remplaçant une relation algébrique simple par une équation différentielle, une loi locale ou un bilan intégral. Cela permet de passer d'une explication qualitative à une prédiction quantitative.
\end{approfondissementbox}

\begin{exercicebox}[Exercice guidé]
À quel moment la fréquence apparente change-t-elle le plus rapidement ?
\begin{enumerate}
  \item Identifier la grandeur principale.
  \item Donner son unité.
  \item Écrire la loi physique pertinente.
  \item Tester l'homogénéité de la relation.
  \item Interpréter le résultat dans le contexte.
\end{enumerate}
\end{exercicebox}

\begin{corrigebox}[Indication]
Au passage proche de l'observateur.
\end{corrigebox}


% ------------------------------------------------------------
\subsection{64. Balle de tennis liftée}
% ------------------------------------------------------------

\begin{exemplebox}[Observation concrète]
Une balle liftée plonge plus vite dans le court.
\end{exemplebox}

\begin{methodebox}[Description qualitative]
La rotation crée une force de Magnus vers le bas pour un topspin.
La stratégie de compréhension est toujours la même : isoler la grandeur qui varie, identifier la loi physique, puis vérifier le sens du phénomène par un ordre de grandeur.
\end{methodebox}

\begin{theorembox}[Mise en équation niveau Terminale]
La relation centrale utilisée ici est :
\[
\vec F_M\propto \vec\omega\times \vec v
\]
Cette formule n'est pas seulement un outil de calcul : elle permet de prévoir comment le phénomène change lorsqu'on modifie une grandeur expérimentale.
\end{theorembox}

\begin{odgbox}[Ordre de grandeur]
Le lift permet de frapper plus fort en gardant la balle dans le terrain.
\end{odgbox}

\begin{attentionbox}[Limites du modèle]
Le modèle présenté repose sur des hypothèses simplificatrices : milieu homogène ou localement homogène, grandeurs moyennes, géométrie idéalisée, absence de couplages secondaires. En situation réelle, les frottements, gradients, turbulences, imperfections géométriques ou effets non linéaires peuvent modifier le résultat.
\end{attentionbox}

\begin{approfondissementbox}[Niveau prépa / Bac+3]
L'écoulement autour de la balle impose une portance dépendant du nombre de Reynolds.
On peut souvent améliorer le modèle en remplaçant une relation algébrique simple par une équation différentielle, une loi locale ou un bilan intégral. Cela permet de passer d'une explication qualitative à une prédiction quantitative.
\end{approfondissementbox}

\begin{exercicebox}[Exercice guidé]
Pourquoi le lift sécurise-t-il un coup puissant ?
\begin{enumerate}
  \item Identifier la grandeur principale.
  \item Donner son unité.
  \item Écrire la loi physique pertinente.
  \item Tester l'homogénéité de la relation.
  \item Interpréter le résultat dans le contexte.
\end{enumerate}
\end{exercicebox}

\begin{corrigebox}[Indication]
La force vers le bas augmente la courbure de la trajectoire.
\end{corrigebox}


% ------------------------------------------------------------
\subsection{65. Effet Magnus}
% ------------------------------------------------------------

\begin{exemplebox}[Observation concrète]
Un objet en rotation dans un fluide subit une force latérale.
\end{exemplebox}

\begin{methodebox}[Description qualitative]
La rotation modifie les vitesses relatives de l'air autour de l'objet.
La stratégie de compréhension est toujours la même : isoler la grandeur qui varie, identifier la loi physique, puis vérifier le sens du phénomène par un ordre de grandeur.
\end{methodebox}

\begin{theorembox}[Mise en équation niveau Terminale]
La relation centrale utilisée ici est :
\[
\vec F_M\sim C\rho R^3\,\vec\omega\times\vec v
\]
Cette formule n'est pas seulement un outil de calcul : elle permet de prévoir comment le phénomène change lorsqu'on modifie une grandeur expérimentale.
\end{theorembox}

\begin{odgbox}[Ordre de grandeur]
Très visible en tennis, football, ping-pong.
\end{odgbox}

\begin{attentionbox}[Limites du modèle]
Le modèle présenté repose sur des hypothèses simplificatrices : milieu homogène ou localement homogène, grandeurs moyennes, géométrie idéalisée, absence de couplages secondaires. En situation réelle, les frottements, gradients, turbulences, imperfections géométriques ou effets non linéaires peuvent modifier le résultat.
\end{attentionbox}

\begin{approfondissementbox}[Niveau prépa / Bac+3]
La force dépend de la circulation du fluide : théorème de Kutta-Joukowski.
On peut souvent améliorer le modèle en remplaçant une relation algébrique simple par une équation différentielle, une loi locale ou un bilan intégral. Cela permet de passer d'une explication qualitative à une prédiction quantitative.
\end{approfondissementbox}

\begin{exercicebox}[Exercice guidé]
Dans quel sens courbe un ballon avec rotation horaire ?
\begin{enumerate}
  \item Identifier la grandeur principale.
  \item Donner son unité.
  \item Écrire la loi physique pertinente.
  \item Tester l'homogénéité de la relation.
  \item Interpréter le résultat dans le contexte.
\end{enumerate}
\end{exercicebox}

\begin{corrigebox}[Indication]
Le sens se déduit de $
ec\omega	imes
ec v$.
\end{corrigebox}


% ------------------------------------------------------------
\subsection{66. Coup franc brossé}
% ------------------------------------------------------------

\begin{exemplebox}[Observation concrète]
Le ballon contourne le mur.
\end{exemplebox}

\begin{methodebox}[Description qualitative]
La rotation latérale produit une force de Magnus horizontale.
La stratégie de compréhension est toujours la même : isoler la grandeur qui varie, identifier la loi physique, puis vérifier le sens du phénomène par un ordre de grandeur.
\end{methodebox}

\begin{theorembox}[Mise en équation niveau Terminale]
La relation centrale utilisée ici est :
\[
a_{\perp}=\frac{F_M}{m}
\]
Cette formule n'est pas seulement un outil de calcul : elle permet de prévoir comment le phénomène change lorsqu'on modifie une grandeur expérimentale.
\end{theorembox}

\begin{odgbox}[Ordre de grandeur]
Une rotation de plusieurs tours par seconde peut courber fortement la trajectoire.
\end{odgbox}

\begin{attentionbox}[Limites du modèle]
Le modèle présenté repose sur des hypothèses simplificatrices : milieu homogène ou localement homogène, grandeurs moyennes, géométrie idéalisée, absence de couplages secondaires. En situation réelle, les frottements, gradients, turbulences, imperfections géométriques ou effets non linéaires peuvent modifier le résultat.
\end{attentionbox}

\begin{approfondissementbox}[Niveau prépa / Bac+3]
Traînée et portance se couplent dans une équation différentielle vectorielle.
On peut souvent améliorer le modèle en remplaçant une relation algébrique simple par une équation différentielle, une loi locale ou un bilan intégral. Cela permet de passer d'une explication qualitative à une prédiction quantitative.
\end{approfondissementbox}

\begin{exercicebox}[Exercice guidé]
Pourquoi une frappe sans rotation va-t-elle plus droit ?
\begin{enumerate}
  \item Identifier la grandeur principale.
  \item Donner son unité.
  \item Écrire la loi physique pertinente.
  \item Tester l'homogénéité de la relation.
  \item Interpréter le résultat dans le contexte.
\end{enumerate}
\end{exercicebox}

\begin{corrigebox}[Indication]
La force de Magnus est faible ou nulle.
\end{corrigebox}


% ------------------------------------------------------------
\subsection{67. Rebond d'une balle}
% ------------------------------------------------------------

\begin{exemplebox}[Observation concrète]
Une balle rebondit moins haut que sa hauteur de chute.
\end{exemplebox}

\begin{methodebox}[Description qualitative]
Une partie de l'énergie est dissipée lors du choc.
La stratégie de compréhension est toujours la même : isoler la grandeur qui varie, identifier la loi physique, puis vérifier le sens du phénomène par un ordre de grandeur.
\end{methodebox}

\begin{theorembox}[Mise en équation niveau Terminale]
La relation centrale utilisée ici est :
\[
e=\frac{v_{\mathrm{après}}}{v_{\mathrm{avant}}}
\]
Cette formule n'est pas seulement un outil de calcul : elle permet de prévoir comment le phénomène change lorsqu'on modifie une grandeur expérimentale.
\end{theorembox}

\begin{odgbox}[Ordre de grandeur]
Le coefficient de restitution est inférieur à 1.
\end{odgbox}

\begin{attentionbox}[Limites du modèle]
Le modèle présenté repose sur des hypothèses simplificatrices : milieu homogène ou localement homogène, grandeurs moyennes, géométrie idéalisée, absence de couplages secondaires. En situation réelle, les frottements, gradients, turbulences, imperfections géométriques ou effets non linéaires peuvent modifier le résultat.
\end{attentionbox}

\begin{approfondissementbox}[Niveau prépa / Bac+3]
Le contact réel implique déformation, viscoélasticité et frottement.
On peut souvent améliorer le modèle en remplaçant une relation algébrique simple par une équation différentielle, une loi locale ou un bilan intégral. Cela permet de passer d'une explication qualitative à une prédiction quantitative.
\end{approfondissementbox}

\begin{exercicebox}[Exercice guidé]
Si $e=0{,}8$, quelle fraction d'énergie verticale reste ?
\begin{enumerate}
  \item Identifier la grandeur principale.
  \item Donner son unité.
  \item Écrire la loi physique pertinente.
  \item Tester l'homogénéité de la relation.
  \item Interpréter le résultat dans le contexte.
\end{enumerate}
\end{exercicebox}

\begin{corrigebox}[Indication]
$e^2=0{,}64$, donc 64 %.
\end{corrigebox}


% ============================================================
\section{Partie — Mécanique}
% ============================================================


% ------------------------------------------------------------
\subsection{68. Frottements solides}
% ------------------------------------------------------------

\begin{exemplebox}[Observation concrète]
Un objet glisse puis s'arrête.
\end{exemplebox}

\begin{methodebox}[Description qualitative]
La force de frottement s'oppose au mouvement.
La stratégie de compréhension est toujours la même : isoler la grandeur qui varie, identifier la loi physique, puis vérifier le sens du phénomène par un ordre de grandeur.
\end{methodebox}

\begin{theorembox}[Mise en équation niveau Terminale]
La relation centrale utilisée ici est :
\[
F_f=\mu N
\]
Cette formule n'est pas seulement un outil de calcul : elle permet de prévoir comment le phénomène change lorsqu'on modifie une grandeur expérimentale.
\end{theorembox}

\begin{odgbox}[Ordre de grandeur]
Le coefficient $\mu$ dépend des matériaux.
\end{odgbox}

\begin{attentionbox}[Limites du modèle]
Le modèle présenté repose sur des hypothèses simplificatrices : milieu homogène ou localement homogène, grandeurs moyennes, géométrie idéalisée, absence de couplages secondaires. En situation réelle, les frottements, gradients, turbulences, imperfections géométriques ou effets non linéaires peuvent modifier le résultat.
\end{attentionbox}

\begin{approfondissementbox}[Niveau prépa / Bac+3]
La loi de Coulomb est empirique et échoue à l'échelle microscopique fine.
On peut souvent améliorer le modèle en remplaçant une relation algébrique simple par une équation différentielle, une loi locale ou un bilan intégral. Cela permet de passer d'une explication qualitative à une prédiction quantitative.
\end{approfondissementbox}

\begin{exercicebox}[Exercice guidé]
Pourquoi un objet plus lourd frotte-t-il davantage ?
\begin{enumerate}
  \item Identifier la grandeur principale.
  \item Donner son unité.
  \item Écrire la loi physique pertinente.
  \item Tester l'homogénéité de la relation.
  \item Interpréter le résultat dans le contexte.
\end{enumerate}
\end{exercicebox}

\begin{corrigebox}[Indication]
La réaction normale augmente.
\end{corrigebox}


% ============================================================
\section{Partie — Mécanique des fluides}
% ============================================================


% ------------------------------------------------------------
\subsection{69. Frottements fluides}
% ------------------------------------------------------------

\begin{exemplebox}[Observation concrète]
Un objet dans l'air ou l'eau subit une traînée.
\end{exemplebox}

\begin{methodebox}[Description qualitative]
Le fluide doit être déplacé et dissipé.
La stratégie de compréhension est toujours la même : isoler la grandeur qui varie, identifier la loi physique, puis vérifier le sens du phénomène par un ordre de grandeur.
\end{methodebox}

\begin{theorembox}[Mise en équation niveau Terminale]
La relation centrale utilisée ici est :
\[
F_D=\frac12\rho C_D S v^2
\]
Cette formule n'est pas seulement un outil de calcul : elle permet de prévoir comment le phénomène change lorsqu'on modifie une grandeur expérimentale.
\end{theorembox}

\begin{odgbox}[Ordre de grandeur]
La traînée croît souvent comme $v^2$ à grande vitesse.
\end{odgbox}

\begin{attentionbox}[Limites du modèle]
Le modèle présenté repose sur des hypothèses simplificatrices : milieu homogène ou localement homogène, grandeurs moyennes, géométrie idéalisée, absence de couplages secondaires. En situation réelle, les frottements, gradients, turbulences, imperfections géométriques ou effets non linéaires peuvent modifier le résultat.
\end{attentionbox}

\begin{approfondissementbox}[Niveau prépa / Bac+3]
À faible Reynolds, la loi de Stokes donne $F=6\pi\eta Rv$.
On peut souvent améliorer le modèle en remplaçant une relation algébrique simple par une équation différentielle, une loi locale ou un bilan intégral. Cela permet de passer d'une explication qualitative à une prédiction quantitative.
\end{approfondissementbox}

\begin{exercicebox}[Exercice guidé]
Pourquoi rouler vite consomme beaucoup plus ?
\begin{enumerate}
  \item Identifier la grandeur principale.
  \item Donner son unité.
  \item Écrire la loi physique pertinente.
  \item Tester l'homogénéité de la relation.
  \item Interpréter le résultat dans le contexte.
\end{enumerate}
\end{exercicebox}

\begin{corrigebox}[Indication]
La puissance de traînée varie comme $v^3$.
\end{corrigebox}


% ============================================================
\section{Partie — Transport}
% ============================================================


% ------------------------------------------------------------
\subsection{70. Aquaplaning}
% ------------------------------------------------------------

\begin{exemplebox}[Observation concrète]
Un pneu perd l'adhérence sur une couche d'eau.
\end{exemplebox}

\begin{methodebox}[Description qualitative]
L'eau ne s'évacue plus assez vite et sépare pneu et route.
La stratégie de compréhension est toujours la même : isoler la grandeur qui varie, identifier la loi physique, puis vérifier le sens du phénomène par un ordre de grandeur.
\end{methodebox}

\begin{theorembox}[Mise en équation niveau Terminale]
La relation centrale utilisée ici est :
\[
\text{pression hydrodynamique}\sim\rho v^2
\]
Cette formule n'est pas seulement un outil de calcul : elle permet de prévoir comment le phénomène change lorsqu'on modifie une grandeur expérimentale.
\end{theorembox}

\begin{odgbox}[Ordre de grandeur]
Le risque augmente fortement avec la vitesse.
\end{odgbox}

\begin{attentionbox}[Limites du modèle]
Le modèle présenté repose sur des hypothèses simplificatrices : milieu homogène ou localement homogène, grandeurs moyennes, géométrie idéalisée, absence de couplages secondaires. En situation réelle, les frottements, gradients, turbulences, imperfections géométriques ou effets non linéaires peuvent modifier le résultat.
\end{attentionbox}

\begin{approfondissementbox}[Niveau prépa / Bac+3]
La stabilité dépend de l'écoulement, de la sculpture du pneu et de la charge.
On peut souvent améliorer le modèle en remplaçant une relation algébrique simple par une équation différentielle, une loi locale ou un bilan intégral. Cela permet de passer d'une explication qualitative à une prédiction quantitative.
\end{approfondissementbox}

\begin{exercicebox}[Exercice guidé]
Pourquoi réduire la vitesse diminue-t-il le risque ?
\begin{enumerate}
  \item Identifier la grandeur principale.
  \item Donner son unité.
  \item Écrire la loi physique pertinente.
  \item Tester l'homogénéité de la relation.
  \item Interpréter le résultat dans le contexte.
\end{enumerate}
\end{exercicebox}

\begin{corrigebox}[Indication]
La pression dynamique diminue comme $v^2$.
\end{corrigebox}


% ============================================================
\section{Partie — Aéronautique}
% ============================================================


% ------------------------------------------------------------
\subsection{71. Portance d'un avion}
% ------------------------------------------------------------

\begin{exemplebox}[Observation concrète]
Une aile maintient l'avion en l'air.
\end{exemplebox}

\begin{methodebox}[Description qualitative]
Elle dévie l'air et crée une différence de pression.
La stratégie de compréhension est toujours la même : isoler la grandeur qui varie, identifier la loi physique, puis vérifier le sens du phénomène par un ordre de grandeur.
\end{methodebox}

\begin{theorembox}[Mise en équation niveau Terminale]
La relation centrale utilisée ici est :
\[
L=\frac12\rho v^2SC_L
\]
Cette formule n'est pas seulement un outil de calcul : elle permet de prévoir comment le phénomène change lorsqu'on modifie une grandeur expérimentale.
\end{theorembox}

\begin{odgbox}[Ordre de grandeur]
La portance augmente comme $v^2$.
\end{odgbox}

\begin{attentionbox}[Limites du modèle]
Le modèle présenté repose sur des hypothèses simplificatrices : milieu homogène ou localement homogène, grandeurs moyennes, géométrie idéalisée, absence de couplages secondaires. En situation réelle, les frottements, gradients, turbulences, imperfections géométriques ou effets non linéaires peuvent modifier le résultat.
\end{attentionbox}

\begin{approfondissementbox}[Niveau prépa / Bac+3]
Le profil d'aile se modélise par circulation et théorème de Kutta-Joukowski.
On peut souvent améliorer le modèle en remplaçant une relation algébrique simple par une équation différentielle, une loi locale ou un bilan intégral. Cela permet de passer d'une explication qualitative à une prédiction quantitative.
\end{approfondissementbox}

\begin{exercicebox}[Exercice guidé]
Pourquoi un avion doit-il accélérer au décollage ?
\begin{enumerate}
  \item Identifier la grandeur principale.
  \item Donner son unité.
  \item Écrire la loi physique pertinente.
  \item Tester l'homogénéité de la relation.
  \item Interpréter le résultat dans le contexte.
\end{enumerate}
\end{exercicebox}

\begin{corrigebox}[Indication]
Il faut atteindre une portance au moins égale au poids.
\end{corrigebox}


% ------------------------------------------------------------
\subsection{72. Traînée aérodynamique}
% ------------------------------------------------------------

\begin{exemplebox}[Observation concrète]
Un véhicule résiste de plus en plus à grande vitesse.
\end{exemplebox}

\begin{methodebox}[Description qualitative]
La traînée est liée au transfert de quantité de mouvement à l'air.
La stratégie de compréhension est toujours la même : isoler la grandeur qui varie, identifier la loi physique, puis vérifier le sens du phénomène par un ordre de grandeur.
\end{methodebox}

\begin{theorembox}[Mise en équation niveau Terminale]
La relation centrale utilisée ici est :
\[
D=\frac12\rho v^2SC_D
\]
Cette formule n'est pas seulement un outil de calcul : elle permet de prévoir comment le phénomène change lorsqu'on modifie une grandeur expérimentale.
\end{theorembox}

\begin{odgbox}[Ordre de grandeur]
À vitesse élevée, la puissance nécessaire croît comme $v^3$.
\end{odgbox}

\begin{attentionbox}[Limites du modèle]
Le modèle présenté repose sur des hypothèses simplificatrices : milieu homogène ou localement homogène, grandeurs moyennes, géométrie idéalisée, absence de couplages secondaires. En situation réelle, les frottements, gradients, turbulences, imperfections géométriques ou effets non linéaires peuvent modifier le résultat.
\end{attentionbox}

\begin{approfondissementbox}[Niveau prépa / Bac+3]
La couche limite peut être laminaire ou turbulente.
On peut souvent améliorer le modèle en remplaçant une relation algébrique simple par une équation différentielle, une loi locale ou un bilan intégral. Cela permet de passer d'une explication qualitative à une prédiction quantitative.
\end{approfondissementbox}

\begin{exercicebox}[Exercice guidé]
Pourquoi les voitures rapides sont-elles profilées ?
\begin{enumerate}
  \item Identifier la grandeur principale.
  \item Donner son unité.
  \item Écrire la loi physique pertinente.
  \item Tester l'homogénéité de la relation.
  \item Interpréter le résultat dans le contexte.
\end{enumerate}
\end{exercicebox}

\begin{corrigebox}[Indication]
Pour réduire $C_D$ et donc la traînée.
\end{corrigebox}


% ============================================================
\section{Partie — Mécanique}
% ============================================================


% ------------------------------------------------------------
\subsection{73. Parachute}
% ------------------------------------------------------------

\begin{exemplebox}[Observation concrète]
Un parachute limite la vitesse de chute.
\end{exemplebox}

\begin{methodebox}[Description qualitative]
La traînée augmente jusqu'à équilibrer le poids.
La stratégie de compréhension est toujours la même : isoler la grandeur qui varie, identifier la loi physique, puis vérifier le sens du phénomène par un ordre de grandeur.
\end{methodebox}

\begin{theorembox}[Mise en équation niveau Terminale]
La relation centrale utilisée ici est :
\[
mg=\frac12\rho C_D S v_T^2
\]
Cette formule n'est pas seulement un outil de calcul : elle permet de prévoir comment le phénomène change lorsqu'on modifie une grandeur expérimentale.
\end{theorembox}

\begin{odgbox}[Ordre de grandeur]
La vitesse terminale diminue quand la surface augmente.
\end{odgbox}

\begin{attentionbox}[Limites du modèle]
Le modèle présenté repose sur des hypothèses simplificatrices : milieu homogène ou localement homogène, grandeurs moyennes, géométrie idéalisée, absence de couplages secondaires. En situation réelle, les frottements, gradients, turbulences, imperfections géométriques ou effets non linéaires peuvent modifier le résultat.
\end{attentionbox}

\begin{approfondissementbox}[Niveau prépa / Bac+3]
La vitesse vérifie une équation différentielle non linéaire.
On peut souvent améliorer le modèle en remplaçant une relation algébrique simple par une équation différentielle, une loi locale ou un bilan intégral. Cela permet de passer d'une explication qualitative à une prédiction quantitative.
\end{approfondissementbox}

\begin{exercicebox}[Exercice guidé]
Pourquoi ouvrir le parachute produit-il une forte décélération ?
\begin{enumerate}
  \item Identifier la grandeur principale.
  \item Donner son unité.
  \item Écrire la loi physique pertinente.
  \item Tester l'homogénéité de la relation.
  \item Interpréter le résultat dans le contexte.
\end{enumerate}
\end{exercicebox}

\begin{corrigebox}[Indication]
La surface $S$ augmente brutalement.
\end{corrigebox}


% ============================================================
\section{Partie — Aérospatial}
% ============================================================


% ------------------------------------------------------------
\subsection{74. Fusée et propulsion}
% ------------------------------------------------------------

\begin{exemplebox}[Observation concrète]
Une fusée avance en expulsant des gaz.
\end{exemplebox}

\begin{methodebox}[Description qualitative]
La quantité de mouvement totale se conserve.
La stratégie de compréhension est toujours la même : isoler la grandeur qui varie, identifier la loi physique, puis vérifier le sens du phénomène par un ordre de grandeur.
\end{methodebox}

\begin{theorembox}[Mise en équation niveau Terminale]
La relation centrale utilisée ici est :
\[
F\simeq \dot m v_e
\]
Cette formule n'est pas seulement un outil de calcul : elle permet de prévoir comment le phénomène change lorsqu'on modifie une grandeur expérimentale.
\end{theorembox}

\begin{odgbox}[Ordre de grandeur]
La poussée dépend du débit massique et de la vitesse d'éjection.
\end{odgbox}

\begin{attentionbox}[Limites du modèle]
Le modèle présenté repose sur des hypothèses simplificatrices : milieu homogène ou localement homogène, grandeurs moyennes, géométrie idéalisée, absence de couplages secondaires. En situation réelle, les frottements, gradients, turbulences, imperfections géométriques ou effets non linéaires peuvent modifier le résultat.
\end{attentionbox}

\begin{approfondissementbox}[Niveau prépa / Bac+3]
L'équation de Tsiolkovski donne $\Delta v=v_e\ln(m_0/m_f)$.
On peut souvent améliorer le modèle en remplaçant une relation algébrique simple par une équation différentielle, une loi locale ou un bilan intégral. Cela permet de passer d'une explication qualitative à une prédiction quantitative.
\end{approfondissementbox}

\begin{exercicebox}[Exercice guidé]
Pourquoi une fusée fonctionne-t-elle dans le vide ?
\begin{enumerate}
  \item Identifier la grandeur principale.
  \item Donner son unité.
  \item Écrire la loi physique pertinente.
  \item Tester l'homogénéité de la relation.
  \item Interpréter le résultat dans le contexte.
\end{enumerate}
\end{exercicebox}

\begin{corrigebox}[Indication]
Elle pousse sur ses propres gaz, pas sur l'air.
\end{corrigebox}


% ============================================================
\section{Partie — Gravitation}
% ============================================================


% ------------------------------------------------------------
\subsection{75. Satellite en orbite}
% ------------------------------------------------------------

\begin{exemplebox}[Observation concrète]
Un satellite tombe sans toucher le sol.
\end{exemplebox}

\begin{methodebox}[Description qualitative]
Sa vitesse horizontale transforme la chute en orbite.
La stratégie de compréhension est toujours la même : isoler la grandeur qui varie, identifier la loi physique, puis vérifier le sens du phénomène par un ordre de grandeur.
\end{methodebox}

\begin{theorembox}[Mise en équation niveau Terminale]
La relation centrale utilisée ici est :
\[
v=\sqrt{\frac{GM}{r}}
\]
Cette formule n'est pas seulement un outil de calcul : elle permet de prévoir comment le phénomène change lorsqu'on modifie une grandeur expérimentale.
\end{theorembox}

\begin{odgbox}[Ordre de grandeur]
En orbite basse, $v$ vaut environ $\SI{7.8}{km.s^{-1}}$.
\end{odgbox}

\begin{attentionbox}[Limites du modèle]
Le modèle présenté repose sur des hypothèses simplificatrices : milieu homogène ou localement homogène, grandeurs moyennes, géométrie idéalisée, absence de couplages secondaires. En situation réelle, les frottements, gradients, turbulences, imperfections géométriques ou effets non linéaires peuvent modifier le résultat.
\end{attentionbox}

\begin{approfondissementbox}[Niveau prépa / Bac+3]
Les orbites képlériennes sont des coniques dans un potentiel $1/r$.
On peut souvent améliorer le modèle en remplaçant une relation algébrique simple par une équation différentielle, une loi locale ou un bilan intégral. Cela permet de passer d'une explication qualitative à une prédiction quantitative.
\end{approfondissementbox}

\begin{exercicebox}[Exercice guidé]
Pourquoi un satellite n'a-t-il pas besoin de moteur permanent ?
\begin{enumerate}
  \item Identifier la grandeur principale.
  \item Donner son unité.
  \item Écrire la loi physique pertinente.
  \item Tester l'homogénéité de la relation.
  \item Interpréter le résultat dans le contexte.
\end{enumerate}
\end{exercicebox}

\begin{corrigebox}[Indication]
Sans frottements, l'inertie et la gravité suffisent.
\end{corrigebox}


% ------------------------------------------------------------
\subsection{76. Marées}
% ------------------------------------------------------------

\begin{exemplebox}[Observation concrète]
La mer monte et descend périodiquement.
\end{exemplebox}

\begin{methodebox}[Description qualitative]
La force gravitationnelle lunaire varie d'un point à l'autre de la Terre.
La stratégie de compréhension est toujours la même : isoler la grandeur qui varie, identifier la loi physique, puis vérifier le sens du phénomène par un ordre de grandeur.
\end{methodebox}

\begin{theorembox}[Mise en équation niveau Terminale]
La relation centrale utilisée ici est :
\[
F\propto\frac1{r^2}
\]
Cette formule n'est pas seulement un outil de calcul : elle permet de prévoir comment le phénomène change lorsqu'on modifie une grandeur expérimentale.
\end{theorembox}

\begin{odgbox}[Ordre de grandeur]
La période dominante est d'environ 12 h 25.
\end{odgbox}

\begin{attentionbox}[Limites du modèle]
Le modèle présenté repose sur des hypothèses simplificatrices : milieu homogène ou localement homogène, grandeurs moyennes, géométrie idéalisée, absence de couplages secondaires. En situation réelle, les frottements, gradients, turbulences, imperfections géométriques ou effets non linéaires peuvent modifier le résultat.
\end{attentionbox}

\begin{approfondissementbox}[Niveau prépa / Bac+3]
Le potentiel de marée est lié au gradient du champ gravitationnel.
On peut souvent améliorer le modèle en remplaçant une relation algébrique simple par une équation différentielle, une loi locale ou un bilan intégral. Cela permet de passer d'une explication qualitative à une prédiction quantitative.
\end{approfondissementbox}

\begin{exercicebox}[Exercice guidé]
Pourquoi la Lune crée-t-elle deux bourrelets ?
\begin{enumerate}
  \item Identifier la grandeur principale.
  \item Donner son unité.
  \item Écrire la loi physique pertinente.
  \item Tester l'homogénéité de la relation.
  \item Interpréter le résultat dans le contexte.
\end{enumerate}
\end{exercicebox}

\begin{corrigebox}[Indication]
C'est un effet différentiel de gravitation et d'inertie.
\end{corrigebox}


% ============================================================
\section{Partie — Mécanique}
% ============================================================


% ------------------------------------------------------------
\subsection{77. Pendule}
% ------------------------------------------------------------

\begin{exemplebox}[Observation concrète]
Un pendule oscille autour d'une position d'équilibre.
\end{exemplebox}

\begin{methodebox}[Description qualitative]
Le poids crée un couple de rappel.
La stratégie de compréhension est toujours la même : isoler la grandeur qui varie, identifier la loi physique, puis vérifier le sens du phénomène par un ordre de grandeur.
\end{methodebox}

\begin{theorembox}[Mise en équation niveau Terminale]
La relation centrale utilisée ici est :
\[
T\simeq2\pi\sqrt{\frac{\ell}{g}}
\]
Cette formule n'est pas seulement un outil de calcul : elle permet de prévoir comment le phénomène change lorsqu'on modifie une grandeur expérimentale.
\end{theorembox}

\begin{odgbox}[Ordre de grandeur]
Un pendule d'un mètre a une période d'environ 2 s.
\end{odgbox}

\begin{attentionbox}[Limites du modèle]
Le modèle présenté repose sur des hypothèses simplificatrices : milieu homogène ou localement homogène, grandeurs moyennes, géométrie idéalisée, absence de couplages secondaires. En situation réelle, les frottements, gradients, turbulences, imperfections géométriques ou effets non linéaires peuvent modifier le résultat.
\end{attentionbox}

\begin{approfondissementbox}[Niveau prépa / Bac+3]
Pour grands angles, l'équation est non linéaire : $	heta''+(g/\ell)\sin	heta=0$.
On peut souvent améliorer le modèle en remplaçant une relation algébrique simple par une équation différentielle, une loi locale ou un bilan intégral. Cela permet de passer d'une explication qualitative à une prédiction quantitative.
\end{approfondissementbox}

\begin{exercicebox}[Exercice guidé]
Pourquoi la période dépend-elle peu de la masse ?
\begin{enumerate}
  \item Identifier la grandeur principale.
  \item Donner son unité.
  \item Écrire la loi physique pertinente.
  \item Tester l'homogénéité de la relation.
  \item Interpréter le résultat dans le contexte.
\end{enumerate}
\end{exercicebox}

\begin{corrigebox}[Indication]
La masse se simplifie dans l'équation du mouvement.
\end{corrigebox}


% ============================================================
\section{Partie — Rotation}
% ============================================================


% ------------------------------------------------------------
\subsection{78. Gyroscope}
% ------------------------------------------------------------

\begin{exemplebox}[Observation concrète]
Une roue en rotation garde son axe presque fixe.
\end{exemplebox}

\begin{methodebox}[Description qualitative]
Le moment cinétique résiste aux changements d'orientation.
La stratégie de compréhension est toujours la même : isoler la grandeur qui varie, identifier la loi physique, puis vérifier le sens du phénomène par un ordre de grandeur.
\end{methodebox}

\begin{theorembox}[Mise en équation niveau Terminale]
La relation centrale utilisée ici est :
\[
\vec L=I\vec\omega
\]
Cette formule n'est pas seulement un outil de calcul : elle permet de prévoir comment le phénomène change lorsqu'on modifie une grandeur expérimentale.
\end{theorembox}

\begin{odgbox}[Ordre de grandeur]
Un gyroscope rapide est plus stable.
\end{odgbox}

\begin{attentionbox}[Limites du modèle]
Le modèle présenté repose sur des hypothèses simplificatrices : milieu homogène ou localement homogène, grandeurs moyennes, géométrie idéalisée, absence de couplages secondaires. En situation réelle, les frottements, gradients, turbulences, imperfections géométriques ou effets non linéaires peuvent modifier le résultat.
\end{attentionbox}

\begin{approfondissementbox}[Niveau prépa / Bac+3]
La précession vérifie $
ec 	au=\dd
ec L/\dd t$.
On peut souvent améliorer le modèle en remplaçant une relation algébrique simple par une équation différentielle, une loi locale ou un bilan intégral. Cela permet de passer d'une explication qualitative à une prédiction quantitative.
\end{approfondissementbox}

\begin{exercicebox}[Exercice guidé]
Pourquoi une roue de vélo aide-t-elle à stabiliser ?
\begin{enumerate}
  \item Identifier la grandeur principale.
  \item Donner son unité.
  \item Écrire la loi physique pertinente.
  \item Tester l'homogénéité de la relation.
  \item Interpréter le résultat dans le contexte.
\end{enumerate}
\end{exercicebox}

\begin{corrigebox}[Indication]
Son moment cinétique s'oppose aux changements brusques.
\end{corrigebox}


% ------------------------------------------------------------
\subsection{79. Toupie}
% ------------------------------------------------------------

\begin{exemplebox}[Observation concrète]
Une toupie reste debout en tournant.
\end{exemplebox}

\begin{methodebox}[Description qualitative]
La rotation transforme la chute en précession.
La stratégie de compréhension est toujours la même : isoler la grandeur qui varie, identifier la loi physique, puis vérifier le sens du phénomène par un ordre de grandeur.
\end{methodebox}

\begin{theorembox}[Mise en équation niveau Terminale]
La relation centrale utilisée ici est :
\[
\Omega\simeq\frac{mgr}{L}
\]
Cette formule n'est pas seulement un outil de calcul : elle permet de prévoir comment le phénomène change lorsqu'on modifie une grandeur expérimentale.
\end{theorembox}

\begin{odgbox}[Ordre de grandeur]
Quand la toupie ralentit, elle finit par tomber.
\end{odgbox}

\begin{attentionbox}[Limites du modèle]
Le modèle présenté repose sur des hypothèses simplificatrices : milieu homogène ou localement homogène, grandeurs moyennes, géométrie idéalisée, absence de couplages secondaires. En situation réelle, les frottements, gradients, turbulences, imperfections géométriques ou effets non linéaires peuvent modifier le résultat.
\end{attentionbox}

\begin{approfondissementbox}[Niveau prépa / Bac+3]
Le mouvement complet d'une toupie lourde est un problème de mécanique du solide.
On peut souvent améliorer le modèle en remplaçant une relation algébrique simple par une équation différentielle, une loi locale ou un bilan intégral. Cela permet de passer d'une explication qualitative à une prédiction quantitative.
\end{approfondissementbox}

\begin{exercicebox}[Exercice guidé]
Pourquoi une toupie immobile tombe-elle ?
\begin{enumerate}
  \item Identifier la grandeur principale.
  \item Donner son unité.
  \item Écrire la loi physique pertinente.
  \item Tester l'homogénéité de la relation.
  \item Interpréter le résultat dans le contexte.
\end{enumerate}
\end{exercicebox}

\begin{corrigebox}[Indication]
Sans moment cinétique, le couple du poids la bascule.
\end{corrigebox}


% ============================================================
\section{Partie — Transport}
% ============================================================


% ------------------------------------------------------------
\subsection{80. Stabilité d'un vélo}
% ------------------------------------------------------------

\begin{exemplebox}[Observation concrète]
Un vélo en mouvement est plus stable qu'à l'arrêt.
\end{exemplebox}

\begin{methodebox}[Description qualitative]
Géométrie de direction, effet gyroscopique et contrôle du cycliste participent.
La stratégie de compréhension est toujours la même : isoler la grandeur qui varie, identifier la loi physique, puis vérifier le sens du phénomène par un ordre de grandeur.
\end{methodebox}

\begin{theorembox}[Mise en équation niveau Terminale]
La relation centrale utilisée ici est :
\[
\tau=\frac{\dd L}{\dd t}
\]
Cette formule n'est pas seulement un outil de calcul : elle permet de prévoir comment le phénomène change lorsqu'on modifie une grandeur expérimentale.
\end{theorembox}

\begin{odgbox}[Ordre de grandeur]
L'effet gyroscopique seul n'explique pas toute la stabilité.
\end{odgbox}

\begin{attentionbox}[Limites du modèle]
Le modèle présenté repose sur des hypothèses simplificatrices : milieu homogène ou localement homogène, grandeurs moyennes, géométrie idéalisée, absence de couplages secondaires. En situation réelle, les frottements, gradients, turbulences, imperfections géométriques ou effets non linéaires peuvent modifier le résultat.
\end{attentionbox}

\begin{approfondissementbox}[Niveau prépa / Bac+3]
Les équations linéarisées du vélo couplent roulis, lacet et direction.
On peut souvent améliorer le modèle en remplaçant une relation algébrique simple par une équation différentielle, une loi locale ou un bilan intégral. Cela permet de passer d'une explication qualitative à une prédiction quantitative.
\end{approfondissementbox}

\begin{exercicebox}[Exercice guidé]
Pourquoi faut-il tourner légèrement pour rattraper une chute ?
\begin{enumerate}
  \item Identifier la grandeur principale.
  \item Donner son unité.
  \item Écrire la loi physique pertinente.
  \item Tester l'homogénéité de la relation.
  \item Interpréter le résultat dans le contexte.
\end{enumerate}
\end{exercicebox}

\begin{corrigebox}[Indication]
Pour replacer le point de contact sous le centre de masse.
\end{corrigebox}


% ============================================================
\section{Partie X — Gagner grâce à la physique : jeux, sports et stratégie}
% ============================================================


% ------------------------------------------------------------
\subsection{Jeu/Sport 1. Billard — Choisir l'angle de tir}
% ------------------------------------------------------------

\begin{exemplebox}[Observation]
Sans effet, la réflexion sur la bande vérifie approximativement $i=r$.
\end{exemplebox}

\begin{theorembox}[Loi physique utile]
\[
m\vec v_1+m\vec v_2=m\vec v_1'+m\vec v_2'
\]
La loi physique ne donne pas une recette magique, mais elle permet de réduire l'incertitude : choisir un angle, prévoir un rebond, anticiper une courbure, ou comprendre une perte d'énergie.
\end{theorembox}

\begin{methodebox}[Comment gagner grâce à la physique]
Viser une bande revient à utiliser la symétrie géométrique ; les effets modifient ensuite la trajectoire.
Il faut ensuite adapter ce principe aux frottements, à la rotation, à l'imprécision du geste et aux conditions réelles du jeu.
\end{methodebox}

\begin{approfondissementbox}[Niveau prépa / Bac+3]
Une analyse complète introduit les degrés de liberté de translation et de rotation. L'énergie cinétique totale peut s'écrire
\[
E_c=\frac12mv^2+\frac12I\omega^2.
\]
Lors d'un choc ou d'un rebond, une partie de cette énergie est transférée, une autre dissipée. La stratégie consiste à contrôler non seulement la vitesse du centre de masse, mais aussi la rotation.
\end{approfondissementbox}

\begin{exercicebox}[Mini-problème]
\begin{enumerate}
  \item Décrire la situation de jeu en termes de forces et de vitesse.
  \item Identifier le rôle de la rotation.
  \item Dire quelle grandeur le joueur peut contrôler directement.
  \item Donner une conséquence stratégique.
  \item Citer une limite du modèle idéal.
\end{enumerate}
\end{exercicebox}

\begin{corrigebox}[Indication]
La bonne réponse doit distinguer ce que prédit la loi idéale et ce que modifient les pertes réelles : frottements, déformation, rotation parasite, qualité du support, précision du geste.
\end{corrigebox}


% ------------------------------------------------------------
\subsection{Jeu/Sport 2. Bowling — Faire courber la boule}
% ------------------------------------------------------------

\begin{exemplebox}[Observation]
La rotation et le frottement avec la piste font passer la boule d'une phase de glissement à une phase de roulement.
\end{exemplebox}

\begin{theorembox}[Loi physique utile]
\[
\vec\tau=I\vec\alpha
\]
La loi physique ne donne pas une recette magique, mais elle permet de réduire l'incertitude : choisir un angle, prévoir un rebond, anticiper une courbure, ou comprendre une perte d'énergie.
\end{theorembox}

\begin{methodebox}[Comment gagner grâce à la physique]
Chercher un angle d'entrée de l'ordre de quelques degrés augmente les chances de strike.
Il faut ensuite adapter ce principe aux frottements, à la rotation, à l'imprécision du geste et aux conditions réelles du jeu.
\end{methodebox}

\begin{approfondissementbox}[Niveau prépa / Bac+3]
Une analyse complète introduit les degrés de liberté de translation et de rotation. L'énergie cinétique totale peut s'écrire
\[
E_c=\frac12mv^2+\frac12I\omega^2.
\]
Lors d'un choc ou d'un rebond, une partie de cette énergie est transférée, une autre dissipée. La stratégie consiste à contrôler non seulement la vitesse du centre de masse, mais aussi la rotation.
\end{approfondissementbox}

\begin{exercicebox}[Mini-problème]
\begin{enumerate}
  \item Décrire la situation de jeu en termes de forces et de vitesse.
  \item Identifier le rôle de la rotation.
  \item Dire quelle grandeur le joueur peut contrôler directement.
  \item Donner une conséquence stratégique.
  \item Citer une limite du modèle idéal.
\end{enumerate}
\end{exercicebox}

\begin{corrigebox}[Indication]
La bonne réponse doit distinguer ce que prédit la loi idéale et ce que modifient les pertes réelles : frottements, déformation, rotation parasite, qualité du support, précision du geste.
\end{corrigebox}


% ------------------------------------------------------------
\subsection{Jeu/Sport 3. Fléchettes — Stabiliser le projectile}
% ------------------------------------------------------------

\begin{exemplebox}[Observation]
Les ailettes déplacent le centre de pression et stabilisent l'orientation.
\end{exemplebox}

\begin{theorembox}[Loi physique utile]
\[
y=x\tan\theta-\frac{gx^2}{2v_0^2\cos^2\theta}
\]
La loi physique ne donne pas une recette magique, mais elle permet de réduire l'incertitude : choisir un angle, prévoir un rebond, anticiper une courbure, ou comprendre une perte d'énergie.
\end{theorembox}

\begin{methodebox}[Comment gagner grâce à la physique]
Une vitesse initiale suffisante réduit la chute et la sensibilité aux erreurs d'angle.
Il faut ensuite adapter ce principe aux frottements, à la rotation, à l'imprécision du geste et aux conditions réelles du jeu.
\end{methodebox}

\begin{approfondissementbox}[Niveau prépa / Bac+3]
Une analyse complète introduit les degrés de liberté de translation et de rotation. L'énergie cinétique totale peut s'écrire
\[
E_c=\frac12mv^2+\frac12I\omega^2.
\]
Lors d'un choc ou d'un rebond, une partie de cette énergie est transférée, une autre dissipée. La stratégie consiste à contrôler non seulement la vitesse du centre de masse, mais aussi la rotation.
\end{approfondissementbox}

\begin{exercicebox}[Mini-problème]
\begin{enumerate}
  \item Décrire la situation de jeu en termes de forces et de vitesse.
  \item Identifier le rôle de la rotation.
  \item Dire quelle grandeur le joueur peut contrôler directement.
  \item Donner une conséquence stratégique.
  \item Citer une limite du modèle idéal.
\end{enumerate}
\end{exercicebox}

\begin{corrigebox}[Indication]
La bonne réponse doit distinguer ce que prédit la loi idéale et ce que modifient les pertes réelles : frottements, déformation, rotation parasite, qualité du support, précision du geste.
\end{corrigebox}


% ------------------------------------------------------------
\subsection{Jeu/Sport 4. Pétanque — Pointer ou tirer}
% ------------------------------------------------------------

\begin{exemplebox}[Observation]
Le choix dépend de la vitesse, de l'angle d'arrivée, du frottement et des chocs.
\end{exemplebox}

\begin{theorembox}[Loi physique utile]
\[
F_f=\mu mg
\]
La loi physique ne donne pas une recette magique, mais elle permet de réduire l'incertitude : choisir un angle, prévoir un rebond, anticiper une courbure, ou comprendre une perte d'énergie.
\end{theorembox}

\begin{methodebox}[Comment gagner grâce à la physique]
Pour pointer, on contrôle surtout la vitesse résiduelle après frottement.
Il faut ensuite adapter ce principe aux frottements, à la rotation, à l'imprécision du geste et aux conditions réelles du jeu.
\end{methodebox}

\begin{approfondissementbox}[Niveau prépa / Bac+3]
Une analyse complète introduit les degrés de liberté de translation et de rotation. L'énergie cinétique totale peut s'écrire
\[
E_c=\frac12mv^2+\frac12I\omega^2.
\]
Lors d'un choc ou d'un rebond, une partie de cette énergie est transférée, une autre dissipée. La stratégie consiste à contrôler non seulement la vitesse du centre de masse, mais aussi la rotation.
\end{approfondissementbox}

\begin{exercicebox}[Mini-problème]
\begin{enumerate}
  \item Décrire la situation de jeu en termes de forces et de vitesse.
  \item Identifier le rôle de la rotation.
  \item Dire quelle grandeur le joueur peut contrôler directement.
  \item Donner une conséquence stratégique.
  \item Citer une limite du modèle idéal.
\end{enumerate}
\end{exercicebox}

\begin{corrigebox}[Indication]
La bonne réponse doit distinguer ce que prédit la loi idéale et ce que modifient les pertes réelles : frottements, déformation, rotation parasite, qualité du support, précision du geste.
\end{corrigebox}


% ------------------------------------------------------------
\subsection{Jeu/Sport 5. Tennis — Topspin gagnant}
% ------------------------------------------------------------

\begin{exemplebox}[Observation]
Le lift crée une force de Magnus vers le bas et permet de frapper fort.
\end{exemplebox}

\begin{theorembox}[Loi physique utile]
\[
\vec F_M\propto \vec\omega\times\vec v
\]
La loi physique ne donne pas une recette magique, mais elle permet de réduire l'incertitude : choisir un angle, prévoir un rebond, anticiper une courbure, ou comprendre une perte d'énergie.
\end{theorembox}

\begin{methodebox}[Comment gagner grâce à la physique]
Un topspin augmente la marge au-dessus du filet tout en ramenant la balle dans le court.
Il faut ensuite adapter ce principe aux frottements, à la rotation, à l'imprécision du geste et aux conditions réelles du jeu.
\end{methodebox}

\begin{approfondissementbox}[Niveau prépa / Bac+3]
Une analyse complète introduit les degrés de liberté de translation et de rotation. L'énergie cinétique totale peut s'écrire
\[
E_c=\frac12mv^2+\frac12I\omega^2.
\]
Lors d'un choc ou d'un rebond, une partie de cette énergie est transférée, une autre dissipée. La stratégie consiste à contrôler non seulement la vitesse du centre de masse, mais aussi la rotation.
\end{approfondissementbox}

\begin{exercicebox}[Mini-problème]
\begin{enumerate}
  \item Décrire la situation de jeu en termes de forces et de vitesse.
  \item Identifier le rôle de la rotation.
  \item Dire quelle grandeur le joueur peut contrôler directement.
  \item Donner une conséquence stratégique.
  \item Citer une limite du modèle idéal.
\end{enumerate}
\end{exercicebox}

\begin{corrigebox}[Indication]
La bonne réponse doit distinguer ce que prédit la loi idéale et ce que modifient les pertes réelles : frottements, déformation, rotation parasite, qualité du support, précision du geste.
\end{corrigebox}


% ------------------------------------------------------------
\subsection{Jeu/Sport 6. Football — Coup franc brossé}
% ------------------------------------------------------------

\begin{exemplebox}[Observation]
La rotation latérale crée une force latérale.
\end{exemplebox}

\begin{theorembox}[Loi physique utile]
\[
\vec F_M\propto \vec\omega\times\vec v
\]
La loi physique ne donne pas une recette magique, mais elle permet de réduire l'incertitude : choisir un angle, prévoir un rebond, anticiper une courbure, ou comprendre une perte d'énergie.
\end{theorembox}

\begin{methodebox}[Comment gagner grâce à la physique]
Plus la vitesse et la rotation sont grandes, plus la courbure est forte, jusqu'aux limites de traînée.
Il faut ensuite adapter ce principe aux frottements, à la rotation, à l'imprécision du geste et aux conditions réelles du jeu.
\end{methodebox}

\begin{approfondissementbox}[Niveau prépa / Bac+3]
Une analyse complète introduit les degrés de liberté de translation et de rotation. L'énergie cinétique totale peut s'écrire
\[
E_c=\frac12mv^2+\frac12I\omega^2.
\]
Lors d'un choc ou d'un rebond, une partie de cette énergie est transférée, une autre dissipée. La stratégie consiste à contrôler non seulement la vitesse du centre de masse, mais aussi la rotation.
\end{approfondissementbox}

\begin{exercicebox}[Mini-problème]
\begin{enumerate}
  \item Décrire la situation de jeu en termes de forces et de vitesse.
  \item Identifier le rôle de la rotation.
  \item Dire quelle grandeur le joueur peut contrôler directement.
  \item Donner une conséquence stratégique.
  \item Citer une limite du modèle idéal.
\end{enumerate}
\end{exercicebox}

\begin{corrigebox}[Indication]
La bonne réponse doit distinguer ce que prédit la loi idéale et ce que modifient les pertes réelles : frottements, déformation, rotation parasite, qualité du support, précision du geste.
\end{corrigebox}


% ------------------------------------------------------------
\subsection{Jeu/Sport 7. Basketball — Angle de tir}
% ------------------------------------------------------------

\begin{exemplebox}[Observation]
Un tir plus arqué augmente la taille apparente du panier.
\end{exemplebox}

\begin{theorembox}[Loi physique utile]
\[
y=x\tan\theta-\frac{gx^2}{2v_0^2\cos^2\theta}
\]
La loi physique ne donne pas une recette magique, mais elle permet de réduire l'incertitude : choisir un angle, prévoir un rebond, anticiper une courbure, ou comprendre une perte d'énergie.
\end{theorembox}

\begin{methodebox}[Comment gagner grâce à la physique]
La rotation arrière stabilise le contact avec l'arceau ou la planche.
Il faut ensuite adapter ce principe aux frottements, à la rotation, à l'imprécision du geste et aux conditions réelles du jeu.
\end{methodebox}

\begin{approfondissementbox}[Niveau prépa / Bac+3]
Une analyse complète introduit les degrés de liberté de translation et de rotation. L'énergie cinétique totale peut s'écrire
\[
E_c=\frac12mv^2+\frac12I\omega^2.
\]
Lors d'un choc ou d'un rebond, une partie de cette énergie est transférée, une autre dissipée. La stratégie consiste à contrôler non seulement la vitesse du centre de masse, mais aussi la rotation.
\end{approfondissementbox}

\begin{exercicebox}[Mini-problème]
\begin{enumerate}
  \item Décrire la situation de jeu en termes de forces et de vitesse.
  \item Identifier le rôle de la rotation.
  \item Dire quelle grandeur le joueur peut contrôler directement.
  \item Donner une conséquence stratégique.
  \item Citer une limite du modèle idéal.
\end{enumerate}
\end{exercicebox}

\begin{corrigebox}[Indication]
La bonne réponse doit distinguer ce que prédit la loi idéale et ce que modifient les pertes réelles : frottements, déformation, rotation parasite, qualité du support, précision du geste.
\end{corrigebox}


% ------------------------------------------------------------
\subsection{Jeu/Sport 8. Golf — Backspin et alvéoles}
% ------------------------------------------------------------

\begin{exemplebox}[Observation]
Le backspin crée de la portance et les alvéoles réduisent la traînée de pression.
\end{exemplebox}

\begin{theorembox}[Loi physique utile]
\[
L=\frac12\rho v^2SC_L
\]
La loi physique ne donne pas une recette magique, mais elle permet de réduire l'incertitude : choisir un angle, prévoir un rebond, anticiper une courbure, ou comprendre une perte d'énergie.
\end{theorembox}

\begin{methodebox}[Comment gagner grâce à la physique]
Un bon compromis portance-traînée maximise la portée.
Il faut ensuite adapter ce principe aux frottements, à la rotation, à l'imprécision du geste et aux conditions réelles du jeu.
\end{methodebox}

\begin{approfondissementbox}[Niveau prépa / Bac+3]
Une analyse complète introduit les degrés de liberté de translation et de rotation. L'énergie cinétique totale peut s'écrire
\[
E_c=\frac12mv^2+\frac12I\omega^2.
\]
Lors d'un choc ou d'un rebond, une partie de cette énergie est transférée, une autre dissipée. La stratégie consiste à contrôler non seulement la vitesse du centre de masse, mais aussi la rotation.
\end{approfondissementbox}

\begin{exercicebox}[Mini-problème]
\begin{enumerate}
  \item Décrire la situation de jeu en termes de forces et de vitesse.
  \item Identifier le rôle de la rotation.
  \item Dire quelle grandeur le joueur peut contrôler directement.
  \item Donner une conséquence stratégique.
  \item Citer une limite du modèle idéal.
\end{enumerate}
\end{exercicebox}

\begin{corrigebox}[Indication]
La bonne réponse doit distinguer ce que prédit la loi idéale et ce que modifient les pertes réelles : frottements, déformation, rotation parasite, qualité du support, précision du geste.
\end{corrigebox}


% ------------------------------------------------------------
\subsection{Jeu/Sport 9. Ping-pong — Lire l'effet adverse}
% ------------------------------------------------------------

\begin{exemplebox}[Observation]
La rotation change le rebond et la courbure.
\end{exemplebox}

\begin{theorembox}[Loi physique utile]
\[
\vec F_M\propto \vec\omega\times\vec v
\]
La loi physique ne donne pas une recette magique, mais elle permet de réduire l'incertitude : choisir un angle, prévoir un rebond, anticiper une courbure, ou comprendre une perte d'énergie.
\end{theorembox}

\begin{methodebox}[Comment gagner grâce à la physique]
Observer le geste adverse permet d'anticiper topspin, backspin ou effet latéral.
Il faut ensuite adapter ce principe aux frottements, à la rotation, à l'imprécision du geste et aux conditions réelles du jeu.
\end{methodebox}

\begin{approfondissementbox}[Niveau prépa / Bac+3]
Une analyse complète introduit les degrés de liberté de translation et de rotation. L'énergie cinétique totale peut s'écrire
\[
E_c=\frac12mv^2+\frac12I\omega^2.
\]
Lors d'un choc ou d'un rebond, une partie de cette énergie est transférée, une autre dissipée. La stratégie consiste à contrôler non seulement la vitesse du centre de masse, mais aussi la rotation.
\end{approfondissementbox}

\begin{exercicebox}[Mini-problème]
\begin{enumerate}
  \item Décrire la situation de jeu en termes de forces et de vitesse.
  \item Identifier le rôle de la rotation.
  \item Dire quelle grandeur le joueur peut contrôler directement.
  \item Donner une conséquence stratégique.
  \item Citer une limite du modèle idéal.
\end{enumerate}
\end{exercicebox}

\begin{corrigebox}[Indication]
La bonne réponse doit distinguer ce que prédit la loi idéale et ce que modifient les pertes réelles : frottements, déformation, rotation parasite, qualité du support, précision du geste.
\end{corrigebox}


% ------------------------------------------------------------
\subsection{Jeu/Sport 10. Kerbal Space Program — Optimiser une mise en orbite}
% ------------------------------------------------------------

\begin{exemplebox}[Observation]
Une orbite se gagne en vitesse horizontale plus qu'en altitude pure.
\end{exemplebox}

\begin{theorembox}[Loi physique utile]
\[
v_{\mathrm{orb}}=\sqrt{\frac{GM}{r}}
\]
La loi physique ne donne pas une recette magique, mais elle permet de réduire l'incertitude : choisir un angle, prévoir un rebond, anticiper une courbure, ou comprendre une perte d'énergie.
\end{theorembox}

\begin{methodebox}[Comment gagner grâce à la physique]
Faire un gravity turn limite les pertes gravitationnelles et aérodynamiques.
Il faut ensuite adapter ce principe aux frottements, à la rotation, à l'imprécision du geste et aux conditions réelles du jeu.
\end{methodebox}

\begin{approfondissementbox}[Niveau prépa / Bac+3]
Une analyse complète introduit les degrés de liberté de translation et de rotation. L'énergie cinétique totale peut s'écrire
\[
E_c=\frac12mv^2+\frac12I\omega^2.
\]
Lors d'un choc ou d'un rebond, une partie de cette énergie est transférée, une autre dissipée. La stratégie consiste à contrôler non seulement la vitesse du centre de masse, mais aussi la rotation.
\end{approfondissementbox}

\begin{exercicebox}[Mini-problème]
\begin{enumerate}
  \item Décrire la situation de jeu en termes de forces et de vitesse.
  \item Identifier le rôle de la rotation.
  \item Dire quelle grandeur le joueur peut contrôler directement.
  \item Donner une conséquence stratégique.
  \item Citer une limite du modèle idéal.
\end{enumerate}
\end{exercicebox}

\begin{corrigebox}[Indication]
La bonne réponse doit distinguer ce que prédit la loi idéale et ce que modifient les pertes réelles : frottements, déformation, rotation parasite, qualité du support, précision du geste.
\end{corrigebox}


% ============================================================
\section{Schémas TikZ simples}
% ============================================================

\subsection{Réfraction air-eau}
\begin{center}
\begin{tikzpicture}[scale=1.0]
  \draw[thick] (-3,0)--(3,0);
  \node at (2.6,0.25) {interface};
  \draw[dashed] (0,-2)--(0,2);
  \draw[->,very thick,blue] (-2,1.5)--(0,0);
  \draw[->,very thick,blue] (0,0)--(1.0,-1.8);
  \draw (0.4,0) arc (0:55:0.4);
  \node at (0.75,0.35) {$i$};
  \draw (0,-0.45) arc (-90:-61:0.45);
  \node at (0.25,-0.8) {$r$};
  \node at (-2.4,1.2) {air $n_1$};
  \node at (2.2,-1.2) {eau $n_2$};
\end{tikzpicture}
\end{center}

\subsection{Fibre optique}
\begin{center}
\begin{tikzpicture}[scale=1.0]
  \draw[thick] (0,0)--(8,0);
  \draw[thick] (0,1)--(8,1);
  \draw[thick] (0,-1)--(8,-1);
  \node at (7.3,0.25) {cœur $n_c$};
  \node at (7.3,1.25) {gaine $n_g$};
  \draw[->,very thick,red] (0.2,0)--(1.5,0.7)--(2.8,-0.7)--(4.1,0.7)--(5.4,-0.7)--(6.7,0.4);
  \node at (3.9,-1.35) {réflexions totales internes};
\end{tikzpicture}
\end{center}

\subsection{Diffraction par une fente}
\begin{center}
\begin{tikzpicture}[scale=1.0]
  \draw[thick] (0,-2)--(0,-0.3);
  \draw[thick] (0,0.3)--(0,2);
  \node at (-0.25,0) {fente};
  \draw[->,blue] (-3,1)--(0,0.2);
  \draw[->,blue] (-3,0)--(0,0);
  \draw[->,blue] (-3,-1)--(0,-0.2);
  \draw[->,red,thick] (0,0)--(3,1.2);
  \draw[->,red,thick] (0,0)--(3,0);
  \draw[->,red,thick] (0,0)--(3,-1.2);
  \draw[dashed] (3,-2)--(3,2);
  \node at (3.3,0) {écran};
\end{tikzpicture}
\end{center}


% ============================================================
\section{Partie XI — Exercices progressifs}
% ============================================================

\subsection{Exercices courts niveau Terminale}

\begin{exercicebox}[Exercice court 1]
On considère un phénomène étudié dans le document. 
\begin{enumerate}
  \item Identifier la grandeur physique principale.
  \item Donner son unité SI.
  \item Écrire la loi simple associée.
  \item Vérifier l'homogénéité.
  \item Donner une phrase d'interprétation.
\end{enumerate}
\end{exercicebox}

\begin{corrigebox}[Indication 1]
La correction doit toujours contenir une loi, une unité, un ordre de grandeur et une phrase de sens physique.
\end{corrigebox}

\begin{exercicebox}[Exercice court 2]
On considère un phénomène étudié dans le document. 
\begin{enumerate}
  \item Identifier la grandeur physique principale.
  \item Donner son unité SI.
  \item Écrire la loi simple associée.
  \item Vérifier l'homogénéité.
  \item Donner une phrase d'interprétation.
\end{enumerate}
\end{exercicebox}

\begin{corrigebox}[Indication 2]
La correction doit toujours contenir une loi, une unité, un ordre de grandeur et une phrase de sens physique.
\end{corrigebox}

\begin{exercicebox}[Exercice court 3]
On considère un phénomène étudié dans le document. 
\begin{enumerate}
  \item Identifier la grandeur physique principale.
  \item Donner son unité SI.
  \item Écrire la loi simple associée.
  \item Vérifier l'homogénéité.
  \item Donner une phrase d'interprétation.
\end{enumerate}
\end{exercicebox}

\begin{corrigebox}[Indication 3]
La correction doit toujours contenir une loi, une unité, un ordre de grandeur et une phrase de sens physique.
\end{corrigebox}

\begin{exercicebox}[Exercice court 4]
On considère un phénomène étudié dans le document. 
\begin{enumerate}
  \item Identifier la grandeur physique principale.
  \item Donner son unité SI.
  \item Écrire la loi simple associée.
  \item Vérifier l'homogénéité.
  \item Donner une phrase d'interprétation.
\end{enumerate}
\end{exercicebox}

\begin{corrigebox}[Indication 4]
La correction doit toujours contenir une loi, une unité, un ordre de grandeur et une phrase de sens physique.
\end{corrigebox}

\begin{exercicebox}[Exercice court 5]
On considère un phénomène étudié dans le document. 
\begin{enumerate}
  \item Identifier la grandeur physique principale.
  \item Donner son unité SI.
  \item Écrire la loi simple associée.
  \item Vérifier l'homogénéité.
  \item Donner une phrase d'interprétation.
\end{enumerate}
\end{exercicebox}

\begin{corrigebox}[Indication 5]
La correction doit toujours contenir une loi, une unité, un ordre de grandeur et une phrase de sens physique.
\end{corrigebox}

\begin{exercicebox}[Exercice court 6]
On considère un phénomène étudié dans le document. 
\begin{enumerate}
  \item Identifier la grandeur physique principale.
  \item Donner son unité SI.
  \item Écrire la loi simple associée.
  \item Vérifier l'homogénéité.
  \item Donner une phrase d'interprétation.
\end{enumerate}
\end{exercicebox}

\begin{corrigebox}[Indication 6]
La correction doit toujours contenir une loi, une unité, un ordre de grandeur et une phrase de sens physique.
\end{corrigebox}

\begin{exercicebox}[Exercice court 7]
On considère un phénomène étudié dans le document. 
\begin{enumerate}
  \item Identifier la grandeur physique principale.
  \item Donner son unité SI.
  \item Écrire la loi simple associée.
  \item Vérifier l'homogénéité.
  \item Donner une phrase d'interprétation.
\end{enumerate}
\end{exercicebox}

\begin{corrigebox}[Indication 7]
La correction doit toujours contenir une loi, une unité, un ordre de grandeur et une phrase de sens physique.
\end{corrigebox}

\begin{exercicebox}[Exercice court 8]
On considère un phénomène étudié dans le document. 
\begin{enumerate}
  \item Identifier la grandeur physique principale.
  \item Donner son unité SI.
  \item Écrire la loi simple associée.
  \item Vérifier l'homogénéité.
  \item Donner une phrase d'interprétation.
\end{enumerate}
\end{exercicebox}

\begin{corrigebox}[Indication 8]
La correction doit toujours contenir une loi, une unité, un ordre de grandeur et une phrase de sens physique.
\end{corrigebox}

\begin{exercicebox}[Exercice court 9]
On considère un phénomène étudié dans le document. 
\begin{enumerate}
  \item Identifier la grandeur physique principale.
  \item Donner son unité SI.
  \item Écrire la loi simple associée.
  \item Vérifier l'homogénéité.
  \item Donner une phrase d'interprétation.
\end{enumerate}
\end{exercicebox}

\begin{corrigebox}[Indication 9]
La correction doit toujours contenir une loi, une unité, un ordre de grandeur et une phrase de sens physique.
\end{corrigebox}

\begin{exercicebox}[Exercice court 10]
On considère un phénomène étudié dans le document. 
\begin{enumerate}
  \item Identifier la grandeur physique principale.
  \item Donner son unité SI.
  \item Écrire la loi simple associée.
  \item Vérifier l'homogénéité.
  \item Donner une phrase d'interprétation.
\end{enumerate}
\end{exercicebox}

\begin{corrigebox}[Indication 10]
La correction doit toujours contenir une loi, une unité, un ordre de grandeur et une phrase de sens physique.
\end{corrigebox}

\begin{exercicebox}[Exercice court 11]
On considère un phénomène étudié dans le document. 
\begin{enumerate}
  \item Identifier la grandeur physique principale.
  \item Donner son unité SI.
  \item Écrire la loi simple associée.
  \item Vérifier l'homogénéité.
  \item Donner une phrase d'interprétation.
\end{enumerate}
\end{exercicebox}

\begin{corrigebox}[Indication 11]
La correction doit toujours contenir une loi, une unité, un ordre de grandeur et une phrase de sens physique.
\end{corrigebox}

\begin{exercicebox}[Exercice court 12]
On considère un phénomène étudié dans le document. 
\begin{enumerate}
  \item Identifier la grandeur physique principale.
  \item Donner son unité SI.
  \item Écrire la loi simple associée.
  \item Vérifier l'homogénéité.
  \item Donner une phrase d'interprétation.
\end{enumerate}
\end{exercicebox}

\begin{corrigebox}[Indication 12]
La correction doit toujours contenir une loi, une unité, un ordre de grandeur et une phrase de sens physique.
\end{corrigebox}

\begin{exercicebox}[Exercice court 13]
On considère un phénomène étudié dans le document. 
\begin{enumerate}
  \item Identifier la grandeur physique principale.
  \item Donner son unité SI.
  \item Écrire la loi simple associée.
  \item Vérifier l'homogénéité.
  \item Donner une phrase d'interprétation.
\end{enumerate}
\end{exercicebox}

\begin{corrigebox}[Indication 13]
La correction doit toujours contenir une loi, une unité, un ordre de grandeur et une phrase de sens physique.
\end{corrigebox}

\begin{exercicebox}[Exercice court 14]
On considère un phénomène étudié dans le document. 
\begin{enumerate}
  \item Identifier la grandeur physique principale.
  \item Donner son unité SI.
  \item Écrire la loi simple associée.
  \item Vérifier l'homogénéité.
  \item Donner une phrase d'interprétation.
\end{enumerate}
\end{exercicebox}

\begin{corrigebox}[Indication 14]
La correction doit toujours contenir une loi, une unité, un ordre de grandeur et une phrase de sens physique.
\end{corrigebox}

\begin{exercicebox}[Exercice court 15]
On considère un phénomène étudié dans le document. 
\begin{enumerate}
  \item Identifier la grandeur physique principale.
  \item Donner son unité SI.
  \item Écrire la loi simple associée.
  \item Vérifier l'homogénéité.
  \item Donner une phrase d'interprétation.
\end{enumerate}
\end{exercicebox}

\begin{corrigebox}[Indication 15]
La correction doit toujours contenir une loi, une unité, un ordre de grandeur et une phrase de sens physique.
\end{corrigebox}

\begin{exercicebox}[Exercice court 16]
On considère un phénomène étudié dans le document. 
\begin{enumerate}
  \item Identifier la grandeur physique principale.
  \item Donner son unité SI.
  \item Écrire la loi simple associée.
  \item Vérifier l'homogénéité.
  \item Donner une phrase d'interprétation.
\end{enumerate}
\end{exercicebox}

\begin{corrigebox}[Indication 16]
La correction doit toujours contenir une loi, une unité, un ordre de grandeur et une phrase de sens physique.
\end{corrigebox}

\begin{exercicebox}[Exercice court 17]
On considère un phénomène étudié dans le document. 
\begin{enumerate}
  \item Identifier la grandeur physique principale.
  \item Donner son unité SI.
  \item Écrire la loi simple associée.
  \item Vérifier l'homogénéité.
  \item Donner une phrase d'interprétation.
\end{enumerate}
\end{exercicebox}

\begin{corrigebox}[Indication 17]
La correction doit toujours contenir une loi, une unité, un ordre de grandeur et une phrase de sens physique.
\end{corrigebox}

\begin{exercicebox}[Exercice court 18]
On considère un phénomène étudié dans le document. 
\begin{enumerate}
  \item Identifier la grandeur physique principale.
  \item Donner son unité SI.
  \item Écrire la loi simple associée.
  \item Vérifier l'homogénéité.
  \item Donner une phrase d'interprétation.
\end{enumerate}
\end{exercicebox}

\begin{corrigebox}[Indication 18]
La correction doit toujours contenir une loi, une unité, un ordre de grandeur et une phrase de sens physique.
\end{corrigebox}

\begin{exercicebox}[Exercice court 19]
On considère un phénomène étudié dans le document. 
\begin{enumerate}
  \item Identifier la grandeur physique principale.
  \item Donner son unité SI.
  \item Écrire la loi simple associée.
  \item Vérifier l'homogénéité.
  \item Donner une phrase d'interprétation.
\end{enumerate}
\end{exercicebox}

\begin{corrigebox}[Indication 19]
La correction doit toujours contenir une loi, une unité, un ordre de grandeur et une phrase de sens physique.
\end{corrigebox}

\begin{exercicebox}[Exercice court 20]
On considère un phénomène étudié dans le document. 
\begin{enumerate}
  \item Identifier la grandeur physique principale.
  \item Donner son unité SI.
  \item Écrire la loi simple associée.
  \item Vérifier l'homogénéité.
  \item Donner une phrase d'interprétation.
\end{enumerate}
\end{exercicebox}

\begin{corrigebox}[Indication 20]
La correction doit toujours contenir une loi, une unité, un ordre de grandeur et une phrase de sens physique.
\end{corrigebox}

\begin{exercicebox}[Exercice court 21]
On considère un phénomène étudié dans le document. 
\begin{enumerate}
  \item Identifier la grandeur physique principale.
  \item Donner son unité SI.
  \item Écrire la loi simple associée.
  \item Vérifier l'homogénéité.
  \item Donner une phrase d'interprétation.
\end{enumerate}
\end{exercicebox}

\begin{corrigebox}[Indication 21]
La correction doit toujours contenir une loi, une unité, un ordre de grandeur et une phrase de sens physique.
\end{corrigebox}

\begin{exercicebox}[Exercice court 22]
On considère un phénomène étudié dans le document. 
\begin{enumerate}
  \item Identifier la grandeur physique principale.
  \item Donner son unité SI.
  \item Écrire la loi simple associée.
  \item Vérifier l'homogénéité.
  \item Donner une phrase d'interprétation.
\end{enumerate}
\end{exercicebox}

\begin{corrigebox}[Indication 22]
La correction doit toujours contenir une loi, une unité, un ordre de grandeur et une phrase de sens physique.
\end{corrigebox}

\begin{exercicebox}[Exercice court 23]
On considère un phénomène étudié dans le document. 
\begin{enumerate}
  \item Identifier la grandeur physique principale.
  \item Donner son unité SI.
  \item Écrire la loi simple associée.
  \item Vérifier l'homogénéité.
  \item Donner une phrase d'interprétation.
\end{enumerate}
\end{exercicebox}

\begin{corrigebox}[Indication 23]
La correction doit toujours contenir une loi, une unité, un ordre de grandeur et une phrase de sens physique.
\end{corrigebox}

\begin{exercicebox}[Exercice court 24]
On considère un phénomène étudié dans le document. 
\begin{enumerate}
  \item Identifier la grandeur physique principale.
  \item Donner son unité SI.
  \item Écrire la loi simple associée.
  \item Vérifier l'homogénéité.
  \item Donner une phrase d'interprétation.
\end{enumerate}
\end{exercicebox}

\begin{corrigebox}[Indication 24]
La correction doit toujours contenir une loi, une unité, un ordre de grandeur et une phrase de sens physique.
\end{corrigebox}

\begin{exercicebox}[Exercice court 25]
On considère un phénomène étudié dans le document. 
\begin{enumerate}
  \item Identifier la grandeur physique principale.
  \item Donner son unité SI.
  \item Écrire la loi simple associée.
  \item Vérifier l'homogénéité.
  \item Donner une phrase d'interprétation.
\end{enumerate}
\end{exercicebox}

\begin{corrigebox}[Indication 25]
La correction doit toujours contenir une loi, une unité, un ordre de grandeur et une phrase de sens physique.
\end{corrigebox}

\begin{exercicebox}[Exercice court 26]
On considère un phénomène étudié dans le document. 
\begin{enumerate}
  \item Identifier la grandeur physique principale.
  \item Donner son unité SI.
  \item Écrire la loi simple associée.
  \item Vérifier l'homogénéité.
  \item Donner une phrase d'interprétation.
\end{enumerate}
\end{exercicebox}

\begin{corrigebox}[Indication 26]
La correction doit toujours contenir une loi, une unité, un ordre de grandeur et une phrase de sens physique.
\end{corrigebox}

\begin{exercicebox}[Exercice court 27]
On considère un phénomène étudié dans le document. 
\begin{enumerate}
  \item Identifier la grandeur physique principale.
  \item Donner son unité SI.
  \item Écrire la loi simple associée.
  \item Vérifier l'homogénéité.
  \item Donner une phrase d'interprétation.
\end{enumerate}
\end{exercicebox}

\begin{corrigebox}[Indication 27]
La correction doit toujours contenir une loi, une unité, un ordre de grandeur et une phrase de sens physique.
\end{corrigebox}

\begin{exercicebox}[Exercice court 28]
On considère un phénomène étudié dans le document. 
\begin{enumerate}
  \item Identifier la grandeur physique principale.
  \item Donner son unité SI.
  \item Écrire la loi simple associée.
  \item Vérifier l'homogénéité.
  \item Donner une phrase d'interprétation.
\end{enumerate}
\end{exercicebox}

\begin{corrigebox}[Indication 28]
La correction doit toujours contenir une loi, une unité, un ordre de grandeur et une phrase de sens physique.
\end{corrigebox}

\begin{exercicebox}[Exercice court 29]
On considère un phénomène étudié dans le document. 
\begin{enumerate}
  \item Identifier la grandeur physique principale.
  \item Donner son unité SI.
  \item Écrire la loi simple associée.
  \item Vérifier l'homogénéité.
  \item Donner une phrase d'interprétation.
\end{enumerate}
\end{exercicebox}

\begin{corrigebox}[Indication 29]
La correction doit toujours contenir une loi, une unité, un ordre de grandeur et une phrase de sens physique.
\end{corrigebox}

\begin{exercicebox}[Exercice court 30]
On considère un phénomène étudié dans le document. 
\begin{enumerate}
  \item Identifier la grandeur physique principale.
  \item Donner son unité SI.
  \item Écrire la loi simple associée.
  \item Vérifier l'homogénéité.
  \item Donner une phrase d'interprétation.
\end{enumerate}
\end{exercicebox}

\begin{corrigebox}[Indication 30]
La correction doit toujours contenir une loi, une unité, un ordre de grandeur et une phrase de sens physique.
\end{corrigebox}

\subsection{Exercices intermédiaires}

\begin{exercicebox}[Exercice intermédiaire 1]
Un phénomène est modélisé par une grandeur $X(t)$ évoluant vers un équilibre $X_{\mathrm{eq}}$.
\[
\frac{\dd X}{\dd t}=-\frac1\tau(X-X_{\mathrm{eq}}).
\]
\begin{enumerate}
  \item Identifier le type d'évolution.
  \item Résoudre l'équation différentielle.
  \item Interpréter $\tau$.
  \item Déterminer le temps nécessaire pour atteindre 95 \% de l'écart à l'équilibre.
  \item Donner une application physique possible.
\end{enumerate}
\end{exercicebox}

\begin{corrigebox}[Correction 1]
La solution est $X(t)=X_{\mathrm{eq}}+(X(0)-X_{\mathrm{eq}})\e^{-t/\tau}$. Pour atteindre 95 \%, il faut résoudre $\e^{-t/\tau}=0{,}05$, soit $t\simeq3\tau$.
\end{corrigebox}

\begin{exercicebox}[Exercice intermédiaire 2]
Un phénomène est modélisé par une grandeur $X(t)$ évoluant vers un équilibre $X_{\mathrm{eq}}$.
\[
\frac{\dd X}{\dd t}=-\frac1\tau(X-X_{\mathrm{eq}}).
\]
\begin{enumerate}
  \item Identifier le type d'évolution.
  \item Résoudre l'équation différentielle.
  \item Interpréter $\tau$.
  \item Déterminer le temps nécessaire pour atteindre 95 \% de l'écart à l'équilibre.
  \item Donner une application physique possible.
\end{enumerate}
\end{exercicebox}

\begin{corrigebox}[Correction 2]
La solution est $X(t)=X_{\mathrm{eq}}+(X(0)-X_{\mathrm{eq}})\e^{-t/\tau}$. Pour atteindre 95 \%, il faut résoudre $\e^{-t/\tau}=0{,}05$, soit $t\simeq3\tau$.
\end{corrigebox}

\begin{exercicebox}[Exercice intermédiaire 3]
Un phénomène est modélisé par une grandeur $X(t)$ évoluant vers un équilibre $X_{\mathrm{eq}}$.
\[
\frac{\dd X}{\dd t}=-\frac1\tau(X-X_{\mathrm{eq}}).
\]
\begin{enumerate}
  \item Identifier le type d'évolution.
  \item Résoudre l'équation différentielle.
  \item Interpréter $\tau$.
  \item Déterminer le temps nécessaire pour atteindre 95 \% de l'écart à l'équilibre.
  \item Donner une application physique possible.
\end{enumerate}
\end{exercicebox}

\begin{corrigebox}[Correction 3]
La solution est $X(t)=X_{\mathrm{eq}}+(X(0)-X_{\mathrm{eq}})\e^{-t/\tau}$. Pour atteindre 95 \%, il faut résoudre $\e^{-t/\tau}=0{,}05$, soit $t\simeq3\tau$.
\end{corrigebox}

\begin{exercicebox}[Exercice intermédiaire 4]
Un phénomène est modélisé par une grandeur $X(t)$ évoluant vers un équilibre $X_{\mathrm{eq}}$.
\[
\frac{\dd X}{\dd t}=-\frac1\tau(X-X_{\mathrm{eq}}).
\]
\begin{enumerate}
  \item Identifier le type d'évolution.
  \item Résoudre l'équation différentielle.
  \item Interpréter $\tau$.
  \item Déterminer le temps nécessaire pour atteindre 95 \% de l'écart à l'équilibre.
  \item Donner une application physique possible.
\end{enumerate}
\end{exercicebox}

\begin{corrigebox}[Correction 4]
La solution est $X(t)=X_{\mathrm{eq}}+(X(0)-X_{\mathrm{eq}})\e^{-t/\tau}$. Pour atteindre 95 \%, il faut résoudre $\e^{-t/\tau}=0{,}05$, soit $t\simeq3\tau$.
\end{corrigebox}

\begin{exercicebox}[Exercice intermédiaire 5]
Un phénomène est modélisé par une grandeur $X(t)$ évoluant vers un équilibre $X_{\mathrm{eq}}$.
\[
\frac{\dd X}{\dd t}=-\frac1\tau(X-X_{\mathrm{eq}}).
\]
\begin{enumerate}
  \item Identifier le type d'évolution.
  \item Résoudre l'équation différentielle.
  \item Interpréter $\tau$.
  \item Déterminer le temps nécessaire pour atteindre 95 \% de l'écart à l'équilibre.
  \item Donner une application physique possible.
\end{enumerate}
\end{exercicebox}

\begin{corrigebox}[Correction 5]
La solution est $X(t)=X_{\mathrm{eq}}+(X(0)-X_{\mathrm{eq}})\e^{-t/\tau}$. Pour atteindre 95 \%, il faut résoudre $\e^{-t/\tau}=0{,}05$, soit $t\simeq3\tau$.
\end{corrigebox}

\begin{exercicebox}[Exercice intermédiaire 6]
Un phénomène est modélisé par une grandeur $X(t)$ évoluant vers un équilibre $X_{\mathrm{eq}}$.
\[
\frac{\dd X}{\dd t}=-\frac1\tau(X-X_{\mathrm{eq}}).
\]
\begin{enumerate}
  \item Identifier le type d'évolution.
  \item Résoudre l'équation différentielle.
  \item Interpréter $\tau$.
  \item Déterminer le temps nécessaire pour atteindre 95 \% de l'écart à l'équilibre.
  \item Donner une application physique possible.
\end{enumerate}
\end{exercicebox}

\begin{corrigebox}[Correction 6]
La solution est $X(t)=X_{\mathrm{eq}}+(X(0)-X_{\mathrm{eq}})\e^{-t/\tau}$. Pour atteindre 95 \%, il faut résoudre $\e^{-t/\tau}=0{,}05$, soit $t\simeq3\tau$.
\end{corrigebox}

\begin{exercicebox}[Exercice intermédiaire 7]
Un phénomène est modélisé par une grandeur $X(t)$ évoluant vers un équilibre $X_{\mathrm{eq}}$.
\[
\frac{\dd X}{\dd t}=-\frac1\tau(X-X_{\mathrm{eq}}).
\]
\begin{enumerate}
  \item Identifier le type d'évolution.
  \item Résoudre l'équation différentielle.
  \item Interpréter $\tau$.
  \item Déterminer le temps nécessaire pour atteindre 95 \% de l'écart à l'équilibre.
  \item Donner une application physique possible.
\end{enumerate}
\end{exercicebox}

\begin{corrigebox}[Correction 7]
La solution est $X(t)=X_{\mathrm{eq}}+(X(0)-X_{\mathrm{eq}})\e^{-t/\tau}$. Pour atteindre 95 \%, il faut résoudre $\e^{-t/\tau}=0{,}05$, soit $t\simeq3\tau$.
\end{corrigebox}

\begin{exercicebox}[Exercice intermédiaire 8]
Un phénomène est modélisé par une grandeur $X(t)$ évoluant vers un équilibre $X_{\mathrm{eq}}$.
\[
\frac{\dd X}{\dd t}=-\frac1\tau(X-X_{\mathrm{eq}}).
\]
\begin{enumerate}
  \item Identifier le type d'évolution.
  \item Résoudre l'équation différentielle.
  \item Interpréter $\tau$.
  \item Déterminer le temps nécessaire pour atteindre 95 \% de l'écart à l'équilibre.
  \item Donner une application physique possible.
\end{enumerate}
\end{exercicebox}

\begin{corrigebox}[Correction 8]
La solution est $X(t)=X_{\mathrm{eq}}+(X(0)-X_{\mathrm{eq}})\e^{-t/\tau}$. Pour atteindre 95 \%, il faut résoudre $\e^{-t/\tau}=0{,}05$, soit $t\simeq3\tau$.
\end{corrigebox}

\begin{exercicebox}[Exercice intermédiaire 9]
Un phénomène est modélisé par une grandeur $X(t)$ évoluant vers un équilibre $X_{\mathrm{eq}}$.
\[
\frac{\dd X}{\dd t}=-\frac1\tau(X-X_{\mathrm{eq}}).
\]
\begin{enumerate}
  \item Identifier le type d'évolution.
  \item Résoudre l'équation différentielle.
  \item Interpréter $\tau$.
  \item Déterminer le temps nécessaire pour atteindre 95 \% de l'écart à l'équilibre.
  \item Donner une application physique possible.
\end{enumerate}
\end{exercicebox}

\begin{corrigebox}[Correction 9]
La solution est $X(t)=X_{\mathrm{eq}}+(X(0)-X_{\mathrm{eq}})\e^{-t/\tau}$. Pour atteindre 95 \%, il faut résoudre $\e^{-t/\tau}=0{,}05$, soit $t\simeq3\tau$.
\end{corrigebox}

\begin{exercicebox}[Exercice intermédiaire 10]
Un phénomène est modélisé par une grandeur $X(t)$ évoluant vers un équilibre $X_{\mathrm{eq}}$.
\[
\frac{\dd X}{\dd t}=-\frac1\tau(X-X_{\mathrm{eq}}).
\]
\begin{enumerate}
  \item Identifier le type d'évolution.
  \item Résoudre l'équation différentielle.
  \item Interpréter $\tau$.
  \item Déterminer le temps nécessaire pour atteindre 95 \% de l'écart à l'équilibre.
  \item Donner une application physique possible.
\end{enumerate}
\end{exercicebox}

\begin{corrigebox}[Correction 10]
La solution est $X(t)=X_{\mathrm{eq}}+(X(0)-X_{\mathrm{eq}})\e^{-t/\tau}$. Pour atteindre 95 \%, il faut résoudre $\e^{-t/\tau}=0{,}05$, soit $t\simeq3\tau$.
\end{corrigebox}

\begin{exercicebox}[Exercice intermédiaire 11]
Un phénomène est modélisé par une grandeur $X(t)$ évoluant vers un équilibre $X_{\mathrm{eq}}$.
\[
\frac{\dd X}{\dd t}=-\frac1\tau(X-X_{\mathrm{eq}}).
\]
\begin{enumerate}
  \item Identifier le type d'évolution.
  \item Résoudre l'équation différentielle.
  \item Interpréter $\tau$.
  \item Déterminer le temps nécessaire pour atteindre 95 \% de l'écart à l'équilibre.
  \item Donner une application physique possible.
\end{enumerate}
\end{exercicebox}

\begin{corrigebox}[Correction 11]
La solution est $X(t)=X_{\mathrm{eq}}+(X(0)-X_{\mathrm{eq}})\e^{-t/\tau}$. Pour atteindre 95 \%, il faut résoudre $\e^{-t/\tau}=0{,}05$, soit $t\simeq3\tau$.
\end{corrigebox}

\begin{exercicebox}[Exercice intermédiaire 12]
Un phénomène est modélisé par une grandeur $X(t)$ évoluant vers un équilibre $X_{\mathrm{eq}}$.
\[
\frac{\dd X}{\dd t}=-\frac1\tau(X-X_{\mathrm{eq}}).
\]
\begin{enumerate}
  \item Identifier le type d'évolution.
  \item Résoudre l'équation différentielle.
  \item Interpréter $\tau$.
  \item Déterminer le temps nécessaire pour atteindre 95 \% de l'écart à l'équilibre.
  \item Donner une application physique possible.
\end{enumerate}
\end{exercicebox}

\begin{corrigebox}[Correction 12]
La solution est $X(t)=X_{\mathrm{eq}}+(X(0)-X_{\mathrm{eq}})\e^{-t/\tau}$. Pour atteindre 95 \%, il faut résoudre $\e^{-t/\tau}=0{,}05$, soit $t\simeq3\tau$.
\end{corrigebox}

\begin{exercicebox}[Exercice intermédiaire 13]
Un phénomène est modélisé par une grandeur $X(t)$ évoluant vers un équilibre $X_{\mathrm{eq}}$.
\[
\frac{\dd X}{\dd t}=-\frac1\tau(X-X_{\mathrm{eq}}).
\]
\begin{enumerate}
  \item Identifier le type d'évolution.
  \item Résoudre l'équation différentielle.
  \item Interpréter $\tau$.
  \item Déterminer le temps nécessaire pour atteindre 95 \% de l'écart à l'équilibre.
  \item Donner une application physique possible.
\end{enumerate}
\end{exercicebox}

\begin{corrigebox}[Correction 13]
La solution est $X(t)=X_{\mathrm{eq}}+(X(0)-X_{\mathrm{eq}})\e^{-t/\tau}$. Pour atteindre 95 \%, il faut résoudre $\e^{-t/\tau}=0{,}05$, soit $t\simeq3\tau$.
\end{corrigebox}

\begin{exercicebox}[Exercice intermédiaire 14]
Un phénomène est modélisé par une grandeur $X(t)$ évoluant vers un équilibre $X_{\mathrm{eq}}$.
\[
\frac{\dd X}{\dd t}=-\frac1\tau(X-X_{\mathrm{eq}}).
\]
\begin{enumerate}
  \item Identifier le type d'évolution.
  \item Résoudre l'équation différentielle.
  \item Interpréter $\tau$.
  \item Déterminer le temps nécessaire pour atteindre 95 \% de l'écart à l'équilibre.
  \item Donner une application physique possible.
\end{enumerate}
\end{exercicebox}

\begin{corrigebox}[Correction 14]
La solution est $X(t)=X_{\mathrm{eq}}+(X(0)-X_{\mathrm{eq}})\e^{-t/\tau}$. Pour atteindre 95 \%, il faut résoudre $\e^{-t/\tau}=0{,}05$, soit $t\simeq3\tau$.
\end{corrigebox}

\begin{exercicebox}[Exercice intermédiaire 15]
Un phénomène est modélisé par une grandeur $X(t)$ évoluant vers un équilibre $X_{\mathrm{eq}}$.
\[
\frac{\dd X}{\dd t}=-\frac1\tau(X-X_{\mathrm{eq}}).
\]
\begin{enumerate}
  \item Identifier le type d'évolution.
  \item Résoudre l'équation différentielle.
  \item Interpréter $\tau$.
  \item Déterminer le temps nécessaire pour atteindre 95 \% de l'écart à l'équilibre.
  \item Donner une application physique possible.
\end{enumerate}
\end{exercicebox}

\begin{corrigebox}[Correction 15]
La solution est $X(t)=X_{\mathrm{eq}}+(X(0)-X_{\mathrm{eq}})\e^{-t/\tau}$. Pour atteindre 95 \%, il faut résoudre $\e^{-t/\tau}=0{,}05$, soit $t\simeq3\tau$.
\end{corrigebox}


\section{Partie XII — Problèmes longs niveau Terminale / prépa / Bac+3}
% ============================================================

\begin{attentionbox}[Format]
Chaque problème est conçu comme un sujet long : contexte, modélisation, calcul, interprétation et synthèse. Il peut être donné en devoir maison ou en oral préparé.
\end{attentionbox}

\subsection{Problème 1 — Mirage sur une route chaude}

\begin{exercicebox}[Énoncé]
Une route noire chauffée par le Soleil crée un gradient de température dans l'air proche du sol. On modélise l'atmosphère par trois couches horizontales d'indices $n_1>n_2>n_3$, la couche la plus basse ayant l'indice le plus faible.
\end{exercicebox}

\textbf{Partie A — Modèle en couches}
\begin{enumerate}
  \item Expliquer pourquoi l'indice de l'air diminue quand la température augmente.
  \item Écrire la loi de Snell-Descartes entre deux couches.
  \item Montrer qualitativement que le rayon s'éloigne de la normale en descendant.
  \item Expliquer pourquoi la trajectoire devient courbe dans un modèle continu.
  \item Interpréter l'image virtuelle observée par le conducteur.
\end{enumerate}

\textbf{Partie B — Passage au modèle continu}
\begin{enumerate}
  \item On suppose $n=n(z)$. Expliquer ce que représente $z$.
  \item Justifier l'idée qu'une quantité proche de $n(z)\sin i(z)$ se conserve.
  \item Décrire le sens de courbure du rayon.
  \item Expliquer pourquoi il n'y a pas de vraie flaque.
  \item Donner deux limites du modèle.
\end{enumerate}

\textbf{Partie C — Synthèse}
\begin{enumerate}
  \item Comparer mirage et réflexion sur l'eau.
  \item Expliquer le rôle du gradient thermique.
  \item Dire pourquoi l'effet est visible à grande distance.
  \item Proposer une expérience de laboratoire.
  \item Rédiger une explication de niveau Grand Oral.
\end{enumerate}

\begin{corrigebox}
L'idée centrale est que le rayon est continûment réfracté vers les zones d'indice plus élevé. L'œil prolonge rectilignement le dernier rayon reçu, ce qui crée une image virtuelle du ciel au niveau de la route.
\end{corrigebox}

\subsection{Problème 2 — Arc-en-ciel primaire}

\textbf{Partie A — Trajet dans une goutte}
\begin{enumerate}
  \item Décrire les trois étapes du trajet : réfraction, réflexion, réfraction.
  \item Écrire $\sin i=n\sin r$.
  \item Donner l'expression $D(i)=\pi+2i-4r$.
  \item Expliquer ce que représente la déviation $D$.
  \item Dire pourquoi $n$ dépend de la couleur.
\end{enumerate}

\textbf{Partie B — Extremum de déviation}
\begin{enumerate}
  \item Expliquer pourquoi on cherche un extremum de $D(i)$.
  \item Dériver implicitement la relation de Snell-Descartes.
  \item Relier $\dd r/\dd i$ à $i$ et $r$.
  \item Écrire la condition $\dd D/\dd i=0$.
  \item Interpréter la concentration de lumière.
\end{enumerate}

\textbf{Partie C — Couleurs}
\begin{enumerate}
  \item Comparer l'indice pour le rouge et le violet.
  \item Dire quelle couleur est la plus déviée.
  \item Expliquer l'ordre des couleurs.
  \item Expliquer pourquoi l'arc est circulaire.
  \item Rédiger une synthèse.
\end{enumerate}

\begin{corrigebox}
L'arc visible correspond aux rayons émergents autour d'une déviation stationnaire. La dispersion sépare les angles selon les longueurs d'onde.
\end{corrigebox}

\subsection{Problème 3 — Fibre optique de télécommunication}

\textbf{Partie A — Guidage}
\begin{enumerate}
  \item Décrire cœur et gaine.
  \item Écrire la condition $n_c>n_g$.
  \item Définir l'angle limite.
  \item Calculer $i_c$ pour $n_c=1{,}48$ et $n_g=1{,}46$.
  \item Interpréter.
\end{enumerate}

\textbf{Partie B — Temps et atténuation}
\begin{enumerate}
  \item Calculer le temps de propagation dans 10 km.
  \item Comparer à la propagation dans le vide.
  \item Définir l'atténuation en dB.
  \item Calculer une puissance de sortie pour une perte de 3 dB.
  \item Discuter l'intérêt pour Internet.
\end{enumerate}

\textbf{Partie C — Limites}
\begin{enumerate}
  \item Expliquer la dispersion modale.
  \item Expliquer la dispersion chromatique.
  \item Comparer fibre multimode et monomode.
  \item Donner une limite technologique.
  \item Conclure.
\end{enumerate}

\begin{corrigebox}
La fibre guide la lumière grâce à la réflexion totale. La vitesse est $c/n$, donc le temps vaut environ $nL/c$. Une perte de 3 dB correspond à une puissance divisée environ par deux.
\end{corrigebox}

\subsection{Problème 4 — Cheveu et diffraction}

\textbf{Partie A — Relation}
\begin{enumerate}
  \item Rappeler $\theta\simeq\lambda/a$.
  \item Relier $\theta$ à $L$ et $D$.
  \item En déduire $a\simeq2\lambda D/L$.
  \item Vérifier les unités.
  \item Expliquer le principe de Babinet.
\end{enumerate}

\textbf{Partie B — Calcul}
\begin{enumerate}
  \item Prendre $\lambda=\SI{650}{nm}$, $D=\SI{2.0}{m}$, $L=\SI{3.0}{cm}$.
  \item Calculer $a$.
  \item Donner le résultat en micromètres.
  \item Vérifier l'ordre de grandeur.
  \item Discuter l'incertitude sur $L$.
\end{enumerate}

\textbf{Partie C — Critique}
\begin{enumerate}
  \item Pourquoi faut-il un laser ?
  \item Pourquoi la pièce doit-elle être sombre ?
  \item Pourquoi mesurer plusieurs franges ?
  \item Comment améliorer la précision ?
  \item Conclure.
\end{enumerate}

\begin{corrigebox}
On trouve une taille de l'ordre de quelques dizaines de micromètres, cohérente avec l'épaisseur d'un cheveu humain.
\end{corrigebox}

\subsection{Problème 5 — Ciel bleu et coucher rouge}

\textbf{Partie A — Diffusion Rayleigh}
\begin{enumerate}
  \item Écrire $I\propto1/\lambda^4$.
  \item Comparer $\lambda_b=\SI{450}{nm}$ et $\lambda_r=\SI{650}{nm}$.
  \item Calculer le rapport des intensités diffusées.
  \item Interpréter.
  \item Expliquer pourquoi le ciel n'est pas violet.
\end{enumerate}

\textbf{Partie B — Coucher de soleil}
\begin{enumerate}
  \item Expliquer le trajet atmosphérique plus long.
  \item Écrire $I=I_0\e^{-\alpha L}$.
  \item Dire comment $\alpha$ dépend de $\lambda$.
  \item Expliquer pourquoi le rouge domine.
  \item Citer l'effet des aérosols.
\end{enumerate}

\textbf{Partie C — Synthèse}
\begin{enumerate}
  \item Relier diffusion et transmission.
  \item Comparer ciel bleu et soleil rouge.
  \item Donner un ordre de grandeur.
  \item Rédiger une réponse niveau Terminale.
  \item Rédiger une réponse niveau prépa.
\end{enumerate}

\begin{corrigebox}
Le bleu est fortement diffusé hors du faisceau solaire direct ; au coucher, le trajet plus long retire davantage les courtes longueurs d'onde, laissant une lumière transmise enrichie en rouge.
\end{corrigebox}


% ============================================================
\section{Partie XIII — Annexes et formulaires}
% ============================================================

\subsection{Formulaire d'optique géométrique}

\begin{methodebox}[Relations essentielles]
\[
n=\frac{c}{v},
\qquad
n_1\sin i_1=n_2\sin i_2,
\qquad
\sin i_c=\frac{n_2}{n_1}.
\]
\[
\frac{1}{OA'}-\frac{1}{OA}=\frac{1}{f'},
\qquad
\gamma=\frac{A'B'}{AB}.
\]
\end{methodebox}

\subsection{Formulaire d'optique ondulatoire}

\begin{methodebox}[Relations essentielles]
\[
v=\lambda f,
\qquad
\delta=k\lambda,
\qquad
\delta=\left(k+\frac12\right)\lambda.
\]
\[
i=\frac{\lambda D}{a},
\qquad
\theta\simeq\frac{\lambda}{a},
\qquad
I=I_0\cos^2\theta.
\]
\end{methodebox}

\subsection{Formulaire de mécanique}

\begin{methodebox}[Relations essentielles]
\[
\sum \vec F=m\vec a,
\qquad
E_c=\frac12mv^2,
\qquad
E_p=mgh,
\qquad
P=\frac{E}{\Delta t}.
\]
\[
F_D=\frac12\rho C_DSv^2,
\qquad
L=\frac12\rho C_LSv^2,
\qquad
v_{\mathrm{orb}}=\sqrt{\frac{GM}{r}}.
\]
\end{methodebox}

\subsection{Méthode pour rédiger une explication physique}

\begin{methodebox}[Structure obligatoire]
\begin{enumerate}
  \item Décrire l'observation.
  \item Identifier les grandeurs physiques.
  \item Choisir un modèle.
  \item Écrire la loi.
  \item Vérifier l'homogénéité.
  \item Calculer un ordre de grandeur.
  \item Interpréter.
  \item Discuter les limites.
\end{enumerate}
\end{methodebox}


% ============================================================
\section{Partie XIV — Grands exercices corrigés : de la question naïve au modèle physique}
% ============================================================

\begin{attentionbox}[Principe de cette partie]
Cette partie transforme les phénomènes physiques en \textbf{longs exercices corrigés}. L'objectif n'est pas seulement de donner une formule, mais de partir d'une question concrète, de construire un modèle, de faire les calculs, puis de discuter les limites. 

Chaque exercice suit la structure :
\begin{enumerate}
  \item question de départ ;
  \item modélisation physique ;
  \item mise en équation ;
  \item calcul numérique ;
  \item interprétation ;
  \item critique du modèle ;
  \item correction détaillée.
\end{enumerate}
\end{attentionbox}

% ------------------------------------------------------------
\subsection{Exercice corrigé 1 — Combien de temps faut-il se frotter les mains pour réchauffer une pièce de \(1^\circ\mathrm C\) ?}
% ------------------------------------------------------------

\begin{exercicebox}[Question de départ]
On imagine une personne enfermée dans une pièce froide. Elle décide de se frotter les mains pour produire de la chaleur. On veut répondre sérieusement à la question :

\begin{center}
\textbf{Combien de temps faudrait-il se frotter les mains pour réchauffer l'air d'une pièce de \(1^\circ\mathrm C\) ?}
\end{center}

On considère une pièce de dimensions \(4{,}0\ \mathrm m \times 3{,}0\ \mathrm m \times 2{,}5\ \mathrm m\). On suppose que la pièce est fermée et que toute l'énergie mécanique dissipée par frottement des mains finit par chauffer l'air de la pièce.

Données :
\[
\rho_{\mathrm{air}}=1{,}2\ \mathrm{kg.m^{-3}},
\qquad
c_{\mathrm{air}}=1000\ \mathrm{J.kg^{-1}.K^{-1}}.
\]
On estime qu'une personne qui se frotte les mains fortement peut dissiper une puissance mécanique moyenne de l'ordre de
\[
P=20\ \mathrm W.
\]
\end{exercicebox}

\subsubsection*{Partie A — Comprendre la question}

\begin{enumerate}
  \item Quelle grandeur physique mesure la chaleur nécessaire pour réchauffer l'air ?
  \item Pourquoi une élévation de \(1^\circ\mathrm C\) est-elle équivalente à une élévation de \(1\ \mathrm K\) ?
  \item Calculer le volume de la pièce.
  \item Calculer la masse d'air contenue dans la pièce.
  \item Donner la relation entre l'énergie thermique reçue par l'air et son élévation de température.
\end{enumerate}

\subsubsection*{Partie B — Mise en équation}

On note :
\[
V
\]
le volume de la pièce,
\[
m
\]
la masse d'air dans la pièce,
\[
\Delta T
\]
l'élévation de température souhaitée,
\[
Q
\]
l'énergie thermique nécessaire.

\begin{enumerate}
  \item Exprimer \(m\) en fonction de \(\rho_{\mathrm{air}}\) et \(V\).
  \item Écrire la relation :
  \[
  Q=mc_{\mathrm{air}}\Delta T.
  \]
  \item Remplacer \(m\) par son expression.
  \item En déduire :
  \[
  Q=\rho_{\mathrm{air}}Vc_{\mathrm{air}}\Delta T.
  \]
  \item Vérifier l'homogénéité de la formule.
\end{enumerate}

\subsubsection*{Partie C — Calcul numérique}

\begin{enumerate}
  \item Calculer :
  \[
  V=4{,}0\times 3{,}0\times 2{,}5.
  \]
  \item En déduire la masse d'air de la pièce.
  \item Calculer l'énergie nécessaire pour augmenter la température de \(1\ \mathrm K\).
  \item Convertir cette énergie en kilojoules.
  \item Si la puissance produite par frottement des mains est \(P=20\ \mathrm W\), déterminer le temps nécessaire :
  \[
  t=\frac{Q}{P}.
  \]
  \item Convertir ce temps en minutes.
\end{enumerate}

\subsubsection*{Partie D — Interprétation physique}

\begin{enumerate}
  \item Le résultat est-il raisonnable à l'échelle humaine ?
  \item Pourquoi a-t-on l'impression que se frotter les mains chauffe vite ?
  \item Quelle partie chauffe réellement en premier ?
  \item Pourquoi chauffer toute une pièce est beaucoup plus difficile ?
  \item Comment le résultat changerait-il si la pièce était deux fois plus grande ?
\end{enumerate}

\subsubsection*{Partie E — Limites du modèle}

Le modèle précédent est très favorable, car il suppose que toute l'énergie produite va dans l'air.

\begin{enumerate}
  \item Dans la réalité, où part une partie de l'énergie produite ?
  \item Pourquoi les murs, le sol, les meubles et les vitres changent-ils le résultat ?
  \item Pourquoi les pertes thermiques vers l'extérieur rendent-elles le chauffage encore plus difficile ?
  \item Expliquer pourquoi on sous-estime fortement le temps réel.
  \item Proposer un modèle plus réaliste.
\end{enumerate}

\subsubsection*{Correction détaillée}

\begin{corrigebox}[Correction de la partie A]
La grandeur utile est l'énergie thermique nécessaire pour élever la température de l'air. Pour une variation de température, une différence de \(1^\circ\mathrm C\) est égale à une différence de \(1\ \mathrm K\). Les échelles Celsius et Kelvin ont le même pas.

Le volume de la pièce vaut :
\[
V=4{,}0\times3{,}0\times2{,}5=30\ \mathrm{m^3}.
\]

La masse d'air est :
\[
m=\rho_{\mathrm{air}}V=1{,}2\times30=36\ \mathrm{kg}.
\]

Pour réchauffer une masse \(m\) d'un matériau de capacité thermique massique \(c\), on utilise :
\[
Q=mc\Delta T.
\]
\end{corrigebox}

\begin{corrigebox}[Correction de la partie B]
On a :
\[
m=\rho_{\mathrm{air}}V.
\]
Donc :
\[
Q=mc_{\mathrm{air}}\Delta T
=\rho_{\mathrm{air}}Vc_{\mathrm{air}}\Delta T.
\]

Vérification dimensionnelle :
\[
[\rho V c\Delta T]
=
\left(\mathrm{kg.m^{-3}}\right)
\left(\mathrm{m^3}\right)
\left(\mathrm{J.kg^{-1}.K^{-1}}\right)
\left(\mathrm K\right)
=\mathrm J.
\]
La formule est homogène.
\end{corrigebox}

\begin{corrigebox}[Correction de la partie C]
On calcule :
\[
Q=36\times1000\times1=36000\ \mathrm J.
\]
Donc :
\[
Q=36\ \mathrm{kJ}.
\]

Si la puissance mécanique transformée en chaleur vaut :
\[
P=20\ \mathrm W=20\ \mathrm{J.s^{-1}},
\]
alors :
\[
t=\frac{Q}{P}
=
\frac{36000}{20}
=
1800\ \mathrm s.
\]
Or :
\[
1800\ \mathrm s = 30\ \mathrm{min}.
\]

Dans ce modèle très optimiste, il faudrait donc environ :
\[
\boxed{30\ \mathrm{minutes}}
\]
de frottement continu pour augmenter de \(1^\circ\mathrm C\) la température de l'air seul dans la pièce.
\end{corrigebox}

\begin{corrigebox}[Correction de la partie D]
Le résultat peut surprendre : trente minutes semblent beaucoup. Pourtant, il faut chauffer non pas seulement les mains, mais environ \(36\ \mathrm{kg}\) d'air. 

Quand on se frotte les mains, la chaleur produite est concentrée dans une petite masse de peau. La température locale augmente donc vite. Pour une pièce, l'énergie se répartit dans un volume beaucoup plus grand.

Si la pièce était deux fois plus grande, la masse d'air serait deux fois plus grande, donc l'énergie nécessaire serait deux fois plus grande. À puissance identique, le temps serait aussi multiplié par deux.
\end{corrigebox}

\begin{corrigebox}[Correction de la partie E]
Le calcul précédent est très favorable. Dans la réalité :
\begin{itemize}
  \item une partie de l'énergie chauffe les mains elles-mêmes ;
  \item une partie est transmise au corps ;
  \item une partie chauffe localement l'air proche ;
  \item une partie est perdue vers les murs, les meubles, le sol, le plafond ;
  \item une partie est perdue par renouvellement d'air ou par conduction à travers les parois.
\end{itemize}

Un modèle plus réaliste devrait inclure une capacité thermique effective beaucoup plus grande que celle de l'air seul :
\[
C_{\mathrm{eff}}=m_{\mathrm{air}}c_{\mathrm{air}}+C_{\mathrm{murs}}+C_{\mathrm{meubles}}+\cdots
\]
et des pertes :
\[
\frac{\dd T}{\dd t}
=
\frac{P}{C_{\mathrm{eff}}}
-
\frac{1}{\tau}(T-T_{\mathrm{ext}}).
\]
Ce modèle montre qu'il serait en pratique quasiment impossible de chauffer significativement une pièce en se frottant simplement les mains.
\end{corrigebox}

\begin{approfondissementbox}[Niveau prépa / Bac+3]
On peut modéliser la pièce comme un système thermique de capacité \(C_{\mathrm{eff}}\) échangeant avec l'extérieur. Si la personne fournit une puissance constante \(P\), alors :
\[
C_{\mathrm{eff}}\frac{\dd T}{\dd t}=P-G(T-T_{\mathrm{ext}}),
\]
où \(G\) est la conductance thermique globale de la pièce vers l'extérieur.

En posant :
\[
\tau=\frac{C_{\mathrm{eff}}}{G},
\]
on obtient :
\[
\frac{\dd T}{\dd t}
=
\frac{P}{C_{\mathrm{eff}}}
-
\frac{1}{\tau}(T-T_{\mathrm{ext}}).
\]
La température tend vers :
\[
T_{\infty}=T_{\mathrm{ext}}+\frac{P}{G}.
\]
Cela signifie qu'une faible puissance ne permet pas de faire monter indéfiniment la température : les pertes finissent par compenser exactement l'apport de chaleur.
\end{approfondissementbox}

% ------------------------------------------------------------
\subsection{Exercice corrigé 2 — Combien de personnes faut-il dans une salle pour la réchauffer de \(1^\circ\mathrm C\) ?}
% ------------------------------------------------------------

\begin{exercicebox}[Question de départ]
Une salle de classe contient initialement de l'air à \(18^\circ\mathrm C\). On se demande combien de personnes il faudrait placer dans cette salle pour réchauffer l'air de \(1^\circ\mathrm C\) en dix minutes, uniquement grâce à la chaleur dégagée par leur corps.

On reprend une salle de volume :
\[
V=150\ \mathrm{m^3}.
\]
On suppose que chaque personne dégage une puissance thermique moyenne :
\[
P_{\mathrm{pers}}=100\ \mathrm W.
\]
On néglige d'abord les pertes et on ne chauffe que l'air.
\end{exercicebox}

\subsubsection*{Partie A — Énergie nécessaire}

\begin{enumerate}
  \item Calculer la masse d'air contenue dans la salle.
  \item Calculer l'énergie nécessaire pour augmenter la température de l'air de \(1\ \mathrm K\).
  \item Convertir dix minutes en secondes.
  \item Déterminer la puissance totale nécessaire.
  \item En déduire le nombre de personnes nécessaires.
\end{enumerate}

\subsubsection*{Correction}

\begin{corrigebox}
La masse d'air vaut :
\[
m=\rho V=1{,}2\times150=180\ \mathrm{kg}.
\]
L'énergie nécessaire vaut :
\[
Q=mc\Delta T=180\times1000\times1=180000\ \mathrm J.
\]
Dix minutes valent :
\[
\Delta t=600\ \mathrm s.
\]
La puissance totale nécessaire est :
\[
P_{\mathrm{tot}}=\frac{Q}{\Delta t}=\frac{180000}{600}=300\ \mathrm W.
\]
Avec \(100\ \mathrm W\) par personne :
\[
N=\frac{300}{100}=3.
\]
Dans ce modèle très optimiste, trois personnes suffiraient à réchauffer l'air de \(1^\circ\mathrm C\) en dix minutes.
\end{corrigebox}

\subsubsection*{Discussion}

\begin{attentionbox}[Pourquoi le résultat réel est différent]
Le résultat est volontairement optimiste. Dans une vraie salle, les personnes ne chauffent pas seulement l'air : elles chauffent aussi les murs, les tables, les chaises, et il existe des pertes vers l'extérieur. La capacité thermique effective de la salle est donc beaucoup plus grande que celle de l'air seul.
\end{attentionbox}

\begin{approfondissementbox}[Niveau prépa / Bac+3]
Si la salle contient des murs et meubles de capacité thermique totale \(C_s\), l'énergie nécessaire devient :
\[
Q=(m_{\mathrm{air}}c_{\mathrm{air}}+C_s)\Delta T.
\]
Comme \(C_s\) peut être très supérieur à \(m_{\mathrm{air}}c_{\mathrm{air}}\), le nombre réel de personnes nécessaires peut être beaucoup plus grand que le calcul naïf.
\end{approfondissementbox}

% ------------------------------------------------------------
\subsection{Exercice corrigé 3 — Peut-on faire bouillir un verre d'eau en criant ?}
% ------------------------------------------------------------

\begin{exercicebox}[Question de départ]
On considère un verre contenant \(200\ \mathrm{mL}\) d'eau initialement à \(20^\circ\mathrm C\). Une personne crie à côté du verre. On veut savoir s'il est possible de porter l'eau à ébullition grâce à l'énergie sonore du cri.

Données :
\[
m_{\mathrm{eau}}=0{,}200\ \mathrm{kg},
\qquad
c_{\mathrm{eau}}=4180\ \mathrm{J.kg^{-1}.K^{-1}},
\qquad
\Delta T=80\ \mathrm K.
\]
On suppose très généreusement que la puissance sonore reçue par l'eau est :
\[
P=1\ \mathrm W.
\]
\end{exercicebox}

\subsubsection*{Partie A — Calcul de l'énergie thermique}

\begin{enumerate}
  \item Calculer l'énergie nécessaire pour chauffer l'eau de \(20^\circ\mathrm C\) à \(100^\circ\mathrm C\).
  \item Calculer le temps nécessaire si la puissance reçue vaut \(1\ \mathrm W\).
  \item Convertir ce temps en heures.
  \item Discuter la valeur choisie pour \(P\).
  \item Conclure.
\end{enumerate}

\begin{corrigebox}[Correction]
L'énergie nécessaire vaut :
\[
Q=mc\Delta T
=
0{,}200\times4180\times80
=
66880\ \mathrm J.
\]
Si la puissance reçue vaut \(1\ \mathrm W\), alors :
\[
t=\frac{Q}{P}=66880\ \mathrm s.
\]
En heures :
\[
t=\frac{66880}{3600}\simeq18{,}6\ \mathrm h.
\]
Même avec une hypothèse extrêmement favorable, il faudrait presque 19 heures. En réalité, la puissance sonore réellement absorbée par l'eau serait bien plus faible que \(1\ \mathrm W\). Il est donc impossible de faire bouillir un verre d'eau en criant.
\end{corrigebox}

\begin{approfondissementbox}[Niveau prépa / Bac+3]
Le son transporte une intensité acoustique :
\[
I=\frac{P}{S}.
\]
Le niveau sonore en décibels est :
\[
L=10\log_{10}\left(\frac{I}{I_0}\right),
\qquad
I_0=10^{-12}\ \mathrm{W.m^{-2}}.
\]
Un son très intense peut correspondre à une puissance surfacique non négligeable, mais l'absorption effective par un petit volume d'eau reste très faible. Le problème réel relève du couplage entre onde acoustique, absorption, réflexion et conversion en chaleur.
\end{approfondissementbox}

% ------------------------------------------------------------
\subsection{Exercice corrigé 4 — Combien d'énergie faut-il pour évaporer la sueur d'un coureur ?}
% ------------------------------------------------------------

\begin{exercicebox}[Question de départ]
Un coureur transpire pendant un effort. On veut comprendre pourquoi l'évaporation de la sueur refroidit efficacement le corps.

On suppose qu'une masse :
\[
m=0{,}50\ \mathrm{kg}
\]
de sueur s'évapore pendant une séance. La chaleur latente de vaporisation de l'eau vaut :
\[
L_v=2{,}3\times10^6\ \mathrm{J.kg^{-1}}.
\]
\end{exercicebox}

\subsubsection*{Partie A — Énergie évacuée}

\begin{enumerate}
  \item Écrire la relation entre énergie évacuée et masse évaporée.
  \item Calculer l'énergie évacuée.
  \item Convertir en kilojoules.
  \item Comparer à l'énergie nécessaire pour chauffer \(1\ \mathrm{kg}\) d'eau de \(1^\circ\mathrm C\).
  \item Expliquer pourquoi la transpiration est efficace.
\end{enumerate}

\begin{corrigebox}
L'énergie évacuée par évaporation est :
\[
Q=mL_v.
\]
Donc :
\[
Q=0{,}50\times2{,}3\times10^6
=
1{,}15\times10^6\ \mathrm J.
\]
Soit :
\[
Q=1150\ \mathrm{kJ}.
\]
C'est énorme à l'échelle du corps humain. Pour chauffer \(1\ \mathrm{kg}\) d'eau de \(1^\circ\mathrm C\), il faut environ \(4180\ \mathrm J\). L'évaporation de seulement \(0{,}50\ \mathrm{kg}\) d'eau retire donc autant d'énergie que le refroidissement de plusieurs centaines de kilogrammes d'eau de \(1^\circ\mathrm C\).
\end{corrigebox}

\begin{approfondissementbox}[Niveau prépa / Bac+3]
Le bilan thermique du corps peut s'écrire :
\[
C\frac{\dd T}{\dd t}=P_{\mathrm{métab}}-P_{\mathrm{conv}}-P_{\mathrm{rad}}-\dot m L_v.
\]
Le terme \(\dot m L_v\) représente la puissance évacuée par évaporation. Par temps humide, \(\dot m\) diminue, donc le refroidissement devient moins efficace.
\end{approfondissementbox}

% ------------------------------------------------------------
\subsection{Exercice corrigé 5 — Pourquoi une voiture consomme beaucoup plus à grande vitesse ?}
% ------------------------------------------------------------

\begin{exercicebox}[Question de départ]
On veut expliquer pourquoi rouler à \(130\ \mathrm{km.h^{-1}}\) consomme beaucoup plus que rouler à \(90\ \mathrm{km.h^{-1}}\), même sur route plate.

On modélise la force de traînée aérodynamique par :
\[
F_D=\frac12\rho C_DSv^2.
\]
La puissance nécessaire pour compenser cette force vaut :
\[
P=F_Dv.
\]
\end{exercicebox}

\subsubsection*{Partie A — Dépendance en vitesse}

\begin{enumerate}
  \item Montrer que la force de traînée est proportionnelle à \(v^2\).
  \item Montrer que la puissance de traînée est proportionnelle à \(v^3\).
  \item Convertir \(90\ \mathrm{km.h^{-1}}\) et \(130\ \mathrm{km.h^{-1}}\) en \(\mathrm{m.s^{-1}}\).
  \item Calculer le rapport :
  \[
  \left(\frac{130}{90}\right)^3.
  \]
  \item Interpréter.
\end{enumerate}

\begin{corrigebox}
Comme :
\[
F_D=\frac12\rho C_DSv^2,
\]
on a :
\[
F_D\propto v^2.
\]
La puissance est :
\[
P=F_Dv,
\]
donc :
\[
P\propto v^3.
\]

Le rapport des puissances vaut :
\[
\left(\frac{130}{90}\right)^3\simeq3{,}0.
\]
La puissance nécessaire pour vaincre la traînée est donc environ trois fois plus grande à \(130\ \mathrm{km.h^{-1}}\) qu'à \(90\ \mathrm{km.h^{-1}}\). C'est une des raisons principales de la forte augmentation de consommation sur autoroute.
\end{corrigebox}

\begin{approfondissementbox}[Niveau prépa / Bac+3]
La traînée quadratique vient d'un transfert de quantité de mouvement au fluide. À grand nombre de Reynolds, les effets inertiels dominent et la force caractéristique est :
\[
F\sim \rho S v^2.
\]
Le coefficient \(C_D\) encode la géométrie, la turbulence, la séparation de couche limite et l'écoulement autour du véhicule.
\end{approfondissementbox}

% ------------------------------------------------------------
\subsection{Exercice corrigé 6 — À quelle vitesse faut-il lancer une balle pour atteindre une cible ?}
% ------------------------------------------------------------

\begin{exercicebox}[Question de départ]
On lance une balle à une distance horizontale \(L=20\ \mathrm m\), avec un angle \(\theta=45^\circ\). On néglige les frottements de l'air. On cherche la vitesse initiale nécessaire pour que la balle retombe à la même hauteur que son point de départ au niveau de la cible.
\end{exercicebox}

\subsubsection*{Partie A — Équations du mouvement}

\begin{enumerate}
  \item Écrire les équations horaires :
  \[
  x(t)=v_0\cos\theta\,t,
  \qquad
  y(t)=v_0\sin\theta\,t-\frac12gt^2.
  \]
  \item Déterminer le temps de vol lorsque \(y=0\).
  \item En déduire la portée :
  \[
  R=\frac{v_0^2\sin(2\theta)}{g}.
  \]
  \item Résoudre \(R=L\).
  \item Calculer \(v_0\).
\end{enumerate}

\begin{corrigebox}
Pour une retombée à même hauteur :
\[
R=\frac{v_0^2\sin(2\theta)}{g}.
\]
Avec \(\theta=45^\circ\), on a :
\[
\sin(2\theta)=\sin(90^\circ)=1.
\]
Donc :
\[
L=\frac{v_0^2}{g},
\]
et :
\[
v_0=\sqrt{gL}.
\]
Numériquement :
\[
v_0=\sqrt{9{,}81\times20}\simeq14\ \mathrm{m.s^{-1}}.
\]
\end{corrigebox}

\begin{approfondissementbox}[Niveau prépa / Bac+3]
Avec frottement quadratique, les équations deviennent :
\[
m\frac{\dd \vec v}{\dd t}=m\vec g-\frac12\rho C_DS v\vec v.
\]
Elles ne donnent plus une parabole exacte. La portée optimale n'est plus obtenue à \(45^\circ\), car la traînée pénalise fortement les trajectoires longues et hautes.
\end{approfondissementbox}

% ------------------------------------------------------------
\subsection{Exercice corrigé 7 — Pourquoi une balle liftée retombe plus vite ?}
% ------------------------------------------------------------

\begin{exercicebox}[Question de départ]
Au tennis, une balle frappée avec un fort lift peut passer haut au-dessus du filet tout en retombant dans le terrain. On veut expliquer ce phénomène par l'effet Magnus.
\end{exercicebox}

\subsubsection*{Partie A — Modélisation}

\begin{enumerate}
  \item Identifier les forces principales sur la balle.
  \item Expliquer qualitativement l'effet de la rotation.
  \item Écrire une force de Magnus sous la forme :
  \[
  \vec F_M\propto \vec\omega\times\vec v.
  \]
  \item Déterminer le sens de la force pour un topspin.
  \item Expliquer l'intérêt tactique du lift.
\end{enumerate}

\begin{corrigebox}
Pour une balle en topspin, la rotation est telle que l'effet Magnus ajoute une force vers le bas. La balle a donc une trajectoire plus courbée que sous l'effet du poids seul. Elle peut être frappée plus vite, passer au-dessus du filet, puis redescendre dans le court.
\end{corrigebox}

\begin{approfondissementbox}[Niveau prépa / Bac+3]
Une modélisation plus complète introduit :
\[
m\frac{\dd \vec v}{\dd t}
=
m\vec g
-\frac12\rho C_D S v\vec v
+
C_M\rho R^3\vec\omega\times\vec v.
\]
Cette équation différentielle vectorielle couple la traînée et la portance de Magnus. Elle doit généralement être résolue numériquement.
\end{approfondissementbox}

% ------------------------------------------------------------
\subsection{Exercice corrigé 8 — Combien de temps un glaçon met-il à fondre ?}
% ------------------------------------------------------------

\begin{exercicebox}[Question de départ]
Un glaçon de masse \(m=30\ \mathrm g\) est placé dans une pièce chaude. On suppose qu'il est déjà à \(0^\circ\mathrm C\). La chaleur latente de fusion de la glace vaut :
\[
L_f=3{,}34\times10^5\ \mathrm{J.kg^{-1}}.
\]
On suppose que le glaçon reçoit une puissance thermique moyenne \(P=5\ \mathrm W\).
\end{exercicebox}

\subsubsection*{Partie A — Fusion}

\begin{enumerate}
  \item Convertir la masse en kilogrammes.
  \item Calculer l'énergie nécessaire pour faire fondre le glaçon.
  \item Calculer le temps de fusion.
  \item Convertir en minutes.
  \item Discuter la valeur de la puissance reçue.
\end{enumerate}

\begin{corrigebox}
La masse vaut :
\[
m=0{,}030\ \mathrm{kg}.
\]
L'énergie nécessaire vaut :
\[
Q=mL_f=0{,}030\times3{,}34\times10^5=10020\ \mathrm J.
\]
Le temps vaut :
\[
t=\frac{Q}{P}=\frac{10020}{5}=2004\ \mathrm s.
\]
Soit environ :
\[
33\ \mathrm{min}.
\]
\end{corrigebox}

\begin{approfondissementbox}[Niveau prépa / Bac+3]
La puissance reçue par le glaçon peut être modélisée comme la somme d'un flux convectif et d'un flux radiatif :
\[
P=hS(T_{\mathrm{air}}-T_{\mathrm{glace}})+\varepsilon\sigma S(T_{\mathrm{parois}}^4-T_{\mathrm{glace}}^4).
\]
Comme la surface du glaçon évolue pendant la fonte, la puissance n'est pas constante.
\end{approfondissementbox}

% ------------------------------------------------------------
\subsection{Exercice corrigé 9 — Pourquoi une fibre optique transmet-elle si vite Internet ?}
% ------------------------------------------------------------

\begin{exercicebox}[Question de départ]
Une information lumineuse parcourt une fibre optique de longueur \(L=100\ \mathrm{km}\). L'indice du cœur vaut \(n=1{,}5\). On veut déterminer le temps de propagation du signal.
\end{exercicebox}

\subsubsection*{Partie A — Temps de propagation}

\begin{enumerate}
  \item Rappeler la relation entre vitesse de la lumière dans le milieu et indice.
  \item Calculer la vitesse de la lumière dans la fibre.
  \item Calculer le temps de propagation sur \(100\ \mathrm{km}\).
  \item Comparer à la propagation dans le vide.
  \item Discuter l'importance des latences réseau.
\end{enumerate}

\begin{corrigebox}
Dans un milieu d'indice \(n\), la vitesse vaut :
\[
v=\frac{c}{n}.
\]
Donc :
\[
v=\frac{3{,}0\times10^8}{1{,}5}=2{,}0\times10^8\ \mathrm{m.s^{-1}}.
\]
Pour \(L=100\ \mathrm{km}=10^5\ \mathrm m\),
\[
t=\frac{L}{v}
=
\frac{10^5}{2{,}0\times10^8}
=
5{,}0\times10^{-4}\ \mathrm s.
\]
Donc :
\[
t=0{,}50\ \mathrm{ms}.
\]
\end{corrigebox}

\begin{approfondissementbox}[Niveau prépa / Bac+3]
La vitesse pertinente pour un paquet d'information n'est pas toujours la vitesse de phase, mais la vitesse de groupe :
\[
v_g=\frac{\dd\omega}{\dd k}.
\]
Dans une fibre dispersive, \(v_g\) dépend de la longueur d'onde, ce qui peut élargir les impulsions lumineuses.
\end{approfondissementbox}

% ------------------------------------------------------------
\subsection{Exercice corrigé 10 — Pourquoi le ciel bleu devient-il rouge au coucher du Soleil ?}
% ------------------------------------------------------------

\begin{exercicebox}[Question de départ]
On veut expliquer avec un calcul simple pourquoi le ciel est bleu le jour et pourquoi le Soleil devient rouge au coucher.
\end{exercicebox}

\subsubsection*{Partie A — Diffusion Rayleigh}

\begin{enumerate}
  \item On admet que l'intensité diffusée vérifie :
  \[
  I\propto \frac{1}{\lambda^4}.
  \]
  \item Comparer la diffusion d'une lumière bleue de longueur d'onde \(450\ \mathrm{nm}\) et d'une lumière rouge de longueur d'onde \(650\ \mathrm{nm}\).
  \item Calculer :
  \[
  \left(\frac{650}{450}\right)^4.
  \]
  \item Interpréter le résultat.
  \item Expliquer pourquoi le ciel n'est pas violet.
\end{enumerate}

\begin{corrigebox}
Le rapport vaut :
\[
\frac{I_{\mathrm{bleu}}}{I_{\mathrm{rouge}}}
=
\left(\frac{650}{450}\right)^4
\simeq 4{,}4.
\]
La lumière bleue est donc diffusée plusieurs fois plus efficacement que la lumière rouge. Le ciel apparaît bleu car cette lumière bleue est redistribuée dans toutes les directions.

Le ciel n'est pas violet car le spectre solaire contient moins de violet, l'atmosphère absorbe une partie du proche ultraviolet et l'œil humain est moins sensible au violet.
\end{corrigebox}

\begin{approfondissementbox}[Niveau prépa / Bac+3]
La diffusion Rayleigh provient de l'oscillation de dipôles induits par le champ électrique lumineux. Pour des particules beaucoup plus petites que la longueur d'onde, la section efficace varie comme :
\[
\sigma\propto\frac{1}{\lambda^4}.
\]
Ce résultat relie directement la couleur du ciel à la nature ondulatoire de la lumière et à la taille moléculaire des diffuseurs.
\end{approfondissementbox}


% ============================================================
\section{Partie XV — Tous les phénomènes sous forme de longs problèmes corrigés}
% ============================================================

\begin{attentionbox}[Nouvelle structure]
Dans cette partie, tous les phénomènes du document sont repris sous forme de \textbf{problèmes corrigés longs}. Chaque problème contient :
\begin{itemize}
  \item une question concrète ;
  \item une modélisation niveau Terminale ;
  \item une mise en équation ;
  \item une réponse rédigée niveau Terminale ;
  \item une réponse approfondie niveau Bac+3 ;
  \item une critique du modèle.
\end{itemize}
\end{attentionbox}

% ------------------------------------------------------------
\subsection{Problème corrigé 1 — Mains et chauffage d'une pièce}
% ------------------------------------------------------------

\begin{exercicebox}[Question-problème]
\textbf{Observation.} Une personne se frotte les mains pour tenter de réchauffer une pièce.

\medskip
\textbf{Question.} Calculer le temps nécessaire pour augmenter la température de l'air d'une pièce de 1 degré.

\medskip
\textbf{Thèmes mobilisés :} Thermique, énergie, puissance, capacité thermique.

\medskip
\textbf{Relation ou idée centrale :}
\[
Q=mc\Delta T,\qquad t=\frac{Q}{P}
\]
\end{exercicebox}

\subsubsection*{Partie A — Modélisation niveau Terminale}

\begin{enumerate}
  \item Décrire précisément le phénomène observé.
  \item Identifier le système physique étudié.
  \item Lister les grandeurs importantes et leurs unités.
  \item Écrire la loi physique simple qui permet de commencer le raisonnement.
  \item Expliquer les hypothèses simplificatrices du modèle Terminale.
\end{enumerate}

\subsubsection*{Partie B — Mise en équation et calcul}

\begin{enumerate}
  \item Définir les variables utiles.
  \item Écrire la relation mathématique principale.
  \item Vérifier l'homogénéité de la formule.
  \item Choisir un ordre de grandeur réaliste pour les données manquantes.
  \item Effectuer le calcul ou établir le raisonnement qualitatif demandé.
  \item Interpréter le résultat en une phrase claire.
\end{enumerate}

\subsubsection*{Partie C — Réponse accessible niveau Terminale}

\begin{corrigebox}[Réponse Terminale]
On part de l'observation concrète : Une personne se frotte les mains pour tenter de réchauffer une pièce.

Au niveau Terminale, on cherche d'abord la grandeur qui contrôle le phénomène. Ici, l'idée centrale est donnée par :
\[
Q=mc\Delta T,\qquad t=\frac{Q}{P}
\]
Cette relation permet de relier la cause physique à l'effet observé. On vérifie les unités, puis on insère des ordres de grandeur réalistes.

La conclusion attendue est la suivante : Calculer le temps nécessaire pour augmenter la température de l'air d'une pièce de 1 degré. La réponse doit toujours préciser les hypothèses utilisées. Si le phénomène dépend d'un milieu réel, d'une géométrie complexe, de frottements, de pertes ou de gradients, le résultat obtenu est un ordre de grandeur et non une valeur exacte.
\end{corrigebox}

\subsubsection*{Partie D — Approfondissement niveau Bac+3}

\begin{approfondissementbox}[Réponse Bac+3]
Au niveau Bac+3, on ne se contente plus d'une formule globale. On remplace le modèle simple par une loi locale, différentielle ou énergétique.

Une forme générale possible est :
\[
\frac{\mathrm d X}{\mathrm dt}=\text{apport}-\text{pertes},
\]
ou, pour un phénomène propagatif :
\[
\frac{\partial^2 \psi}{\partial t^2}=c^2\Delta\psi.
\]
Selon le cas, on peut aussi utiliser un principe variationnel, par exemple le principe de Fermat en optique :
\[
\delta \int n\,\mathrm ds=0.
\]

Pour ce phénomène, le modèle Bac+3 consiste à partir de la relation centrale
\[
Q=mc\Delta T,\qquad t=\frac{Q}{P}
\]
puis à introduire les corrections négligées : dépendance spatiale, dépendance temporelle, non-linéarités, pertes, conditions aux limites ou fluctuations. Cela permet de passer d'une réponse d'ordre de grandeur à une prédiction plus robuste.
\end{approfondissementbox}

\subsubsection*{Partie E — Limites, critique et prolongement}

\begin{enumerate}
  \item Citer au moins deux limites du modèle Terminale.
  \item Dire quelles grandeurs seraient difficiles à mesurer expérimentalement.
  \item Proposer une expérience simple permettant de tester le modèle.
  \item Expliquer ce qui changerait si l'on doublait une grandeur clé.
  \item Rédiger une conclusion courte, accessible à un élève de Terminale.
\end{enumerate}

\begin{corrigebox}[Synthèse rédigée]
Le phénomène « Mains et chauffage d'une pièce » s'explique par un petit nombre de lois physiques simples, mais sa manifestation réelle dépend souvent de nombreux paramètres. Le modèle Terminale donne l'intuition et les ordres de grandeur ; le modèle Bac+3 introduit les variations locales, les pertes, les équations différentielles ou les principes variationnels nécessaires pour décrire finement la réalité.
\end{corrigebox}

% ------------------------------------------------------------
\subsection{Problème corrigé 2 — Mirage inférieur sur une route chaude}
% ------------------------------------------------------------

\begin{exercicebox}[Question-problème]
\textbf{Observation.} Sur une route chauffée, on croit voir une flaque au loin.

\medskip
\textbf{Question.} Expliquer pourquoi le rayon se courbe et pourquoi l'image est virtuelle.

\medskip
\textbf{Thèmes mobilisés :} Réfraction, gradient d'indice, optique géométrique.

\medskip
\textbf{Relation ou idée centrale :}
\[
n_k\sin i_k=n_{k+1}\sin i_{k+1}
\]
\end{exercicebox}

\subsubsection*{Partie A — Modélisation niveau Terminale}

\begin{enumerate}
  \item Décrire précisément le phénomène observé.
  \item Identifier le système physique étudié.
  \item Lister les grandeurs importantes et leurs unités.
  \item Écrire la loi physique simple qui permet de commencer le raisonnement.
  \item Expliquer les hypothèses simplificatrices du modèle Terminale.
\end{enumerate}

\subsubsection*{Partie B — Mise en équation et calcul}

\begin{enumerate}
  \item Définir les variables utiles.
  \item Écrire la relation mathématique principale.
  \item Vérifier l'homogénéité de la formule.
  \item Choisir un ordre de grandeur réaliste pour les données manquantes.
  \item Effectuer le calcul ou établir le raisonnement qualitatif demandé.
  \item Interpréter le résultat en une phrase claire.
\end{enumerate}

\subsubsection*{Partie C — Réponse accessible niveau Terminale}

\begin{corrigebox}[Réponse Terminale]
On part de l'observation concrète : Sur une route chauffée, on croit voir une flaque au loin.

Au niveau Terminale, on cherche d'abord la grandeur qui contrôle le phénomène. Ici, l'idée centrale est donnée par :
\[
n_k\sin i_k=n_{k+1}\sin i_{k+1}
\]
Cette relation permet de relier la cause physique à l'effet observé. On vérifie les unités, puis on insère des ordres de grandeur réalistes.

La conclusion attendue est la suivante : Expliquer pourquoi le rayon se courbe et pourquoi l'image est virtuelle. La réponse doit toujours préciser les hypothèses utilisées. Si le phénomène dépend d'un milieu réel, d'une géométrie complexe, de frottements, de pertes ou de gradients, le résultat obtenu est un ordre de grandeur et non une valeur exacte.
\end{corrigebox}

\subsubsection*{Partie D — Approfondissement niveau Bac+3}

\begin{approfondissementbox}[Réponse Bac+3]
Au niveau Bac+3, on ne se contente plus d'une formule globale. On remplace le modèle simple par une loi locale, différentielle ou énergétique.

Une forme générale possible est :
\[
\frac{\mathrm d X}{\mathrm dt}=\text{apport}-\text{pertes},
\]
ou, pour un phénomène propagatif :
\[
\frac{\partial^2 \psi}{\partial t^2}=c^2\Delta\psi.
\]
Selon le cas, on peut aussi utiliser un principe variationnel, par exemple le principe de Fermat en optique :
\[
\delta \int n\,\mathrm ds=0.
\]

Pour ce phénomène, le modèle Bac+3 consiste à partir de la relation centrale
\[
n_k\sin i_k=n_{k+1}\sin i_{k+1}
\]
puis à introduire les corrections négligées : dépendance spatiale, dépendance temporelle, non-linéarités, pertes, conditions aux limites ou fluctuations. Cela permet de passer d'une réponse d'ordre de grandeur à une prédiction plus robuste.
\end{approfondissementbox}

\subsubsection*{Partie E — Limites, critique et prolongement}

\begin{enumerate}
  \item Citer au moins deux limites du modèle Terminale.
  \item Dire quelles grandeurs seraient difficiles à mesurer expérimentalement.
  \item Proposer une expérience simple permettant de tester le modèle.
  \item Expliquer ce qui changerait si l'on doublait une grandeur clé.
  \item Rédiger une conclusion courte, accessible à un élève de Terminale.
\end{enumerate}

\begin{corrigebox}[Synthèse rédigée]
Le phénomène « Mirage inférieur sur une route chaude » s'explique par un petit nombre de lois physiques simples, mais sa manifestation réelle dépend souvent de nombreux paramètres. Le modèle Terminale donne l'intuition et les ordres de grandeur ; le modèle Bac+3 introduit les variations locales, les pertes, les équations différentielles ou les principes variationnels nécessaires pour décrire finement la réalité.
\end{corrigebox}

% ------------------------------------------------------------
\subsection{Problème corrigé 3 — Mirage supérieur}
% ------------------------------------------------------------

\begin{exercicebox}[Question-problème]
\textbf{Observation.} Un bateau ou un paysage semble flotter au-dessus de l'horizon.

\medskip
\textbf{Question.} Comprendre comment un gradient inverse peut rendre visible un objet lointain.

\medskip
\textbf{Thèmes mobilisés :} Réfraction atmosphérique, inversion thermique.

\medskip
\textbf{Relation ou idée centrale :}
\[
n=n(z)
\]
\end{exercicebox}

\subsubsection*{Partie A — Modélisation niveau Terminale}

\begin{enumerate}
  \item Décrire précisément le phénomène observé.
  \item Identifier le système physique étudié.
  \item Lister les grandeurs importantes et leurs unités.
  \item Écrire la loi physique simple qui permet de commencer le raisonnement.
  \item Expliquer les hypothèses simplificatrices du modèle Terminale.
\end{enumerate}

\subsubsection*{Partie B — Mise en équation et calcul}

\begin{enumerate}
  \item Définir les variables utiles.
  \item Écrire la relation mathématique principale.
  \item Vérifier l'homogénéité de la formule.
  \item Choisir un ordre de grandeur réaliste pour les données manquantes.
  \item Effectuer le calcul ou établir le raisonnement qualitatif demandé.
  \item Interpréter le résultat en une phrase claire.
\end{enumerate}

\subsubsection*{Partie C — Réponse accessible niveau Terminale}

\begin{corrigebox}[Réponse Terminale]
On part de l'observation concrète : Un bateau ou un paysage semble flotter au-dessus de l'horizon.

Au niveau Terminale, on cherche d'abord la grandeur qui contrôle le phénomène. Ici, l'idée centrale est donnée par :
\[
n=n(z)
\]
Cette relation permet de relier la cause physique à l'effet observé. On vérifie les unités, puis on insère des ordres de grandeur réalistes.

La conclusion attendue est la suivante : Comprendre comment un gradient inverse peut rendre visible un objet lointain. La réponse doit toujours préciser les hypothèses utilisées. Si le phénomène dépend d'un milieu réel, d'une géométrie complexe, de frottements, de pertes ou de gradients, le résultat obtenu est un ordre de grandeur et non une valeur exacte.
\end{corrigebox}

\subsubsection*{Partie D — Approfondissement niveau Bac+3}

\begin{approfondissementbox}[Réponse Bac+3]
Au niveau Bac+3, on ne se contente plus d'une formule globale. On remplace le modèle simple par une loi locale, différentielle ou énergétique.

Une forme générale possible est :
\[
\frac{\mathrm d X}{\mathrm dt}=\text{apport}-\text{pertes},
\]
ou, pour un phénomène propagatif :
\[
\frac{\partial^2 \psi}{\partial t^2}=c^2\Delta\psi.
\]
Selon le cas, on peut aussi utiliser un principe variationnel, par exemple le principe de Fermat en optique :
\[
\delta \int n\,\mathrm ds=0.
\]

Pour ce phénomène, le modèle Bac+3 consiste à partir de la relation centrale
\[
n=n(z)
\]
puis à introduire les corrections négligées : dépendance spatiale, dépendance temporelle, non-linéarités, pertes, conditions aux limites ou fluctuations. Cela permet de passer d'une réponse d'ordre de grandeur à une prédiction plus robuste.
\end{approfondissementbox}

\subsubsection*{Partie E — Limites, critique et prolongement}

\begin{enumerate}
  \item Citer au moins deux limites du modèle Terminale.
  \item Dire quelles grandeurs seraient difficiles à mesurer expérimentalement.
  \item Proposer une expérience simple permettant de tester le modèle.
  \item Expliquer ce qui changerait si l'on doublait une grandeur clé.
  \item Rédiger une conclusion courte, accessible à un élève de Terminale.
\end{enumerate}

\begin{corrigebox}[Synthèse rédigée]
Le phénomène « Mirage supérieur » s'explique par un petit nombre de lois physiques simples, mais sa manifestation réelle dépend souvent de nombreux paramètres. Le modèle Terminale donne l'intuition et les ordres de grandeur ; le modèle Bac+3 introduit les variations locales, les pertes, les équations différentielles ou les principes variationnels nécessaires pour décrire finement la réalité.
\end{corrigebox}

% ------------------------------------------------------------
\subsection{Problème corrigé 4 — Fata Morgana}
% ------------------------------------------------------------

\begin{exercicebox}[Question-problème]
\textbf{Observation.} Des objets lointains sont étirés ou empilés verticalement.

\medskip
\textbf{Question.} Expliquer la formation d'images multiples par couches d'air.

\medskip
\textbf{Thèmes mobilisés :} Stratification atmosphérique, images multiples.

\medskip
\textbf{Relation ou idée centrale :}
\[
\text{plusieurs chemins optiques}
\]
\end{exercicebox}

\subsubsection*{Partie A — Modélisation niveau Terminale}

\begin{enumerate}
  \item Décrire précisément le phénomène observé.
  \item Identifier le système physique étudié.
  \item Lister les grandeurs importantes et leurs unités.
  \item Écrire la loi physique simple qui permet de commencer le raisonnement.
  \item Expliquer les hypothèses simplificatrices du modèle Terminale.
\end{enumerate}

\subsubsection*{Partie B — Mise en équation et calcul}

\begin{enumerate}
  \item Définir les variables utiles.
  \item Écrire la relation mathématique principale.
  \item Vérifier l'homogénéité de la formule.
  \item Choisir un ordre de grandeur réaliste pour les données manquantes.
  \item Effectuer le calcul ou établir le raisonnement qualitatif demandé.
  \item Interpréter le résultat en une phrase claire.
\end{enumerate}

\subsubsection*{Partie C — Réponse accessible niveau Terminale}

\begin{corrigebox}[Réponse Terminale]
On part de l'observation concrète : Des objets lointains sont étirés ou empilés verticalement.

Au niveau Terminale, on cherche d'abord la grandeur qui contrôle le phénomène. Ici, l'idée centrale est donnée par :
\[
\text{plusieurs chemins optiques}
\]
Cette relation permet de relier la cause physique à l'effet observé. On vérifie les unités, puis on insère des ordres de grandeur réalistes.

La conclusion attendue est la suivante : Expliquer la formation d'images multiples par couches d'air. La réponse doit toujours préciser les hypothèses utilisées. Si le phénomène dépend d'un milieu réel, d'une géométrie complexe, de frottements, de pertes ou de gradients, le résultat obtenu est un ordre de grandeur et non une valeur exacte.
\end{corrigebox}

\subsubsection*{Partie D — Approfondissement niveau Bac+3}

\begin{approfondissementbox}[Réponse Bac+3]
Au niveau Bac+3, on ne se contente plus d'une formule globale. On remplace le modèle simple par une loi locale, différentielle ou énergétique.

Une forme générale possible est :
\[
\frac{\mathrm d X}{\mathrm dt}=\text{apport}-\text{pertes},
\]
ou, pour un phénomène propagatif :
\[
\frac{\partial^2 \psi}{\partial t^2}=c^2\Delta\psi.
\]
Selon le cas, on peut aussi utiliser un principe variationnel, par exemple le principe de Fermat en optique :
\[
\delta \int n\,\mathrm ds=0.
\]

Pour ce phénomène, le modèle Bac+3 consiste à partir de la relation centrale
\[
\text{plusieurs chemins optiques}
\]
puis à introduire les corrections négligées : dépendance spatiale, dépendance temporelle, non-linéarités, pertes, conditions aux limites ou fluctuations. Cela permet de passer d'une réponse d'ordre de grandeur à une prédiction plus robuste.
\end{approfondissementbox}

\subsubsection*{Partie E — Limites, critique et prolongement}

\begin{enumerate}
  \item Citer au moins deux limites du modèle Terminale.
  \item Dire quelles grandeurs seraient difficiles à mesurer expérimentalement.
  \item Proposer une expérience simple permettant de tester le modèle.
  \item Expliquer ce qui changerait si l'on doublait une grandeur clé.
  \item Rédiger une conclusion courte, accessible à un élève de Terminale.
\end{enumerate}

\begin{corrigebox}[Synthèse rédigée]
Le phénomène « Fata Morgana » s'explique par un petit nombre de lois physiques simples, mais sa manifestation réelle dépend souvent de nombreux paramètres. Le modèle Terminale donne l'intuition et les ordres de grandeur ; le modèle Bac+3 introduit les variations locales, les pertes, les équations différentielles ou les principes variationnels nécessaires pour décrire finement la réalité.
\end{corrigebox}

% ------------------------------------------------------------
\subsection{Problème corrigé 5 — Arc-en-ciel primaire}
% ------------------------------------------------------------

\begin{exercicebox}[Question-problème]
\textbf{Observation.} Un arc coloré apparaît face au Soleil après la pluie.

\medskip
\textbf{Question.} Montrer pourquoi l'arc apparaît autour de 42 degrés.

\medskip
\textbf{Thèmes mobilisés :} Réfraction, réflexion interne, dispersion.

\medskip
\textbf{Relation ou idée centrale :}
\[
D(i)=\pi+2i-4r,\qquad \sin i=n\sin r
\]
\end{exercicebox}

\subsubsection*{Partie A — Modélisation niveau Terminale}

\begin{enumerate}
  \item Décrire précisément le phénomène observé.
  \item Identifier le système physique étudié.
  \item Lister les grandeurs importantes et leurs unités.
  \item Écrire la loi physique simple qui permet de commencer le raisonnement.
  \item Expliquer les hypothèses simplificatrices du modèle Terminale.
\end{enumerate}

\subsubsection*{Partie B — Mise en équation et calcul}

\begin{enumerate}
  \item Définir les variables utiles.
  \item Écrire la relation mathématique principale.
  \item Vérifier l'homogénéité de la formule.
  \item Choisir un ordre de grandeur réaliste pour les données manquantes.
  \item Effectuer le calcul ou établir le raisonnement qualitatif demandé.
  \item Interpréter le résultat en une phrase claire.
\end{enumerate}

\subsubsection*{Partie C — Réponse accessible niveau Terminale}

\begin{corrigebox}[Réponse Terminale]
On part de l'observation concrète : Un arc coloré apparaît face au Soleil après la pluie.

Au niveau Terminale, on cherche d'abord la grandeur qui contrôle le phénomène. Ici, l'idée centrale est donnée par :
\[
D(i)=\pi+2i-4r,\qquad \sin i=n\sin r
\]
Cette relation permet de relier la cause physique à l'effet observé. On vérifie les unités, puis on insère des ordres de grandeur réalistes.

La conclusion attendue est la suivante : Montrer pourquoi l'arc apparaît autour de 42 degrés. La réponse doit toujours préciser les hypothèses utilisées. Si le phénomène dépend d'un milieu réel, d'une géométrie complexe, de frottements, de pertes ou de gradients, le résultat obtenu est un ordre de grandeur et non une valeur exacte.
\end{corrigebox}

\subsubsection*{Partie D — Approfondissement niveau Bac+3}

\begin{approfondissementbox}[Réponse Bac+3]
Au niveau Bac+3, on ne se contente plus d'une formule globale. On remplace le modèle simple par une loi locale, différentielle ou énergétique.

Une forme générale possible est :
\[
\frac{\mathrm d X}{\mathrm dt}=\text{apport}-\text{pertes},
\]
ou, pour un phénomène propagatif :
\[
\frac{\partial^2 \psi}{\partial t^2}=c^2\Delta\psi.
\]
Selon le cas, on peut aussi utiliser un principe variationnel, par exemple le principe de Fermat en optique :
\[
\delta \int n\,\mathrm ds=0.
\]

Pour ce phénomène, le modèle Bac+3 consiste à partir de la relation centrale
\[
D(i)=\pi+2i-4r,\qquad \sin i=n\sin r
\]
puis à introduire les corrections négligées : dépendance spatiale, dépendance temporelle, non-linéarités, pertes, conditions aux limites ou fluctuations. Cela permet de passer d'une réponse d'ordre de grandeur à une prédiction plus robuste.
\end{approfondissementbox}

\subsubsection*{Partie E — Limites, critique et prolongement}

\begin{enumerate}
  \item Citer au moins deux limites du modèle Terminale.
  \item Dire quelles grandeurs seraient difficiles à mesurer expérimentalement.
  \item Proposer une expérience simple permettant de tester le modèle.
  \item Expliquer ce qui changerait si l'on doublait une grandeur clé.
  \item Rédiger une conclusion courte, accessible à un élève de Terminale.
\end{enumerate}

\begin{corrigebox}[Synthèse rédigée]
Le phénomène « Arc-en-ciel primaire » s'explique par un petit nombre de lois physiques simples, mais sa manifestation réelle dépend souvent de nombreux paramètres. Le modèle Terminale donne l'intuition et les ordres de grandeur ; le modèle Bac+3 introduit les variations locales, les pertes, les équations différentielles ou les principes variationnels nécessaires pour décrire finement la réalité.
\end{corrigebox}

% ------------------------------------------------------------
\subsection{Problème corrigé 6 — Arc-en-ciel secondaire}
% ------------------------------------------------------------

\begin{exercicebox}[Question-problème]
\textbf{Observation.} Un deuxième arc plus faible apparaît parfois avec couleurs inversées.

\medskip
\textbf{Question.} Comparer l'arc primaire et l'arc secondaire.

\medskip
\textbf{Thèmes mobilisés :} Deux réflexions internes, dispersion.

\medskip
\textbf{Relation ou idée centrale :}
\[
D_2(i)=2\pi+2i-6r
\]
\end{exercicebox}

\subsubsection*{Partie A — Modélisation niveau Terminale}

\begin{enumerate}
  \item Décrire précisément le phénomène observé.
  \item Identifier le système physique étudié.
  \item Lister les grandeurs importantes et leurs unités.
  \item Écrire la loi physique simple qui permet de commencer le raisonnement.
  \item Expliquer les hypothèses simplificatrices du modèle Terminale.
\end{enumerate}

\subsubsection*{Partie B — Mise en équation et calcul}

\begin{enumerate}
  \item Définir les variables utiles.
  \item Écrire la relation mathématique principale.
  \item Vérifier l'homogénéité de la formule.
  \item Choisir un ordre de grandeur réaliste pour les données manquantes.
  \item Effectuer le calcul ou établir le raisonnement qualitatif demandé.
  \item Interpréter le résultat en une phrase claire.
\end{enumerate}

\subsubsection*{Partie C — Réponse accessible niveau Terminale}

\begin{corrigebox}[Réponse Terminale]
On part de l'observation concrète : Un deuxième arc plus faible apparaît parfois avec couleurs inversées.

Au niveau Terminale, on cherche d'abord la grandeur qui contrôle le phénomène. Ici, l'idée centrale est donnée par :
\[
D_2(i)=2\pi+2i-6r
\]
Cette relation permet de relier la cause physique à l'effet observé. On vérifie les unités, puis on insère des ordres de grandeur réalistes.

La conclusion attendue est la suivante : Comparer l'arc primaire et l'arc secondaire. La réponse doit toujours préciser les hypothèses utilisées. Si le phénomène dépend d'un milieu réel, d'une géométrie complexe, de frottements, de pertes ou de gradients, le résultat obtenu est un ordre de grandeur et non une valeur exacte.
\end{corrigebox}

\subsubsection*{Partie D — Approfondissement niveau Bac+3}

\begin{approfondissementbox}[Réponse Bac+3]
Au niveau Bac+3, on ne se contente plus d'une formule globale. On remplace le modèle simple par une loi locale, différentielle ou énergétique.

Une forme générale possible est :
\[
\frac{\mathrm d X}{\mathrm dt}=\text{apport}-\text{pertes},
\]
ou, pour un phénomène propagatif :
\[
\frac{\partial^2 \psi}{\partial t^2}=c^2\Delta\psi.
\]
Selon le cas, on peut aussi utiliser un principe variationnel, par exemple le principe de Fermat en optique :
\[
\delta \int n\,\mathrm ds=0.
\]

Pour ce phénomène, le modèle Bac+3 consiste à partir de la relation centrale
\[
D_2(i)=2\pi+2i-6r
\]
puis à introduire les corrections négligées : dépendance spatiale, dépendance temporelle, non-linéarités, pertes, conditions aux limites ou fluctuations. Cela permet de passer d'une réponse d'ordre de grandeur à une prédiction plus robuste.
\end{approfondissementbox}

\subsubsection*{Partie E — Limites, critique et prolongement}

\begin{enumerate}
  \item Citer au moins deux limites du modèle Terminale.
  \item Dire quelles grandeurs seraient difficiles à mesurer expérimentalement.
  \item Proposer une expérience simple permettant de tester le modèle.
  \item Expliquer ce qui changerait si l'on doublait une grandeur clé.
  \item Rédiger une conclusion courte, accessible à un élève de Terminale.
\end{enumerate}

\begin{corrigebox}[Synthèse rédigée]
Le phénomène « Arc-en-ciel secondaire » s'explique par un petit nombre de lois physiques simples, mais sa manifestation réelle dépend souvent de nombreux paramètres. Le modèle Terminale donne l'intuition et les ordres de grandeur ; le modèle Bac+3 introduit les variations locales, les pertes, les équations différentielles ou les principes variationnels nécessaires pour décrire finement la réalité.
\end{corrigebox}

% ------------------------------------------------------------
\subsection{Problème corrigé 7 — Bande sombre d'Alexandre}
% ------------------------------------------------------------

\begin{exercicebox}[Question-problème]
\textbf{Observation.} La zone entre deux arcs-en-ciel semble sombre.

\medskip
\textbf{Question.} Expliquer l'absence relative de rayons dans certaines directions.

\medskip
\textbf{Thèmes mobilisés :} Distribution angulaire des rayons.

\medskip
\textbf{Relation ou idée centrale :}
\[
\theta_1\simeq42^\circ,\qquad \theta_2\simeq51^\circ
\]
\end{exercicebox}

\subsubsection*{Partie A — Modélisation niveau Terminale}

\begin{enumerate}
  \item Décrire précisément le phénomène observé.
  \item Identifier le système physique étudié.
  \item Lister les grandeurs importantes et leurs unités.
  \item Écrire la loi physique simple qui permet de commencer le raisonnement.
  \item Expliquer les hypothèses simplificatrices du modèle Terminale.
\end{enumerate}

\subsubsection*{Partie B — Mise en équation et calcul}

\begin{enumerate}
  \item Définir les variables utiles.
  \item Écrire la relation mathématique principale.
  \item Vérifier l'homogénéité de la formule.
  \item Choisir un ordre de grandeur réaliste pour les données manquantes.
  \item Effectuer le calcul ou établir le raisonnement qualitatif demandé.
  \item Interpréter le résultat en une phrase claire.
\end{enumerate}

\subsubsection*{Partie C — Réponse accessible niveau Terminale}

\begin{corrigebox}[Réponse Terminale]
On part de l'observation concrète : La zone entre deux arcs-en-ciel semble sombre.

Au niveau Terminale, on cherche d'abord la grandeur qui contrôle le phénomène. Ici, l'idée centrale est donnée par :
\[
\theta_1\simeq42^\circ,\qquad \theta_2\simeq51^\circ
\]
Cette relation permet de relier la cause physique à l'effet observé. On vérifie les unités, puis on insère des ordres de grandeur réalistes.

La conclusion attendue est la suivante : Expliquer l'absence relative de rayons dans certaines directions. La réponse doit toujours préciser les hypothèses utilisées. Si le phénomène dépend d'un milieu réel, d'une géométrie complexe, de frottements, de pertes ou de gradients, le résultat obtenu est un ordre de grandeur et non une valeur exacte.
\end{corrigebox}

\subsubsection*{Partie D — Approfondissement niveau Bac+3}

\begin{approfondissementbox}[Réponse Bac+3]
Au niveau Bac+3, on ne se contente plus d'une formule globale. On remplace le modèle simple par une loi locale, différentielle ou énergétique.

Une forme générale possible est :
\[
\frac{\mathrm d X}{\mathrm dt}=\text{apport}-\text{pertes},
\]
ou, pour un phénomène propagatif :
\[
\frac{\partial^2 \psi}{\partial t^2}=c^2\Delta\psi.
\]
Selon le cas, on peut aussi utiliser un principe variationnel, par exemple le principe de Fermat en optique :
\[
\delta \int n\,\mathrm ds=0.
\]

Pour ce phénomène, le modèle Bac+3 consiste à partir de la relation centrale
\[
\theta_1\simeq42^\circ,\qquad \theta_2\simeq51^\circ
\]
puis à introduire les corrections négligées : dépendance spatiale, dépendance temporelle, non-linéarités, pertes, conditions aux limites ou fluctuations. Cela permet de passer d'une réponse d'ordre de grandeur à une prédiction plus robuste.
\end{approfondissementbox}

\subsubsection*{Partie E — Limites, critique et prolongement}

\begin{enumerate}
  \item Citer au moins deux limites du modèle Terminale.
  \item Dire quelles grandeurs seraient difficiles à mesurer expérimentalement.
  \item Proposer une expérience simple permettant de tester le modèle.
  \item Expliquer ce qui changerait si l'on doublait une grandeur clé.
  \item Rédiger une conclusion courte, accessible à un élève de Terminale.
\end{enumerate}

\begin{corrigebox}[Synthèse rédigée]
Le phénomène « Bande sombre d'Alexandre » s'explique par un petit nombre de lois physiques simples, mais sa manifestation réelle dépend souvent de nombreux paramètres. Le modèle Terminale donne l'intuition et les ordres de grandeur ; le modèle Bac+3 introduit les variations locales, les pertes, les équations différentielles ou les principes variationnels nécessaires pour décrire finement la réalité.
\end{corrigebox}

% ------------------------------------------------------------
\subsection{Problème corrigé 8 — Halo solaire}
% ------------------------------------------------------------

\begin{exercicebox}[Question-problème]
\textbf{Observation.} Un cercle lumineux apparaît autour du Soleil.

\medskip
\textbf{Question.} Relier le halo à la géométrie des cristaux.

\medskip
\textbf{Thèmes mobilisés :} Réfraction par cristaux de glace.

\medskip
\textbf{Relation ou idée centrale :}
\[
\text{déviation minimale dans un prisme}
\]
\end{exercicebox}

\subsubsection*{Partie A — Modélisation niveau Terminale}

\begin{enumerate}
  \item Décrire précisément le phénomène observé.
  \item Identifier le système physique étudié.
  \item Lister les grandeurs importantes et leurs unités.
  \item Écrire la loi physique simple qui permet de commencer le raisonnement.
  \item Expliquer les hypothèses simplificatrices du modèle Terminale.
\end{enumerate}

\subsubsection*{Partie B — Mise en équation et calcul}

\begin{enumerate}
  \item Définir les variables utiles.
  \item Écrire la relation mathématique principale.
  \item Vérifier l'homogénéité de la formule.
  \item Choisir un ordre de grandeur réaliste pour les données manquantes.
  \item Effectuer le calcul ou établir le raisonnement qualitatif demandé.
  \item Interpréter le résultat en une phrase claire.
\end{enumerate}

\subsubsection*{Partie C — Réponse accessible niveau Terminale}

\begin{corrigebox}[Réponse Terminale]
On part de l'observation concrète : Un cercle lumineux apparaît autour du Soleil.

Au niveau Terminale, on cherche d'abord la grandeur qui contrôle le phénomène. Ici, l'idée centrale est donnée par :
\[
\text{déviation minimale dans un prisme}
\]
Cette relation permet de relier la cause physique à l'effet observé. On vérifie les unités, puis on insère des ordres de grandeur réalistes.

La conclusion attendue est la suivante : Relier le halo à la géométrie des cristaux. La réponse doit toujours préciser les hypothèses utilisées. Si le phénomène dépend d'un milieu réel, d'une géométrie complexe, de frottements, de pertes ou de gradients, le résultat obtenu est un ordre de grandeur et non une valeur exacte.
\end{corrigebox}

\subsubsection*{Partie D — Approfondissement niveau Bac+3}

\begin{approfondissementbox}[Réponse Bac+3]
Au niveau Bac+3, on ne se contente plus d'une formule globale. On remplace le modèle simple par une loi locale, différentielle ou énergétique.

Une forme générale possible est :
\[
\frac{\mathrm d X}{\mathrm dt}=\text{apport}-\text{pertes},
\]
ou, pour un phénomène propagatif :
\[
\frac{\partial^2 \psi}{\partial t^2}=c^2\Delta\psi.
\]
Selon le cas, on peut aussi utiliser un principe variationnel, par exemple le principe de Fermat en optique :
\[
\delta \int n\,\mathrm ds=0.
\]

Pour ce phénomène, le modèle Bac+3 consiste à partir de la relation centrale
\[
\text{déviation minimale dans un prisme}
\]
puis à introduire les corrections négligées : dépendance spatiale, dépendance temporelle, non-linéarités, pertes, conditions aux limites ou fluctuations. Cela permet de passer d'une réponse d'ordre de grandeur à une prédiction plus robuste.
\end{approfondissementbox}

\subsubsection*{Partie E — Limites, critique et prolongement}

\begin{enumerate}
  \item Citer au moins deux limites du modèle Terminale.
  \item Dire quelles grandeurs seraient difficiles à mesurer expérimentalement.
  \item Proposer une expérience simple permettant de tester le modèle.
  \item Expliquer ce qui changerait si l'on doublait une grandeur clé.
  \item Rédiger une conclusion courte, accessible à un élève de Terminale.
\end{enumerate}

\begin{corrigebox}[Synthèse rédigée]
Le phénomène « Halo solaire » s'explique par un petit nombre de lois physiques simples, mais sa manifestation réelle dépend souvent de nombreux paramètres. Le modèle Terminale donne l'intuition et les ordres de grandeur ; le modèle Bac+3 introduit les variations locales, les pertes, les équations différentielles ou les principes variationnels nécessaires pour décrire finement la réalité.
\end{corrigebox}

% ------------------------------------------------------------
\subsection{Problème corrigé 9 — Parhélies}
% ------------------------------------------------------------

\begin{exercicebox}[Question-problème]
\textbf{Observation.} Deux taches brillantes apparaissent de part et d'autre du Soleil.

\medskip
\textbf{Question.} Expliquer pourquoi les parhélies ne forment pas toujours un cercle complet.

\medskip
\textbf{Thèmes mobilisés :} Cristaux orientés, réfraction sélective.

\medskip
\textbf{Relation ou idée centrale :}
\[
\text{directions privilégiées}
\]
\end{exercicebox}

\subsubsection*{Partie A — Modélisation niveau Terminale}

\begin{enumerate}
  \item Décrire précisément le phénomène observé.
  \item Identifier le système physique étudié.
  \item Lister les grandeurs importantes et leurs unités.
  \item Écrire la loi physique simple qui permet de commencer le raisonnement.
  \item Expliquer les hypothèses simplificatrices du modèle Terminale.
\end{enumerate}

\subsubsection*{Partie B — Mise en équation et calcul}

\begin{enumerate}
  \item Définir les variables utiles.
  \item Écrire la relation mathématique principale.
  \item Vérifier l'homogénéité de la formule.
  \item Choisir un ordre de grandeur réaliste pour les données manquantes.
  \item Effectuer le calcul ou établir le raisonnement qualitatif demandé.
  \item Interpréter le résultat en une phrase claire.
\end{enumerate}

\subsubsection*{Partie C — Réponse accessible niveau Terminale}

\begin{corrigebox}[Réponse Terminale]
On part de l'observation concrète : Deux taches brillantes apparaissent de part et d'autre du Soleil.

Au niveau Terminale, on cherche d'abord la grandeur qui contrôle le phénomène. Ici, l'idée centrale est donnée par :
\[
\text{directions privilégiées}
\]
Cette relation permet de relier la cause physique à l'effet observé. On vérifie les unités, puis on insère des ordres de grandeur réalistes.

La conclusion attendue est la suivante : Expliquer pourquoi les parhélies ne forment pas toujours un cercle complet. La réponse doit toujours préciser les hypothèses utilisées. Si le phénomène dépend d'un milieu réel, d'une géométrie complexe, de frottements, de pertes ou de gradients, le résultat obtenu est un ordre de grandeur et non une valeur exacte.
\end{corrigebox}

\subsubsection*{Partie D — Approfondissement niveau Bac+3}

\begin{approfondissementbox}[Réponse Bac+3]
Au niveau Bac+3, on ne se contente plus d'une formule globale. On remplace le modèle simple par une loi locale, différentielle ou énergétique.

Une forme générale possible est :
\[
\frac{\mathrm d X}{\mathrm dt}=\text{apport}-\text{pertes},
\]
ou, pour un phénomène propagatif :
\[
\frac{\partial^2 \psi}{\partial t^2}=c^2\Delta\psi.
\]
Selon le cas, on peut aussi utiliser un principe variationnel, par exemple le principe de Fermat en optique :
\[
\delta \int n\,\mathrm ds=0.
\]

Pour ce phénomène, le modèle Bac+3 consiste à partir de la relation centrale
\[
\text{directions privilégiées}
\]
puis à introduire les corrections négligées : dépendance spatiale, dépendance temporelle, non-linéarités, pertes, conditions aux limites ou fluctuations. Cela permet de passer d'une réponse d'ordre de grandeur à une prédiction plus robuste.
\end{approfondissementbox}

\subsubsection*{Partie E — Limites, critique et prolongement}

\begin{enumerate}
  \item Citer au moins deux limites du modèle Terminale.
  \item Dire quelles grandeurs seraient difficiles à mesurer expérimentalement.
  \item Proposer une expérience simple permettant de tester le modèle.
  \item Expliquer ce qui changerait si l'on doublait une grandeur clé.
  \item Rédiger une conclusion courte, accessible à un élève de Terminale.
\end{enumerate}

\begin{corrigebox}[Synthèse rédigée]
Le phénomène « Parhélies » s'explique par un petit nombre de lois physiques simples, mais sa manifestation réelle dépend souvent de nombreux paramètres. Le modèle Terminale donne l'intuition et les ordres de grandeur ; le modèle Bac+3 introduit les variations locales, les pertes, les équations différentielles ou les principes variationnels nécessaires pour décrire finement la réalité.
\end{corrigebox}

% ------------------------------------------------------------
\subsection{Problème corrigé 10 — Ciel bleu}
% ------------------------------------------------------------

\begin{exercicebox}[Question-problème]
\textbf{Observation.} Le ciel est bleu loin du Soleil.

\medskip
\textbf{Question.} Comparer la diffusion du bleu et du rouge.

\medskip
\textbf{Thèmes mobilisés :} Diffusion Rayleigh.

\medskip
\textbf{Relation ou idée centrale :}
\[
I\propto\frac1{\lambda^4}
\]
\end{exercicebox}

\subsubsection*{Partie A — Modélisation niveau Terminale}

\begin{enumerate}
  \item Décrire précisément le phénomène observé.
  \item Identifier le système physique étudié.
  \item Lister les grandeurs importantes et leurs unités.
  \item Écrire la loi physique simple qui permet de commencer le raisonnement.
  \item Expliquer les hypothèses simplificatrices du modèle Terminale.
\end{enumerate}

\subsubsection*{Partie B — Mise en équation et calcul}

\begin{enumerate}
  \item Définir les variables utiles.
  \item Écrire la relation mathématique principale.
  \item Vérifier l'homogénéité de la formule.
  \item Choisir un ordre de grandeur réaliste pour les données manquantes.
  \item Effectuer le calcul ou établir le raisonnement qualitatif demandé.
  \item Interpréter le résultat en une phrase claire.
\end{enumerate}

\subsubsection*{Partie C — Réponse accessible niveau Terminale}

\begin{corrigebox}[Réponse Terminale]
On part de l'observation concrète : Le ciel est bleu loin du Soleil.

Au niveau Terminale, on cherche d'abord la grandeur qui contrôle le phénomène. Ici, l'idée centrale est donnée par :
\[
I\propto\frac1{\lambda^4}
\]
Cette relation permet de relier la cause physique à l'effet observé. On vérifie les unités, puis on insère des ordres de grandeur réalistes.

La conclusion attendue est la suivante : Comparer la diffusion du bleu et du rouge. La réponse doit toujours préciser les hypothèses utilisées. Si le phénomène dépend d'un milieu réel, d'une géométrie complexe, de frottements, de pertes ou de gradients, le résultat obtenu est un ordre de grandeur et non une valeur exacte.
\end{corrigebox}

\subsubsection*{Partie D — Approfondissement niveau Bac+3}

\begin{approfondissementbox}[Réponse Bac+3]
Au niveau Bac+3, on ne se contente plus d'une formule globale. On remplace le modèle simple par une loi locale, différentielle ou énergétique.

Une forme générale possible est :
\[
\frac{\mathrm d X}{\mathrm dt}=\text{apport}-\text{pertes},
\]
ou, pour un phénomène propagatif :
\[
\frac{\partial^2 \psi}{\partial t^2}=c^2\Delta\psi.
\]
Selon le cas, on peut aussi utiliser un principe variationnel, par exemple le principe de Fermat en optique :
\[
\delta \int n\,\mathrm ds=0.
\]

Pour ce phénomène, le modèle Bac+3 consiste à partir de la relation centrale
\[
I\propto\frac1{\lambda^4}
\]
puis à introduire les corrections négligées : dépendance spatiale, dépendance temporelle, non-linéarités, pertes, conditions aux limites ou fluctuations. Cela permet de passer d'une réponse d'ordre de grandeur à une prédiction plus robuste.
\end{approfondissementbox}

\subsubsection*{Partie E — Limites, critique et prolongement}

\begin{enumerate}
  \item Citer au moins deux limites du modèle Terminale.
  \item Dire quelles grandeurs seraient difficiles à mesurer expérimentalement.
  \item Proposer une expérience simple permettant de tester le modèle.
  \item Expliquer ce qui changerait si l'on doublait une grandeur clé.
  \item Rédiger une conclusion courte, accessible à un élève de Terminale.
\end{enumerate}

\begin{corrigebox}[Synthèse rédigée]
Le phénomène « Ciel bleu » s'explique par un petit nombre de lois physiques simples, mais sa manifestation réelle dépend souvent de nombreux paramètres. Le modèle Terminale donne l'intuition et les ordres de grandeur ; le modèle Bac+3 introduit les variations locales, les pertes, les équations différentielles ou les principes variationnels nécessaires pour décrire finement la réalité.
\end{corrigebox}

% ------------------------------------------------------------
\subsection{Problème corrigé 11 — Coucher de soleil rouge}
% ------------------------------------------------------------

\begin{exercicebox}[Question-problème]
\textbf{Observation.} Le Soleil devient rouge près de l'horizon.

\medskip
\textbf{Question.} Expliquer pourquoi le rouge domine au coucher.

\medskip
\textbf{Thèmes mobilisés :} Diffusion et transmission atmosphérique.

\medskip
\textbf{Relation ou idée centrale :}
\[
I=I_0e^{-\alpha(\lambda)L}
\]
\end{exercicebox}

\subsubsection*{Partie A — Modélisation niveau Terminale}

\begin{enumerate}
  \item Décrire précisément le phénomène observé.
  \item Identifier le système physique étudié.
  \item Lister les grandeurs importantes et leurs unités.
  \item Écrire la loi physique simple qui permet de commencer le raisonnement.
  \item Expliquer les hypothèses simplificatrices du modèle Terminale.
\end{enumerate}

\subsubsection*{Partie B — Mise en équation et calcul}

\begin{enumerate}
  \item Définir les variables utiles.
  \item Écrire la relation mathématique principale.
  \item Vérifier l'homogénéité de la formule.
  \item Choisir un ordre de grandeur réaliste pour les données manquantes.
  \item Effectuer le calcul ou établir le raisonnement qualitatif demandé.
  \item Interpréter le résultat en une phrase claire.
\end{enumerate}

\subsubsection*{Partie C — Réponse accessible niveau Terminale}

\begin{corrigebox}[Réponse Terminale]
On part de l'observation concrète : Le Soleil devient rouge près de l'horizon.

Au niveau Terminale, on cherche d'abord la grandeur qui contrôle le phénomène. Ici, l'idée centrale est donnée par :
\[
I=I_0e^{-\alpha(\lambda)L}
\]
Cette relation permet de relier la cause physique à l'effet observé. On vérifie les unités, puis on insère des ordres de grandeur réalistes.

La conclusion attendue est la suivante : Expliquer pourquoi le rouge domine au coucher. La réponse doit toujours préciser les hypothèses utilisées. Si le phénomène dépend d'un milieu réel, d'une géométrie complexe, de frottements, de pertes ou de gradients, le résultat obtenu est un ordre de grandeur et non une valeur exacte.
\end{corrigebox}

\subsubsection*{Partie D — Approfondissement niveau Bac+3}

\begin{approfondissementbox}[Réponse Bac+3]
Au niveau Bac+3, on ne se contente plus d'une formule globale. On remplace le modèle simple par une loi locale, différentielle ou énergétique.

Une forme générale possible est :
\[
\frac{\mathrm d X}{\mathrm dt}=\text{apport}-\text{pertes},
\]
ou, pour un phénomène propagatif :
\[
\frac{\partial^2 \psi}{\partial t^2}=c^2\Delta\psi.
\]
Selon le cas, on peut aussi utiliser un principe variationnel, par exemple le principe de Fermat en optique :
\[
\delta \int n\,\mathrm ds=0.
\]

Pour ce phénomène, le modèle Bac+3 consiste à partir de la relation centrale
\[
I=I_0e^{-\alpha(\lambda)L}
\]
puis à introduire les corrections négligées : dépendance spatiale, dépendance temporelle, non-linéarités, pertes, conditions aux limites ou fluctuations. Cela permet de passer d'une réponse d'ordre de grandeur à une prédiction plus robuste.
\end{approfondissementbox}

\subsubsection*{Partie E — Limites, critique et prolongement}

\begin{enumerate}
  \item Citer au moins deux limites du modèle Terminale.
  \item Dire quelles grandeurs seraient difficiles à mesurer expérimentalement.
  \item Proposer une expérience simple permettant de tester le modèle.
  \item Expliquer ce qui changerait si l'on doublait une grandeur clé.
  \item Rédiger une conclusion courte, accessible à un élève de Terminale.
\end{enumerate}

\begin{corrigebox}[Synthèse rédigée]
Le phénomène « Coucher de soleil rouge » s'explique par un petit nombre de lois physiques simples, mais sa manifestation réelle dépend souvent de nombreux paramètres. Le modèle Terminale donne l'intuition et les ordres de grandeur ; le modèle Bac+3 introduit les variations locales, les pertes, les équations différentielles ou les principes variationnels nécessaires pour décrire finement la réalité.
\end{corrigebox}

% ------------------------------------------------------------
\subsection{Problème corrigé 12 — Nuages blancs}
% ------------------------------------------------------------

\begin{exercicebox}[Question-problème]
\textbf{Observation.} Les nuages apparaissent blancs ou gris.

\medskip
\textbf{Question.} Comprendre pourquoi toutes les couleurs sont diffusées.

\medskip
\textbf{Thèmes mobilisés :} Diffusion de Mie.

\medskip
\textbf{Relation ou idée centrale :}
\[
a\gg\lambda
\]
\end{exercicebox}

\subsubsection*{Partie A — Modélisation niveau Terminale}

\begin{enumerate}
  \item Décrire précisément le phénomène observé.
  \item Identifier le système physique étudié.
  \item Lister les grandeurs importantes et leurs unités.
  \item Écrire la loi physique simple qui permet de commencer le raisonnement.
  \item Expliquer les hypothèses simplificatrices du modèle Terminale.
\end{enumerate}

\subsubsection*{Partie B — Mise en équation et calcul}

\begin{enumerate}
  \item Définir les variables utiles.
  \item Écrire la relation mathématique principale.
  \item Vérifier l'homogénéité de la formule.
  \item Choisir un ordre de grandeur réaliste pour les données manquantes.
  \item Effectuer le calcul ou établir le raisonnement qualitatif demandé.
  \item Interpréter le résultat en une phrase claire.
\end{enumerate}

\subsubsection*{Partie C — Réponse accessible niveau Terminale}

\begin{corrigebox}[Réponse Terminale]
On part de l'observation concrète : Les nuages apparaissent blancs ou gris.

Au niveau Terminale, on cherche d'abord la grandeur qui contrôle le phénomène. Ici, l'idée centrale est donnée par :
\[
a\gg\lambda
\]
Cette relation permet de relier la cause physique à l'effet observé. On vérifie les unités, puis on insère des ordres de grandeur réalistes.

La conclusion attendue est la suivante : Comprendre pourquoi toutes les couleurs sont diffusées. La réponse doit toujours préciser les hypothèses utilisées. Si le phénomène dépend d'un milieu réel, d'une géométrie complexe, de frottements, de pertes ou de gradients, le résultat obtenu est un ordre de grandeur et non une valeur exacte.
\end{corrigebox}

\subsubsection*{Partie D — Approfondissement niveau Bac+3}

\begin{approfondissementbox}[Réponse Bac+3]
Au niveau Bac+3, on ne se contente plus d'une formule globale. On remplace le modèle simple par une loi locale, différentielle ou énergétique.

Une forme générale possible est :
\[
\frac{\mathrm d X}{\mathrm dt}=\text{apport}-\text{pertes},
\]
ou, pour un phénomène propagatif :
\[
\frac{\partial^2 \psi}{\partial t^2}=c^2\Delta\psi.
\]
Selon le cas, on peut aussi utiliser un principe variationnel, par exemple le principe de Fermat en optique :
\[
\delta \int n\,\mathrm ds=0.
\]

Pour ce phénomène, le modèle Bac+3 consiste à partir de la relation centrale
\[
a\gg\lambda
\]
puis à introduire les corrections négligées : dépendance spatiale, dépendance temporelle, non-linéarités, pertes, conditions aux limites ou fluctuations. Cela permet de passer d'une réponse d'ordre de grandeur à une prédiction plus robuste.
\end{approfondissementbox}

\subsubsection*{Partie E — Limites, critique et prolongement}

\begin{enumerate}
  \item Citer au moins deux limites du modèle Terminale.
  \item Dire quelles grandeurs seraient difficiles à mesurer expérimentalement.
  \item Proposer une expérience simple permettant de tester le modèle.
  \item Expliquer ce qui changerait si l'on doublait une grandeur clé.
  \item Rédiger une conclusion courte, accessible à un élève de Terminale.
\end{enumerate}

\begin{corrigebox}[Synthèse rédigée]
Le phénomène « Nuages blancs » s'explique par un petit nombre de lois physiques simples, mais sa manifestation réelle dépend souvent de nombreux paramètres. Le modèle Terminale donne l'intuition et les ordres de grandeur ; le modèle Bac+3 introduit les variations locales, les pertes, les équations différentielles ou les principes variationnels nécessaires pour décrire finement la réalité.
\end{corrigebox}

% ------------------------------------------------------------
\subsection{Problème corrigé 13 — Brouillard}
% ------------------------------------------------------------

\begin{exercicebox}[Question-problème]
\textbf{Observation.} Le brouillard réduit fortement la visibilité.

\medskip
\textbf{Question.} Relier visibilité et atténuation exponentielle.

\medskip
\textbf{Thèmes mobilisés :} Diffusion multiple, atténuation.

\medskip
\textbf{Relation ou idée centrale :}
\[
I(L)=I_0e^{-\beta L}
\]
\end{exercicebox}

\subsubsection*{Partie A — Modélisation niveau Terminale}

\begin{enumerate}
  \item Décrire précisément le phénomène observé.
  \item Identifier le système physique étudié.
  \item Lister les grandeurs importantes et leurs unités.
  \item Écrire la loi physique simple qui permet de commencer le raisonnement.
  \item Expliquer les hypothèses simplificatrices du modèle Terminale.
\end{enumerate}

\subsubsection*{Partie B — Mise en équation et calcul}

\begin{enumerate}
  \item Définir les variables utiles.
  \item Écrire la relation mathématique principale.
  \item Vérifier l'homogénéité de la formule.
  \item Choisir un ordre de grandeur réaliste pour les données manquantes.
  \item Effectuer le calcul ou établir le raisonnement qualitatif demandé.
  \item Interpréter le résultat en une phrase claire.
\end{enumerate}

\subsubsection*{Partie C — Réponse accessible niveau Terminale}

\begin{corrigebox}[Réponse Terminale]
On part de l'observation concrète : Le brouillard réduit fortement la visibilité.

Au niveau Terminale, on cherche d'abord la grandeur qui contrôle le phénomène. Ici, l'idée centrale est donnée par :
\[
I(L)=I_0e^{-\beta L}
\]
Cette relation permet de relier la cause physique à l'effet observé. On vérifie les unités, puis on insère des ordres de grandeur réalistes.

La conclusion attendue est la suivante : Relier visibilité et atténuation exponentielle. La réponse doit toujours préciser les hypothèses utilisées. Si le phénomène dépend d'un milieu réel, d'une géométrie complexe, de frottements, de pertes ou de gradients, le résultat obtenu est un ordre de grandeur et non une valeur exacte.
\end{corrigebox}

\subsubsection*{Partie D — Approfondissement niveau Bac+3}

\begin{approfondissementbox}[Réponse Bac+3]
Au niveau Bac+3, on ne se contente plus d'une formule globale. On remplace le modèle simple par une loi locale, différentielle ou énergétique.

Une forme générale possible est :
\[
\frac{\mathrm d X}{\mathrm dt}=\text{apport}-\text{pertes},
\]
ou, pour un phénomène propagatif :
\[
\frac{\partial^2 \psi}{\partial t^2}=c^2\Delta\psi.
\]
Selon le cas, on peut aussi utiliser un principe variationnel, par exemple le principe de Fermat en optique :
\[
\delta \int n\,\mathrm ds=0.
\]

Pour ce phénomène, le modèle Bac+3 consiste à partir de la relation centrale
\[
I(L)=I_0e^{-\beta L}
\]
puis à introduire les corrections négligées : dépendance spatiale, dépendance temporelle, non-linéarités, pertes, conditions aux limites ou fluctuations. Cela permet de passer d'une réponse d'ordre de grandeur à une prédiction plus robuste.
\end{approfondissementbox}

\subsubsection*{Partie E — Limites, critique et prolongement}

\begin{enumerate}
  \item Citer au moins deux limites du modèle Terminale.
  \item Dire quelles grandeurs seraient difficiles à mesurer expérimentalement.
  \item Proposer une expérience simple permettant de tester le modèle.
  \item Expliquer ce qui changerait si l'on doublait une grandeur clé.
  \item Rédiger une conclusion courte, accessible à un élève de Terminale.
\end{enumerate}

\begin{corrigebox}[Synthèse rédigée]
Le phénomène « Brouillard » s'explique par un petit nombre de lois physiques simples, mais sa manifestation réelle dépend souvent de nombreux paramètres. Le modèle Terminale donne l'intuition et les ordres de grandeur ; le modèle Bac+3 introduit les variations locales, les pertes, les équations différentielles ou les principes variationnels nécessaires pour décrire finement la réalité.
\end{corrigebox}

% ------------------------------------------------------------
\subsection{Problème corrigé 14 — Diffusion atmosphérique}
% ------------------------------------------------------------

\begin{exercicebox}[Question-problème]
\textbf{Observation.} La lumière solaire est redistribuée dans l'atmosphère.

\medskip
\textbf{Question.} Classer les régimes de diffusion selon la taille des particules.

\medskip
\textbf{Thèmes mobilisés :} Rayleigh, Mie, aérosols.

\medskip
\textbf{Relation ou idée centrale :}
\[
I_{\mathrm{diff}}=I_R+I_M
\]
\end{exercicebox}

\subsubsection*{Partie A — Modélisation niveau Terminale}

\begin{enumerate}
  \item Décrire précisément le phénomène observé.
  \item Identifier le système physique étudié.
  \item Lister les grandeurs importantes et leurs unités.
  \item Écrire la loi physique simple qui permet de commencer le raisonnement.
  \item Expliquer les hypothèses simplificatrices du modèle Terminale.
\end{enumerate}

\subsubsection*{Partie B — Mise en équation et calcul}

\begin{enumerate}
  \item Définir les variables utiles.
  \item Écrire la relation mathématique principale.
  \item Vérifier l'homogénéité de la formule.
  \item Choisir un ordre de grandeur réaliste pour les données manquantes.
  \item Effectuer le calcul ou établir le raisonnement qualitatif demandé.
  \item Interpréter le résultat en une phrase claire.
\end{enumerate}

\subsubsection*{Partie C — Réponse accessible niveau Terminale}

\begin{corrigebox}[Réponse Terminale]
On part de l'observation concrète : La lumière solaire est redistribuée dans l'atmosphère.

Au niveau Terminale, on cherche d'abord la grandeur qui contrôle le phénomène. Ici, l'idée centrale est donnée par :
\[
I_{\mathrm{diff}}=I_R+I_M
\]
Cette relation permet de relier la cause physique à l'effet observé. On vérifie les unités, puis on insère des ordres de grandeur réalistes.

La conclusion attendue est la suivante : Classer les régimes de diffusion selon la taille des particules. La réponse doit toujours préciser les hypothèses utilisées. Si le phénomène dépend d'un milieu réel, d'une géométrie complexe, de frottements, de pertes ou de gradients, le résultat obtenu est un ordre de grandeur et non une valeur exacte.
\end{corrigebox}

\subsubsection*{Partie D — Approfondissement niveau Bac+3}

\begin{approfondissementbox}[Réponse Bac+3]
Au niveau Bac+3, on ne se contente plus d'une formule globale. On remplace le modèle simple par une loi locale, différentielle ou énergétique.

Une forme générale possible est :
\[
\frac{\mathrm d X}{\mathrm dt}=\text{apport}-\text{pertes},
\]
ou, pour un phénomène propagatif :
\[
\frac{\partial^2 \psi}{\partial t^2}=c^2\Delta\psi.
\]
Selon le cas, on peut aussi utiliser un principe variationnel, par exemple le principe de Fermat en optique :
\[
\delta \int n\,\mathrm ds=0.
\]

Pour ce phénomène, le modèle Bac+3 consiste à partir de la relation centrale
\[
I_{\mathrm{diff}}=I_R+I_M
\]
puis à introduire les corrections négligées : dépendance spatiale, dépendance temporelle, non-linéarités, pertes, conditions aux limites ou fluctuations. Cela permet de passer d'une réponse d'ordre de grandeur à une prédiction plus robuste.
\end{approfondissementbox}

\subsubsection*{Partie E — Limites, critique et prolongement}

\begin{enumerate}
  \item Citer au moins deux limites du modèle Terminale.
  \item Dire quelles grandeurs seraient difficiles à mesurer expérimentalement.
  \item Proposer une expérience simple permettant de tester le modèle.
  \item Expliquer ce qui changerait si l'on doublait une grandeur clé.
  \item Rédiger une conclusion courte, accessible à un élève de Terminale.
\end{enumerate}

\begin{corrigebox}[Synthèse rédigée]
Le phénomène « Diffusion atmosphérique » s'explique par un petit nombre de lois physiques simples, mais sa manifestation réelle dépend souvent de nombreux paramètres. Le modèle Terminale donne l'intuition et les ordres de grandeur ; le modèle Bac+3 introduit les variations locales, les pertes, les équations différentielles ou les principes variationnels nécessaires pour décrire finement la réalité.
\end{corrigebox}

% ------------------------------------------------------------
\subsection{Problème corrigé 15 — Éclair}
% ------------------------------------------------------------

\begin{exercicebox}[Question-problème]
\textbf{Observation.} Une décharge lumineuse traverse l'air.

\medskip
\textbf{Question.} Estimer l'ordre de grandeur d'une tension de claquage.

\medskip
\textbf{Thèmes mobilisés :} Champ électrique, ionisation, plasma.

\medskip
\textbf{Relation ou idée centrale :}
\[
E_{\mathrm{rupture}}\sim 3\times10^6\,\mathrm{V.m^{-1}}
\]
\end{exercicebox}

\subsubsection*{Partie A — Modélisation niveau Terminale}

\begin{enumerate}
  \item Décrire précisément le phénomène observé.
  \item Identifier le système physique étudié.
  \item Lister les grandeurs importantes et leurs unités.
  \item Écrire la loi physique simple qui permet de commencer le raisonnement.
  \item Expliquer les hypothèses simplificatrices du modèle Terminale.
\end{enumerate}

\subsubsection*{Partie B — Mise en équation et calcul}

\begin{enumerate}
  \item Définir les variables utiles.
  \item Écrire la relation mathématique principale.
  \item Vérifier l'homogénéité de la formule.
  \item Choisir un ordre de grandeur réaliste pour les données manquantes.
  \item Effectuer le calcul ou établir le raisonnement qualitatif demandé.
  \item Interpréter le résultat en une phrase claire.
\end{enumerate}

\subsubsection*{Partie C — Réponse accessible niveau Terminale}

\begin{corrigebox}[Réponse Terminale]
On part de l'observation concrète : Une décharge lumineuse traverse l'air.

Au niveau Terminale, on cherche d'abord la grandeur qui contrôle le phénomène. Ici, l'idée centrale est donnée par :
\[
E_{\mathrm{rupture}}\sim 3\times10^6\,\mathrm{V.m^{-1}}
\]
Cette relation permet de relier la cause physique à l'effet observé. On vérifie les unités, puis on insère des ordres de grandeur réalistes.

La conclusion attendue est la suivante : Estimer l'ordre de grandeur d'une tension de claquage. La réponse doit toujours préciser les hypothèses utilisées. Si le phénomène dépend d'un milieu réel, d'une géométrie complexe, de frottements, de pertes ou de gradients, le résultat obtenu est un ordre de grandeur et non une valeur exacte.
\end{corrigebox}

\subsubsection*{Partie D — Approfondissement niveau Bac+3}

\begin{approfondissementbox}[Réponse Bac+3]
Au niveau Bac+3, on ne se contente plus d'une formule globale. On remplace le modèle simple par une loi locale, différentielle ou énergétique.

Une forme générale possible est :
\[
\frac{\mathrm d X}{\mathrm dt}=\text{apport}-\text{pertes},
\]
ou, pour un phénomène propagatif :
\[
\frac{\partial^2 \psi}{\partial t^2}=c^2\Delta\psi.
\]
Selon le cas, on peut aussi utiliser un principe variationnel, par exemple le principe de Fermat en optique :
\[
\delta \int n\,\mathrm ds=0.
\]

Pour ce phénomène, le modèle Bac+3 consiste à partir de la relation centrale
\[
E_{\mathrm{rupture}}\sim 3\times10^6\,\mathrm{V.m^{-1}}
\]
puis à introduire les corrections négligées : dépendance spatiale, dépendance temporelle, non-linéarités, pertes, conditions aux limites ou fluctuations. Cela permet de passer d'une réponse d'ordre de grandeur à une prédiction plus robuste.
\end{approfondissementbox}

\subsubsection*{Partie E — Limites, critique et prolongement}

\begin{enumerate}
  \item Citer au moins deux limites du modèle Terminale.
  \item Dire quelles grandeurs seraient difficiles à mesurer expérimentalement.
  \item Proposer une expérience simple permettant de tester le modèle.
  \item Expliquer ce qui changerait si l'on doublait une grandeur clé.
  \item Rédiger une conclusion courte, accessible à un élève de Terminale.
\end{enumerate}

\begin{corrigebox}[Synthèse rédigée]
Le phénomène « Éclair » s'explique par un petit nombre de lois physiques simples, mais sa manifestation réelle dépend souvent de nombreux paramètres. Le modèle Terminale donne l'intuition et les ordres de grandeur ; le modèle Bac+3 introduit les variations locales, les pertes, les équations différentielles ou les principes variationnels nécessaires pour décrire finement la réalité.
\end{corrigebox}

% ------------------------------------------------------------
\subsection{Problème corrigé 16 — Tonnerre}
% ------------------------------------------------------------

\begin{exercicebox}[Question-problème]
\textbf{Observation.} Un bruit violent suit l'éclair.

\medskip
\textbf{Question.} Relier délai éclair-tonnerre et distance.

\medskip
\textbf{Thèmes mobilisés :} Onde de choc, propagation sonore.

\medskip
\textbf{Relation ou idée centrale :}
\[
d\simeq c_s\Delta t
\]
\end{exercicebox}

\subsubsection*{Partie A — Modélisation niveau Terminale}

\begin{enumerate}
  \item Décrire précisément le phénomène observé.
  \item Identifier le système physique étudié.
  \item Lister les grandeurs importantes et leurs unités.
  \item Écrire la loi physique simple qui permet de commencer le raisonnement.
  \item Expliquer les hypothèses simplificatrices du modèle Terminale.
\end{enumerate}

\subsubsection*{Partie B — Mise en équation et calcul}

\begin{enumerate}
  \item Définir les variables utiles.
  \item Écrire la relation mathématique principale.
  \item Vérifier l'homogénéité de la formule.
  \item Choisir un ordre de grandeur réaliste pour les données manquantes.
  \item Effectuer le calcul ou établir le raisonnement qualitatif demandé.
  \item Interpréter le résultat en une phrase claire.
\end{enumerate}

\subsubsection*{Partie C — Réponse accessible niveau Terminale}

\begin{corrigebox}[Réponse Terminale]
On part de l'observation concrète : Un bruit violent suit l'éclair.

Au niveau Terminale, on cherche d'abord la grandeur qui contrôle le phénomène. Ici, l'idée centrale est donnée par :
\[
d\simeq c_s\Delta t
\]
Cette relation permet de relier la cause physique à l'effet observé. On vérifie les unités, puis on insère des ordres de grandeur réalistes.

La conclusion attendue est la suivante : Relier délai éclair-tonnerre et distance. La réponse doit toujours préciser les hypothèses utilisées. Si le phénomène dépend d'un milieu réel, d'une géométrie complexe, de frottements, de pertes ou de gradients, le résultat obtenu est un ordre de grandeur et non une valeur exacte.
\end{corrigebox}

\subsubsection*{Partie D — Approfondissement niveau Bac+3}

\begin{approfondissementbox}[Réponse Bac+3]
Au niveau Bac+3, on ne se contente plus d'une formule globale. On remplace le modèle simple par une loi locale, différentielle ou énergétique.

Une forme générale possible est :
\[
\frac{\mathrm d X}{\mathrm dt}=\text{apport}-\text{pertes},
\]
ou, pour un phénomène propagatif :
\[
\frac{\partial^2 \psi}{\partial t^2}=c^2\Delta\psi.
\]
Selon le cas, on peut aussi utiliser un principe variationnel, par exemple le principe de Fermat en optique :
\[
\delta \int n\,\mathrm ds=0.
\]

Pour ce phénomène, le modèle Bac+3 consiste à partir de la relation centrale
\[
d\simeq c_s\Delta t
\]
puis à introduire les corrections négligées : dépendance spatiale, dépendance temporelle, non-linéarités, pertes, conditions aux limites ou fluctuations. Cela permet de passer d'une réponse d'ordre de grandeur à une prédiction plus robuste.
\end{approfondissementbox}

\subsubsection*{Partie E — Limites, critique et prolongement}

\begin{enumerate}
  \item Citer au moins deux limites du modèle Terminale.
  \item Dire quelles grandeurs seraient difficiles à mesurer expérimentalement.
  \item Proposer une expérience simple permettant de tester le modèle.
  \item Expliquer ce qui changerait si l'on doublait une grandeur clé.
  \item Rédiger une conclusion courte, accessible à un élève de Terminale.
\end{enumerate}

\begin{corrigebox}[Synthèse rédigée]
Le phénomène « Tonnerre » s'explique par un petit nombre de lois physiques simples, mais sa manifestation réelle dépend souvent de nombreux paramètres. Le modèle Terminale donne l'intuition et les ordres de grandeur ; le modèle Bac+3 introduit les variations locales, les pertes, les équations différentielles ou les principes variationnels nécessaires pour décrire finement la réalité.
\end{corrigebox}

% ------------------------------------------------------------
\subsection{Problème corrigé 17 — Distance d'un orage}
% ------------------------------------------------------------

\begin{exercicebox}[Question-problème]
\textbf{Observation.} On compte les secondes entre l'éclair et le tonnerre.

\medskip
\textbf{Question.} Obtenir une méthode rapide d'estimation.

\medskip
\textbf{Thèmes mobilisés :} Propagation lumière/son.

\medskip
\textbf{Relation ou idée centrale :}
\[
d=c_s\Delta t
\]
\end{exercicebox}

\subsubsection*{Partie A — Modélisation niveau Terminale}

\begin{enumerate}
  \item Décrire précisément le phénomène observé.
  \item Identifier le système physique étudié.
  \item Lister les grandeurs importantes et leurs unités.
  \item Écrire la loi physique simple qui permet de commencer le raisonnement.
  \item Expliquer les hypothèses simplificatrices du modèle Terminale.
\end{enumerate}

\subsubsection*{Partie B — Mise en équation et calcul}

\begin{enumerate}
  \item Définir les variables utiles.
  \item Écrire la relation mathématique principale.
  \item Vérifier l'homogénéité de la formule.
  \item Choisir un ordre de grandeur réaliste pour les données manquantes.
  \item Effectuer le calcul ou établir le raisonnement qualitatif demandé.
  \item Interpréter le résultat en une phrase claire.
\end{enumerate}

\subsubsection*{Partie C — Réponse accessible niveau Terminale}

\begin{corrigebox}[Réponse Terminale]
On part de l'observation concrète : On compte les secondes entre l'éclair et le tonnerre.

Au niveau Terminale, on cherche d'abord la grandeur qui contrôle le phénomène. Ici, l'idée centrale est donnée par :
\[
d=c_s\Delta t
\]
Cette relation permet de relier la cause physique à l'effet observé. On vérifie les unités, puis on insère des ordres de grandeur réalistes.

La conclusion attendue est la suivante : Obtenir une méthode rapide d'estimation. La réponse doit toujours préciser les hypothèses utilisées. Si le phénomène dépend d'un milieu réel, d'une géométrie complexe, de frottements, de pertes ou de gradients, le résultat obtenu est un ordre de grandeur et non une valeur exacte.
\end{corrigebox}

\subsubsection*{Partie D — Approfondissement niveau Bac+3}

\begin{approfondissementbox}[Réponse Bac+3]
Au niveau Bac+3, on ne se contente plus d'une formule globale. On remplace le modèle simple par une loi locale, différentielle ou énergétique.

Une forme générale possible est :
\[
\frac{\mathrm d X}{\mathrm dt}=\text{apport}-\text{pertes},
\]
ou, pour un phénomène propagatif :
\[
\frac{\partial^2 \psi}{\partial t^2}=c^2\Delta\psi.
\]
Selon le cas, on peut aussi utiliser un principe variationnel, par exemple le principe de Fermat en optique :
\[
\delta \int n\,\mathrm ds=0.
\]

Pour ce phénomène, le modèle Bac+3 consiste à partir de la relation centrale
\[
d=c_s\Delta t
\]
puis à introduire les corrections négligées : dépendance spatiale, dépendance temporelle, non-linéarités, pertes, conditions aux limites ou fluctuations. Cela permet de passer d'une réponse d'ordre de grandeur à une prédiction plus robuste.
\end{approfondissementbox}

\subsubsection*{Partie E — Limites, critique et prolongement}

\begin{enumerate}
  \item Citer au moins deux limites du modèle Terminale.
  \item Dire quelles grandeurs seraient difficiles à mesurer expérimentalement.
  \item Proposer une expérience simple permettant de tester le modèle.
  \item Expliquer ce qui changerait si l'on doublait une grandeur clé.
  \item Rédiger une conclusion courte, accessible à un élève de Terminale.
\end{enumerate}

\begin{corrigebox}[Synthèse rédigée]
Le phénomène « Distance d'un orage » s'explique par un petit nombre de lois physiques simples, mais sa manifestation réelle dépend souvent de nombreux paramètres. Le modèle Terminale donne l'intuition et les ordres de grandeur ; le modèle Bac+3 introduit les variations locales, les pertes, les équations différentielles ou les principes variationnels nécessaires pour décrire finement la réalité.
\end{corrigebox}

% ------------------------------------------------------------
\subsection{Problème corrigé 18 — Givre}
% ------------------------------------------------------------

\begin{exercicebox}[Question-problème]
\textbf{Observation.} Une couche blanche apparaît sur les surfaces froides.

\medskip
\textbf{Question.} Expliquer pourquoi le givre apparaît par ciel clair.

\medskip
\textbf{Thèmes mobilisés :} Changement d'état, sublimation inverse.

\medskip
\textbf{Relation ou idée centrale :}
\[
T_{\mathrm{surface}}<0^\circ\mathrm C
\]
\end{exercicebox}

\subsubsection*{Partie A — Modélisation niveau Terminale}

\begin{enumerate}
  \item Décrire précisément le phénomène observé.
  \item Identifier le système physique étudié.
  \item Lister les grandeurs importantes et leurs unités.
  \item Écrire la loi physique simple qui permet de commencer le raisonnement.
  \item Expliquer les hypothèses simplificatrices du modèle Terminale.
\end{enumerate}

\subsubsection*{Partie B — Mise en équation et calcul}

\begin{enumerate}
  \item Définir les variables utiles.
  \item Écrire la relation mathématique principale.
  \item Vérifier l'homogénéité de la formule.
  \item Choisir un ordre de grandeur réaliste pour les données manquantes.
  \item Effectuer le calcul ou établir le raisonnement qualitatif demandé.
  \item Interpréter le résultat en une phrase claire.
\end{enumerate}

\subsubsection*{Partie C — Réponse accessible niveau Terminale}

\begin{corrigebox}[Réponse Terminale]
On part de l'observation concrète : Une couche blanche apparaît sur les surfaces froides.

Au niveau Terminale, on cherche d'abord la grandeur qui contrôle le phénomène. Ici, l'idée centrale est donnée par :
\[
T_{\mathrm{surface}}<0^\circ\mathrm C
\]
Cette relation permet de relier la cause physique à l'effet observé. On vérifie les unités, puis on insère des ordres de grandeur réalistes.

La conclusion attendue est la suivante : Expliquer pourquoi le givre apparaît par ciel clair. La réponse doit toujours préciser les hypothèses utilisées. Si le phénomène dépend d'un milieu réel, d'une géométrie complexe, de frottements, de pertes ou de gradients, le résultat obtenu est un ordre de grandeur et non une valeur exacte.
\end{corrigebox}

\subsubsection*{Partie D — Approfondissement niveau Bac+3}

\begin{approfondissementbox}[Réponse Bac+3]
Au niveau Bac+3, on ne se contente plus d'une formule globale. On remplace le modèle simple par une loi locale, différentielle ou énergétique.

Une forme générale possible est :
\[
\frac{\mathrm d X}{\mathrm dt}=\text{apport}-\text{pertes},
\]
ou, pour un phénomène propagatif :
\[
\frac{\partial^2 \psi}{\partial t^2}=c^2\Delta\psi.
\]
Selon le cas, on peut aussi utiliser un principe variationnel, par exemple le principe de Fermat en optique :
\[
\delta \int n\,\mathrm ds=0.
\]

Pour ce phénomène, le modèle Bac+3 consiste à partir de la relation centrale
\[
T_{\mathrm{surface}}<0^\circ\mathrm C
\]
puis à introduire les corrections négligées : dépendance spatiale, dépendance temporelle, non-linéarités, pertes, conditions aux limites ou fluctuations. Cela permet de passer d'une réponse d'ordre de grandeur à une prédiction plus robuste.
\end{approfondissementbox}

\subsubsection*{Partie E — Limites, critique et prolongement}

\begin{enumerate}
  \item Citer au moins deux limites du modèle Terminale.
  \item Dire quelles grandeurs seraient difficiles à mesurer expérimentalement.
  \item Proposer une expérience simple permettant de tester le modèle.
  \item Expliquer ce qui changerait si l'on doublait une grandeur clé.
  \item Rédiger une conclusion courte, accessible à un élève de Terminale.
\end{enumerate}

\begin{corrigebox}[Synthèse rédigée]
Le phénomène « Givre » s'explique par un petit nombre de lois physiques simples, mais sa manifestation réelle dépend souvent de nombreux paramètres. Le modèle Terminale donne l'intuition et les ordres de grandeur ; le modèle Bac+3 introduit les variations locales, les pertes, les équations différentielles ou les principes variationnels nécessaires pour décrire finement la réalité.
\end{corrigebox}

% ------------------------------------------------------------
\subsection{Problème corrigé 19 — Rosée}
% ------------------------------------------------------------

\begin{exercicebox}[Question-problème]
\textbf{Observation.} Des gouttes apparaissent sur l'herbe le matin.

\medskip
\textbf{Question.} Relier humidité, refroidissement et condensation.

\medskip
\textbf{Thèmes mobilisés :} Condensation, point de rosée.

\medskip
\textbf{Relation ou idée centrale :}
\[
T_{\mathrm{surface}}<T_{\mathrm{rosée}}
\]
\end{exercicebox}

\subsubsection*{Partie A — Modélisation niveau Terminale}

\begin{enumerate}
  \item Décrire précisément le phénomène observé.
  \item Identifier le système physique étudié.
  \item Lister les grandeurs importantes et leurs unités.
  \item Écrire la loi physique simple qui permet de commencer le raisonnement.
  \item Expliquer les hypothèses simplificatrices du modèle Terminale.
\end{enumerate}

\subsubsection*{Partie B — Mise en équation et calcul}

\begin{enumerate}
  \item Définir les variables utiles.
  \item Écrire la relation mathématique principale.
  \item Vérifier l'homogénéité de la formule.
  \item Choisir un ordre de grandeur réaliste pour les données manquantes.
  \item Effectuer le calcul ou établir le raisonnement qualitatif demandé.
  \item Interpréter le résultat en une phrase claire.
\end{enumerate}

\subsubsection*{Partie C — Réponse accessible niveau Terminale}

\begin{corrigebox}[Réponse Terminale]
On part de l'observation concrète : Des gouttes apparaissent sur l'herbe le matin.

Au niveau Terminale, on cherche d'abord la grandeur qui contrôle le phénomène. Ici, l'idée centrale est donnée par :
\[
T_{\mathrm{surface}}<T_{\mathrm{rosée}}
\]
Cette relation permet de relier la cause physique à l'effet observé. On vérifie les unités, puis on insère des ordres de grandeur réalistes.

La conclusion attendue est la suivante : Relier humidité, refroidissement et condensation. La réponse doit toujours préciser les hypothèses utilisées. Si le phénomène dépend d'un milieu réel, d'une géométrie complexe, de frottements, de pertes ou de gradients, le résultat obtenu est un ordre de grandeur et non une valeur exacte.
\end{corrigebox}

\subsubsection*{Partie D — Approfondissement niveau Bac+3}

\begin{approfondissementbox}[Réponse Bac+3]
Au niveau Bac+3, on ne se contente plus d'une formule globale. On remplace le modèle simple par une loi locale, différentielle ou énergétique.

Une forme générale possible est :
\[
\frac{\mathrm d X}{\mathrm dt}=\text{apport}-\text{pertes},
\]
ou, pour un phénomène propagatif :
\[
\frac{\partial^2 \psi}{\partial t^2}=c^2\Delta\psi.
\]
Selon le cas, on peut aussi utiliser un principe variationnel, par exemple le principe de Fermat en optique :
\[
\delta \int n\,\mathrm ds=0.
\]

Pour ce phénomène, le modèle Bac+3 consiste à partir de la relation centrale
\[
T_{\mathrm{surface}}<T_{\mathrm{rosée}}
\]
puis à introduire les corrections négligées : dépendance spatiale, dépendance temporelle, non-linéarités, pertes, conditions aux limites ou fluctuations. Cela permet de passer d'une réponse d'ordre de grandeur à une prédiction plus robuste.
\end{approfondissementbox}

\subsubsection*{Partie E — Limites, critique et prolongement}

\begin{enumerate}
  \item Citer au moins deux limites du modèle Terminale.
  \item Dire quelles grandeurs seraient difficiles à mesurer expérimentalement.
  \item Proposer une expérience simple permettant de tester le modèle.
  \item Expliquer ce qui changerait si l'on doublait une grandeur clé.
  \item Rédiger une conclusion courte, accessible à un élève de Terminale.
\end{enumerate}

\begin{corrigebox}[Synthèse rédigée]
Le phénomène « Rosée » s'explique par un petit nombre de lois physiques simples, mais sa manifestation réelle dépend souvent de nombreux paramètres. Le modèle Terminale donne l'intuition et les ordres de grandeur ; le modèle Bac+3 introduit les variations locales, les pertes, les équations différentielles ou les principes variationnels nécessaires pour décrire finement la réalité.
\end{corrigebox}

% ------------------------------------------------------------
\subsection{Problème corrigé 20 — Buée sur une vitre}
% ------------------------------------------------------------

\begin{exercicebox}[Question-problème]
\textbf{Observation.} Une vitre froide se couvre de gouttelettes.

\medskip
\textbf{Question.} Expliquer pourquoi souffler sur une vitre peut créer de la buée.

\medskip
\textbf{Thèmes mobilisés :} Condensation sur surface froide.

\medskip
\textbf{Relation ou idée centrale :}
\[
\mathrm{HR}=100\%
\]
\end{exercicebox}

\subsubsection*{Partie A — Modélisation niveau Terminale}

\begin{enumerate}
  \item Décrire précisément le phénomène observé.
  \item Identifier le système physique étudié.
  \item Lister les grandeurs importantes et leurs unités.
  \item Écrire la loi physique simple qui permet de commencer le raisonnement.
  \item Expliquer les hypothèses simplificatrices du modèle Terminale.
\end{enumerate}

\subsubsection*{Partie B — Mise en équation et calcul}

\begin{enumerate}
  \item Définir les variables utiles.
  \item Écrire la relation mathématique principale.
  \item Vérifier l'homogénéité de la formule.
  \item Choisir un ordre de grandeur réaliste pour les données manquantes.
  \item Effectuer le calcul ou établir le raisonnement qualitatif demandé.
  \item Interpréter le résultat en une phrase claire.
\end{enumerate}

\subsubsection*{Partie C — Réponse accessible niveau Terminale}

\begin{corrigebox}[Réponse Terminale]
On part de l'observation concrète : Une vitre froide se couvre de gouttelettes.

Au niveau Terminale, on cherche d'abord la grandeur qui contrôle le phénomène. Ici, l'idée centrale est donnée par :
\[
\mathrm{HR}=100\%
\]
Cette relation permet de relier la cause physique à l'effet observé. On vérifie les unités, puis on insère des ordres de grandeur réalistes.

La conclusion attendue est la suivante : Expliquer pourquoi souffler sur une vitre peut créer de la buée. La réponse doit toujours préciser les hypothèses utilisées. Si le phénomène dépend d'un milieu réel, d'une géométrie complexe, de frottements, de pertes ou de gradients, le résultat obtenu est un ordre de grandeur et non une valeur exacte.
\end{corrigebox}

\subsubsection*{Partie D — Approfondissement niveau Bac+3}

\begin{approfondissementbox}[Réponse Bac+3]
Au niveau Bac+3, on ne se contente plus d'une formule globale. On remplace le modèle simple par une loi locale, différentielle ou énergétique.

Une forme générale possible est :
\[
\frac{\mathrm d X}{\mathrm dt}=\text{apport}-\text{pertes},
\]
ou, pour un phénomène propagatif :
\[
\frac{\partial^2 \psi}{\partial t^2}=c^2\Delta\psi.
\]
Selon le cas, on peut aussi utiliser un principe variationnel, par exemple le principe de Fermat en optique :
\[
\delta \int n\,\mathrm ds=0.
\]

Pour ce phénomène, le modèle Bac+3 consiste à partir de la relation centrale
\[
\mathrm{HR}=100\%
\]
puis à introduire les corrections négligées : dépendance spatiale, dépendance temporelle, non-linéarités, pertes, conditions aux limites ou fluctuations. Cela permet de passer d'une réponse d'ordre de grandeur à une prédiction plus robuste.
\end{approfondissementbox}

\subsubsection*{Partie E — Limites, critique et prolongement}

\begin{enumerate}
  \item Citer au moins deux limites du modèle Terminale.
  \item Dire quelles grandeurs seraient difficiles à mesurer expérimentalement.
  \item Proposer une expérience simple permettant de tester le modèle.
  \item Expliquer ce qui changerait si l'on doublait une grandeur clé.
  \item Rédiger une conclusion courte, accessible à un élève de Terminale.
\end{enumerate}

\begin{corrigebox}[Synthèse rédigée]
Le phénomène « Buée sur une vitre » s'explique par un petit nombre de lois physiques simples, mais sa manifestation réelle dépend souvent de nombreux paramètres. Le modèle Terminale donne l'intuition et les ordres de grandeur ; le modèle Bac+3 introduit les variations locales, les pertes, les équations différentielles ou les principes variationnels nécessaires pour décrire finement la réalité.
\end{corrigebox}

% ------------------------------------------------------------
\subsection{Problème corrigé 21 — Vapeur visible}
% ------------------------------------------------------------

\begin{exercicebox}[Question-problème]
\textbf{Observation.} Un panache blanc apparaît au-dessus d'une casserole.

\medskip
\textbf{Question.} Comprendre pourquoi la vapeur d'eau pure est invisible.

\medskip
\textbf{Thèmes mobilisés :} Condensation en microgouttelettes.

\medskip
\textbf{Relation ou idée centrale :}
\[
\text{vapeur invisible}\to\text{gouttelettes visibles}
\]
\end{exercicebox}

\subsubsection*{Partie A — Modélisation niveau Terminale}

\begin{enumerate}
  \item Décrire précisément le phénomène observé.
  \item Identifier le système physique étudié.
  \item Lister les grandeurs importantes et leurs unités.
  \item Écrire la loi physique simple qui permet de commencer le raisonnement.
  \item Expliquer les hypothèses simplificatrices du modèle Terminale.
\end{enumerate}

\subsubsection*{Partie B — Mise en équation et calcul}

\begin{enumerate}
  \item Définir les variables utiles.
  \item Écrire la relation mathématique principale.
  \item Vérifier l'homogénéité de la formule.
  \item Choisir un ordre de grandeur réaliste pour les données manquantes.
  \item Effectuer le calcul ou établir le raisonnement qualitatif demandé.
  \item Interpréter le résultat en une phrase claire.
\end{enumerate}

\subsubsection*{Partie C — Réponse accessible niveau Terminale}

\begin{corrigebox}[Réponse Terminale]
On part de l'observation concrète : Un panache blanc apparaît au-dessus d'une casserole.

Au niveau Terminale, on cherche d'abord la grandeur qui contrôle le phénomène. Ici, l'idée centrale est donnée par :
\[
\text{vapeur invisible}\to\text{gouttelettes visibles}
\]
Cette relation permet de relier la cause physique à l'effet observé. On vérifie les unités, puis on insère des ordres de grandeur réalistes.

La conclusion attendue est la suivante : Comprendre pourquoi la vapeur d'eau pure est invisible. La réponse doit toujours préciser les hypothèses utilisées. Si le phénomène dépend d'un milieu réel, d'une géométrie complexe, de frottements, de pertes ou de gradients, le résultat obtenu est un ordre de grandeur et non une valeur exacte.
\end{corrigebox}

\subsubsection*{Partie D — Approfondissement niveau Bac+3}

\begin{approfondissementbox}[Réponse Bac+3]
Au niveau Bac+3, on ne se contente plus d'une formule globale. On remplace le modèle simple par une loi locale, différentielle ou énergétique.

Une forme générale possible est :
\[
\frac{\mathrm d X}{\mathrm dt}=\text{apport}-\text{pertes},
\]
ou, pour un phénomène propagatif :
\[
\frac{\partial^2 \psi}{\partial t^2}=c^2\Delta\psi.
\]
Selon le cas, on peut aussi utiliser un principe variationnel, par exemple le principe de Fermat en optique :
\[
\delta \int n\,\mathrm ds=0.
\]

Pour ce phénomène, le modèle Bac+3 consiste à partir de la relation centrale
\[
\text{vapeur invisible}\to\text{gouttelettes visibles}
\]
puis à introduire les corrections négligées : dépendance spatiale, dépendance temporelle, non-linéarités, pertes, conditions aux limites ou fluctuations. Cela permet de passer d'une réponse d'ordre de grandeur à une prédiction plus robuste.
\end{approfondissementbox}

\subsubsection*{Partie E — Limites, critique et prolongement}

\begin{enumerate}
  \item Citer au moins deux limites du modèle Terminale.
  \item Dire quelles grandeurs seraient difficiles à mesurer expérimentalement.
  \item Proposer une expérience simple permettant de tester le modèle.
  \item Expliquer ce qui changerait si l'on doublait une grandeur clé.
  \item Rédiger une conclusion courte, accessible à un élève de Terminale.
\end{enumerate}

\begin{corrigebox}[Synthèse rédigée]
Le phénomène « Vapeur visible » s'explique par un petit nombre de lois physiques simples, mais sa manifestation réelle dépend souvent de nombreux paramètres. Le modèle Terminale donne l'intuition et les ordres de grandeur ; le modèle Bac+3 introduit les variations locales, les pertes, les équations différentielles ou les principes variationnels nécessaires pour décrire finement la réalité.
\end{corrigebox}

% ------------------------------------------------------------
\subsection{Problème corrigé 22 — Ébullition de l'eau}
% ------------------------------------------------------------

\begin{exercicebox}[Question-problème]
\textbf{Observation.} Des bulles se forment dans tout le liquide.

\medskip
\textbf{Question.} Définir la condition d'ébullition.

\medskip
\textbf{Thèmes mobilisés :} Pression de vapeur saturante.

\medskip
\textbf{Relation ou idée centrale :}
\[
P_{\mathrm{sat}}(T)=P_{\mathrm{ext}}
\]
\end{exercicebox}

\subsubsection*{Partie A — Modélisation niveau Terminale}

\begin{enumerate}
  \item Décrire précisément le phénomène observé.
  \item Identifier le système physique étudié.
  \item Lister les grandeurs importantes et leurs unités.
  \item Écrire la loi physique simple qui permet de commencer le raisonnement.
  \item Expliquer les hypothèses simplificatrices du modèle Terminale.
\end{enumerate}

\subsubsection*{Partie B — Mise en équation et calcul}

\begin{enumerate}
  \item Définir les variables utiles.
  \item Écrire la relation mathématique principale.
  \item Vérifier l'homogénéité de la formule.
  \item Choisir un ordre de grandeur réaliste pour les données manquantes.
  \item Effectuer le calcul ou établir le raisonnement qualitatif demandé.
  \item Interpréter le résultat en une phrase claire.
\end{enumerate}

\subsubsection*{Partie C — Réponse accessible niveau Terminale}

\begin{corrigebox}[Réponse Terminale]
On part de l'observation concrète : Des bulles se forment dans tout le liquide.

Au niveau Terminale, on cherche d'abord la grandeur qui contrôle le phénomène. Ici, l'idée centrale est donnée par :
\[
P_{\mathrm{sat}}(T)=P_{\mathrm{ext}}
\]
Cette relation permet de relier la cause physique à l'effet observé. On vérifie les unités, puis on insère des ordres de grandeur réalistes.

La conclusion attendue est la suivante : Définir la condition d'ébullition. La réponse doit toujours préciser les hypothèses utilisées. Si le phénomène dépend d'un milieu réel, d'une géométrie complexe, de frottements, de pertes ou de gradients, le résultat obtenu est un ordre de grandeur et non une valeur exacte.
\end{corrigebox}

\subsubsection*{Partie D — Approfondissement niveau Bac+3}

\begin{approfondissementbox}[Réponse Bac+3]
Au niveau Bac+3, on ne se contente plus d'une formule globale. On remplace le modèle simple par une loi locale, différentielle ou énergétique.

Une forme générale possible est :
\[
\frac{\mathrm d X}{\mathrm dt}=\text{apport}-\text{pertes},
\]
ou, pour un phénomène propagatif :
\[
\frac{\partial^2 \psi}{\partial t^2}=c^2\Delta\psi.
\]
Selon le cas, on peut aussi utiliser un principe variationnel, par exemple le principe de Fermat en optique :
\[
\delta \int n\,\mathrm ds=0.
\]

Pour ce phénomène, le modèle Bac+3 consiste à partir de la relation centrale
\[
P_{\mathrm{sat}}(T)=P_{\mathrm{ext}}
\]
puis à introduire les corrections négligées : dépendance spatiale, dépendance temporelle, non-linéarités, pertes, conditions aux limites ou fluctuations. Cela permet de passer d'une réponse d'ordre de grandeur à une prédiction plus robuste.
\end{approfondissementbox}

\subsubsection*{Partie E — Limites, critique et prolongement}

\begin{enumerate}
  \item Citer au moins deux limites du modèle Terminale.
  \item Dire quelles grandeurs seraient difficiles à mesurer expérimentalement.
  \item Proposer une expérience simple permettant de tester le modèle.
  \item Expliquer ce qui changerait si l'on doublait une grandeur clé.
  \item Rédiger une conclusion courte, accessible à un élève de Terminale.
\end{enumerate}

\begin{corrigebox}[Synthèse rédigée]
Le phénomène « Ébullition de l'eau » s'explique par un petit nombre de lois physiques simples, mais sa manifestation réelle dépend souvent de nombreux paramètres. Le modèle Terminale donne l'intuition et les ordres de grandeur ; le modèle Bac+3 introduit les variations locales, les pertes, les équations différentielles ou les principes variationnels nécessaires pour décrire finement la réalité.
\end{corrigebox}

% ------------------------------------------------------------
\subsection{Problème corrigé 23 — Ébullition en altitude}
% ------------------------------------------------------------

\begin{exercicebox}[Question-problème]
\textbf{Observation.} L'eau bout plus bas en altitude.

\medskip
\textbf{Question.} Expliquer pourquoi la cuisson change en montagne.

\medskip
\textbf{Thèmes mobilisés :} Pression atmosphérique décroissante.

\medskip
\textbf{Relation ou idée centrale :}
\[
P(z)=P_0e^{-z/H}
\]
\end{exercicebox}

\subsubsection*{Partie A — Modélisation niveau Terminale}

\begin{enumerate}
  \item Décrire précisément le phénomène observé.
  \item Identifier le système physique étudié.
  \item Lister les grandeurs importantes et leurs unités.
  \item Écrire la loi physique simple qui permet de commencer le raisonnement.
  \item Expliquer les hypothèses simplificatrices du modèle Terminale.
\end{enumerate}

\subsubsection*{Partie B — Mise en équation et calcul}

\begin{enumerate}
  \item Définir les variables utiles.
  \item Écrire la relation mathématique principale.
  \item Vérifier l'homogénéité de la formule.
  \item Choisir un ordre de grandeur réaliste pour les données manquantes.
  \item Effectuer le calcul ou établir le raisonnement qualitatif demandé.
  \item Interpréter le résultat en une phrase claire.
\end{enumerate}

\subsubsection*{Partie C — Réponse accessible niveau Terminale}

\begin{corrigebox}[Réponse Terminale]
On part de l'observation concrète : L'eau bout plus bas en altitude.

Au niveau Terminale, on cherche d'abord la grandeur qui contrôle le phénomène. Ici, l'idée centrale est donnée par :
\[
P(z)=P_0e^{-z/H}
\]
Cette relation permet de relier la cause physique à l'effet observé. On vérifie les unités, puis on insère des ordres de grandeur réalistes.

La conclusion attendue est la suivante : Expliquer pourquoi la cuisson change en montagne. La réponse doit toujours préciser les hypothèses utilisées. Si le phénomène dépend d'un milieu réel, d'une géométrie complexe, de frottements, de pertes ou de gradients, le résultat obtenu est un ordre de grandeur et non une valeur exacte.
\end{corrigebox}

\subsubsection*{Partie D — Approfondissement niveau Bac+3}

\begin{approfondissementbox}[Réponse Bac+3]
Au niveau Bac+3, on ne se contente plus d'une formule globale. On remplace le modèle simple par une loi locale, différentielle ou énergétique.

Une forme générale possible est :
\[
\frac{\mathrm d X}{\mathrm dt}=\text{apport}-\text{pertes},
\]
ou, pour un phénomène propagatif :
\[
\frac{\partial^2 \psi}{\partial t^2}=c^2\Delta\psi.
\]
Selon le cas, on peut aussi utiliser un principe variationnel, par exemple le principe de Fermat en optique :
\[
\delta \int n\,\mathrm ds=0.
\]

Pour ce phénomène, le modèle Bac+3 consiste à partir de la relation centrale
\[
P(z)=P_0e^{-z/H}
\]
puis à introduire les corrections négligées : dépendance spatiale, dépendance temporelle, non-linéarités, pertes, conditions aux limites ou fluctuations. Cela permet de passer d'une réponse d'ordre de grandeur à une prédiction plus robuste.
\end{approfondissementbox}

\subsubsection*{Partie E — Limites, critique et prolongement}

\begin{enumerate}
  \item Citer au moins deux limites du modèle Terminale.
  \item Dire quelles grandeurs seraient difficiles à mesurer expérimentalement.
  \item Proposer une expérience simple permettant de tester le modèle.
  \item Expliquer ce qui changerait si l'on doublait une grandeur clé.
  \item Rédiger une conclusion courte, accessible à un élève de Terminale.
\end{enumerate}

\begin{corrigebox}[Synthèse rédigée]
Le phénomène « Ébullition en altitude » s'explique par un petit nombre de lois physiques simples, mais sa manifestation réelle dépend souvent de nombreux paramètres. Le modèle Terminale donne l'intuition et les ordres de grandeur ; le modèle Bac+3 introduit les variations locales, les pertes, les équations différentielles ou les principes variationnels nécessaires pour décrire finement la réalité.
\end{corrigebox}

% ------------------------------------------------------------
\subsection{Problème corrigé 24 — Cocotte-minute}
% ------------------------------------------------------------

\begin{exercicebox}[Question-problème]
\textbf{Observation.} Les aliments cuisent plus vite sous pression.

\medskip
\textbf{Question.} Relier pression et température d'ébullition.

\medskip
\textbf{Thèmes mobilisés :} Ébullition à pression élevée.

\medskip
\textbf{Relation ou idée centrale :}
\[
P_{\mathrm{int}}>P_{\mathrm{atm}}
\]
\end{exercicebox}

\subsubsection*{Partie A — Modélisation niveau Terminale}

\begin{enumerate}
  \item Décrire précisément le phénomène observé.
  \item Identifier le système physique étudié.
  \item Lister les grandeurs importantes et leurs unités.
  \item Écrire la loi physique simple qui permet de commencer le raisonnement.
  \item Expliquer les hypothèses simplificatrices du modèle Terminale.
\end{enumerate}

\subsubsection*{Partie B — Mise en équation et calcul}

\begin{enumerate}
  \item Définir les variables utiles.
  \item Écrire la relation mathématique principale.
  \item Vérifier l'homogénéité de la formule.
  \item Choisir un ordre de grandeur réaliste pour les données manquantes.
  \item Effectuer le calcul ou établir le raisonnement qualitatif demandé.
  \item Interpréter le résultat en une phrase claire.
\end{enumerate}

\subsubsection*{Partie C — Réponse accessible niveau Terminale}

\begin{corrigebox}[Réponse Terminale]
On part de l'observation concrète : Les aliments cuisent plus vite sous pression.

Au niveau Terminale, on cherche d'abord la grandeur qui contrôle le phénomène. Ici, l'idée centrale est donnée par :
\[
P_{\mathrm{int}}>P_{\mathrm{atm}}
\]
Cette relation permet de relier la cause physique à l'effet observé. On vérifie les unités, puis on insère des ordres de grandeur réalistes.

La conclusion attendue est la suivante : Relier pression et température d'ébullition. La réponse doit toujours préciser les hypothèses utilisées. Si le phénomène dépend d'un milieu réel, d'une géométrie complexe, de frottements, de pertes ou de gradients, le résultat obtenu est un ordre de grandeur et non une valeur exacte.
\end{corrigebox}

\subsubsection*{Partie D — Approfondissement niveau Bac+3}

\begin{approfondissementbox}[Réponse Bac+3]
Au niveau Bac+3, on ne se contente plus d'une formule globale. On remplace le modèle simple par une loi locale, différentielle ou énergétique.

Une forme générale possible est :
\[
\frac{\mathrm d X}{\mathrm dt}=\text{apport}-\text{pertes},
\]
ou, pour un phénomène propagatif :
\[
\frac{\partial^2 \psi}{\partial t^2}=c^2\Delta\psi.
\]
Selon le cas, on peut aussi utiliser un principe variationnel, par exemple le principe de Fermat en optique :
\[
\delta \int n\,\mathrm ds=0.
\]

Pour ce phénomène, le modèle Bac+3 consiste à partir de la relation centrale
\[
P_{\mathrm{int}}>P_{\mathrm{atm}}
\]
puis à introduire les corrections négligées : dépendance spatiale, dépendance temporelle, non-linéarités, pertes, conditions aux limites ou fluctuations. Cela permet de passer d'une réponse d'ordre de grandeur à une prédiction plus robuste.
\end{approfondissementbox}

\subsubsection*{Partie E — Limites, critique et prolongement}

\begin{enumerate}
  \item Citer au moins deux limites du modèle Terminale.
  \item Dire quelles grandeurs seraient difficiles à mesurer expérimentalement.
  \item Proposer une expérience simple permettant de tester le modèle.
  \item Expliquer ce qui changerait si l'on doublait une grandeur clé.
  \item Rédiger une conclusion courte, accessible à un élève de Terminale.
\end{enumerate}

\begin{corrigebox}[Synthèse rédigée]
Le phénomène « Cocotte-minute » s'explique par un petit nombre de lois physiques simples, mais sa manifestation réelle dépend souvent de nombreux paramètres. Le modèle Terminale donne l'intuition et les ordres de grandeur ; le modèle Bac+3 introduit les variations locales, les pertes, les équations différentielles ou les principes variationnels nécessaires pour décrire finement la réalité.
\end{corrigebox}

% ------------------------------------------------------------
\subsection{Problème corrigé 25 — Effet de serre}
% ------------------------------------------------------------

\begin{exercicebox}[Question-problème]
\textbf{Observation.} L'atmosphère réchauffe la surface terrestre.

\medskip
\textbf{Question.} Construire un modèle radiatif simple.

\medskip
\textbf{Thèmes mobilisés :} Bilan radiatif, infrarouge.

\medskip
\textbf{Relation ou idée centrale :}
\[
P=\sigma S T^4
\]
\end{exercicebox}

\subsubsection*{Partie A — Modélisation niveau Terminale}

\begin{enumerate}
  \item Décrire précisément le phénomène observé.
  \item Identifier le système physique étudié.
  \item Lister les grandeurs importantes et leurs unités.
  \item Écrire la loi physique simple qui permet de commencer le raisonnement.
  \item Expliquer les hypothèses simplificatrices du modèle Terminale.
\end{enumerate}

\subsubsection*{Partie B — Mise en équation et calcul}

\begin{enumerate}
  \item Définir les variables utiles.
  \item Écrire la relation mathématique principale.
  \item Vérifier l'homogénéité de la formule.
  \item Choisir un ordre de grandeur réaliste pour les données manquantes.
  \item Effectuer le calcul ou établir le raisonnement qualitatif demandé.
  \item Interpréter le résultat en une phrase claire.
\end{enumerate}

\subsubsection*{Partie C — Réponse accessible niveau Terminale}

\begin{corrigebox}[Réponse Terminale]
On part de l'observation concrète : L'atmosphère réchauffe la surface terrestre.

Au niveau Terminale, on cherche d'abord la grandeur qui contrôle le phénomène. Ici, l'idée centrale est donnée par :
\[
P=\sigma S T^4
\]
Cette relation permet de relier la cause physique à l'effet observé. On vérifie les unités, puis on insère des ordres de grandeur réalistes.

La conclusion attendue est la suivante : Construire un modèle radiatif simple. La réponse doit toujours préciser les hypothèses utilisées. Si le phénomène dépend d'un milieu réel, d'une géométrie complexe, de frottements, de pertes ou de gradients, le résultat obtenu est un ordre de grandeur et non une valeur exacte.
\end{corrigebox}

\subsubsection*{Partie D — Approfondissement niveau Bac+3}

\begin{approfondissementbox}[Réponse Bac+3]
Au niveau Bac+3, on ne se contente plus d'une formule globale. On remplace le modèle simple par une loi locale, différentielle ou énergétique.

Une forme générale possible est :
\[
\frac{\mathrm d X}{\mathrm dt}=\text{apport}-\text{pertes},
\]
ou, pour un phénomène propagatif :
\[
\frac{\partial^2 \psi}{\partial t^2}=c^2\Delta\psi.
\]
Selon le cas, on peut aussi utiliser un principe variationnel, par exemple le principe de Fermat en optique :
\[
\delta \int n\,\mathrm ds=0.
\]

Pour ce phénomène, le modèle Bac+3 consiste à partir de la relation centrale
\[
P=\sigma S T^4
\]
puis à introduire les corrections négligées : dépendance spatiale, dépendance temporelle, non-linéarités, pertes, conditions aux limites ou fluctuations. Cela permet de passer d'une réponse d'ordre de grandeur à une prédiction plus robuste.
\end{approfondissementbox}

\subsubsection*{Partie E — Limites, critique et prolongement}

\begin{enumerate}
  \item Citer au moins deux limites du modèle Terminale.
  \item Dire quelles grandeurs seraient difficiles à mesurer expérimentalement.
  \item Proposer une expérience simple permettant de tester le modèle.
  \item Expliquer ce qui changerait si l'on doublait une grandeur clé.
  \item Rédiger une conclusion courte, accessible à un élève de Terminale.
\end{enumerate}

\begin{corrigebox}[Synthèse rédigée]
Le phénomène « Effet de serre » s'explique par un petit nombre de lois physiques simples, mais sa manifestation réelle dépend souvent de nombreux paramètres. Le modèle Terminale donne l'intuition et les ordres de grandeur ; le modèle Bac+3 introduit les variations locales, les pertes, les équations différentielles ou les principes variationnels nécessaires pour décrire finement la réalité.
\end{corrigebox}

% ------------------------------------------------------------
\subsection{Problème corrigé 26 — Route au soleil}
% ------------------------------------------------------------

\begin{exercicebox}[Question-problème]
\textbf{Observation.} L'asphalte devient brûlant.

\medskip
\textbf{Question.} Expliquer pourquoi les surfaces sombres chauffent.

\medskip
\textbf{Thèmes mobilisés :} Absorption solaire, albédo.

\medskip
\textbf{Relation ou idée centrale :}
\[
Q=mc\Delta T
\]
\end{exercicebox}

\subsubsection*{Partie A — Modélisation niveau Terminale}

\begin{enumerate}
  \item Décrire précisément le phénomène observé.
  \item Identifier le système physique étudié.
  \item Lister les grandeurs importantes et leurs unités.
  \item Écrire la loi physique simple qui permet de commencer le raisonnement.
  \item Expliquer les hypothèses simplificatrices du modèle Terminale.
\end{enumerate}

\subsubsection*{Partie B — Mise en équation et calcul}

\begin{enumerate}
  \item Définir les variables utiles.
  \item Écrire la relation mathématique principale.
  \item Vérifier l'homogénéité de la formule.
  \item Choisir un ordre de grandeur réaliste pour les données manquantes.
  \item Effectuer le calcul ou établir le raisonnement qualitatif demandé.
  \item Interpréter le résultat en une phrase claire.
\end{enumerate}

\subsubsection*{Partie C — Réponse accessible niveau Terminale}

\begin{corrigebox}[Réponse Terminale]
On part de l'observation concrète : L'asphalte devient brûlant.

Au niveau Terminale, on cherche d'abord la grandeur qui contrôle le phénomène. Ici, l'idée centrale est donnée par :
\[
Q=mc\Delta T
\]
Cette relation permet de relier la cause physique à l'effet observé. On vérifie les unités, puis on insère des ordres de grandeur réalistes.

La conclusion attendue est la suivante : Expliquer pourquoi les surfaces sombres chauffent. La réponse doit toujours préciser les hypothèses utilisées. Si le phénomène dépend d'un milieu réel, d'une géométrie complexe, de frottements, de pertes ou de gradients, le résultat obtenu est un ordre de grandeur et non une valeur exacte.
\end{corrigebox}

\subsubsection*{Partie D — Approfondissement niveau Bac+3}

\begin{approfondissementbox}[Réponse Bac+3]
Au niveau Bac+3, on ne se contente plus d'une formule globale. On remplace le modèle simple par une loi locale, différentielle ou énergétique.

Une forme générale possible est :
\[
\frac{\mathrm d X}{\mathrm dt}=\text{apport}-\text{pertes},
\]
ou, pour un phénomène propagatif :
\[
\frac{\partial^2 \psi}{\partial t^2}=c^2\Delta\psi.
\]
Selon le cas, on peut aussi utiliser un principe variationnel, par exemple le principe de Fermat en optique :
\[
\delta \int n\,\mathrm ds=0.
\]

Pour ce phénomène, le modèle Bac+3 consiste à partir de la relation centrale
\[
Q=mc\Delta T
\]
puis à introduire les corrections négligées : dépendance spatiale, dépendance temporelle, non-linéarités, pertes, conditions aux limites ou fluctuations. Cela permet de passer d'une réponse d'ordre de grandeur à une prédiction plus robuste.
\end{approfondissementbox}

\subsubsection*{Partie E — Limites, critique et prolongement}

\begin{enumerate}
  \item Citer au moins deux limites du modèle Terminale.
  \item Dire quelles grandeurs seraient difficiles à mesurer expérimentalement.
  \item Proposer une expérience simple permettant de tester le modèle.
  \item Expliquer ce qui changerait si l'on doublait une grandeur clé.
  \item Rédiger une conclusion courte, accessible à un élève de Terminale.
\end{enumerate}

\begin{corrigebox}[Synthèse rédigée]
Le phénomène « Route au soleil » s'explique par un petit nombre de lois physiques simples, mais sa manifestation réelle dépend souvent de nombreux paramètres. Le modèle Terminale donne l'intuition et les ordres de grandeur ; le modèle Bac+3 introduit les variations locales, les pertes, les équations différentielles ou les principes variationnels nécessaires pour décrire finement la réalité.
\end{corrigebox}

% ------------------------------------------------------------
\subsection{Problème corrigé 27 — Sable brûlant}
% ------------------------------------------------------------

\begin{exercicebox}[Question-problème]
\textbf{Observation.} Le sable brûle les pieds en été.

\medskip
\textbf{Question.} Relier surface chaude et faible diffusion thermique.

\medskip
\textbf{Thèmes mobilisés :} Absorption, faible conduction.

\medskip
\textbf{Relation ou idée centrale :}
\[
\Phi=-\lambda_{\mathrm{th}}\frac{\Delta T}{e}
\]
\end{exercicebox}

\subsubsection*{Partie A — Modélisation niveau Terminale}

\begin{enumerate}
  \item Décrire précisément le phénomène observé.
  \item Identifier le système physique étudié.
  \item Lister les grandeurs importantes et leurs unités.
  \item Écrire la loi physique simple qui permet de commencer le raisonnement.
  \item Expliquer les hypothèses simplificatrices du modèle Terminale.
\end{enumerate}

\subsubsection*{Partie B — Mise en équation et calcul}

\begin{enumerate}
  \item Définir les variables utiles.
  \item Écrire la relation mathématique principale.
  \item Vérifier l'homogénéité de la formule.
  \item Choisir un ordre de grandeur réaliste pour les données manquantes.
  \item Effectuer le calcul ou établir le raisonnement qualitatif demandé.
  \item Interpréter le résultat en une phrase claire.
\end{enumerate}

\subsubsection*{Partie C — Réponse accessible niveau Terminale}

\begin{corrigebox}[Réponse Terminale]
On part de l'observation concrète : Le sable brûle les pieds en été.

Au niveau Terminale, on cherche d'abord la grandeur qui contrôle le phénomène. Ici, l'idée centrale est donnée par :
\[
\Phi=-\lambda_{\mathrm{th}}\frac{\Delta T}{e}
\]
Cette relation permet de relier la cause physique à l'effet observé. On vérifie les unités, puis on insère des ordres de grandeur réalistes.

La conclusion attendue est la suivante : Relier surface chaude et faible diffusion thermique. La réponse doit toujours préciser les hypothèses utilisées. Si le phénomène dépend d'un milieu réel, d'une géométrie complexe, de frottements, de pertes ou de gradients, le résultat obtenu est un ordre de grandeur et non une valeur exacte.
\end{corrigebox}

\subsubsection*{Partie D — Approfondissement niveau Bac+3}

\begin{approfondissementbox}[Réponse Bac+3]
Au niveau Bac+3, on ne se contente plus d'une formule globale. On remplace le modèle simple par une loi locale, différentielle ou énergétique.

Une forme générale possible est :
\[
\frac{\mathrm d X}{\mathrm dt}=\text{apport}-\text{pertes},
\]
ou, pour un phénomène propagatif :
\[
\frac{\partial^2 \psi}{\partial t^2}=c^2\Delta\psi.
\]
Selon le cas, on peut aussi utiliser un principe variationnel, par exemple le principe de Fermat en optique :
\[
\delta \int n\,\mathrm ds=0.
\]

Pour ce phénomène, le modèle Bac+3 consiste à partir de la relation centrale
\[
\Phi=-\lambda_{\mathrm{th}}\frac{\Delta T}{e}
\]
puis à introduire les corrections négligées : dépendance spatiale, dépendance temporelle, non-linéarités, pertes, conditions aux limites ou fluctuations. Cela permet de passer d'une réponse d'ordre de grandeur à une prédiction plus robuste.
\end{approfondissementbox}

\subsubsection*{Partie E — Limites, critique et prolongement}

\begin{enumerate}
  \item Citer au moins deux limites du modèle Terminale.
  \item Dire quelles grandeurs seraient difficiles à mesurer expérimentalement.
  \item Proposer une expérience simple permettant de tester le modèle.
  \item Expliquer ce qui changerait si l'on doublait une grandeur clé.
  \item Rédiger une conclusion courte, accessible à un élève de Terminale.
\end{enumerate}

\begin{corrigebox}[Synthèse rédigée]
Le phénomène « Sable brûlant » s'explique par un petit nombre de lois physiques simples, mais sa manifestation réelle dépend souvent de nombreux paramètres. Le modèle Terminale donne l'intuition et les ordres de grandeur ; le modèle Bac+3 introduit les variations locales, les pertes, les équations différentielles ou les principes variationnels nécessaires pour décrire finement la réalité.
\end{corrigebox}

% ------------------------------------------------------------
\subsection{Problème corrigé 28 — Métal plus froid que bois}
% ------------------------------------------------------------

\begin{exercicebox}[Question-problème]
\textbf{Observation.} Le métal semble plus froid que le bois.

\medskip
\textbf{Question.} Comprendre la sensation thermique.

\medskip
\textbf{Thèmes mobilisés :} Conductivité thermique.

\medskip
\textbf{Relation ou idée centrale :}
\[
\Phi\propto\lambda_{\mathrm{th}}\Delta T
\]
\end{exercicebox}

\subsubsection*{Partie A — Modélisation niveau Terminale}

\begin{enumerate}
  \item Décrire précisément le phénomène observé.
  \item Identifier le système physique étudié.
  \item Lister les grandeurs importantes et leurs unités.
  \item Écrire la loi physique simple qui permet de commencer le raisonnement.
  \item Expliquer les hypothèses simplificatrices du modèle Terminale.
\end{enumerate}

\subsubsection*{Partie B — Mise en équation et calcul}

\begin{enumerate}
  \item Définir les variables utiles.
  \item Écrire la relation mathématique principale.
  \item Vérifier l'homogénéité de la formule.
  \item Choisir un ordre de grandeur réaliste pour les données manquantes.
  \item Effectuer le calcul ou établir le raisonnement qualitatif demandé.
  \item Interpréter le résultat en une phrase claire.
\end{enumerate}

\subsubsection*{Partie C — Réponse accessible niveau Terminale}

\begin{corrigebox}[Réponse Terminale]
On part de l'observation concrète : Le métal semble plus froid que le bois.

Au niveau Terminale, on cherche d'abord la grandeur qui contrôle le phénomène. Ici, l'idée centrale est donnée par :
\[
\Phi\propto\lambda_{\mathrm{th}}\Delta T
\]
Cette relation permet de relier la cause physique à l'effet observé. On vérifie les unités, puis on insère des ordres de grandeur réalistes.

La conclusion attendue est la suivante : Comprendre la sensation thermique. La réponse doit toujours préciser les hypothèses utilisées. Si le phénomène dépend d'un milieu réel, d'une géométrie complexe, de frottements, de pertes ou de gradients, le résultat obtenu est un ordre de grandeur et non une valeur exacte.
\end{corrigebox}

\subsubsection*{Partie D — Approfondissement niveau Bac+3}

\begin{approfondissementbox}[Réponse Bac+3]
Au niveau Bac+3, on ne se contente plus d'une formule globale. On remplace le modèle simple par une loi locale, différentielle ou énergétique.

Une forme générale possible est :
\[
\frac{\mathrm d X}{\mathrm dt}=\text{apport}-\text{pertes},
\]
ou, pour un phénomène propagatif :
\[
\frac{\partial^2 \psi}{\partial t^2}=c^2\Delta\psi.
\]
Selon le cas, on peut aussi utiliser un principe variationnel, par exemple le principe de Fermat en optique :
\[
\delta \int n\,\mathrm ds=0.
\]

Pour ce phénomène, le modèle Bac+3 consiste à partir de la relation centrale
\[
\Phi\propto\lambda_{\mathrm{th}}\Delta T
\]
puis à introduire les corrections négligées : dépendance spatiale, dépendance temporelle, non-linéarités, pertes, conditions aux limites ou fluctuations. Cela permet de passer d'une réponse d'ordre de grandeur à une prédiction plus robuste.
\end{approfondissementbox}

\subsubsection*{Partie E — Limites, critique et prolongement}

\begin{enumerate}
  \item Citer au moins deux limites du modèle Terminale.
  \item Dire quelles grandeurs seraient difficiles à mesurer expérimentalement.
  \item Proposer une expérience simple permettant de tester le modèle.
  \item Expliquer ce qui changerait si l'on doublait une grandeur clé.
  \item Rédiger une conclusion courte, accessible à un élève de Terminale.
\end{enumerate}

\begin{corrigebox}[Synthèse rédigée]
Le phénomène « Métal plus froid que bois » s'explique par un petit nombre de lois physiques simples, mais sa manifestation réelle dépend souvent de nombreux paramètres. Le modèle Terminale donne l'intuition et les ordres de grandeur ; le modèle Bac+3 introduit les variations locales, les pertes, les équations différentielles ou les principes variationnels nécessaires pour décrire finement la réalité.
\end{corrigebox}

% ------------------------------------------------------------
\subsection{Problème corrigé 29 — Refroidissement par évaporation}
% ------------------------------------------------------------

\begin{exercicebox}[Question-problème]
\textbf{Observation.} Un linge humide refroidit en séchant.

\medskip
\textbf{Question.} Calculer l'énergie emportée par évaporation.

\medskip
\textbf{Thèmes mobilisés :} Chaleur latente.

\medskip
\textbf{Relation ou idée centrale :}
\[
Q=mL_v
\]
\end{exercicebox}

\subsubsection*{Partie A — Modélisation niveau Terminale}

\begin{enumerate}
  \item Décrire précisément le phénomène observé.
  \item Identifier le système physique étudié.
  \item Lister les grandeurs importantes et leurs unités.
  \item Écrire la loi physique simple qui permet de commencer le raisonnement.
  \item Expliquer les hypothèses simplificatrices du modèle Terminale.
\end{enumerate}

\subsubsection*{Partie B — Mise en équation et calcul}

\begin{enumerate}
  \item Définir les variables utiles.
  \item Écrire la relation mathématique principale.
  \item Vérifier l'homogénéité de la formule.
  \item Choisir un ordre de grandeur réaliste pour les données manquantes.
  \item Effectuer le calcul ou établir le raisonnement qualitatif demandé.
  \item Interpréter le résultat en une phrase claire.
\end{enumerate}

\subsubsection*{Partie C — Réponse accessible niveau Terminale}

\begin{corrigebox}[Réponse Terminale]
On part de l'observation concrète : Un linge humide refroidit en séchant.

Au niveau Terminale, on cherche d'abord la grandeur qui contrôle le phénomène. Ici, l'idée centrale est donnée par :
\[
Q=mL_v
\]
Cette relation permet de relier la cause physique à l'effet observé. On vérifie les unités, puis on insère des ordres de grandeur réalistes.

La conclusion attendue est la suivante : Calculer l'énergie emportée par évaporation. La réponse doit toujours préciser les hypothèses utilisées. Si le phénomène dépend d'un milieu réel, d'une géométrie complexe, de frottements, de pertes ou de gradients, le résultat obtenu est un ordre de grandeur et non une valeur exacte.
\end{corrigebox}

\subsubsection*{Partie D — Approfondissement niveau Bac+3}

\begin{approfondissementbox}[Réponse Bac+3]
Au niveau Bac+3, on ne se contente plus d'une formule globale. On remplace le modèle simple par une loi locale, différentielle ou énergétique.

Une forme générale possible est :
\[
\frac{\mathrm d X}{\mathrm dt}=\text{apport}-\text{pertes},
\]
ou, pour un phénomène propagatif :
\[
\frac{\partial^2 \psi}{\partial t^2}=c^2\Delta\psi.
\]
Selon le cas, on peut aussi utiliser un principe variationnel, par exemple le principe de Fermat en optique :
\[
\delta \int n\,\mathrm ds=0.
\]

Pour ce phénomène, le modèle Bac+3 consiste à partir de la relation centrale
\[
Q=mL_v
\]
puis à introduire les corrections négligées : dépendance spatiale, dépendance temporelle, non-linéarités, pertes, conditions aux limites ou fluctuations. Cela permet de passer d'une réponse d'ordre de grandeur à une prédiction plus robuste.
\end{approfondissementbox}

\subsubsection*{Partie E — Limites, critique et prolongement}

\begin{enumerate}
  \item Citer au moins deux limites du modèle Terminale.
  \item Dire quelles grandeurs seraient difficiles à mesurer expérimentalement.
  \item Proposer une expérience simple permettant de tester le modèle.
  \item Expliquer ce qui changerait si l'on doublait une grandeur clé.
  \item Rédiger une conclusion courte, accessible à un élève de Terminale.
\end{enumerate}

\begin{corrigebox}[Synthèse rédigée]
Le phénomène « Refroidissement par évaporation » s'explique par un petit nombre de lois physiques simples, mais sa manifestation réelle dépend souvent de nombreux paramètres. Le modèle Terminale donne l'intuition et les ordres de grandeur ; le modèle Bac+3 introduit les variations locales, les pertes, les équations différentielles ou les principes variationnels nécessaires pour décrire finement la réalité.
\end{corrigebox}

% ------------------------------------------------------------
\subsection{Problème corrigé 30 — Transpiration}
% ------------------------------------------------------------

\begin{exercicebox}[Question-problème]
\textbf{Observation.} Le corps se refroidit en transpirant.

\medskip
\textbf{Question.} Relier débit de sueur et puissance de refroidissement.

\medskip
\textbf{Thèmes mobilisés :} Évaporation, bilan thermique.

\medskip
\textbf{Relation ou idée centrale :}
\[
P=\dot mL_v
\]
\end{exercicebox}

\subsubsection*{Partie A — Modélisation niveau Terminale}

\begin{enumerate}
  \item Décrire précisément le phénomène observé.
  \item Identifier le système physique étudié.
  \item Lister les grandeurs importantes et leurs unités.
  \item Écrire la loi physique simple qui permet de commencer le raisonnement.
  \item Expliquer les hypothèses simplificatrices du modèle Terminale.
\end{enumerate}

\subsubsection*{Partie B — Mise en équation et calcul}

\begin{enumerate}
  \item Définir les variables utiles.
  \item Écrire la relation mathématique principale.
  \item Vérifier l'homogénéité de la formule.
  \item Choisir un ordre de grandeur réaliste pour les données manquantes.
  \item Effectuer le calcul ou établir le raisonnement qualitatif demandé.
  \item Interpréter le résultat en une phrase claire.
\end{enumerate}

\subsubsection*{Partie C — Réponse accessible niveau Terminale}

\begin{corrigebox}[Réponse Terminale]
On part de l'observation concrète : Le corps se refroidit en transpirant.

Au niveau Terminale, on cherche d'abord la grandeur qui contrôle le phénomène. Ici, l'idée centrale est donnée par :
\[
P=\dot mL_v
\]
Cette relation permet de relier la cause physique à l'effet observé. On vérifie les unités, puis on insère des ordres de grandeur réalistes.

La conclusion attendue est la suivante : Relier débit de sueur et puissance de refroidissement. La réponse doit toujours préciser les hypothèses utilisées. Si le phénomène dépend d'un milieu réel, d'une géométrie complexe, de frottements, de pertes ou de gradients, le résultat obtenu est un ordre de grandeur et non une valeur exacte.
\end{corrigebox}

\subsubsection*{Partie D — Approfondissement niveau Bac+3}

\begin{approfondissementbox}[Réponse Bac+3]
Au niveau Bac+3, on ne se contente plus d'une formule globale. On remplace le modèle simple par une loi locale, différentielle ou énergétique.

Une forme générale possible est :
\[
\frac{\mathrm d X}{\mathrm dt}=\text{apport}-\text{pertes},
\]
ou, pour un phénomène propagatif :
\[
\frac{\partial^2 \psi}{\partial t^2}=c^2\Delta\psi.
\]
Selon le cas, on peut aussi utiliser un principe variationnel, par exemple le principe de Fermat en optique :
\[
\delta \int n\,\mathrm ds=0.
\]

Pour ce phénomène, le modèle Bac+3 consiste à partir de la relation centrale
\[
P=\dot mL_v
\]
puis à introduire les corrections négligées : dépendance spatiale, dépendance temporelle, non-linéarités, pertes, conditions aux limites ou fluctuations. Cela permet de passer d'une réponse d'ordre de grandeur à une prédiction plus robuste.
\end{approfondissementbox}

\subsubsection*{Partie E — Limites, critique et prolongement}

\begin{enumerate}
  \item Citer au moins deux limites du modèle Terminale.
  \item Dire quelles grandeurs seraient difficiles à mesurer expérimentalement.
  \item Proposer une expérience simple permettant de tester le modèle.
  \item Expliquer ce qui changerait si l'on doublait une grandeur clé.
  \item Rédiger une conclusion courte, accessible à un élève de Terminale.
\end{enumerate}

\begin{corrigebox}[Synthèse rédigée]
Le phénomène « Transpiration » s'explique par un petit nombre de lois physiques simples, mais sa manifestation réelle dépend souvent de nombreux paramètres. Le modèle Terminale donne l'intuition et les ordres de grandeur ; le modèle Bac+3 introduit les variations locales, les pertes, les équations différentielles ou les principes variationnels nécessaires pour décrire finement la réalité.
\end{corrigebox}

% ------------------------------------------------------------
\subsection{Problème corrigé 31 — Isolation d'une maison}
% ------------------------------------------------------------

\begin{exercicebox}[Question-problème]
\textbf{Observation.} Une maison isolée conserve mieux la chaleur.

\medskip
\textbf{Question.} Relier isolation et puissance de chauffage.

\medskip
\textbf{Thèmes mobilisés :} Résistance thermique.

\medskip
\textbf{Relation ou idée centrale :}
\[
\Phi=\frac{\Delta T}{R_{\mathrm{th}}}
\]
\end{exercicebox}

\subsubsection*{Partie A — Modélisation niveau Terminale}

\begin{enumerate}
  \item Décrire précisément le phénomène observé.
  \item Identifier le système physique étudié.
  \item Lister les grandeurs importantes et leurs unités.
  \item Écrire la loi physique simple qui permet de commencer le raisonnement.
  \item Expliquer les hypothèses simplificatrices du modèle Terminale.
\end{enumerate}

\subsubsection*{Partie B — Mise en équation et calcul}

\begin{enumerate}
  \item Définir les variables utiles.
  \item Écrire la relation mathématique principale.
  \item Vérifier l'homogénéité de la formule.
  \item Choisir un ordre de grandeur réaliste pour les données manquantes.
  \item Effectuer le calcul ou établir le raisonnement qualitatif demandé.
  \item Interpréter le résultat en une phrase claire.
\end{enumerate}

\subsubsection*{Partie C — Réponse accessible niveau Terminale}

\begin{corrigebox}[Réponse Terminale]
On part de l'observation concrète : Une maison isolée conserve mieux la chaleur.

Au niveau Terminale, on cherche d'abord la grandeur qui contrôle le phénomène. Ici, l'idée centrale est donnée par :
\[
\Phi=\frac{\Delta T}{R_{\mathrm{th}}}
\]
Cette relation permet de relier la cause physique à l'effet observé. On vérifie les unités, puis on insère des ordres de grandeur réalistes.

La conclusion attendue est la suivante : Relier isolation et puissance de chauffage. La réponse doit toujours préciser les hypothèses utilisées. Si le phénomène dépend d'un milieu réel, d'une géométrie complexe, de frottements, de pertes ou de gradients, le résultat obtenu est un ordre de grandeur et non une valeur exacte.
\end{corrigebox}

\subsubsection*{Partie D — Approfondissement niveau Bac+3}

\begin{approfondissementbox}[Réponse Bac+3]
Au niveau Bac+3, on ne se contente plus d'une formule globale. On remplace le modèle simple par une loi locale, différentielle ou énergétique.

Une forme générale possible est :
\[
\frac{\mathrm d X}{\mathrm dt}=\text{apport}-\text{pertes},
\]
ou, pour un phénomène propagatif :
\[
\frac{\partial^2 \psi}{\partial t^2}=c^2\Delta\psi.
\]
Selon le cas, on peut aussi utiliser un principe variationnel, par exemple le principe de Fermat en optique :
\[
\delta \int n\,\mathrm ds=0.
\]

Pour ce phénomène, le modèle Bac+3 consiste à partir de la relation centrale
\[
\Phi=\frac{\Delta T}{R_{\mathrm{th}}}
\]
puis à introduire les corrections négligées : dépendance spatiale, dépendance temporelle, non-linéarités, pertes, conditions aux limites ou fluctuations. Cela permet de passer d'une réponse d'ordre de grandeur à une prédiction plus robuste.
\end{approfondissementbox}

\subsubsection*{Partie E — Limites, critique et prolongement}

\begin{enumerate}
  \item Citer au moins deux limites du modèle Terminale.
  \item Dire quelles grandeurs seraient difficiles à mesurer expérimentalement.
  \item Proposer une expérience simple permettant de tester le modèle.
  \item Expliquer ce qui changerait si l'on doublait une grandeur clé.
  \item Rédiger une conclusion courte, accessible à un élève de Terminale.
\end{enumerate}

\begin{corrigebox}[Synthèse rédigée]
Le phénomène « Isolation d'une maison » s'explique par un petit nombre de lois physiques simples, mais sa manifestation réelle dépend souvent de nombreux paramètres. Le modèle Terminale donne l'intuition et les ordres de grandeur ; le modèle Bac+3 introduit les variations locales, les pertes, les équations différentielles ou les principes variationnels nécessaires pour décrire finement la réalité.
\end{corrigebox}

% ------------------------------------------------------------
\subsection{Problème corrigé 32 — Double vitrage}
% ------------------------------------------------------------

\begin{exercicebox}[Question-problème]
\textbf{Observation.} Une fenêtre double limite les pertes.

\medskip
\textbf{Question.} Expliquer l'intérêt de la lame d'air ou de gaz.

\medskip
\textbf{Thèmes mobilisés :} Conduction, convection, résistance thermique.

\medskip
\textbf{Relation ou idée centrale :}
\[
R_{\mathrm{tot}}=R_1+R_2+R_{\mathrm{gaz}}
\]
\end{exercicebox}

\subsubsection*{Partie A — Modélisation niveau Terminale}

\begin{enumerate}
  \item Décrire précisément le phénomène observé.
  \item Identifier le système physique étudié.
  \item Lister les grandeurs importantes et leurs unités.
  \item Écrire la loi physique simple qui permet de commencer le raisonnement.
  \item Expliquer les hypothèses simplificatrices du modèle Terminale.
\end{enumerate}

\subsubsection*{Partie B — Mise en équation et calcul}

\begin{enumerate}
  \item Définir les variables utiles.
  \item Écrire la relation mathématique principale.
  \item Vérifier l'homogénéité de la formule.
  \item Choisir un ordre de grandeur réaliste pour les données manquantes.
  \item Effectuer le calcul ou établir le raisonnement qualitatif demandé.
  \item Interpréter le résultat en une phrase claire.
\end{enumerate}

\subsubsection*{Partie C — Réponse accessible niveau Terminale}

\begin{corrigebox}[Réponse Terminale]
On part de l'observation concrète : Une fenêtre double limite les pertes.

Au niveau Terminale, on cherche d'abord la grandeur qui contrôle le phénomène. Ici, l'idée centrale est donnée par :
\[
R_{\mathrm{tot}}=R_1+R_2+R_{\mathrm{gaz}}
\]
Cette relation permet de relier la cause physique à l'effet observé. On vérifie les unités, puis on insère des ordres de grandeur réalistes.

La conclusion attendue est la suivante : Expliquer l'intérêt de la lame d'air ou de gaz. La réponse doit toujours préciser les hypothèses utilisées. Si le phénomène dépend d'un milieu réel, d'une géométrie complexe, de frottements, de pertes ou de gradients, le résultat obtenu est un ordre de grandeur et non une valeur exacte.
\end{corrigebox}

\subsubsection*{Partie D — Approfondissement niveau Bac+3}

\begin{approfondissementbox}[Réponse Bac+3]
Au niveau Bac+3, on ne se contente plus d'une formule globale. On remplace le modèle simple par une loi locale, différentielle ou énergétique.

Une forme générale possible est :
\[
\frac{\mathrm d X}{\mathrm dt}=\text{apport}-\text{pertes},
\]
ou, pour un phénomène propagatif :
\[
\frac{\partial^2 \psi}{\partial t^2}=c^2\Delta\psi.
\]
Selon le cas, on peut aussi utiliser un principe variationnel, par exemple le principe de Fermat en optique :
\[
\delta \int n\,\mathrm ds=0.
\]

Pour ce phénomène, le modèle Bac+3 consiste à partir de la relation centrale
\[
R_{\mathrm{tot}}=R_1+R_2+R_{\mathrm{gaz}}
\]
puis à introduire les corrections négligées : dépendance spatiale, dépendance temporelle, non-linéarités, pertes, conditions aux limites ou fluctuations. Cela permet de passer d'une réponse d'ordre de grandeur à une prédiction plus robuste.
\end{approfondissementbox}

\subsubsection*{Partie E — Limites, critique et prolongement}

\begin{enumerate}
  \item Citer au moins deux limites du modèle Terminale.
  \item Dire quelles grandeurs seraient difficiles à mesurer expérimentalement.
  \item Proposer une expérience simple permettant de tester le modèle.
  \item Expliquer ce qui changerait si l'on doublait une grandeur clé.
  \item Rédiger une conclusion courte, accessible à un élève de Terminale.
\end{enumerate}

\begin{corrigebox}[Synthèse rédigée]
Le phénomène « Double vitrage » s'explique par un petit nombre de lois physiques simples, mais sa manifestation réelle dépend souvent de nombreux paramètres. Le modèle Terminale donne l'intuition et les ordres de grandeur ; le modèle Bac+3 introduit les variations locales, les pertes, les équations différentielles ou les principes variationnels nécessaires pour décrire finement la réalité.
\end{corrigebox}

% ------------------------------------------------------------
\subsection{Problème corrigé 33 — Fibre optique}
% ------------------------------------------------------------

\begin{exercicebox}[Question-problème]
\textbf{Observation.} La lumière se propage dans un fil de verre.

\medskip
\textbf{Question.} Calculer l'angle limite et le temps de propagation.

\medskip
\textbf{Thèmes mobilisés :} Réflexion totale, indice.

\medskip
\textbf{Relation ou idée centrale :}
\[
\sin i_c=\frac{n_g}{n_c}
\]
\end{exercicebox}

\subsubsection*{Partie A — Modélisation niveau Terminale}

\begin{enumerate}
  \item Décrire précisément le phénomène observé.
  \item Identifier le système physique étudié.
  \item Lister les grandeurs importantes et leurs unités.
  \item Écrire la loi physique simple qui permet de commencer le raisonnement.
  \item Expliquer les hypothèses simplificatrices du modèle Terminale.
\end{enumerate}

\subsubsection*{Partie B — Mise en équation et calcul}

\begin{enumerate}
  \item Définir les variables utiles.
  \item Écrire la relation mathématique principale.
  \item Vérifier l'homogénéité de la formule.
  \item Choisir un ordre de grandeur réaliste pour les données manquantes.
  \item Effectuer le calcul ou établir le raisonnement qualitatif demandé.
  \item Interpréter le résultat en une phrase claire.
\end{enumerate}

\subsubsection*{Partie C — Réponse accessible niveau Terminale}

\begin{corrigebox}[Réponse Terminale]
On part de l'observation concrète : La lumière se propage dans un fil de verre.

Au niveau Terminale, on cherche d'abord la grandeur qui contrôle le phénomène. Ici, l'idée centrale est donnée par :
\[
\sin i_c=\frac{n_g}{n_c}
\]
Cette relation permet de relier la cause physique à l'effet observé. On vérifie les unités, puis on insère des ordres de grandeur réalistes.

La conclusion attendue est la suivante : Calculer l'angle limite et le temps de propagation. La réponse doit toujours préciser les hypothèses utilisées. Si le phénomène dépend d'un milieu réel, d'une géométrie complexe, de frottements, de pertes ou de gradients, le résultat obtenu est un ordre de grandeur et non une valeur exacte.
\end{corrigebox}

\subsubsection*{Partie D — Approfondissement niveau Bac+3}

\begin{approfondissementbox}[Réponse Bac+3]
Au niveau Bac+3, on ne se contente plus d'une formule globale. On remplace le modèle simple par une loi locale, différentielle ou énergétique.

Une forme générale possible est :
\[
\frac{\mathrm d X}{\mathrm dt}=\text{apport}-\text{pertes},
\]
ou, pour un phénomène propagatif :
\[
\frac{\partial^2 \psi}{\partial t^2}=c^2\Delta\psi.
\]
Selon le cas, on peut aussi utiliser un principe variationnel, par exemple le principe de Fermat en optique :
\[
\delta \int n\,\mathrm ds=0.
\]

Pour ce phénomène, le modèle Bac+3 consiste à partir de la relation centrale
\[
\sin i_c=\frac{n_g}{n_c}
\]
puis à introduire les corrections négligées : dépendance spatiale, dépendance temporelle, non-linéarités, pertes, conditions aux limites ou fluctuations. Cela permet de passer d'une réponse d'ordre de grandeur à une prédiction plus robuste.
\end{approfondissementbox}

\subsubsection*{Partie E — Limites, critique et prolongement}

\begin{enumerate}
  \item Citer au moins deux limites du modèle Terminale.
  \item Dire quelles grandeurs seraient difficiles à mesurer expérimentalement.
  \item Proposer une expérience simple permettant de tester le modèle.
  \item Expliquer ce qui changerait si l'on doublait une grandeur clé.
  \item Rédiger une conclusion courte, accessible à un élève de Terminale.
\end{enumerate}

\begin{corrigebox}[Synthèse rédigée]
Le phénomène « Fibre optique » s'explique par un petit nombre de lois physiques simples, mais sa manifestation réelle dépend souvent de nombreux paramètres. Le modèle Terminale donne l'intuition et les ordres de grandeur ; le modèle Bac+3 introduit les variations locales, les pertes, les équations différentielles ou les principes variationnels nécessaires pour décrire finement la réalité.
\end{corrigebox}

% ------------------------------------------------------------
\subsection{Problème corrigé 34 — Réflexion totale interne}
% ------------------------------------------------------------

\begin{exercicebox}[Question-problème]
\textbf{Observation.} Un rayon peut être totalement réfléchi.

\medskip
\textbf{Question.} Déterminer quand la réfraction devient impossible.

\medskip
\textbf{Thèmes mobilisés :} Angle limite, indice.

\medskip
\textbf{Relation ou idée centrale :}
\[
\sin i_c=\frac{n_2}{n_1}
\]
\end{exercicebox}

\subsubsection*{Partie A — Modélisation niveau Terminale}

\begin{enumerate}
  \item Décrire précisément le phénomène observé.
  \item Identifier le système physique étudié.
  \item Lister les grandeurs importantes et leurs unités.
  \item Écrire la loi physique simple qui permet de commencer le raisonnement.
  \item Expliquer les hypothèses simplificatrices du modèle Terminale.
\end{enumerate}

\subsubsection*{Partie B — Mise en équation et calcul}

\begin{enumerate}
  \item Définir les variables utiles.
  \item Écrire la relation mathématique principale.
  \item Vérifier l'homogénéité de la formule.
  \item Choisir un ordre de grandeur réaliste pour les données manquantes.
  \item Effectuer le calcul ou établir le raisonnement qualitatif demandé.
  \item Interpréter le résultat en une phrase claire.
\end{enumerate}

\subsubsection*{Partie C — Réponse accessible niveau Terminale}

\begin{corrigebox}[Réponse Terminale]
On part de l'observation concrète : Un rayon peut être totalement réfléchi.

Au niveau Terminale, on cherche d'abord la grandeur qui contrôle le phénomène. Ici, l'idée centrale est donnée par :
\[
\sin i_c=\frac{n_2}{n_1}
\]
Cette relation permet de relier la cause physique à l'effet observé. On vérifie les unités, puis on insère des ordres de grandeur réalistes.

La conclusion attendue est la suivante : Déterminer quand la réfraction devient impossible. La réponse doit toujours préciser les hypothèses utilisées. Si le phénomène dépend d'un milieu réel, d'une géométrie complexe, de frottements, de pertes ou de gradients, le résultat obtenu est un ordre de grandeur et non une valeur exacte.
\end{corrigebox}

\subsubsection*{Partie D — Approfondissement niveau Bac+3}

\begin{approfondissementbox}[Réponse Bac+3]
Au niveau Bac+3, on ne se contente plus d'une formule globale. On remplace le modèle simple par une loi locale, différentielle ou énergétique.

Une forme générale possible est :
\[
\frac{\mathrm d X}{\mathrm dt}=\text{apport}-\text{pertes},
\]
ou, pour un phénomène propagatif :
\[
\frac{\partial^2 \psi}{\partial t^2}=c^2\Delta\psi.
\]
Selon le cas, on peut aussi utiliser un principe variationnel, par exemple le principe de Fermat en optique :
\[
\delta \int n\,\mathrm ds=0.
\]

Pour ce phénomène, le modèle Bac+3 consiste à partir de la relation centrale
\[
\sin i_c=\frac{n_2}{n_1}
\]
puis à introduire les corrections négligées : dépendance spatiale, dépendance temporelle, non-linéarités, pertes, conditions aux limites ou fluctuations. Cela permet de passer d'une réponse d'ordre de grandeur à une prédiction plus robuste.
\end{approfondissementbox}

\subsubsection*{Partie E — Limites, critique et prolongement}

\begin{enumerate}
  \item Citer au moins deux limites du modèle Terminale.
  \item Dire quelles grandeurs seraient difficiles à mesurer expérimentalement.
  \item Proposer une expérience simple permettant de tester le modèle.
  \item Expliquer ce qui changerait si l'on doublait une grandeur clé.
  \item Rédiger une conclusion courte, accessible à un élève de Terminale.
\end{enumerate}

\begin{corrigebox}[Synthèse rédigée]
Le phénomène « Réflexion totale interne » s'explique par un petit nombre de lois physiques simples, mais sa manifestation réelle dépend souvent de nombreux paramètres. Le modèle Terminale donne l'intuition et les ordres de grandeur ; le modèle Bac+3 introduit les variations locales, les pertes, les équations différentielles ou les principes variationnels nécessaires pour décrire finement la réalité.
\end{corrigebox}

% ------------------------------------------------------------
\subsection{Problème corrigé 35 — Lentilles}
% ------------------------------------------------------------

\begin{exercicebox}[Question-problème]
\textbf{Observation.} Une lentille forme une image.

\medskip
\textbf{Question.} Calculer la position d'une image.

\medskip
\textbf{Thèmes mobilisés :} Conjugaison, foyer, vergence.

\medskip
\textbf{Relation ou idée centrale :}
\[
\frac1{OA'}-\frac1{OA}=\frac1{f'}
\]
\end{exercicebox}

\subsubsection*{Partie A — Modélisation niveau Terminale}

\begin{enumerate}
  \item Décrire précisément le phénomène observé.
  \item Identifier le système physique étudié.
  \item Lister les grandeurs importantes et leurs unités.
  \item Écrire la loi physique simple qui permet de commencer le raisonnement.
  \item Expliquer les hypothèses simplificatrices du modèle Terminale.
\end{enumerate}

\subsubsection*{Partie B — Mise en équation et calcul}

\begin{enumerate}
  \item Définir les variables utiles.
  \item Écrire la relation mathématique principale.
  \item Vérifier l'homogénéité de la formule.
  \item Choisir un ordre de grandeur réaliste pour les données manquantes.
  \item Effectuer le calcul ou établir le raisonnement qualitatif demandé.
  \item Interpréter le résultat en une phrase claire.
\end{enumerate}

\subsubsection*{Partie C — Réponse accessible niveau Terminale}

\begin{corrigebox}[Réponse Terminale]
On part de l'observation concrète : Une lentille forme une image.

Au niveau Terminale, on cherche d'abord la grandeur qui contrôle le phénomène. Ici, l'idée centrale est donnée par :
\[
\frac1{OA'}-\frac1{OA}=\frac1{f'}
\]
Cette relation permet de relier la cause physique à l'effet observé. On vérifie les unités, puis on insère des ordres de grandeur réalistes.

La conclusion attendue est la suivante : Calculer la position d'une image. La réponse doit toujours préciser les hypothèses utilisées. Si le phénomène dépend d'un milieu réel, d'une géométrie complexe, de frottements, de pertes ou de gradients, le résultat obtenu est un ordre de grandeur et non une valeur exacte.
\end{corrigebox}

\subsubsection*{Partie D — Approfondissement niveau Bac+3}

\begin{approfondissementbox}[Réponse Bac+3]
Au niveau Bac+3, on ne se contente plus d'une formule globale. On remplace le modèle simple par une loi locale, différentielle ou énergétique.

Une forme générale possible est :
\[
\frac{\mathrm d X}{\mathrm dt}=\text{apport}-\text{pertes},
\]
ou, pour un phénomène propagatif :
\[
\frac{\partial^2 \psi}{\partial t^2}=c^2\Delta\psi.
\]
Selon le cas, on peut aussi utiliser un principe variationnel, par exemple le principe de Fermat en optique :
\[
\delta \int n\,\mathrm ds=0.
\]

Pour ce phénomène, le modèle Bac+3 consiste à partir de la relation centrale
\[
\frac1{OA'}-\frac1{OA}=\frac1{f'}
\]
puis à introduire les corrections négligées : dépendance spatiale, dépendance temporelle, non-linéarités, pertes, conditions aux limites ou fluctuations. Cela permet de passer d'une réponse d'ordre de grandeur à une prédiction plus robuste.
\end{approfondissementbox}

\subsubsection*{Partie E — Limites, critique et prolongement}

\begin{enumerate}
  \item Citer au moins deux limites du modèle Terminale.
  \item Dire quelles grandeurs seraient difficiles à mesurer expérimentalement.
  \item Proposer une expérience simple permettant de tester le modèle.
  \item Expliquer ce qui changerait si l'on doublait une grandeur clé.
  \item Rédiger une conclusion courte, accessible à un élève de Terminale.
\end{enumerate}

\begin{corrigebox}[Synthèse rédigée]
Le phénomène « Lentilles » s'explique par un petit nombre de lois physiques simples, mais sa manifestation réelle dépend souvent de nombreux paramètres. Le modèle Terminale donne l'intuition et les ordres de grandeur ; le modèle Bac+3 introduit les variations locales, les pertes, les équations différentielles ou les principes variationnels nécessaires pour décrire finement la réalité.
\end{corrigebox}

% ------------------------------------------------------------
\subsection{Problème corrigé 36 — Œil humain}
% ------------------------------------------------------------

\begin{exercicebox}[Question-problème]
\textbf{Observation.} L'œil forme une image sur la rétine.

\medskip
\textbf{Question.} Expliquer le rôle du cristallin.

\medskip
\textbf{Thèmes mobilisés :} Optique physiologique, accommodation.

\medskip
\textbf{Relation ou idée centrale :}
\[
\text{image nette sur la rétine}
\]
\end{exercicebox}

\subsubsection*{Partie A — Modélisation niveau Terminale}

\begin{enumerate}
  \item Décrire précisément le phénomène observé.
  \item Identifier le système physique étudié.
  \item Lister les grandeurs importantes et leurs unités.
  \item Écrire la loi physique simple qui permet de commencer le raisonnement.
  \item Expliquer les hypothèses simplificatrices du modèle Terminale.
\end{enumerate}

\subsubsection*{Partie B — Mise en équation et calcul}

\begin{enumerate}
  \item Définir les variables utiles.
  \item Écrire la relation mathématique principale.
  \item Vérifier l'homogénéité de la formule.
  \item Choisir un ordre de grandeur réaliste pour les données manquantes.
  \item Effectuer le calcul ou établir le raisonnement qualitatif demandé.
  \item Interpréter le résultat en une phrase claire.
\end{enumerate}

\subsubsection*{Partie C — Réponse accessible niveau Terminale}

\begin{corrigebox}[Réponse Terminale]
On part de l'observation concrète : L'œil forme une image sur la rétine.

Au niveau Terminale, on cherche d'abord la grandeur qui contrôle le phénomène. Ici, l'idée centrale est donnée par :
\[
\text{image nette sur la rétine}
\]
Cette relation permet de relier la cause physique à l'effet observé. On vérifie les unités, puis on insère des ordres de grandeur réalistes.

La conclusion attendue est la suivante : Expliquer le rôle du cristallin. La réponse doit toujours préciser les hypothèses utilisées. Si le phénomène dépend d'un milieu réel, d'une géométrie complexe, de frottements, de pertes ou de gradients, le résultat obtenu est un ordre de grandeur et non une valeur exacte.
\end{corrigebox}

\subsubsection*{Partie D — Approfondissement niveau Bac+3}

\begin{approfondissementbox}[Réponse Bac+3]
Au niveau Bac+3, on ne se contente plus d'une formule globale. On remplace le modèle simple par une loi locale, différentielle ou énergétique.

Une forme générale possible est :
\[
\frac{\mathrm d X}{\mathrm dt}=\text{apport}-\text{pertes},
\]
ou, pour un phénomène propagatif :
\[
\frac{\partial^2 \psi}{\partial t^2}=c^2\Delta\psi.
\]
Selon le cas, on peut aussi utiliser un principe variationnel, par exemple le principe de Fermat en optique :
\[
\delta \int n\,\mathrm ds=0.
\]

Pour ce phénomène, le modèle Bac+3 consiste à partir de la relation centrale
\[
\text{image nette sur la rétine}
\]
puis à introduire les corrections négligées : dépendance spatiale, dépendance temporelle, non-linéarités, pertes, conditions aux limites ou fluctuations. Cela permet de passer d'une réponse d'ordre de grandeur à une prédiction plus robuste.
\end{approfondissementbox}

\subsubsection*{Partie E — Limites, critique et prolongement}

\begin{enumerate}
  \item Citer au moins deux limites du modèle Terminale.
  \item Dire quelles grandeurs seraient difficiles à mesurer expérimentalement.
  \item Proposer une expérience simple permettant de tester le modèle.
  \item Expliquer ce qui changerait si l'on doublait une grandeur clé.
  \item Rédiger une conclusion courte, accessible à un élève de Terminale.
\end{enumerate}

\begin{corrigebox}[Synthèse rédigée]
Le phénomène « Œil humain » s'explique par un petit nombre de lois physiques simples, mais sa manifestation réelle dépend souvent de nombreux paramètres. Le modèle Terminale donne l'intuition et les ordres de grandeur ; le modèle Bac+3 introduit les variations locales, les pertes, les équations différentielles ou les principes variationnels nécessaires pour décrire finement la réalité.
\end{corrigebox}

% ------------------------------------------------------------
\subsection{Problème corrigé 37 — Myopie}
% ------------------------------------------------------------

\begin{exercicebox}[Question-problème]
\textbf{Observation.} Un myope voit flou de loin.

\medskip
\textbf{Question.} Expliquer la correction.

\medskip
\textbf{Thèmes mobilisés :} Image avant la rétine.

\medskip
\textbf{Relation ou idée centrale :}
\[
\text{correction par lentille divergente}
\]
\end{exercicebox}

\subsubsection*{Partie A — Modélisation niveau Terminale}

\begin{enumerate}
  \item Décrire précisément le phénomène observé.
  \item Identifier le système physique étudié.
  \item Lister les grandeurs importantes et leurs unités.
  \item Écrire la loi physique simple qui permet de commencer le raisonnement.
  \item Expliquer les hypothèses simplificatrices du modèle Terminale.
\end{enumerate}

\subsubsection*{Partie B — Mise en équation et calcul}

\begin{enumerate}
  \item Définir les variables utiles.
  \item Écrire la relation mathématique principale.
  \item Vérifier l'homogénéité de la formule.
  \item Choisir un ordre de grandeur réaliste pour les données manquantes.
  \item Effectuer le calcul ou établir le raisonnement qualitatif demandé.
  \item Interpréter le résultat en une phrase claire.
\end{enumerate}

\subsubsection*{Partie C — Réponse accessible niveau Terminale}

\begin{corrigebox}[Réponse Terminale]
On part de l'observation concrète : Un myope voit flou de loin.

Au niveau Terminale, on cherche d'abord la grandeur qui contrôle le phénomène. Ici, l'idée centrale est donnée par :
\[
\text{correction par lentille divergente}
\]
Cette relation permet de relier la cause physique à l'effet observé. On vérifie les unités, puis on insère des ordres de grandeur réalistes.

La conclusion attendue est la suivante : Expliquer la correction. La réponse doit toujours préciser les hypothèses utilisées. Si le phénomène dépend d'un milieu réel, d'une géométrie complexe, de frottements, de pertes ou de gradients, le résultat obtenu est un ordre de grandeur et non une valeur exacte.
\end{corrigebox}

\subsubsection*{Partie D — Approfondissement niveau Bac+3}

\begin{approfondissementbox}[Réponse Bac+3]
Au niveau Bac+3, on ne se contente plus d'une formule globale. On remplace le modèle simple par une loi locale, différentielle ou énergétique.

Une forme générale possible est :
\[
\frac{\mathrm d X}{\mathrm dt}=\text{apport}-\text{pertes},
\]
ou, pour un phénomène propagatif :
\[
\frac{\partial^2 \psi}{\partial t^2}=c^2\Delta\psi.
\]
Selon le cas, on peut aussi utiliser un principe variationnel, par exemple le principe de Fermat en optique :
\[
\delta \int n\,\mathrm ds=0.
\]

Pour ce phénomène, le modèle Bac+3 consiste à partir de la relation centrale
\[
\text{correction par lentille divergente}
\]
puis à introduire les corrections négligées : dépendance spatiale, dépendance temporelle, non-linéarités, pertes, conditions aux limites ou fluctuations. Cela permet de passer d'une réponse d'ordre de grandeur à une prédiction plus robuste.
\end{approfondissementbox}

\subsubsection*{Partie E — Limites, critique et prolongement}

\begin{enumerate}
  \item Citer au moins deux limites du modèle Terminale.
  \item Dire quelles grandeurs seraient difficiles à mesurer expérimentalement.
  \item Proposer une expérience simple permettant de tester le modèle.
  \item Expliquer ce qui changerait si l'on doublait une grandeur clé.
  \item Rédiger une conclusion courte, accessible à un élève de Terminale.
\end{enumerate}

\begin{corrigebox}[Synthèse rédigée]
Le phénomène « Myopie » s'explique par un petit nombre de lois physiques simples, mais sa manifestation réelle dépend souvent de nombreux paramètres. Le modèle Terminale donne l'intuition et les ordres de grandeur ; le modèle Bac+3 introduit les variations locales, les pertes, les équations différentielles ou les principes variationnels nécessaires pour décrire finement la réalité.
\end{corrigebox}

% ------------------------------------------------------------
\subsection{Problème corrigé 38 — Hypermétropie}
% ------------------------------------------------------------

\begin{exercicebox}[Question-problème]
\textbf{Observation.} Un hypermétrope fatigue en vision proche.

\medskip
\textbf{Question.} Expliquer la correction.

\medskip
\textbf{Thèmes mobilisés :} Image derrière la rétine.

\medskip
\textbf{Relation ou idée centrale :}
\[
\text{correction par lentille convergente}
\]
\end{exercicebox}

\subsubsection*{Partie A — Modélisation niveau Terminale}

\begin{enumerate}
  \item Décrire précisément le phénomène observé.
  \item Identifier le système physique étudié.
  \item Lister les grandeurs importantes et leurs unités.
  \item Écrire la loi physique simple qui permet de commencer le raisonnement.
  \item Expliquer les hypothèses simplificatrices du modèle Terminale.
\end{enumerate}

\subsubsection*{Partie B — Mise en équation et calcul}

\begin{enumerate}
  \item Définir les variables utiles.
  \item Écrire la relation mathématique principale.
  \item Vérifier l'homogénéité de la formule.
  \item Choisir un ordre de grandeur réaliste pour les données manquantes.
  \item Effectuer le calcul ou établir le raisonnement qualitatif demandé.
  \item Interpréter le résultat en une phrase claire.
\end{enumerate}

\subsubsection*{Partie C — Réponse accessible niveau Terminale}

\begin{corrigebox}[Réponse Terminale]
On part de l'observation concrète : Un hypermétrope fatigue en vision proche.

Au niveau Terminale, on cherche d'abord la grandeur qui contrôle le phénomène. Ici, l'idée centrale est donnée par :
\[
\text{correction par lentille convergente}
\]
Cette relation permet de relier la cause physique à l'effet observé. On vérifie les unités, puis on insère des ordres de grandeur réalistes.

La conclusion attendue est la suivante : Expliquer la correction. La réponse doit toujours préciser les hypothèses utilisées. Si le phénomène dépend d'un milieu réel, d'une géométrie complexe, de frottements, de pertes ou de gradients, le résultat obtenu est un ordre de grandeur et non une valeur exacte.
\end{corrigebox}

\subsubsection*{Partie D — Approfondissement niveau Bac+3}

\begin{approfondissementbox}[Réponse Bac+3]
Au niveau Bac+3, on ne se contente plus d'une formule globale. On remplace le modèle simple par une loi locale, différentielle ou énergétique.

Une forme générale possible est :
\[
\frac{\mathrm d X}{\mathrm dt}=\text{apport}-\text{pertes},
\]
ou, pour un phénomène propagatif :
\[
\frac{\partial^2 \psi}{\partial t^2}=c^2\Delta\psi.
\]
Selon le cas, on peut aussi utiliser un principe variationnel, par exemple le principe de Fermat en optique :
\[
\delta \int n\,\mathrm ds=0.
\]

Pour ce phénomène, le modèle Bac+3 consiste à partir de la relation centrale
\[
\text{correction par lentille convergente}
\]
puis à introduire les corrections négligées : dépendance spatiale, dépendance temporelle, non-linéarités, pertes, conditions aux limites ou fluctuations. Cela permet de passer d'une réponse d'ordre de grandeur à une prédiction plus robuste.
\end{approfondissementbox}

\subsubsection*{Partie E — Limites, critique et prolongement}

\begin{enumerate}
  \item Citer au moins deux limites du modèle Terminale.
  \item Dire quelles grandeurs seraient difficiles à mesurer expérimentalement.
  \item Proposer une expérience simple permettant de tester le modèle.
  \item Expliquer ce qui changerait si l'on doublait une grandeur clé.
  \item Rédiger une conclusion courte, accessible à un élève de Terminale.
\end{enumerate}

\begin{corrigebox}[Synthèse rédigée]
Le phénomène « Hypermétropie » s'explique par un petit nombre de lois physiques simples, mais sa manifestation réelle dépend souvent de nombreux paramètres. Le modèle Terminale donne l'intuition et les ordres de grandeur ; le modèle Bac+3 introduit les variations locales, les pertes, les équations différentielles ou les principes variationnels nécessaires pour décrire finement la réalité.
\end{corrigebox}

% ------------------------------------------------------------
\subsection{Problème corrigé 39 — Appareil photo}
% ------------------------------------------------------------

\begin{exercicebox}[Question-problème]
\textbf{Observation.} Un objectif forme une image sur un capteur.

\medskip
\textbf{Question.} Relier mise au point et conjugaison.

\medskip
\textbf{Thèmes mobilisés :} Lentille, diaphragme, profondeur de champ.

\medskip
\textbf{Relation ou idée centrale :}
\[
\frac1{p'}-\frac1p=\frac1f
\]
\end{exercicebox}

\subsubsection*{Partie A — Modélisation niveau Terminale}

\begin{enumerate}
  \item Décrire précisément le phénomène observé.
  \item Identifier le système physique étudié.
  \item Lister les grandeurs importantes et leurs unités.
  \item Écrire la loi physique simple qui permet de commencer le raisonnement.
  \item Expliquer les hypothèses simplificatrices du modèle Terminale.
\end{enumerate}

\subsubsection*{Partie B — Mise en équation et calcul}

\begin{enumerate}
  \item Définir les variables utiles.
  \item Écrire la relation mathématique principale.
  \item Vérifier l'homogénéité de la formule.
  \item Choisir un ordre de grandeur réaliste pour les données manquantes.
  \item Effectuer le calcul ou établir le raisonnement qualitatif demandé.
  \item Interpréter le résultat en une phrase claire.
\end{enumerate}

\subsubsection*{Partie C — Réponse accessible niveau Terminale}

\begin{corrigebox}[Réponse Terminale]
On part de l'observation concrète : Un objectif forme une image sur un capteur.

Au niveau Terminale, on cherche d'abord la grandeur qui contrôle le phénomène. Ici, l'idée centrale est donnée par :
\[
\frac1{p'}-\frac1p=\frac1f
\]
Cette relation permet de relier la cause physique à l'effet observé. On vérifie les unités, puis on insère des ordres de grandeur réalistes.

La conclusion attendue est la suivante : Relier mise au point et conjugaison. La réponse doit toujours préciser les hypothèses utilisées. Si le phénomène dépend d'un milieu réel, d'une géométrie complexe, de frottements, de pertes ou de gradients, le résultat obtenu est un ordre de grandeur et non une valeur exacte.
\end{corrigebox}

\subsubsection*{Partie D — Approfondissement niveau Bac+3}

\begin{approfondissementbox}[Réponse Bac+3]
Au niveau Bac+3, on ne se contente plus d'une formule globale. On remplace le modèle simple par une loi locale, différentielle ou énergétique.

Une forme générale possible est :
\[
\frac{\mathrm d X}{\mathrm dt}=\text{apport}-\text{pertes},
\]
ou, pour un phénomène propagatif :
\[
\frac{\partial^2 \psi}{\partial t^2}=c^2\Delta\psi.
\]
Selon le cas, on peut aussi utiliser un principe variationnel, par exemple le principe de Fermat en optique :
\[
\delta \int n\,\mathrm ds=0.
\]

Pour ce phénomène, le modèle Bac+3 consiste à partir de la relation centrale
\[
\frac1{p'}-\frac1p=\frac1f
\]
puis à introduire les corrections négligées : dépendance spatiale, dépendance temporelle, non-linéarités, pertes, conditions aux limites ou fluctuations. Cela permet de passer d'une réponse d'ordre de grandeur à une prédiction plus robuste.
\end{approfondissementbox}

\subsubsection*{Partie E — Limites, critique et prolongement}

\begin{enumerate}
  \item Citer au moins deux limites du modèle Terminale.
  \item Dire quelles grandeurs seraient difficiles à mesurer expérimentalement.
  \item Proposer une expérience simple permettant de tester le modèle.
  \item Expliquer ce qui changerait si l'on doublait une grandeur clé.
  \item Rédiger une conclusion courte, accessible à un élève de Terminale.
\end{enumerate}

\begin{corrigebox}[Synthèse rédigée]
Le phénomène « Appareil photo » s'explique par un petit nombre de lois physiques simples, mais sa manifestation réelle dépend souvent de nombreux paramètres. Le modèle Terminale donne l'intuition et les ordres de grandeur ; le modèle Bac+3 introduit les variations locales, les pertes, les équations différentielles ou les principes variationnels nécessaires pour décrire finement la réalité.
\end{corrigebox}

% ------------------------------------------------------------
\subsection{Problème corrigé 40 — Loupe}
% ------------------------------------------------------------

\begin{exercicebox}[Question-problème]
\textbf{Observation.} Une loupe agrandit un objet proche.

\medskip
\textbf{Question.} Comprendre l'image virtuelle agrandie.

\medskip
\textbf{Thèmes mobilisés :} Grossissement angulaire.

\medskip
\textbf{Relation ou idée centrale :}
\[
G\simeq\frac{25\,\mathrm{cm}}{f'}
\]
\end{exercicebox}

\subsubsection*{Partie A — Modélisation niveau Terminale}

\begin{enumerate}
  \item Décrire précisément le phénomène observé.
  \item Identifier le système physique étudié.
  \item Lister les grandeurs importantes et leurs unités.
  \item Écrire la loi physique simple qui permet de commencer le raisonnement.
  \item Expliquer les hypothèses simplificatrices du modèle Terminale.
\end{enumerate}

\subsubsection*{Partie B — Mise en équation et calcul}

\begin{enumerate}
  \item Définir les variables utiles.
  \item Écrire la relation mathématique principale.
  \item Vérifier l'homogénéité de la formule.
  \item Choisir un ordre de grandeur réaliste pour les données manquantes.
  \item Effectuer le calcul ou établir le raisonnement qualitatif demandé.
  \item Interpréter le résultat en une phrase claire.
\end{enumerate}

\subsubsection*{Partie C — Réponse accessible niveau Terminale}

\begin{corrigebox}[Réponse Terminale]
On part de l'observation concrète : Une loupe agrandit un objet proche.

Au niveau Terminale, on cherche d'abord la grandeur qui contrôle le phénomène. Ici, l'idée centrale est donnée par :
\[
G\simeq\frac{25\,\mathrm{cm}}{f'}
\]
Cette relation permet de relier la cause physique à l'effet observé. On vérifie les unités, puis on insère des ordres de grandeur réalistes.

La conclusion attendue est la suivante : Comprendre l'image virtuelle agrandie. La réponse doit toujours préciser les hypothèses utilisées. Si le phénomène dépend d'un milieu réel, d'une géométrie complexe, de frottements, de pertes ou de gradients, le résultat obtenu est un ordre de grandeur et non une valeur exacte.
\end{corrigebox}

\subsubsection*{Partie D — Approfondissement niveau Bac+3}

\begin{approfondissementbox}[Réponse Bac+3]
Au niveau Bac+3, on ne se contente plus d'une formule globale. On remplace le modèle simple par une loi locale, différentielle ou énergétique.

Une forme générale possible est :
\[
\frac{\mathrm d X}{\mathrm dt}=\text{apport}-\text{pertes},
\]
ou, pour un phénomène propagatif :
\[
\frac{\partial^2 \psi}{\partial t^2}=c^2\Delta\psi.
\]
Selon le cas, on peut aussi utiliser un principe variationnel, par exemple le principe de Fermat en optique :
\[
\delta \int n\,\mathrm ds=0.
\]

Pour ce phénomène, le modèle Bac+3 consiste à partir de la relation centrale
\[
G\simeq\frac{25\,\mathrm{cm}}{f'}
\]
puis à introduire les corrections négligées : dépendance spatiale, dépendance temporelle, non-linéarités, pertes, conditions aux limites ou fluctuations. Cela permet de passer d'une réponse d'ordre de grandeur à une prédiction plus robuste.
\end{approfondissementbox}

\subsubsection*{Partie E — Limites, critique et prolongement}

\begin{enumerate}
  \item Citer au moins deux limites du modèle Terminale.
  \item Dire quelles grandeurs seraient difficiles à mesurer expérimentalement.
  \item Proposer une expérience simple permettant de tester le modèle.
  \item Expliquer ce qui changerait si l'on doublait une grandeur clé.
  \item Rédiger une conclusion courte, accessible à un élève de Terminale.
\end{enumerate}

\begin{corrigebox}[Synthèse rédigée]
Le phénomène « Loupe » s'explique par un petit nombre de lois physiques simples, mais sa manifestation réelle dépend souvent de nombreux paramètres. Le modèle Terminale donne l'intuition et les ordres de grandeur ; le modèle Bac+3 introduit les variations locales, les pertes, les équations différentielles ou les principes variationnels nécessaires pour décrire finement la réalité.
\end{corrigebox}

% ------------------------------------------------------------
\subsection{Problème corrigé 41 — Microscope}
% ------------------------------------------------------------

\begin{exercicebox}[Question-problème]
\textbf{Observation.} Un microscope observe des objets très petits.

\medskip
\textbf{Question.} Relier grossissement et résolution.

\medskip
\textbf{Thèmes mobilisés :} Objectif, oculaire, diffraction.

\medskip
\textbf{Relation ou idée centrale :}
\[
G=G_{\mathrm{obj}}G_{\mathrm{oc}}
\]
\end{exercicebox}

\subsubsection*{Partie A — Modélisation niveau Terminale}

\begin{enumerate}
  \item Décrire précisément le phénomène observé.
  \item Identifier le système physique étudié.
  \item Lister les grandeurs importantes et leurs unités.
  \item Écrire la loi physique simple qui permet de commencer le raisonnement.
  \item Expliquer les hypothèses simplificatrices du modèle Terminale.
\end{enumerate}

\subsubsection*{Partie B — Mise en équation et calcul}

\begin{enumerate}
  \item Définir les variables utiles.
  \item Écrire la relation mathématique principale.
  \item Vérifier l'homogénéité de la formule.
  \item Choisir un ordre de grandeur réaliste pour les données manquantes.
  \item Effectuer le calcul ou établir le raisonnement qualitatif demandé.
  \item Interpréter le résultat en une phrase claire.
\end{enumerate}

\subsubsection*{Partie C — Réponse accessible niveau Terminale}

\begin{corrigebox}[Réponse Terminale]
On part de l'observation concrète : Un microscope observe des objets très petits.

Au niveau Terminale, on cherche d'abord la grandeur qui contrôle le phénomène. Ici, l'idée centrale est donnée par :
\[
G=G_{\mathrm{obj}}G_{\mathrm{oc}}
\]
Cette relation permet de relier la cause physique à l'effet observé. On vérifie les unités, puis on insère des ordres de grandeur réalistes.

La conclusion attendue est la suivante : Relier grossissement et résolution. La réponse doit toujours préciser les hypothèses utilisées. Si le phénomène dépend d'un milieu réel, d'une géométrie complexe, de frottements, de pertes ou de gradients, le résultat obtenu est un ordre de grandeur et non une valeur exacte.
\end{corrigebox}

\subsubsection*{Partie D — Approfondissement niveau Bac+3}

\begin{approfondissementbox}[Réponse Bac+3]
Au niveau Bac+3, on ne se contente plus d'une formule globale. On remplace le modèle simple par une loi locale, différentielle ou énergétique.

Une forme générale possible est :
\[
\frac{\mathrm d X}{\mathrm dt}=\text{apport}-\text{pertes},
\]
ou, pour un phénomène propagatif :
\[
\frac{\partial^2 \psi}{\partial t^2}=c^2\Delta\psi.
\]
Selon le cas, on peut aussi utiliser un principe variationnel, par exemple le principe de Fermat en optique :
\[
\delta \int n\,\mathrm ds=0.
\]

Pour ce phénomène, le modèle Bac+3 consiste à partir de la relation centrale
\[
G=G_{\mathrm{obj}}G_{\mathrm{oc}}
\]
puis à introduire les corrections négligées : dépendance spatiale, dépendance temporelle, non-linéarités, pertes, conditions aux limites ou fluctuations. Cela permet de passer d'une réponse d'ordre de grandeur à une prédiction plus robuste.
\end{approfondissementbox}

\subsubsection*{Partie E — Limites, critique et prolongement}

\begin{enumerate}
  \item Citer au moins deux limites du modèle Terminale.
  \item Dire quelles grandeurs seraient difficiles à mesurer expérimentalement.
  \item Proposer une expérience simple permettant de tester le modèle.
  \item Expliquer ce qui changerait si l'on doublait une grandeur clé.
  \item Rédiger une conclusion courte, accessible à un élève de Terminale.
\end{enumerate}

\begin{corrigebox}[Synthèse rédigée]
Le phénomène « Microscope » s'explique par un petit nombre de lois physiques simples, mais sa manifestation réelle dépend souvent de nombreux paramètres. Le modèle Terminale donne l'intuition et les ordres de grandeur ; le modèle Bac+3 introduit les variations locales, les pertes, les équations différentielles ou les principes variationnels nécessaires pour décrire finement la réalité.
\end{corrigebox}

% ------------------------------------------------------------
\subsection{Problème corrigé 42 — Télescope}
% ------------------------------------------------------------

\begin{exercicebox}[Question-problème]
\textbf{Observation.} Un télescope observe des objets lointains.

\medskip
\textbf{Question.} Relier diamètre et résolution.

\medskip
\textbf{Thèmes mobilisés :} Collecte de lumière, résolution.

\medskip
\textbf{Relation ou idée centrale :}
\[
\theta_{\min}\simeq1.22\frac{\lambda}{D}
\]
\end{exercicebox}

\subsubsection*{Partie A — Modélisation niveau Terminale}

\begin{enumerate}
  \item Décrire précisément le phénomène observé.
  \item Identifier le système physique étudié.
  \item Lister les grandeurs importantes et leurs unités.
  \item Écrire la loi physique simple qui permet de commencer le raisonnement.
  \item Expliquer les hypothèses simplificatrices du modèle Terminale.
\end{enumerate}

\subsubsection*{Partie B — Mise en équation et calcul}

\begin{enumerate}
  \item Définir les variables utiles.
  \item Écrire la relation mathématique principale.
  \item Vérifier l'homogénéité de la formule.
  \item Choisir un ordre de grandeur réaliste pour les données manquantes.
  \item Effectuer le calcul ou établir le raisonnement qualitatif demandé.
  \item Interpréter le résultat en une phrase claire.
\end{enumerate}

\subsubsection*{Partie C — Réponse accessible niveau Terminale}

\begin{corrigebox}[Réponse Terminale]
On part de l'observation concrète : Un télescope observe des objets lointains.

Au niveau Terminale, on cherche d'abord la grandeur qui contrôle le phénomène. Ici, l'idée centrale est donnée par :
\[
\theta_{\min}\simeq1.22\frac{\lambda}{D}
\]
Cette relation permet de relier la cause physique à l'effet observé. On vérifie les unités, puis on insère des ordres de grandeur réalistes.

La conclusion attendue est la suivante : Relier diamètre et résolution. La réponse doit toujours préciser les hypothèses utilisées. Si le phénomène dépend d'un milieu réel, d'une géométrie complexe, de frottements, de pertes ou de gradients, le résultat obtenu est un ordre de grandeur et non une valeur exacte.
\end{corrigebox}

\subsubsection*{Partie D — Approfondissement niveau Bac+3}

\begin{approfondissementbox}[Réponse Bac+3]
Au niveau Bac+3, on ne se contente plus d'une formule globale. On remplace le modèle simple par une loi locale, différentielle ou énergétique.

Une forme générale possible est :
\[
\frac{\mathrm d X}{\mathrm dt}=\text{apport}-\text{pertes},
\]
ou, pour un phénomène propagatif :
\[
\frac{\partial^2 \psi}{\partial t^2}=c^2\Delta\psi.
\]
Selon le cas, on peut aussi utiliser un principe variationnel, par exemple le principe de Fermat en optique :
\[
\delta \int n\,\mathrm ds=0.
\]

Pour ce phénomène, le modèle Bac+3 consiste à partir de la relation centrale
\[
\theta_{\min}\simeq1.22\frac{\lambda}{D}
\]
puis à introduire les corrections négligées : dépendance spatiale, dépendance temporelle, non-linéarités, pertes, conditions aux limites ou fluctuations. Cela permet de passer d'une réponse d'ordre de grandeur à une prédiction plus robuste.
\end{approfondissementbox}

\subsubsection*{Partie E — Limites, critique et prolongement}

\begin{enumerate}
  \item Citer au moins deux limites du modèle Terminale.
  \item Dire quelles grandeurs seraient difficiles à mesurer expérimentalement.
  \item Proposer une expérience simple permettant de tester le modèle.
  \item Expliquer ce qui changerait si l'on doublait une grandeur clé.
  \item Rédiger une conclusion courte, accessible à un élève de Terminale.
\end{enumerate}

\begin{corrigebox}[Synthèse rédigée]
Le phénomène « Télescope » s'explique par un petit nombre de lois physiques simples, mais sa manifestation réelle dépend souvent de nombreux paramètres. Le modèle Terminale donne l'intuition et les ordres de grandeur ; le modèle Bac+3 introduit les variations locales, les pertes, les équations différentielles ou les principes variationnels nécessaires pour décrire finement la réalité.
\end{corrigebox}

% ------------------------------------------------------------
\subsection{Problème corrigé 43 — Diffraction par une fente}
% ------------------------------------------------------------

\begin{exercicebox}[Question-problème]
\textbf{Observation.} Une onde s'élargit après une ouverture.

\medskip
\textbf{Question.} Montrer pourquoi une petite fente élargit le faisceau.

\medskip
\textbf{Thèmes mobilisés :} Diffraction.

\medskip
\textbf{Relation ou idée centrale :}
\[
\theta\simeq\frac{\lambda}{a}
\]
\end{exercicebox}

\subsubsection*{Partie A — Modélisation niveau Terminale}

\begin{enumerate}
  \item Décrire précisément le phénomène observé.
  \item Identifier le système physique étudié.
  \item Lister les grandeurs importantes et leurs unités.
  \item Écrire la loi physique simple qui permet de commencer le raisonnement.
  \item Expliquer les hypothèses simplificatrices du modèle Terminale.
\end{enumerate}

\subsubsection*{Partie B — Mise en équation et calcul}

\begin{enumerate}
  \item Définir les variables utiles.
  \item Écrire la relation mathématique principale.
  \item Vérifier l'homogénéité de la formule.
  \item Choisir un ordre de grandeur réaliste pour les données manquantes.
  \item Effectuer le calcul ou établir le raisonnement qualitatif demandé.
  \item Interpréter le résultat en une phrase claire.
\end{enumerate}

\subsubsection*{Partie C — Réponse accessible niveau Terminale}

\begin{corrigebox}[Réponse Terminale]
On part de l'observation concrète : Une onde s'élargit après une ouverture.

Au niveau Terminale, on cherche d'abord la grandeur qui contrôle le phénomène. Ici, l'idée centrale est donnée par :
\[
\theta\simeq\frac{\lambda}{a}
\]
Cette relation permet de relier la cause physique à l'effet observé. On vérifie les unités, puis on insère des ordres de grandeur réalistes.

La conclusion attendue est la suivante : Montrer pourquoi une petite fente élargit le faisceau. La réponse doit toujours préciser les hypothèses utilisées. Si le phénomène dépend d'un milieu réel, d'une géométrie complexe, de frottements, de pertes ou de gradients, le résultat obtenu est un ordre de grandeur et non une valeur exacte.
\end{corrigebox}

\subsubsection*{Partie D — Approfondissement niveau Bac+3}

\begin{approfondissementbox}[Réponse Bac+3]
Au niveau Bac+3, on ne se contente plus d'une formule globale. On remplace le modèle simple par une loi locale, différentielle ou énergétique.

Une forme générale possible est :
\[
\frac{\mathrm d X}{\mathrm dt}=\text{apport}-\text{pertes},
\]
ou, pour un phénomène propagatif :
\[
\frac{\partial^2 \psi}{\partial t^2}=c^2\Delta\psi.
\]
Selon le cas, on peut aussi utiliser un principe variationnel, par exemple le principe de Fermat en optique :
\[
\delta \int n\,\mathrm ds=0.
\]

Pour ce phénomène, le modèle Bac+3 consiste à partir de la relation centrale
\[
\theta\simeq\frac{\lambda}{a}
\]
puis à introduire les corrections négligées : dépendance spatiale, dépendance temporelle, non-linéarités, pertes, conditions aux limites ou fluctuations. Cela permet de passer d'une réponse d'ordre de grandeur à une prédiction plus robuste.
\end{approfondissementbox}

\subsubsection*{Partie E — Limites, critique et prolongement}

\begin{enumerate}
  \item Citer au moins deux limites du modèle Terminale.
  \item Dire quelles grandeurs seraient difficiles à mesurer expérimentalement.
  \item Proposer une expérience simple permettant de tester le modèle.
  \item Expliquer ce qui changerait si l'on doublait une grandeur clé.
  \item Rédiger une conclusion courte, accessible à un élève de Terminale.
\end{enumerate}

\begin{corrigebox}[Synthèse rédigée]
Le phénomène « Diffraction par une fente » s'explique par un petit nombre de lois physiques simples, mais sa manifestation réelle dépend souvent de nombreux paramètres. Le modèle Terminale donne l'intuition et les ordres de grandeur ; le modèle Bac+3 introduit les variations locales, les pertes, les équations différentielles ou les principes variationnels nécessaires pour décrire finement la réalité.
\end{corrigebox}

% ------------------------------------------------------------
\subsection{Problème corrigé 44 — Mesure d'un cheveu par diffraction}
% ------------------------------------------------------------

\begin{exercicebox}[Question-problème]
\textbf{Observation.} Un cheveu produit une figure de diffraction.

\medskip
\textbf{Question.} Mesurer une petite dimension optiquement.

\medskip
\textbf{Thèmes mobilisés :} Babinet, laser, diffraction.

\medskip
\textbf{Relation ou idée centrale :}
\[
a\simeq\frac{2\lambda D}{L}
\]
\end{exercicebox}

\subsubsection*{Partie A — Modélisation niveau Terminale}

\begin{enumerate}
  \item Décrire précisément le phénomène observé.
  \item Identifier le système physique étudié.
  \item Lister les grandeurs importantes et leurs unités.
  \item Écrire la loi physique simple qui permet de commencer le raisonnement.
  \item Expliquer les hypothèses simplificatrices du modèle Terminale.
\end{enumerate}

\subsubsection*{Partie B — Mise en équation et calcul}

\begin{enumerate}
  \item Définir les variables utiles.
  \item Écrire la relation mathématique principale.
  \item Vérifier l'homogénéité de la formule.
  \item Choisir un ordre de grandeur réaliste pour les données manquantes.
  \item Effectuer le calcul ou établir le raisonnement qualitatif demandé.
  \item Interpréter le résultat en une phrase claire.
\end{enumerate}

\subsubsection*{Partie C — Réponse accessible niveau Terminale}

\begin{corrigebox}[Réponse Terminale]
On part de l'observation concrète : Un cheveu produit une figure de diffraction.

Au niveau Terminale, on cherche d'abord la grandeur qui contrôle le phénomène. Ici, l'idée centrale est donnée par :
\[
a\simeq\frac{2\lambda D}{L}
\]
Cette relation permet de relier la cause physique à l'effet observé. On vérifie les unités, puis on insère des ordres de grandeur réalistes.

La conclusion attendue est la suivante : Mesurer une petite dimension optiquement. La réponse doit toujours préciser les hypothèses utilisées. Si le phénomène dépend d'un milieu réel, d'une géométrie complexe, de frottements, de pertes ou de gradients, le résultat obtenu est un ordre de grandeur et non une valeur exacte.
\end{corrigebox}

\subsubsection*{Partie D — Approfondissement niveau Bac+3}

\begin{approfondissementbox}[Réponse Bac+3]
Au niveau Bac+3, on ne se contente plus d'une formule globale. On remplace le modèle simple par une loi locale, différentielle ou énergétique.

Une forme générale possible est :
\[
\frac{\mathrm d X}{\mathrm dt}=\text{apport}-\text{pertes},
\]
ou, pour un phénomène propagatif :
\[
\frac{\partial^2 \psi}{\partial t^2}=c^2\Delta\psi.
\]
Selon le cas, on peut aussi utiliser un principe variationnel, par exemple le principe de Fermat en optique :
\[
\delta \int n\,\mathrm ds=0.
\]

Pour ce phénomène, le modèle Bac+3 consiste à partir de la relation centrale
\[
a\simeq\frac{2\lambda D}{L}
\]
puis à introduire les corrections négligées : dépendance spatiale, dépendance temporelle, non-linéarités, pertes, conditions aux limites ou fluctuations. Cela permet de passer d'une réponse d'ordre de grandeur à une prédiction plus robuste.
\end{approfondissementbox}

\subsubsection*{Partie E — Limites, critique et prolongement}

\begin{enumerate}
  \item Citer au moins deux limites du modèle Terminale.
  \item Dire quelles grandeurs seraient difficiles à mesurer expérimentalement.
  \item Proposer une expérience simple permettant de tester le modèle.
  \item Expliquer ce qui changerait si l'on doublait une grandeur clé.
  \item Rédiger une conclusion courte, accessible à un élève de Terminale.
\end{enumerate}

\begin{corrigebox}[Synthèse rédigée]
Le phénomène « Mesure d'un cheveu par diffraction » s'explique par un petit nombre de lois physiques simples, mais sa manifestation réelle dépend souvent de nombreux paramètres. Le modèle Terminale donne l'intuition et les ordres de grandeur ; le modèle Bac+3 introduit les variations locales, les pertes, les équations différentielles ou les principes variationnels nécessaires pour décrire finement la réalité.
\end{corrigebox}

% ------------------------------------------------------------
\subsection{Problème corrigé 45 — Fentes d'Young}
% ------------------------------------------------------------

\begin{exercicebox}[Question-problème]
\textbf{Observation.} Des franges claires et sombres apparaissent.

\medskip
\textbf{Question.} Calculer un interfrange.

\medskip
\textbf{Thèmes mobilisés :} Interférences.

\medskip
\textbf{Relation ou idée centrale :}
\[
i=\frac{\lambda D}{a}
\]
\end{exercicebox}

\subsubsection*{Partie A — Modélisation niveau Terminale}

\begin{enumerate}
  \item Décrire précisément le phénomène observé.
  \item Identifier le système physique étudié.
  \item Lister les grandeurs importantes et leurs unités.
  \item Écrire la loi physique simple qui permet de commencer le raisonnement.
  \item Expliquer les hypothèses simplificatrices du modèle Terminale.
\end{enumerate}

\subsubsection*{Partie B — Mise en équation et calcul}

\begin{enumerate}
  \item Définir les variables utiles.
  \item Écrire la relation mathématique principale.
  \item Vérifier l'homogénéité de la formule.
  \item Choisir un ordre de grandeur réaliste pour les données manquantes.
  \item Effectuer le calcul ou établir le raisonnement qualitatif demandé.
  \item Interpréter le résultat en une phrase claire.
\end{enumerate}

\subsubsection*{Partie C — Réponse accessible niveau Terminale}

\begin{corrigebox}[Réponse Terminale]
On part de l'observation concrète : Des franges claires et sombres apparaissent.

Au niveau Terminale, on cherche d'abord la grandeur qui contrôle le phénomène. Ici, l'idée centrale est donnée par :
\[
i=\frac{\lambda D}{a}
\]
Cette relation permet de relier la cause physique à l'effet observé. On vérifie les unités, puis on insère des ordres de grandeur réalistes.

La conclusion attendue est la suivante : Calculer un interfrange. La réponse doit toujours préciser les hypothèses utilisées. Si le phénomène dépend d'un milieu réel, d'une géométrie complexe, de frottements, de pertes ou de gradients, le résultat obtenu est un ordre de grandeur et non une valeur exacte.
\end{corrigebox}

\subsubsection*{Partie D — Approfondissement niveau Bac+3}

\begin{approfondissementbox}[Réponse Bac+3]
Au niveau Bac+3, on ne se contente plus d'une formule globale. On remplace le modèle simple par une loi locale, différentielle ou énergétique.

Une forme générale possible est :
\[
\frac{\mathrm d X}{\mathrm dt}=\text{apport}-\text{pertes},
\]
ou, pour un phénomène propagatif :
\[
\frac{\partial^2 \psi}{\partial t^2}=c^2\Delta\psi.
\]
Selon le cas, on peut aussi utiliser un principe variationnel, par exemple le principe de Fermat en optique :
\[
\delta \int n\,\mathrm ds=0.
\]

Pour ce phénomène, le modèle Bac+3 consiste à partir de la relation centrale
\[
i=\frac{\lambda D}{a}
\]
puis à introduire les corrections négligées : dépendance spatiale, dépendance temporelle, non-linéarités, pertes, conditions aux limites ou fluctuations. Cela permet de passer d'une réponse d'ordre de grandeur à une prédiction plus robuste.
\end{approfondissementbox}

\subsubsection*{Partie E — Limites, critique et prolongement}

\begin{enumerate}
  \item Citer au moins deux limites du modèle Terminale.
  \item Dire quelles grandeurs seraient difficiles à mesurer expérimentalement.
  \item Proposer une expérience simple permettant de tester le modèle.
  \item Expliquer ce qui changerait si l'on doublait une grandeur clé.
  \item Rédiger une conclusion courte, accessible à un élève de Terminale.
\end{enumerate}

\begin{corrigebox}[Synthèse rédigée]
Le phénomène « Fentes d'Young » s'explique par un petit nombre de lois physiques simples, mais sa manifestation réelle dépend souvent de nombreux paramètres. Le modèle Terminale donne l'intuition et les ordres de grandeur ; le modèle Bac+3 introduit les variations locales, les pertes, les équations différentielles ou les principes variationnels nécessaires pour décrire finement la réalité.
\end{corrigebox}

% ------------------------------------------------------------
\subsection{Problème corrigé 46 — Bulles de savon}
% ------------------------------------------------------------

\begin{exercicebox}[Question-problème]
\textbf{Observation.} Les bulles présentent des couleurs.

\medskip
\textbf{Question.} Relier couleur et épaisseur.

\medskip
\textbf{Thèmes mobilisés :} Interférences en couche mince.

\medskip
\textbf{Relation ou idée centrale :}
\[
2ne\simeq k\lambda
\]
\end{exercicebox}

\subsubsection*{Partie A — Modélisation niveau Terminale}

\begin{enumerate}
  \item Décrire précisément le phénomène observé.
  \item Identifier le système physique étudié.
  \item Lister les grandeurs importantes et leurs unités.
  \item Écrire la loi physique simple qui permet de commencer le raisonnement.
  \item Expliquer les hypothèses simplificatrices du modèle Terminale.
\end{enumerate}

\subsubsection*{Partie B — Mise en équation et calcul}

\begin{enumerate}
  \item Définir les variables utiles.
  \item Écrire la relation mathématique principale.
  \item Vérifier l'homogénéité de la formule.
  \item Choisir un ordre de grandeur réaliste pour les données manquantes.
  \item Effectuer le calcul ou établir le raisonnement qualitatif demandé.
  \item Interpréter le résultat en une phrase claire.
\end{enumerate}

\subsubsection*{Partie C — Réponse accessible niveau Terminale}

\begin{corrigebox}[Réponse Terminale]
On part de l'observation concrète : Les bulles présentent des couleurs.

Au niveau Terminale, on cherche d'abord la grandeur qui contrôle le phénomène. Ici, l'idée centrale est donnée par :
\[
2ne\simeq k\lambda
\]
Cette relation permet de relier la cause physique à l'effet observé. On vérifie les unités, puis on insère des ordres de grandeur réalistes.

La conclusion attendue est la suivante : Relier couleur et épaisseur. La réponse doit toujours préciser les hypothèses utilisées. Si le phénomène dépend d'un milieu réel, d'une géométrie complexe, de frottements, de pertes ou de gradients, le résultat obtenu est un ordre de grandeur et non une valeur exacte.
\end{corrigebox}

\subsubsection*{Partie D — Approfondissement niveau Bac+3}

\begin{approfondissementbox}[Réponse Bac+3]
Au niveau Bac+3, on ne se contente plus d'une formule globale. On remplace le modèle simple par une loi locale, différentielle ou énergétique.

Une forme générale possible est :
\[
\frac{\mathrm d X}{\mathrm dt}=\text{apport}-\text{pertes},
\]
ou, pour un phénomène propagatif :
\[
\frac{\partial^2 \psi}{\partial t^2}=c^2\Delta\psi.
\]
Selon le cas, on peut aussi utiliser un principe variationnel, par exemple le principe de Fermat en optique :
\[
\delta \int n\,\mathrm ds=0.
\]

Pour ce phénomène, le modèle Bac+3 consiste à partir de la relation centrale
\[
2ne\simeq k\lambda
\]
puis à introduire les corrections négligées : dépendance spatiale, dépendance temporelle, non-linéarités, pertes, conditions aux limites ou fluctuations. Cela permet de passer d'une réponse d'ordre de grandeur à une prédiction plus robuste.
\end{approfondissementbox}

\subsubsection*{Partie E — Limites, critique et prolongement}

\begin{enumerate}
  \item Citer au moins deux limites du modèle Terminale.
  \item Dire quelles grandeurs seraient difficiles à mesurer expérimentalement.
  \item Proposer une expérience simple permettant de tester le modèle.
  \item Expliquer ce qui changerait si l'on doublait une grandeur clé.
  \item Rédiger une conclusion courte, accessible à un élève de Terminale.
\end{enumerate}

\begin{corrigebox}[Synthèse rédigée]
Le phénomène « Bulles de savon » s'explique par un petit nombre de lois physiques simples, mais sa manifestation réelle dépend souvent de nombreux paramètres. Le modèle Terminale donne l'intuition et les ordres de grandeur ; le modèle Bac+3 introduit les variations locales, les pertes, les équations différentielles ou les principes variationnels nécessaires pour décrire finement la réalité.
\end{corrigebox}

% ------------------------------------------------------------
\subsection{Problème corrigé 47 — Taches d'huile}
% ------------------------------------------------------------

\begin{exercicebox}[Question-problème]
\textbf{Observation.} Une tache d'huile est irisée.

\medskip
\textbf{Question.} Expliquer la dépendance angulaire.

\medskip
\textbf{Thèmes mobilisés :} Interférences en film mince.

\medskip
\textbf{Relation ou idée centrale :}
\[
\delta\simeq2ne\cos r
\]
\end{exercicebox}

\subsubsection*{Partie A — Modélisation niveau Terminale}

\begin{enumerate}
  \item Décrire précisément le phénomène observé.
  \item Identifier le système physique étudié.
  \item Lister les grandeurs importantes et leurs unités.
  \item Écrire la loi physique simple qui permet de commencer le raisonnement.
  \item Expliquer les hypothèses simplificatrices du modèle Terminale.
\end{enumerate}

\subsubsection*{Partie B — Mise en équation et calcul}

\begin{enumerate}
  \item Définir les variables utiles.
  \item Écrire la relation mathématique principale.
  \item Vérifier l'homogénéité de la formule.
  \item Choisir un ordre de grandeur réaliste pour les données manquantes.
  \item Effectuer le calcul ou établir le raisonnement qualitatif demandé.
  \item Interpréter le résultat en une phrase claire.
\end{enumerate}

\subsubsection*{Partie C — Réponse accessible niveau Terminale}

\begin{corrigebox}[Réponse Terminale]
On part de l'observation concrète : Une tache d'huile est irisée.

Au niveau Terminale, on cherche d'abord la grandeur qui contrôle le phénomène. Ici, l'idée centrale est donnée par :
\[
\delta\simeq2ne\cos r
\]
Cette relation permet de relier la cause physique à l'effet observé. On vérifie les unités, puis on insère des ordres de grandeur réalistes.

La conclusion attendue est la suivante : Expliquer la dépendance angulaire. La réponse doit toujours préciser les hypothèses utilisées. Si le phénomène dépend d'un milieu réel, d'une géométrie complexe, de frottements, de pertes ou de gradients, le résultat obtenu est un ordre de grandeur et non une valeur exacte.
\end{corrigebox}

\subsubsection*{Partie D — Approfondissement niveau Bac+3}

\begin{approfondissementbox}[Réponse Bac+3]
Au niveau Bac+3, on ne se contente plus d'une formule globale. On remplace le modèle simple par une loi locale, différentielle ou énergétique.

Une forme générale possible est :
\[
\frac{\mathrm d X}{\mathrm dt}=\text{apport}-\text{pertes},
\]
ou, pour un phénomène propagatif :
\[
\frac{\partial^2 \psi}{\partial t^2}=c^2\Delta\psi.
\]
Selon le cas, on peut aussi utiliser un principe variationnel, par exemple le principe de Fermat en optique :
\[
\delta \int n\,\mathrm ds=0.
\]

Pour ce phénomène, le modèle Bac+3 consiste à partir de la relation centrale
\[
\delta\simeq2ne\cos r
\]
puis à introduire les corrections négligées : dépendance spatiale, dépendance temporelle, non-linéarités, pertes, conditions aux limites ou fluctuations. Cela permet de passer d'une réponse d'ordre de grandeur à une prédiction plus robuste.
\end{approfondissementbox}

\subsubsection*{Partie E — Limites, critique et prolongement}

\begin{enumerate}
  \item Citer au moins deux limites du modèle Terminale.
  \item Dire quelles grandeurs seraient difficiles à mesurer expérimentalement.
  \item Proposer une expérience simple permettant de tester le modèle.
  \item Expliquer ce qui changerait si l'on doublait une grandeur clé.
  \item Rédiger une conclusion courte, accessible à un élève de Terminale.
\end{enumerate}

\begin{corrigebox}[Synthèse rédigée]
Le phénomène « Taches d'huile » s'explique par un petit nombre de lois physiques simples, mais sa manifestation réelle dépend souvent de nombreux paramètres. Le modèle Terminale donne l'intuition et les ordres de grandeur ; le modèle Bac+3 introduit les variations locales, les pertes, les équations différentielles ou les principes variationnels nécessaires pour décrire finement la réalité.
\end{corrigebox}

% ------------------------------------------------------------
\subsection{Problème corrigé 48 — Polarisation}
% ------------------------------------------------------------

\begin{exercicebox}[Question-problème]
\textbf{Observation.} Certains filtres bloquent la lumière.

\medskip
\textbf{Question.} Calculer une intensité transmise.

\medskip
\textbf{Thèmes mobilisés :} Onde transverse, loi de Malus.

\medskip
\textbf{Relation ou idée centrale :}
\[
I=I_0\cos^2\theta
\]
\end{exercicebox}

\subsubsection*{Partie A — Modélisation niveau Terminale}

\begin{enumerate}
  \item Décrire précisément le phénomène observé.
  \item Identifier le système physique étudié.
  \item Lister les grandeurs importantes et leurs unités.
  \item Écrire la loi physique simple qui permet de commencer le raisonnement.
  \item Expliquer les hypothèses simplificatrices du modèle Terminale.
\end{enumerate}

\subsubsection*{Partie B — Mise en équation et calcul}

\begin{enumerate}
  \item Définir les variables utiles.
  \item Écrire la relation mathématique principale.
  \item Vérifier l'homogénéité de la formule.
  \item Choisir un ordre de grandeur réaliste pour les données manquantes.
  \item Effectuer le calcul ou établir le raisonnement qualitatif demandé.
  \item Interpréter le résultat en une phrase claire.
\end{enumerate}

\subsubsection*{Partie C — Réponse accessible niveau Terminale}

\begin{corrigebox}[Réponse Terminale]
On part de l'observation concrète : Certains filtres bloquent la lumière.

Au niveau Terminale, on cherche d'abord la grandeur qui contrôle le phénomène. Ici, l'idée centrale est donnée par :
\[
I=I_0\cos^2\theta
\]
Cette relation permet de relier la cause physique à l'effet observé. On vérifie les unités, puis on insère des ordres de grandeur réalistes.

La conclusion attendue est la suivante : Calculer une intensité transmise. La réponse doit toujours préciser les hypothèses utilisées. Si le phénomène dépend d'un milieu réel, d'une géométrie complexe, de frottements, de pertes ou de gradients, le résultat obtenu est un ordre de grandeur et non une valeur exacte.
\end{corrigebox}

\subsubsection*{Partie D — Approfondissement niveau Bac+3}

\begin{approfondissementbox}[Réponse Bac+3]
Au niveau Bac+3, on ne se contente plus d'une formule globale. On remplace le modèle simple par une loi locale, différentielle ou énergétique.

Une forme générale possible est :
\[
\frac{\mathrm d X}{\mathrm dt}=\text{apport}-\text{pertes},
\]
ou, pour un phénomène propagatif :
\[
\frac{\partial^2 \psi}{\partial t^2}=c^2\Delta\psi.
\]
Selon le cas, on peut aussi utiliser un principe variationnel, par exemple le principe de Fermat en optique :
\[
\delta \int n\,\mathrm ds=0.
\]

Pour ce phénomène, le modèle Bac+3 consiste à partir de la relation centrale
\[
I=I_0\cos^2\theta
\]
puis à introduire les corrections négligées : dépendance spatiale, dépendance temporelle, non-linéarités, pertes, conditions aux limites ou fluctuations. Cela permet de passer d'une réponse d'ordre de grandeur à une prédiction plus robuste.
\end{approfondissementbox}

\subsubsection*{Partie E — Limites, critique et prolongement}

\begin{enumerate}
  \item Citer au moins deux limites du modèle Terminale.
  \item Dire quelles grandeurs seraient difficiles à mesurer expérimentalement.
  \item Proposer une expérience simple permettant de tester le modèle.
  \item Expliquer ce qui changerait si l'on doublait une grandeur clé.
  \item Rédiger une conclusion courte, accessible à un élève de Terminale.
\end{enumerate}

\begin{corrigebox}[Synthèse rédigée]
Le phénomène « Polarisation » s'explique par un petit nombre de lois physiques simples, mais sa manifestation réelle dépend souvent de nombreux paramètres. Le modèle Terminale donne l'intuition et les ordres de grandeur ; le modèle Bac+3 introduit les variations locales, les pertes, les équations différentielles ou les principes variationnels nécessaires pour décrire finement la réalité.
\end{corrigebox}

% ------------------------------------------------------------
\subsection{Problème corrigé 49 — Lunettes polarisantes}
% ------------------------------------------------------------

\begin{exercicebox}[Question-problème]
\textbf{Observation.} Les reflets sont atténués.

\medskip
\textbf{Question.} Comprendre l'atténuation des reflets.

\medskip
\textbf{Thèmes mobilisés :} Polarisation par réflexion.

\medskip
\textbf{Relation ou idée centrale :}
\[
I=I_0\cos^2\theta
\]
\end{exercicebox}

\subsubsection*{Partie A — Modélisation niveau Terminale}

\begin{enumerate}
  \item Décrire précisément le phénomène observé.
  \item Identifier le système physique étudié.
  \item Lister les grandeurs importantes et leurs unités.
  \item Écrire la loi physique simple qui permet de commencer le raisonnement.
  \item Expliquer les hypothèses simplificatrices du modèle Terminale.
\end{enumerate}

\subsubsection*{Partie B — Mise en équation et calcul}

\begin{enumerate}
  \item Définir les variables utiles.
  \item Écrire la relation mathématique principale.
  \item Vérifier l'homogénéité de la formule.
  \item Choisir un ordre de grandeur réaliste pour les données manquantes.
  \item Effectuer le calcul ou établir le raisonnement qualitatif demandé.
  \item Interpréter le résultat en une phrase claire.
\end{enumerate}

\subsubsection*{Partie C — Réponse accessible niveau Terminale}

\begin{corrigebox}[Réponse Terminale]
On part de l'observation concrète : Les reflets sont atténués.

Au niveau Terminale, on cherche d'abord la grandeur qui contrôle le phénomène. Ici, l'idée centrale est donnée par :
\[
I=I_0\cos^2\theta
\]
Cette relation permet de relier la cause physique à l'effet observé. On vérifie les unités, puis on insère des ordres de grandeur réalistes.

La conclusion attendue est la suivante : Comprendre l'atténuation des reflets. La réponse doit toujours préciser les hypothèses utilisées. Si le phénomène dépend d'un milieu réel, d'une géométrie complexe, de frottements, de pertes ou de gradients, le résultat obtenu est un ordre de grandeur et non une valeur exacte.
\end{corrigebox}

\subsubsection*{Partie D — Approfondissement niveau Bac+3}

\begin{approfondissementbox}[Réponse Bac+3]
Au niveau Bac+3, on ne se contente plus d'une formule globale. On remplace le modèle simple par une loi locale, différentielle ou énergétique.

Une forme générale possible est :
\[
\frac{\mathrm d X}{\mathrm dt}=\text{apport}-\text{pertes},
\]
ou, pour un phénomène propagatif :
\[
\frac{\partial^2 \psi}{\partial t^2}=c^2\Delta\psi.
\]
Selon le cas, on peut aussi utiliser un principe variationnel, par exemple le principe de Fermat en optique :
\[
\delta \int n\,\mathrm ds=0.
\]

Pour ce phénomène, le modèle Bac+3 consiste à partir de la relation centrale
\[
I=I_0\cos^2\theta
\]
puis à introduire les corrections négligées : dépendance spatiale, dépendance temporelle, non-linéarités, pertes, conditions aux limites ou fluctuations. Cela permet de passer d'une réponse d'ordre de grandeur à une prédiction plus robuste.
\end{approfondissementbox}

\subsubsection*{Partie E — Limites, critique et prolongement}

\begin{enumerate}
  \item Citer au moins deux limites du modèle Terminale.
  \item Dire quelles grandeurs seraient difficiles à mesurer expérimentalement.
  \item Proposer une expérience simple permettant de tester le modèle.
  \item Expliquer ce qui changerait si l'on doublait une grandeur clé.
  \item Rédiger une conclusion courte, accessible à un élève de Terminale.
\end{enumerate}

\begin{corrigebox}[Synthèse rédigée]
Le phénomène « Lunettes polarisantes » s'explique par un petit nombre de lois physiques simples, mais sa manifestation réelle dépend souvent de nombreux paramètres. Le modèle Terminale donne l'intuition et les ordres de grandeur ; le modèle Bac+3 introduit les variations locales, les pertes, les équations différentielles ou les principes variationnels nécessaires pour décrire finement la réalité.
\end{corrigebox}

% ------------------------------------------------------------
\subsection{Problème corrigé 50 — Écrans LCD}
% ------------------------------------------------------------

\begin{exercicebox}[Question-problème]
\textbf{Observation.} Un écran module la lumière.

\medskip
\textbf{Question.} Expliquer comment un pixel change de luminosité.

\medskip
\textbf{Thèmes mobilisés :} Polarisation, cristaux liquides.

\medskip
\textbf{Relation ou idée centrale :}
\[
I=I_0\cos^2\theta
\]
\end{exercicebox}

\subsubsection*{Partie A — Modélisation niveau Terminale}

\begin{enumerate}
  \item Décrire précisément le phénomène observé.
  \item Identifier le système physique étudié.
  \item Lister les grandeurs importantes et leurs unités.
  \item Écrire la loi physique simple qui permet de commencer le raisonnement.
  \item Expliquer les hypothèses simplificatrices du modèle Terminale.
\end{enumerate}

\subsubsection*{Partie B — Mise en équation et calcul}

\begin{enumerate}
  \item Définir les variables utiles.
  \item Écrire la relation mathématique principale.
  \item Vérifier l'homogénéité de la formule.
  \item Choisir un ordre de grandeur réaliste pour les données manquantes.
  \item Effectuer le calcul ou établir le raisonnement qualitatif demandé.
  \item Interpréter le résultat en une phrase claire.
\end{enumerate}

\subsubsection*{Partie C — Réponse accessible niveau Terminale}

\begin{corrigebox}[Réponse Terminale]
On part de l'observation concrète : Un écran module la lumière.

Au niveau Terminale, on cherche d'abord la grandeur qui contrôle le phénomène. Ici, l'idée centrale est donnée par :
\[
I=I_0\cos^2\theta
\]
Cette relation permet de relier la cause physique à l'effet observé. On vérifie les unités, puis on insère des ordres de grandeur réalistes.

La conclusion attendue est la suivante : Expliquer comment un pixel change de luminosité. La réponse doit toujours préciser les hypothèses utilisées. Si le phénomène dépend d'un milieu réel, d'une géométrie complexe, de frottements, de pertes ou de gradients, le résultat obtenu est un ordre de grandeur et non une valeur exacte.
\end{corrigebox}

\subsubsection*{Partie D — Approfondissement niveau Bac+3}

\begin{approfondissementbox}[Réponse Bac+3]
Au niveau Bac+3, on ne se contente plus d'une formule globale. On remplace le modèle simple par une loi locale, différentielle ou énergétique.

Une forme générale possible est :
\[
\frac{\mathrm d X}{\mathrm dt}=\text{apport}-\text{pertes},
\]
ou, pour un phénomène propagatif :
\[
\frac{\partial^2 \psi}{\partial t^2}=c^2\Delta\psi.
\]
Selon le cas, on peut aussi utiliser un principe variationnel, par exemple le principe de Fermat en optique :
\[
\delta \int n\,\mathrm ds=0.
\]

Pour ce phénomène, le modèle Bac+3 consiste à partir de la relation centrale
\[
I=I_0\cos^2\theta
\]
puis à introduire les corrections négligées : dépendance spatiale, dépendance temporelle, non-linéarités, pertes, conditions aux limites ou fluctuations. Cela permet de passer d'une réponse d'ordre de grandeur à une prédiction plus robuste.
\end{approfondissementbox}

\subsubsection*{Partie E — Limites, critique et prolongement}

\begin{enumerate}
  \item Citer au moins deux limites du modèle Terminale.
  \item Dire quelles grandeurs seraient difficiles à mesurer expérimentalement.
  \item Proposer une expérience simple permettant de tester le modèle.
  \item Expliquer ce qui changerait si l'on doublait une grandeur clé.
  \item Rédiger une conclusion courte, accessible à un élève de Terminale.
\end{enumerate}

\begin{corrigebox}[Synthèse rédigée]
Le phénomène « Écrans LCD » s'explique par un petit nombre de lois physiques simples, mais sa manifestation réelle dépend souvent de nombreux paramètres. Le modèle Terminale donne l'intuition et les ordres de grandeur ; le modèle Bac+3 introduit les variations locales, les pertes, les équations différentielles ou les principes variationnels nécessaires pour décrire finement la réalité.
\end{corrigebox}

% ------------------------------------------------------------
\subsection{Problème corrigé 51 — Laser}
% ------------------------------------------------------------

\begin{exercicebox}[Question-problème]
\textbf{Observation.} Un laser produit un faisceau cohérent.

\medskip
\textbf{Question.} Expliquer la monochromaticité et la directivité.

\medskip
\textbf{Thèmes mobilisés :} Émission stimulée, cohérence.

\medskip
\textbf{Relation ou idée centrale :}
\[
E=h\nu
\]
\end{exercicebox}

\subsubsection*{Partie A — Modélisation niveau Terminale}

\begin{enumerate}
  \item Décrire précisément le phénomène observé.
  \item Identifier le système physique étudié.
  \item Lister les grandeurs importantes et leurs unités.
  \item Écrire la loi physique simple qui permet de commencer le raisonnement.
  \item Expliquer les hypothèses simplificatrices du modèle Terminale.
\end{enumerate}

\subsubsection*{Partie B — Mise en équation et calcul}

\begin{enumerate}
  \item Définir les variables utiles.
  \item Écrire la relation mathématique principale.
  \item Vérifier l'homogénéité de la formule.
  \item Choisir un ordre de grandeur réaliste pour les données manquantes.
  \item Effectuer le calcul ou établir le raisonnement qualitatif demandé.
  \item Interpréter le résultat en une phrase claire.
\end{enumerate}

\subsubsection*{Partie C — Réponse accessible niveau Terminale}

\begin{corrigebox}[Réponse Terminale]
On part de l'observation concrète : Un laser produit un faisceau cohérent.

Au niveau Terminale, on cherche d'abord la grandeur qui contrôle le phénomène. Ici, l'idée centrale est donnée par :
\[
E=h\nu
\]
Cette relation permet de relier la cause physique à l'effet observé. On vérifie les unités, puis on insère des ordres de grandeur réalistes.

La conclusion attendue est la suivante : Expliquer la monochromaticité et la directivité. La réponse doit toujours préciser les hypothèses utilisées. Si le phénomène dépend d'un milieu réel, d'une géométrie complexe, de frottements, de pertes ou de gradients, le résultat obtenu est un ordre de grandeur et non une valeur exacte.
\end{corrigebox}

\subsubsection*{Partie D — Approfondissement niveau Bac+3}

\begin{approfondissementbox}[Réponse Bac+3]
Au niveau Bac+3, on ne se contente plus d'une formule globale. On remplace le modèle simple par une loi locale, différentielle ou énergétique.

Une forme générale possible est :
\[
\frac{\mathrm d X}{\mathrm dt}=\text{apport}-\text{pertes},
\]
ou, pour un phénomène propagatif :
\[
\frac{\partial^2 \psi}{\partial t^2}=c^2\Delta\psi.
\]
Selon le cas, on peut aussi utiliser un principe variationnel, par exemple le principe de Fermat en optique :
\[
\delta \int n\,\mathrm ds=0.
\]

Pour ce phénomène, le modèle Bac+3 consiste à partir de la relation centrale
\[
E=h\nu
\]
puis à introduire les corrections négligées : dépendance spatiale, dépendance temporelle, non-linéarités, pertes, conditions aux limites ou fluctuations. Cela permet de passer d'une réponse d'ordre de grandeur à une prédiction plus robuste.
\end{approfondissementbox}

\subsubsection*{Partie E — Limites, critique et prolongement}

\begin{enumerate}
  \item Citer au moins deux limites du modèle Terminale.
  \item Dire quelles grandeurs seraient difficiles à mesurer expérimentalement.
  \item Proposer une expérience simple permettant de tester le modèle.
  \item Expliquer ce qui changerait si l'on doublait une grandeur clé.
  \item Rédiger une conclusion courte, accessible à un élève de Terminale.
\end{enumerate}

\begin{corrigebox}[Synthèse rédigée]
Le phénomène « Laser » s'explique par un petit nombre de lois physiques simples, mais sa manifestation réelle dépend souvent de nombreux paramètres. Le modèle Terminale donne l'intuition et les ordres de grandeur ; le modèle Bac+3 introduit les variations locales, les pertes, les équations différentielles ou les principes variationnels nécessaires pour décrire finement la réalité.
\end{corrigebox}

% ------------------------------------------------------------
\subsection{Problème corrigé 52 — Lidar}
% ------------------------------------------------------------

\begin{exercicebox}[Question-problème]
\textbf{Observation.} Un lidar mesure une distance.

\medskip
\textbf{Question.} Calculer une distance par délai.

\medskip
\textbf{Thèmes mobilisés :} Temps de vol.

\medskip
\textbf{Relation ou idée centrale :}
\[
d=\frac{c\Delta t}{2}
\]
\end{exercicebox}

\subsubsection*{Partie A — Modélisation niveau Terminale}

\begin{enumerate}
  \item Décrire précisément le phénomène observé.
  \item Identifier le système physique étudié.
  \item Lister les grandeurs importantes et leurs unités.
  \item Écrire la loi physique simple qui permet de commencer le raisonnement.
  \item Expliquer les hypothèses simplificatrices du modèle Terminale.
\end{enumerate}

\subsubsection*{Partie B — Mise en équation et calcul}

\begin{enumerate}
  \item Définir les variables utiles.
  \item Écrire la relation mathématique principale.
  \item Vérifier l'homogénéité de la formule.
  \item Choisir un ordre de grandeur réaliste pour les données manquantes.
  \item Effectuer le calcul ou établir le raisonnement qualitatif demandé.
  \item Interpréter le résultat en une phrase claire.
\end{enumerate}

\subsubsection*{Partie C — Réponse accessible niveau Terminale}

\begin{corrigebox}[Réponse Terminale]
On part de l'observation concrète : Un lidar mesure une distance.

Au niveau Terminale, on cherche d'abord la grandeur qui contrôle le phénomène. Ici, l'idée centrale est donnée par :
\[
d=\frac{c\Delta t}{2}
\]
Cette relation permet de relier la cause physique à l'effet observé. On vérifie les unités, puis on insère des ordres de grandeur réalistes.

La conclusion attendue est la suivante : Calculer une distance par délai. La réponse doit toujours préciser les hypothèses utilisées. Si le phénomène dépend d'un milieu réel, d'une géométrie complexe, de frottements, de pertes ou de gradients, le résultat obtenu est un ordre de grandeur et non une valeur exacte.
\end{corrigebox}

\subsubsection*{Partie D — Approfondissement niveau Bac+3}

\begin{approfondissementbox}[Réponse Bac+3]
Au niveau Bac+3, on ne se contente plus d'une formule globale. On remplace le modèle simple par une loi locale, différentielle ou énergétique.

Une forme générale possible est :
\[
\frac{\mathrm d X}{\mathrm dt}=\text{apport}-\text{pertes},
\]
ou, pour un phénomène propagatif :
\[
\frac{\partial^2 \psi}{\partial t^2}=c^2\Delta\psi.
\]
Selon le cas, on peut aussi utiliser un principe variationnel, par exemple le principe de Fermat en optique :
\[
\delta \int n\,\mathrm ds=0.
\]

Pour ce phénomène, le modèle Bac+3 consiste à partir de la relation centrale
\[
d=\frac{c\Delta t}{2}
\]
puis à introduire les corrections négligées : dépendance spatiale, dépendance temporelle, non-linéarités, pertes, conditions aux limites ou fluctuations. Cela permet de passer d'une réponse d'ordre de grandeur à une prédiction plus robuste.
\end{approfondissementbox}

\subsubsection*{Partie E — Limites, critique et prolongement}

\begin{enumerate}
  \item Citer au moins deux limites du modèle Terminale.
  \item Dire quelles grandeurs seraient difficiles à mesurer expérimentalement.
  \item Proposer une expérience simple permettant de tester le modèle.
  \item Expliquer ce qui changerait si l'on doublait une grandeur clé.
  \item Rédiger une conclusion courte, accessible à un élève de Terminale.
\end{enumerate}

\begin{corrigebox}[Synthèse rédigée]
Le phénomène « Lidar » s'explique par un petit nombre de lois physiques simples, mais sa manifestation réelle dépend souvent de nombreux paramètres. Le modèle Terminale donne l'intuition et les ordres de grandeur ; le modèle Bac+3 introduit les variations locales, les pertes, les équations différentielles ou les principes variationnels nécessaires pour décrire finement la réalité.
\end{corrigebox}

% ------------------------------------------------------------
\subsection{Problème corrigé 53 — Spectroscopie}
% ------------------------------------------------------------

\begin{exercicebox}[Question-problème]
\textbf{Observation.} Un spectre révèle une composition.

\medskip
\textbf{Question.} Relier absorbance et concentration.

\medskip
\textbf{Thèmes mobilisés :} Absorption, émission, Beer-Lambert.

\medskip
\textbf{Relation ou idée centrale :}
\[
A=\varepsilon\ell C
\]
\end{exercicebox}

\subsubsection*{Partie A — Modélisation niveau Terminale}

\begin{enumerate}
  \item Décrire précisément le phénomène observé.
  \item Identifier le système physique étudié.
  \item Lister les grandeurs importantes et leurs unités.
  \item Écrire la loi physique simple qui permet de commencer le raisonnement.
  \item Expliquer les hypothèses simplificatrices du modèle Terminale.
\end{enumerate}

\subsubsection*{Partie B — Mise en équation et calcul}

\begin{enumerate}
  \item Définir les variables utiles.
  \item Écrire la relation mathématique principale.
  \item Vérifier l'homogénéité de la formule.
  \item Choisir un ordre de grandeur réaliste pour les données manquantes.
  \item Effectuer le calcul ou établir le raisonnement qualitatif demandé.
  \item Interpréter le résultat en une phrase claire.
\end{enumerate}

\subsubsection*{Partie C — Réponse accessible niveau Terminale}

\begin{corrigebox}[Réponse Terminale]
On part de l'observation concrète : Un spectre révèle une composition.

Au niveau Terminale, on cherche d'abord la grandeur qui contrôle le phénomène. Ici, l'idée centrale est donnée par :
\[
A=\varepsilon\ell C
\]
Cette relation permet de relier la cause physique à l'effet observé. On vérifie les unités, puis on insère des ordres de grandeur réalistes.

La conclusion attendue est la suivante : Relier absorbance et concentration. La réponse doit toujours préciser les hypothèses utilisées. Si le phénomène dépend d'un milieu réel, d'une géométrie complexe, de frottements, de pertes ou de gradients, le résultat obtenu est un ordre de grandeur et non une valeur exacte.
\end{corrigebox}

\subsubsection*{Partie D — Approfondissement niveau Bac+3}

\begin{approfondissementbox}[Réponse Bac+3]
Au niveau Bac+3, on ne se contente plus d'une formule globale. On remplace le modèle simple par une loi locale, différentielle ou énergétique.

Une forme générale possible est :
\[
\frac{\mathrm d X}{\mathrm dt}=\text{apport}-\text{pertes},
\]
ou, pour un phénomène propagatif :
\[
\frac{\partial^2 \psi}{\partial t^2}=c^2\Delta\psi.
\]
Selon le cas, on peut aussi utiliser un principe variationnel, par exemple le principe de Fermat en optique :
\[
\delta \int n\,\mathrm ds=0.
\]

Pour ce phénomène, le modèle Bac+3 consiste à partir de la relation centrale
\[
A=\varepsilon\ell C
\]
puis à introduire les corrections négligées : dépendance spatiale, dépendance temporelle, non-linéarités, pertes, conditions aux limites ou fluctuations. Cela permet de passer d'une réponse d'ordre de grandeur à une prédiction plus robuste.
\end{approfondissementbox}

\subsubsection*{Partie E — Limites, critique et prolongement}

\begin{enumerate}
  \item Citer au moins deux limites du modèle Terminale.
  \item Dire quelles grandeurs seraient difficiles à mesurer expérimentalement.
  \item Proposer une expérience simple permettant de tester le modèle.
  \item Expliquer ce qui changerait si l'on doublait une grandeur clé.
  \item Rédiger une conclusion courte, accessible à un élève de Terminale.
\end{enumerate}

\begin{corrigebox}[Synthèse rédigée]
Le phénomène « Spectroscopie » s'explique par un petit nombre de lois physiques simples, mais sa manifestation réelle dépend souvent de nombreux paramètres. Le modèle Terminale donne l'intuition et les ordres de grandeur ; le modèle Bac+3 introduit les variations locales, les pertes, les équations différentielles ou les principes variationnels nécessaires pour décrire finement la réalité.
\end{corrigebox}

% ------------------------------------------------------------
\subsection{Problème corrigé 54 — Effet Doppler sonore}
% ------------------------------------------------------------

\begin{exercicebox}[Question-problème]
\textbf{Observation.} Une ambulance change de son.

\medskip
\textbf{Question.} Calculer une fréquence perçue.

\medskip
\textbf{Thèmes mobilisés :} Source en mouvement.

\medskip
\textbf{Relation ou idée centrale :}
\[
f'=f\frac{c}{c-v_s}
\]
\end{exercicebox}

\subsubsection*{Partie A — Modélisation niveau Terminale}

\begin{enumerate}
  \item Décrire précisément le phénomène observé.
  \item Identifier le système physique étudié.
  \item Lister les grandeurs importantes et leurs unités.
  \item Écrire la loi physique simple qui permet de commencer le raisonnement.
  \item Expliquer les hypothèses simplificatrices du modèle Terminale.
\end{enumerate}

\subsubsection*{Partie B — Mise en équation et calcul}

\begin{enumerate}
  \item Définir les variables utiles.
  \item Écrire la relation mathématique principale.
  \item Vérifier l'homogénéité de la formule.
  \item Choisir un ordre de grandeur réaliste pour les données manquantes.
  \item Effectuer le calcul ou établir le raisonnement qualitatif demandé.
  \item Interpréter le résultat en une phrase claire.
\end{enumerate}

\subsubsection*{Partie C — Réponse accessible niveau Terminale}

\begin{corrigebox}[Réponse Terminale]
On part de l'observation concrète : Une ambulance change de son.

Au niveau Terminale, on cherche d'abord la grandeur qui contrôle le phénomène. Ici, l'idée centrale est donnée par :
\[
f'=f\frac{c}{c-v_s}
\]
Cette relation permet de relier la cause physique à l'effet observé. On vérifie les unités, puis on insère des ordres de grandeur réalistes.

La conclusion attendue est la suivante : Calculer une fréquence perçue. La réponse doit toujours préciser les hypothèses utilisées. Si le phénomène dépend d'un milieu réel, d'une géométrie complexe, de frottements, de pertes ou de gradients, le résultat obtenu est un ordre de grandeur et non une valeur exacte.
\end{corrigebox}

\subsubsection*{Partie D — Approfondissement niveau Bac+3}

\begin{approfondissementbox}[Réponse Bac+3]
Au niveau Bac+3, on ne se contente plus d'une formule globale. On remplace le modèle simple par une loi locale, différentielle ou énergétique.

Une forme générale possible est :
\[
\frac{\mathrm d X}{\mathrm dt}=\text{apport}-\text{pertes},
\]
ou, pour un phénomène propagatif :
\[
\frac{\partial^2 \psi}{\partial t^2}=c^2\Delta\psi.
\]
Selon le cas, on peut aussi utiliser un principe variationnel, par exemple le principe de Fermat en optique :
\[
\delta \int n\,\mathrm ds=0.
\]

Pour ce phénomène, le modèle Bac+3 consiste à partir de la relation centrale
\[
f'=f\frac{c}{c-v_s}
\]
puis à introduire les corrections négligées : dépendance spatiale, dépendance temporelle, non-linéarités, pertes, conditions aux limites ou fluctuations. Cela permet de passer d'une réponse d'ordre de grandeur à une prédiction plus robuste.
\end{approfondissementbox}

\subsubsection*{Partie E — Limites, critique et prolongement}

\begin{enumerate}
  \item Citer au moins deux limites du modèle Terminale.
  \item Dire quelles grandeurs seraient difficiles à mesurer expérimentalement.
  \item Proposer une expérience simple permettant de tester le modèle.
  \item Expliquer ce qui changerait si l'on doublait une grandeur clé.
  \item Rédiger une conclusion courte, accessible à un élève de Terminale.
\end{enumerate}

\begin{corrigebox}[Synthèse rédigée]
Le phénomène « Effet Doppler sonore » s'explique par un petit nombre de lois physiques simples, mais sa manifestation réelle dépend souvent de nombreux paramètres. Le modèle Terminale donne l'intuition et les ordres de grandeur ; le modèle Bac+3 introduit les variations locales, les pertes, les équations différentielles ou les principes variationnels nécessaires pour décrire finement la réalité.
\end{corrigebox}

% ------------------------------------------------------------
\subsection{Problème corrigé 55 — Effet Doppler lumineux}
% ------------------------------------------------------------

\begin{exercicebox}[Question-problème]
\textbf{Observation.} Une étoile peut être décalée vers le rouge.

\medskip
\textbf{Question.} Relier décalage et vitesse.

\medskip
\textbf{Thèmes mobilisés :} Décalage spectral.

\medskip
\textbf{Relation ou idée centrale :}
\[
\frac{\Delta\lambda}{\lambda}\simeq\frac{v}{c}
\]
\end{exercicebox}

\subsubsection*{Partie A — Modélisation niveau Terminale}

\begin{enumerate}
  \item Décrire précisément le phénomène observé.
  \item Identifier le système physique étudié.
  \item Lister les grandeurs importantes et leurs unités.
  \item Écrire la loi physique simple qui permet de commencer le raisonnement.
  \item Expliquer les hypothèses simplificatrices du modèle Terminale.
\end{enumerate}

\subsubsection*{Partie B — Mise en équation et calcul}

\begin{enumerate}
  \item Définir les variables utiles.
  \item Écrire la relation mathématique principale.
  \item Vérifier l'homogénéité de la formule.
  \item Choisir un ordre de grandeur réaliste pour les données manquantes.
  \item Effectuer le calcul ou établir le raisonnement qualitatif demandé.
  \item Interpréter le résultat en une phrase claire.
\end{enumerate}

\subsubsection*{Partie C — Réponse accessible niveau Terminale}

\begin{corrigebox}[Réponse Terminale]
On part de l'observation concrète : Une étoile peut être décalée vers le rouge.

Au niveau Terminale, on cherche d'abord la grandeur qui contrôle le phénomène. Ici, l'idée centrale est donnée par :
\[
\frac{\Delta\lambda}{\lambda}\simeq\frac{v}{c}
\]
Cette relation permet de relier la cause physique à l'effet observé. On vérifie les unités, puis on insère des ordres de grandeur réalistes.

La conclusion attendue est la suivante : Relier décalage et vitesse. La réponse doit toujours préciser les hypothèses utilisées. Si le phénomène dépend d'un milieu réel, d'une géométrie complexe, de frottements, de pertes ou de gradients, le résultat obtenu est un ordre de grandeur et non une valeur exacte.
\end{corrigebox}

\subsubsection*{Partie D — Approfondissement niveau Bac+3}

\begin{approfondissementbox}[Réponse Bac+3]
Au niveau Bac+3, on ne se contente plus d'une formule globale. On remplace le modèle simple par une loi locale, différentielle ou énergétique.

Une forme générale possible est :
\[
\frac{\mathrm d X}{\mathrm dt}=\text{apport}-\text{pertes},
\]
ou, pour un phénomène propagatif :
\[
\frac{\partial^2 \psi}{\partial t^2}=c^2\Delta\psi.
\]
Selon le cas, on peut aussi utiliser un principe variationnel, par exemple le principe de Fermat en optique :
\[
\delta \int n\,\mathrm ds=0.
\]

Pour ce phénomène, le modèle Bac+3 consiste à partir de la relation centrale
\[
\frac{\Delta\lambda}{\lambda}\simeq\frac{v}{c}
\]
puis à introduire les corrections négligées : dépendance spatiale, dépendance temporelle, non-linéarités, pertes, conditions aux limites ou fluctuations. Cela permet de passer d'une réponse d'ordre de grandeur à une prédiction plus robuste.
\end{approfondissementbox}

\subsubsection*{Partie E — Limites, critique et prolongement}

\begin{enumerate}
  \item Citer au moins deux limites du modèle Terminale.
  \item Dire quelles grandeurs seraient difficiles à mesurer expérimentalement.
  \item Proposer une expérience simple permettant de tester le modèle.
  \item Expliquer ce qui changerait si l'on doublait une grandeur clé.
  \item Rédiger une conclusion courte, accessible à un élève de Terminale.
\end{enumerate}

\begin{corrigebox}[Synthèse rédigée]
Le phénomène « Effet Doppler lumineux » s'explique par un petit nombre de lois physiques simples, mais sa manifestation réelle dépend souvent de nombreux paramètres. Le modèle Terminale donne l'intuition et les ordres de grandeur ; le modèle Bac+3 introduit les variations locales, les pertes, les équations différentielles ou les principes variationnels nécessaires pour décrire finement la réalité.
\end{corrigebox}

% ------------------------------------------------------------
\subsection{Problème corrigé 56 — Mur du son}
% ------------------------------------------------------------

\begin{exercicebox}[Question-problème]
\textbf{Observation.} Un avion atteint Mach 1.

\medskip
\textbf{Question.} Comprendre le régime supersonique.

\medskip
\textbf{Thèmes mobilisés :} Vitesse du son, accumulation d'ondes.

\medskip
\textbf{Relation ou idée centrale :}
\[
M=\frac{v}{c_s}
\]
\end{exercicebox}

\subsubsection*{Partie A — Modélisation niveau Terminale}

\begin{enumerate}
  \item Décrire précisément le phénomène observé.
  \item Identifier le système physique étudié.
  \item Lister les grandeurs importantes et leurs unités.
  \item Écrire la loi physique simple qui permet de commencer le raisonnement.
  \item Expliquer les hypothèses simplificatrices du modèle Terminale.
\end{enumerate}

\subsubsection*{Partie B — Mise en équation et calcul}

\begin{enumerate}
  \item Définir les variables utiles.
  \item Écrire la relation mathématique principale.
  \item Vérifier l'homogénéité de la formule.
  \item Choisir un ordre de grandeur réaliste pour les données manquantes.
  \item Effectuer le calcul ou établir le raisonnement qualitatif demandé.
  \item Interpréter le résultat en une phrase claire.
\end{enumerate}

\subsubsection*{Partie C — Réponse accessible niveau Terminale}

\begin{corrigebox}[Réponse Terminale]
On part de l'observation concrète : Un avion atteint Mach 1.

Au niveau Terminale, on cherche d'abord la grandeur qui contrôle le phénomène. Ici, l'idée centrale est donnée par :
\[
M=\frac{v}{c_s}
\]
Cette relation permet de relier la cause physique à l'effet observé. On vérifie les unités, puis on insère des ordres de grandeur réalistes.

La conclusion attendue est la suivante : Comprendre le régime supersonique. La réponse doit toujours préciser les hypothèses utilisées. Si le phénomène dépend d'un milieu réel, d'une géométrie complexe, de frottements, de pertes ou de gradients, le résultat obtenu est un ordre de grandeur et non une valeur exacte.
\end{corrigebox}

\subsubsection*{Partie D — Approfondissement niveau Bac+3}

\begin{approfondissementbox}[Réponse Bac+3]
Au niveau Bac+3, on ne se contente plus d'une formule globale. On remplace le modèle simple par une loi locale, différentielle ou énergétique.

Une forme générale possible est :
\[
\frac{\mathrm d X}{\mathrm dt}=\text{apport}-\text{pertes},
\]
ou, pour un phénomène propagatif :
\[
\frac{\partial^2 \psi}{\partial t^2}=c^2\Delta\psi.
\]
Selon le cas, on peut aussi utiliser un principe variationnel, par exemple le principe de Fermat en optique :
\[
\delta \int n\,\mathrm ds=0.
\]

Pour ce phénomène, le modèle Bac+3 consiste à partir de la relation centrale
\[
M=\frac{v}{c_s}
\]
puis à introduire les corrections négligées : dépendance spatiale, dépendance temporelle, non-linéarités, pertes, conditions aux limites ou fluctuations. Cela permet de passer d'une réponse d'ordre de grandeur à une prédiction plus robuste.
\end{approfondissementbox}

\subsubsection*{Partie E — Limites, critique et prolongement}

\begin{enumerate}
  \item Citer au moins deux limites du modèle Terminale.
  \item Dire quelles grandeurs seraient difficiles à mesurer expérimentalement.
  \item Proposer une expérience simple permettant de tester le modèle.
  \item Expliquer ce qui changerait si l'on doublait une grandeur clé.
  \item Rédiger une conclusion courte, accessible à un élève de Terminale.
\end{enumerate}

\begin{corrigebox}[Synthèse rédigée]
Le phénomène « Mur du son » s'explique par un petit nombre de lois physiques simples, mais sa manifestation réelle dépend souvent de nombreux paramètres. Le modèle Terminale donne l'intuition et les ordres de grandeur ; le modèle Bac+3 introduit les variations locales, les pertes, les équations différentielles ou les principes variationnels nécessaires pour décrire finement la réalité.
\end{corrigebox}

% ------------------------------------------------------------
\subsection{Problème corrigé 57 — Bang supersonique}
% ------------------------------------------------------------

\begin{exercicebox}[Question-problème]
\textbf{Observation.} Un bang accompagne un avion supersonique.

\medskip
\textbf{Question.} Relier angle et nombre de Mach.

\medskip
\textbf{Thèmes mobilisés :} Cône de Mach.

\medskip
\textbf{Relation ou idée centrale :}
\[
\sin\theta=\frac1M
\]
\end{exercicebox}

\subsubsection*{Partie A — Modélisation niveau Terminale}

\begin{enumerate}
  \item Décrire précisément le phénomène observé.
  \item Identifier le système physique étudié.
  \item Lister les grandeurs importantes et leurs unités.
  \item Écrire la loi physique simple qui permet de commencer le raisonnement.
  \item Expliquer les hypothèses simplificatrices du modèle Terminale.
\end{enumerate}

\subsubsection*{Partie B — Mise en équation et calcul}

\begin{enumerate}
  \item Définir les variables utiles.
  \item Écrire la relation mathématique principale.
  \item Vérifier l'homogénéité de la formule.
  \item Choisir un ordre de grandeur réaliste pour les données manquantes.
  \item Effectuer le calcul ou établir le raisonnement qualitatif demandé.
  \item Interpréter le résultat en une phrase claire.
\end{enumerate}

\subsubsection*{Partie C — Réponse accessible niveau Terminale}

\begin{corrigebox}[Réponse Terminale]
On part de l'observation concrète : Un bang accompagne un avion supersonique.

Au niveau Terminale, on cherche d'abord la grandeur qui contrôle le phénomène. Ici, l'idée centrale est donnée par :
\[
\sin\theta=\frac1M
\]
Cette relation permet de relier la cause physique à l'effet observé. On vérifie les unités, puis on insère des ordres de grandeur réalistes.

La conclusion attendue est la suivante : Relier angle et nombre de Mach. La réponse doit toujours préciser les hypothèses utilisées. Si le phénomène dépend d'un milieu réel, d'une géométrie complexe, de frottements, de pertes ou de gradients, le résultat obtenu est un ordre de grandeur et non une valeur exacte.
\end{corrigebox}

\subsubsection*{Partie D — Approfondissement niveau Bac+3}

\begin{approfondissementbox}[Réponse Bac+3]
Au niveau Bac+3, on ne se contente plus d'une formule globale. On remplace le modèle simple par une loi locale, différentielle ou énergétique.

Une forme générale possible est :
\[
\frac{\mathrm d X}{\mathrm dt}=\text{apport}-\text{pertes},
\]
ou, pour un phénomène propagatif :
\[
\frac{\partial^2 \psi}{\partial t^2}=c^2\Delta\psi.
\]
Selon le cas, on peut aussi utiliser un principe variationnel, par exemple le principe de Fermat en optique :
\[
\delta \int n\,\mathrm ds=0.
\]

Pour ce phénomène, le modèle Bac+3 consiste à partir de la relation centrale
\[
\sin\theta=\frac1M
\]
puis à introduire les corrections négligées : dépendance spatiale, dépendance temporelle, non-linéarités, pertes, conditions aux limites ou fluctuations. Cela permet de passer d'une réponse d'ordre de grandeur à une prédiction plus robuste.
\end{approfondissementbox}

\subsubsection*{Partie E — Limites, critique et prolongement}

\begin{enumerate}
  \item Citer au moins deux limites du modèle Terminale.
  \item Dire quelles grandeurs seraient difficiles à mesurer expérimentalement.
  \item Proposer une expérience simple permettant de tester le modèle.
  \item Expliquer ce qui changerait si l'on doublait une grandeur clé.
  \item Rédiger une conclusion courte, accessible à un élève de Terminale.
\end{enumerate}

\begin{corrigebox}[Synthèse rédigée]
Le phénomène « Bang supersonique » s'explique par un petit nombre de lois physiques simples, mais sa manifestation réelle dépend souvent de nombreux paramètres. Le modèle Terminale donne l'intuition et les ordres de grandeur ; le modèle Bac+3 introduit les variations locales, les pertes, les équations différentielles ou les principes variationnels nécessaires pour décrire finement la réalité.
\end{corrigebox}

% ------------------------------------------------------------
\subsection{Problème corrigé 58 — Résonance d'un verre}
% ------------------------------------------------------------

\begin{exercicebox}[Question-problème]
\textbf{Observation.} Un verre vibre fortement à certaines fréquences.

\medskip
\textbf{Question.} Comprendre l'amplification résonante.

\medskip
\textbf{Thèmes mobilisés :} Oscillateur forcé.

\medskip
\textbf{Relation ou idée centrale :}
\[
\omega_0=\sqrt{\frac{k}{m}}
\]
\end{exercicebox}

\subsubsection*{Partie A — Modélisation niveau Terminale}

\begin{enumerate}
  \item Décrire précisément le phénomène observé.
  \item Identifier le système physique étudié.
  \item Lister les grandeurs importantes et leurs unités.
  \item Écrire la loi physique simple qui permet de commencer le raisonnement.
  \item Expliquer les hypothèses simplificatrices du modèle Terminale.
\end{enumerate}

\subsubsection*{Partie B — Mise en équation et calcul}

\begin{enumerate}
  \item Définir les variables utiles.
  \item Écrire la relation mathématique principale.
  \item Vérifier l'homogénéité de la formule.
  \item Choisir un ordre de grandeur réaliste pour les données manquantes.
  \item Effectuer le calcul ou établir le raisonnement qualitatif demandé.
  \item Interpréter le résultat en une phrase claire.
\end{enumerate}

\subsubsection*{Partie C — Réponse accessible niveau Terminale}

\begin{corrigebox}[Réponse Terminale]
On part de l'observation concrète : Un verre vibre fortement à certaines fréquences.

Au niveau Terminale, on cherche d'abord la grandeur qui contrôle le phénomène. Ici, l'idée centrale est donnée par :
\[
\omega_0=\sqrt{\frac{k}{m}}
\]
Cette relation permet de relier la cause physique à l'effet observé. On vérifie les unités, puis on insère des ordres de grandeur réalistes.

La conclusion attendue est la suivante : Comprendre l'amplification résonante. La réponse doit toujours préciser les hypothèses utilisées. Si le phénomène dépend d'un milieu réel, d'une géométrie complexe, de frottements, de pertes ou de gradients, le résultat obtenu est un ordre de grandeur et non une valeur exacte.
\end{corrigebox}

\subsubsection*{Partie D — Approfondissement niveau Bac+3}

\begin{approfondissementbox}[Réponse Bac+3]
Au niveau Bac+3, on ne se contente plus d'une formule globale. On remplace le modèle simple par une loi locale, différentielle ou énergétique.

Une forme générale possible est :
\[
\frac{\mathrm d X}{\mathrm dt}=\text{apport}-\text{pertes},
\]
ou, pour un phénomène propagatif :
\[
\frac{\partial^2 \psi}{\partial t^2}=c^2\Delta\psi.
\]
Selon le cas, on peut aussi utiliser un principe variationnel, par exemple le principe de Fermat en optique :
\[
\delta \int n\,\mathrm ds=0.
\]

Pour ce phénomène, le modèle Bac+3 consiste à partir de la relation centrale
\[
\omega_0=\sqrt{\frac{k}{m}}
\]
puis à introduire les corrections négligées : dépendance spatiale, dépendance temporelle, non-linéarités, pertes, conditions aux limites ou fluctuations. Cela permet de passer d'une réponse d'ordre de grandeur à une prédiction plus robuste.
\end{approfondissementbox}

\subsubsection*{Partie E — Limites, critique et prolongement}

\begin{enumerate}
  \item Citer au moins deux limites du modèle Terminale.
  \item Dire quelles grandeurs seraient difficiles à mesurer expérimentalement.
  \item Proposer une expérience simple permettant de tester le modèle.
  \item Expliquer ce qui changerait si l'on doublait une grandeur clé.
  \item Rédiger une conclusion courte, accessible à un élève de Terminale.
\end{enumerate}

\begin{corrigebox}[Synthèse rédigée]
Le phénomène « Résonance d'un verre » s'explique par un petit nombre de lois physiques simples, mais sa manifestation réelle dépend souvent de nombreux paramètres. Le modèle Terminale donne l'intuition et les ordres de grandeur ; le modèle Bac+3 introduit les variations locales, les pertes, les équations différentielles ou les principes variationnels nécessaires pour décrire finement la réalité.
\end{corrigebox}

% ------------------------------------------------------------
\subsection{Problème corrigé 59 — Résonance d'un pont}
% ------------------------------------------------------------

\begin{exercicebox}[Question-problème]
\textbf{Observation.} Un pont peut osciller dangereusement.

\medskip
\textbf{Question.} Relier excitation et mode propre.

\medskip
\textbf{Thèmes mobilisés :} Modes propres, excitation.

\medskip
\textbf{Relation ou idée centrale :}
\[
\omega\simeq\omega_0
\]
\end{exercicebox}

\subsubsection*{Partie A — Modélisation niveau Terminale}

\begin{enumerate}
  \item Décrire précisément le phénomène observé.
  \item Identifier le système physique étudié.
  \item Lister les grandeurs importantes et leurs unités.
  \item Écrire la loi physique simple qui permet de commencer le raisonnement.
  \item Expliquer les hypothèses simplificatrices du modèle Terminale.
\end{enumerate}

\subsubsection*{Partie B — Mise en équation et calcul}

\begin{enumerate}
  \item Définir les variables utiles.
  \item Écrire la relation mathématique principale.
  \item Vérifier l'homogénéité de la formule.
  \item Choisir un ordre de grandeur réaliste pour les données manquantes.
  \item Effectuer le calcul ou établir le raisonnement qualitatif demandé.
  \item Interpréter le résultat en une phrase claire.
\end{enumerate}

\subsubsection*{Partie C — Réponse accessible niveau Terminale}

\begin{corrigebox}[Réponse Terminale]
On part de l'observation concrète : Un pont peut osciller dangereusement.

Au niveau Terminale, on cherche d'abord la grandeur qui contrôle le phénomène. Ici, l'idée centrale est donnée par :
\[
\omega\simeq\omega_0
\]
Cette relation permet de relier la cause physique à l'effet observé. On vérifie les unités, puis on insère des ordres de grandeur réalistes.

La conclusion attendue est la suivante : Relier excitation et mode propre. La réponse doit toujours préciser les hypothèses utilisées. Si le phénomène dépend d'un milieu réel, d'une géométrie complexe, de frottements, de pertes ou de gradients, le résultat obtenu est un ordre de grandeur et non une valeur exacte.
\end{corrigebox}

\subsubsection*{Partie D — Approfondissement niveau Bac+3}

\begin{approfondissementbox}[Réponse Bac+3]
Au niveau Bac+3, on ne se contente plus d'une formule globale. On remplace le modèle simple par une loi locale, différentielle ou énergétique.

Une forme générale possible est :
\[
\frac{\mathrm d X}{\mathrm dt}=\text{apport}-\text{pertes},
\]
ou, pour un phénomène propagatif :
\[
\frac{\partial^2 \psi}{\partial t^2}=c^2\Delta\psi.
\]
Selon le cas, on peut aussi utiliser un principe variationnel, par exemple le principe de Fermat en optique :
\[
\delta \int n\,\mathrm ds=0.
\]

Pour ce phénomène, le modèle Bac+3 consiste à partir de la relation centrale
\[
\omega\simeq\omega_0
\]
puis à introduire les corrections négligées : dépendance spatiale, dépendance temporelle, non-linéarités, pertes, conditions aux limites ou fluctuations. Cela permet de passer d'une réponse d'ordre de grandeur à une prédiction plus robuste.
\end{approfondissementbox}

\subsubsection*{Partie E — Limites, critique et prolongement}

\begin{enumerate}
  \item Citer au moins deux limites du modèle Terminale.
  \item Dire quelles grandeurs seraient difficiles à mesurer expérimentalement.
  \item Proposer une expérience simple permettant de tester le modèle.
  \item Expliquer ce qui changerait si l'on doublait une grandeur clé.
  \item Rédiger une conclusion courte, accessible à un élève de Terminale.
\end{enumerate}

\begin{corrigebox}[Synthèse rédigée]
Le phénomène « Résonance d'un pont » s'explique par un petit nombre de lois physiques simples, mais sa manifestation réelle dépend souvent de nombreux paramètres. Le modèle Terminale donne l'intuition et les ordres de grandeur ; le modèle Bac+3 introduit les variations locales, les pertes, les équations différentielles ou les principes variationnels nécessaires pour décrire finement la réalité.
\end{corrigebox}

% ------------------------------------------------------------
\subsection{Problème corrigé 60 — Instrument à corde}
% ------------------------------------------------------------

\begin{exercicebox}[Question-problème]
\textbf{Observation.} Une corde vibre à des harmoniques.

\medskip
\textbf{Question.} Relier longueur et hauteur du son.

\medskip
\textbf{Thèmes mobilisés :} Onde stationnaire.

\medskip
\textbf{Relation ou idée centrale :}
\[
f_n=n\frac{v}{2L}
\]
\end{exercicebox}

\subsubsection*{Partie A — Modélisation niveau Terminale}

\begin{enumerate}
  \item Décrire précisément le phénomène observé.
  \item Identifier le système physique étudié.
  \item Lister les grandeurs importantes et leurs unités.
  \item Écrire la loi physique simple qui permet de commencer le raisonnement.
  \item Expliquer les hypothèses simplificatrices du modèle Terminale.
\end{enumerate}

\subsubsection*{Partie B — Mise en équation et calcul}

\begin{enumerate}
  \item Définir les variables utiles.
  \item Écrire la relation mathématique principale.
  \item Vérifier l'homogénéité de la formule.
  \item Choisir un ordre de grandeur réaliste pour les données manquantes.
  \item Effectuer le calcul ou établir le raisonnement qualitatif demandé.
  \item Interpréter le résultat en une phrase claire.
\end{enumerate}

\subsubsection*{Partie C — Réponse accessible niveau Terminale}

\begin{corrigebox}[Réponse Terminale]
On part de l'observation concrète : Une corde vibre à des harmoniques.

Au niveau Terminale, on cherche d'abord la grandeur qui contrôle le phénomène. Ici, l'idée centrale est donnée par :
\[
f_n=n\frac{v}{2L}
\]
Cette relation permet de relier la cause physique à l'effet observé. On vérifie les unités, puis on insère des ordres de grandeur réalistes.

La conclusion attendue est la suivante : Relier longueur et hauteur du son. La réponse doit toujours préciser les hypothèses utilisées. Si le phénomène dépend d'un milieu réel, d'une géométrie complexe, de frottements, de pertes ou de gradients, le résultat obtenu est un ordre de grandeur et non une valeur exacte.
\end{corrigebox}

\subsubsection*{Partie D — Approfondissement niveau Bac+3}

\begin{approfondissementbox}[Réponse Bac+3]
Au niveau Bac+3, on ne se contente plus d'une formule globale. On remplace le modèle simple par une loi locale, différentielle ou énergétique.

Une forme générale possible est :
\[
\frac{\mathrm d X}{\mathrm dt}=\text{apport}-\text{pertes},
\]
ou, pour un phénomène propagatif :
\[
\frac{\partial^2 \psi}{\partial t^2}=c^2\Delta\psi.
\]
Selon le cas, on peut aussi utiliser un principe variationnel, par exemple le principe de Fermat en optique :
\[
\delta \int n\,\mathrm ds=0.
\]

Pour ce phénomène, le modèle Bac+3 consiste à partir de la relation centrale
\[
f_n=n\frac{v}{2L}
\]
puis à introduire les corrections négligées : dépendance spatiale, dépendance temporelle, non-linéarités, pertes, conditions aux limites ou fluctuations. Cela permet de passer d'une réponse d'ordre de grandeur à une prédiction plus robuste.
\end{approfondissementbox}

\subsubsection*{Partie E — Limites, critique et prolongement}

\begin{enumerate}
  \item Citer au moins deux limites du modèle Terminale.
  \item Dire quelles grandeurs seraient difficiles à mesurer expérimentalement.
  \item Proposer une expérience simple permettant de tester le modèle.
  \item Expliquer ce qui changerait si l'on doublait une grandeur clé.
  \item Rédiger une conclusion courte, accessible à un élève de Terminale.
\end{enumerate}

\begin{corrigebox}[Synthèse rédigée]
Le phénomène « Instrument à corde » s'explique par un petit nombre de lois physiques simples, mais sa manifestation réelle dépend souvent de nombreux paramètres. Le modèle Terminale donne l'intuition et les ordres de grandeur ; le modèle Bac+3 introduit les variations locales, les pertes, les équations différentielles ou les principes variationnels nécessaires pour décrire finement la réalité.
\end{corrigebox}

% ------------------------------------------------------------
\subsection{Problème corrigé 61 — Instrument à vent}
% ------------------------------------------------------------

\begin{exercicebox}[Question-problème]
\textbf{Observation.} Une colonne d'air résonne.

\medskip
\textbf{Question.} Relier longueur du tube et fréquence.

\medskip
\textbf{Thèmes mobilisés :} Modes acoustiques.

\medskip
\textbf{Relation ou idée centrale :}
\[
f_n=n\frac{c}{2L}
\]
\end{exercicebox}

\subsubsection*{Partie A — Modélisation niveau Terminale}

\begin{enumerate}
  \item Décrire précisément le phénomène observé.
  \item Identifier le système physique étudié.
  \item Lister les grandeurs importantes et leurs unités.
  \item Écrire la loi physique simple qui permet de commencer le raisonnement.
  \item Expliquer les hypothèses simplificatrices du modèle Terminale.
\end{enumerate}

\subsubsection*{Partie B — Mise en équation et calcul}

\begin{enumerate}
  \item Définir les variables utiles.
  \item Écrire la relation mathématique principale.
  \item Vérifier l'homogénéité de la formule.
  \item Choisir un ordre de grandeur réaliste pour les données manquantes.
  \item Effectuer le calcul ou établir le raisonnement qualitatif demandé.
  \item Interpréter le résultat en une phrase claire.
\end{enumerate}

\subsubsection*{Partie C — Réponse accessible niveau Terminale}

\begin{corrigebox}[Réponse Terminale]
On part de l'observation concrète : Une colonne d'air résonne.

Au niveau Terminale, on cherche d'abord la grandeur qui contrôle le phénomène. Ici, l'idée centrale est donnée par :
\[
f_n=n\frac{c}{2L}
\]
Cette relation permet de relier la cause physique à l'effet observé. On vérifie les unités, puis on insère des ordres de grandeur réalistes.

La conclusion attendue est la suivante : Relier longueur du tube et fréquence. La réponse doit toujours préciser les hypothèses utilisées. Si le phénomène dépend d'un milieu réel, d'une géométrie complexe, de frottements, de pertes ou de gradients, le résultat obtenu est un ordre de grandeur et non une valeur exacte.
\end{corrigebox}

\subsubsection*{Partie D — Approfondissement niveau Bac+3}

\begin{approfondissementbox}[Réponse Bac+3]
Au niveau Bac+3, on ne se contente plus d'une formule globale. On remplace le modèle simple par une loi locale, différentielle ou énergétique.

Une forme générale possible est :
\[
\frac{\mathrm d X}{\mathrm dt}=\text{apport}-\text{pertes},
\]
ou, pour un phénomène propagatif :
\[
\frac{\partial^2 \psi}{\partial t^2}=c^2\Delta\psi.
\]
Selon le cas, on peut aussi utiliser un principe variationnel, par exemple le principe de Fermat en optique :
\[
\delta \int n\,\mathrm ds=0.
\]

Pour ce phénomène, le modèle Bac+3 consiste à partir de la relation centrale
\[
f_n=n\frac{c}{2L}
\]
puis à introduire les corrections négligées : dépendance spatiale, dépendance temporelle, non-linéarités, pertes, conditions aux limites ou fluctuations. Cela permet de passer d'une réponse d'ordre de grandeur à une prédiction plus robuste.
\end{approfondissementbox}

\subsubsection*{Partie E — Limites, critique et prolongement}

\begin{enumerate}
  \item Citer au moins deux limites du modèle Terminale.
  \item Dire quelles grandeurs seraient difficiles à mesurer expérimentalement.
  \item Proposer une expérience simple permettant de tester le modèle.
  \item Expliquer ce qui changerait si l'on doublait une grandeur clé.
  \item Rédiger une conclusion courte, accessible à un élève de Terminale.
\end{enumerate}

\begin{corrigebox}[Synthèse rédigée]
Le phénomène « Instrument à vent » s'explique par un petit nombre de lois physiques simples, mais sa manifestation réelle dépend souvent de nombreux paramètres. Le modèle Terminale donne l'intuition et les ordres de grandeur ; le modèle Bac+3 introduit les variations locales, les pertes, les équations différentielles ou les principes variationnels nécessaires pour décrire finement la réalité.
\end{corrigebox}

% ------------------------------------------------------------
\subsection{Problème corrigé 62 — Écho}
% ------------------------------------------------------------

\begin{exercicebox}[Question-problème]
\textbf{Observation.} Un son revient après réflexion.

\medskip
\textbf{Question.} Mesurer une distance par écho.

\medskip
\textbf{Thèmes mobilisés :} Aller-retour sonore.

\medskip
\textbf{Relation ou idée centrale :}
\[
d=\frac{c_s\Delta t}{2}
\]
\end{exercicebox}

\subsubsection*{Partie A — Modélisation niveau Terminale}

\begin{enumerate}
  \item Décrire précisément le phénomène observé.
  \item Identifier le système physique étudié.
  \item Lister les grandeurs importantes et leurs unités.
  \item Écrire la loi physique simple qui permet de commencer le raisonnement.
  \item Expliquer les hypothèses simplificatrices du modèle Terminale.
\end{enumerate}

\subsubsection*{Partie B — Mise en équation et calcul}

\begin{enumerate}
  \item Définir les variables utiles.
  \item Écrire la relation mathématique principale.
  \item Vérifier l'homogénéité de la formule.
  \item Choisir un ordre de grandeur réaliste pour les données manquantes.
  \item Effectuer le calcul ou établir le raisonnement qualitatif demandé.
  \item Interpréter le résultat en une phrase claire.
\end{enumerate}

\subsubsection*{Partie C — Réponse accessible niveau Terminale}

\begin{corrigebox}[Réponse Terminale]
On part de l'observation concrète : Un son revient après réflexion.

Au niveau Terminale, on cherche d'abord la grandeur qui contrôle le phénomène. Ici, l'idée centrale est donnée par :
\[
d=\frac{c_s\Delta t}{2}
\]
Cette relation permet de relier la cause physique à l'effet observé. On vérifie les unités, puis on insère des ordres de grandeur réalistes.

La conclusion attendue est la suivante : Mesurer une distance par écho. La réponse doit toujours préciser les hypothèses utilisées. Si le phénomène dépend d'un milieu réel, d'une géométrie complexe, de frottements, de pertes ou de gradients, le résultat obtenu est un ordre de grandeur et non une valeur exacte.
\end{corrigebox}

\subsubsection*{Partie D — Approfondissement niveau Bac+3}

\begin{approfondissementbox}[Réponse Bac+3]
Au niveau Bac+3, on ne se contente plus d'une formule globale. On remplace le modèle simple par une loi locale, différentielle ou énergétique.

Une forme générale possible est :
\[
\frac{\mathrm d X}{\mathrm dt}=\text{apport}-\text{pertes},
\]
ou, pour un phénomène propagatif :
\[
\frac{\partial^2 \psi}{\partial t^2}=c^2\Delta\psi.
\]
Selon le cas, on peut aussi utiliser un principe variationnel, par exemple le principe de Fermat en optique :
\[
\delta \int n\,\mathrm ds=0.
\]

Pour ce phénomène, le modèle Bac+3 consiste à partir de la relation centrale
\[
d=\frac{c_s\Delta t}{2}
\]
puis à introduire les corrections négligées : dépendance spatiale, dépendance temporelle, non-linéarités, pertes, conditions aux limites ou fluctuations. Cela permet de passer d'une réponse d'ordre de grandeur à une prédiction plus robuste.
\end{approfondissementbox}

\subsubsection*{Partie E — Limites, critique et prolongement}

\begin{enumerate}
  \item Citer au moins deux limites du modèle Terminale.
  \item Dire quelles grandeurs seraient difficiles à mesurer expérimentalement.
  \item Proposer une expérience simple permettant de tester le modèle.
  \item Expliquer ce qui changerait si l'on doublait une grandeur clé.
  \item Rédiger une conclusion courte, accessible à un élève de Terminale.
\end{enumerate}

\begin{corrigebox}[Synthèse rédigée]
Le phénomène « Écho » s'explique par un petit nombre de lois physiques simples, mais sa manifestation réelle dépend souvent de nombreux paramètres. Le modèle Terminale donne l'intuition et les ordres de grandeur ; le modèle Bac+3 introduit les variations locales, les pertes, les équations différentielles ou les principes variationnels nécessaires pour décrire finement la réalité.
\end{corrigebox}

% ------------------------------------------------------------
\subsection{Problème corrigé 63 — Acoustique d'une salle}
% ------------------------------------------------------------

\begin{exercicebox}[Question-problème]
\textbf{Observation.} Une salle réverbère plus ou moins.

\medskip
\textbf{Question.} Relier absorption et temps de réverbération.

\medskip
\textbf{Thèmes mobilisés :} Réflexions multiples, absorption.

\medskip
\textbf{Relation ou idée centrale :}
\[
T_R\simeq0.16\frac{V}{A}
\]
\end{exercicebox}

\subsubsection*{Partie A — Modélisation niveau Terminale}

\begin{enumerate}
  \item Décrire précisément le phénomène observé.
  \item Identifier le système physique étudié.
  \item Lister les grandeurs importantes et leurs unités.
  \item Écrire la loi physique simple qui permet de commencer le raisonnement.
  \item Expliquer les hypothèses simplificatrices du modèle Terminale.
\end{enumerate}

\subsubsection*{Partie B — Mise en équation et calcul}

\begin{enumerate}
  \item Définir les variables utiles.
  \item Écrire la relation mathématique principale.
  \item Vérifier l'homogénéité de la formule.
  \item Choisir un ordre de grandeur réaliste pour les données manquantes.
  \item Effectuer le calcul ou établir le raisonnement qualitatif demandé.
  \item Interpréter le résultat en une phrase claire.
\end{enumerate}

\subsubsection*{Partie C — Réponse accessible niveau Terminale}

\begin{corrigebox}[Réponse Terminale]
On part de l'observation concrète : Une salle réverbère plus ou moins.

Au niveau Terminale, on cherche d'abord la grandeur qui contrôle le phénomène. Ici, l'idée centrale est donnée par :
\[
T_R\simeq0.16\frac{V}{A}
\]
Cette relation permet de relier la cause physique à l'effet observé. On vérifie les unités, puis on insère des ordres de grandeur réalistes.

La conclusion attendue est la suivante : Relier absorption et temps de réverbération. La réponse doit toujours préciser les hypothèses utilisées. Si le phénomène dépend d'un milieu réel, d'une géométrie complexe, de frottements, de pertes ou de gradients, le résultat obtenu est un ordre de grandeur et non une valeur exacte.
\end{corrigebox}

\subsubsection*{Partie D — Approfondissement niveau Bac+3}

\begin{approfondissementbox}[Réponse Bac+3]
Au niveau Bac+3, on ne se contente plus d'une formule globale. On remplace le modèle simple par une loi locale, différentielle ou énergétique.

Une forme générale possible est :
\[
\frac{\mathrm d X}{\mathrm dt}=\text{apport}-\text{pertes},
\]
ou, pour un phénomène propagatif :
\[
\frac{\partial^2 \psi}{\partial t^2}=c^2\Delta\psi.
\]
Selon le cas, on peut aussi utiliser un principe variationnel, par exemple le principe de Fermat en optique :
\[
\delta \int n\,\mathrm ds=0.
\]

Pour ce phénomène, le modèle Bac+3 consiste à partir de la relation centrale
\[
T_R\simeq0.16\frac{V}{A}
\]
puis à introduire les corrections négligées : dépendance spatiale, dépendance temporelle, non-linéarités, pertes, conditions aux limites ou fluctuations. Cela permet de passer d'une réponse d'ordre de grandeur à une prédiction plus robuste.
\end{approfondissementbox}

\subsubsection*{Partie E — Limites, critique et prolongement}

\begin{enumerate}
  \item Citer au moins deux limites du modèle Terminale.
  \item Dire quelles grandeurs seraient difficiles à mesurer expérimentalement.
  \item Proposer une expérience simple permettant de tester le modèle.
  \item Expliquer ce qui changerait si l'on doublait une grandeur clé.
  \item Rédiger une conclusion courte, accessible à un élève de Terminale.
\end{enumerate}

\begin{corrigebox}[Synthèse rédigée]
Le phénomène « Acoustique d'une salle » s'explique par un petit nombre de lois physiques simples, mais sa manifestation réelle dépend souvent de nombreux paramètres. Le modèle Terminale donne l'intuition et les ordres de grandeur ; le modèle Bac+3 introduit les variations locales, les pertes, les équations différentielles ou les principes variationnels nécessaires pour décrire finement la réalité.
\end{corrigebox}

% ------------------------------------------------------------
\subsection{Problème corrigé 64 — Balle de tennis liftée}
% ------------------------------------------------------------

\begin{exercicebox}[Question-problème]
\textbf{Observation.} Une balle liftée retombe vite.

\medskip
\textbf{Question.} Expliquer le topspin.

\medskip
\textbf{Thèmes mobilisés :} Effet Magnus.

\medskip
\textbf{Relation ou idée centrale :}
\[
\vec F_M\propto\vec\omega\times\vec v
\]
\end{exercicebox}

\subsubsection*{Partie A — Modélisation niveau Terminale}

\begin{enumerate}
  \item Décrire précisément le phénomène observé.
  \item Identifier le système physique étudié.
  \item Lister les grandeurs importantes et leurs unités.
  \item Écrire la loi physique simple qui permet de commencer le raisonnement.
  \item Expliquer les hypothèses simplificatrices du modèle Terminale.
\end{enumerate}

\subsubsection*{Partie B — Mise en équation et calcul}

\begin{enumerate}
  \item Définir les variables utiles.
  \item Écrire la relation mathématique principale.
  \item Vérifier l'homogénéité de la formule.
  \item Choisir un ordre de grandeur réaliste pour les données manquantes.
  \item Effectuer le calcul ou établir le raisonnement qualitatif demandé.
  \item Interpréter le résultat en une phrase claire.
\end{enumerate}

\subsubsection*{Partie C — Réponse accessible niveau Terminale}

\begin{corrigebox}[Réponse Terminale]
On part de l'observation concrète : Une balle liftée retombe vite.

Au niveau Terminale, on cherche d'abord la grandeur qui contrôle le phénomène. Ici, l'idée centrale est donnée par :
\[
\vec F_M\propto\vec\omega\times\vec v
\]
Cette relation permet de relier la cause physique à l'effet observé. On vérifie les unités, puis on insère des ordres de grandeur réalistes.

La conclusion attendue est la suivante : Expliquer le topspin. La réponse doit toujours préciser les hypothèses utilisées. Si le phénomène dépend d'un milieu réel, d'une géométrie complexe, de frottements, de pertes ou de gradients, le résultat obtenu est un ordre de grandeur et non une valeur exacte.
\end{corrigebox}

\subsubsection*{Partie D — Approfondissement niveau Bac+3}

\begin{approfondissementbox}[Réponse Bac+3]
Au niveau Bac+3, on ne se contente plus d'une formule globale. On remplace le modèle simple par une loi locale, différentielle ou énergétique.

Une forme générale possible est :
\[
\frac{\mathrm d X}{\mathrm dt}=\text{apport}-\text{pertes},
\]
ou, pour un phénomène propagatif :
\[
\frac{\partial^2 \psi}{\partial t^2}=c^2\Delta\psi.
\]
Selon le cas, on peut aussi utiliser un principe variationnel, par exemple le principe de Fermat en optique :
\[
\delta \int n\,\mathrm ds=0.
\]

Pour ce phénomène, le modèle Bac+3 consiste à partir de la relation centrale
\[
\vec F_M\propto\vec\omega\times\vec v
\]
puis à introduire les corrections négligées : dépendance spatiale, dépendance temporelle, non-linéarités, pertes, conditions aux limites ou fluctuations. Cela permet de passer d'une réponse d'ordre de grandeur à une prédiction plus robuste.
\end{approfondissementbox}

\subsubsection*{Partie E — Limites, critique et prolongement}

\begin{enumerate}
  \item Citer au moins deux limites du modèle Terminale.
  \item Dire quelles grandeurs seraient difficiles à mesurer expérimentalement.
  \item Proposer une expérience simple permettant de tester le modèle.
  \item Expliquer ce qui changerait si l'on doublait une grandeur clé.
  \item Rédiger une conclusion courte, accessible à un élève de Terminale.
\end{enumerate}

\begin{corrigebox}[Synthèse rédigée]
Le phénomène « Balle de tennis liftée » s'explique par un petit nombre de lois physiques simples, mais sa manifestation réelle dépend souvent de nombreux paramètres. Le modèle Terminale donne l'intuition et les ordres de grandeur ; le modèle Bac+3 introduit les variations locales, les pertes, les équations différentielles ou les principes variationnels nécessaires pour décrire finement la réalité.
\end{corrigebox}

% ------------------------------------------------------------
\subsection{Problème corrigé 65 — Coup franc brossé}
% ------------------------------------------------------------

\begin{exercicebox}[Question-problème]
\textbf{Observation.} Un ballon contourne un mur.

\medskip
\textbf{Question.} Relier rotation et trajectoire courbe.

\medskip
\textbf{Thèmes mobilisés :} Effet Magnus latéral.

\medskip
\textbf{Relation ou idée centrale :}
\[
a_\perp=\frac{F_M}{m}
\]
\end{exercicebox}

\subsubsection*{Partie A — Modélisation niveau Terminale}

\begin{enumerate}
  \item Décrire précisément le phénomène observé.
  \item Identifier le système physique étudié.
  \item Lister les grandeurs importantes et leurs unités.
  \item Écrire la loi physique simple qui permet de commencer le raisonnement.
  \item Expliquer les hypothèses simplificatrices du modèle Terminale.
\end{enumerate}

\subsubsection*{Partie B — Mise en équation et calcul}

\begin{enumerate}
  \item Définir les variables utiles.
  \item Écrire la relation mathématique principale.
  \item Vérifier l'homogénéité de la formule.
  \item Choisir un ordre de grandeur réaliste pour les données manquantes.
  \item Effectuer le calcul ou établir le raisonnement qualitatif demandé.
  \item Interpréter le résultat en une phrase claire.
\end{enumerate}

\subsubsection*{Partie C — Réponse accessible niveau Terminale}

\begin{corrigebox}[Réponse Terminale]
On part de l'observation concrète : Un ballon contourne un mur.

Au niveau Terminale, on cherche d'abord la grandeur qui contrôle le phénomène. Ici, l'idée centrale est donnée par :
\[
a_\perp=\frac{F_M}{m}
\]
Cette relation permet de relier la cause physique à l'effet observé. On vérifie les unités, puis on insère des ordres de grandeur réalistes.

La conclusion attendue est la suivante : Relier rotation et trajectoire courbe. La réponse doit toujours préciser les hypothèses utilisées. Si le phénomène dépend d'un milieu réel, d'une géométrie complexe, de frottements, de pertes ou de gradients, le résultat obtenu est un ordre de grandeur et non une valeur exacte.
\end{corrigebox}

\subsubsection*{Partie D — Approfondissement niveau Bac+3}

\begin{approfondissementbox}[Réponse Bac+3]
Au niveau Bac+3, on ne se contente plus d'une formule globale. On remplace le modèle simple par une loi locale, différentielle ou énergétique.

Une forme générale possible est :
\[
\frac{\mathrm d X}{\mathrm dt}=\text{apport}-\text{pertes},
\]
ou, pour un phénomène propagatif :
\[
\frac{\partial^2 \psi}{\partial t^2}=c^2\Delta\psi.
\]
Selon le cas, on peut aussi utiliser un principe variationnel, par exemple le principe de Fermat en optique :
\[
\delta \int n\,\mathrm ds=0.
\]

Pour ce phénomène, le modèle Bac+3 consiste à partir de la relation centrale
\[
a_\perp=\frac{F_M}{m}
\]
puis à introduire les corrections négligées : dépendance spatiale, dépendance temporelle, non-linéarités, pertes, conditions aux limites ou fluctuations. Cela permet de passer d'une réponse d'ordre de grandeur à une prédiction plus robuste.
\end{approfondissementbox}

\subsubsection*{Partie E — Limites, critique et prolongement}

\begin{enumerate}
  \item Citer au moins deux limites du modèle Terminale.
  \item Dire quelles grandeurs seraient difficiles à mesurer expérimentalement.
  \item Proposer une expérience simple permettant de tester le modèle.
  \item Expliquer ce qui changerait si l'on doublait une grandeur clé.
  \item Rédiger une conclusion courte, accessible à un élève de Terminale.
\end{enumerate}

\begin{corrigebox}[Synthèse rédigée]
Le phénomène « Coup franc brossé » s'explique par un petit nombre de lois physiques simples, mais sa manifestation réelle dépend souvent de nombreux paramètres. Le modèle Terminale donne l'intuition et les ordres de grandeur ; le modèle Bac+3 introduit les variations locales, les pertes, les équations différentielles ou les principes variationnels nécessaires pour décrire finement la réalité.
\end{corrigebox}

% ------------------------------------------------------------
\subsection{Problème corrigé 66 — Rebond d'une balle}
% ------------------------------------------------------------

\begin{exercicebox}[Question-problème]
\textbf{Observation.} Une balle rebondit moins haut.

\medskip
\textbf{Question.} Relier perte d'énergie et hauteur de rebond.

\medskip
\textbf{Thèmes mobilisés :} Coefficient de restitution.

\medskip
\textbf{Relation ou idée centrale :}
\[
e=\frac{v_{\mathrm{après}}}{v_{\mathrm{avant}}}
\]
\end{exercicebox}

\subsubsection*{Partie A — Modélisation niveau Terminale}

\begin{enumerate}
  \item Décrire précisément le phénomène observé.
  \item Identifier le système physique étudié.
  \item Lister les grandeurs importantes et leurs unités.
  \item Écrire la loi physique simple qui permet de commencer le raisonnement.
  \item Expliquer les hypothèses simplificatrices du modèle Terminale.
\end{enumerate}

\subsubsection*{Partie B — Mise en équation et calcul}

\begin{enumerate}
  \item Définir les variables utiles.
  \item Écrire la relation mathématique principale.
  \item Vérifier l'homogénéité de la formule.
  \item Choisir un ordre de grandeur réaliste pour les données manquantes.
  \item Effectuer le calcul ou établir le raisonnement qualitatif demandé.
  \item Interpréter le résultat en une phrase claire.
\end{enumerate}

\subsubsection*{Partie C — Réponse accessible niveau Terminale}

\begin{corrigebox}[Réponse Terminale]
On part de l'observation concrète : Une balle rebondit moins haut.

Au niveau Terminale, on cherche d'abord la grandeur qui contrôle le phénomène. Ici, l'idée centrale est donnée par :
\[
e=\frac{v_{\mathrm{après}}}{v_{\mathrm{avant}}}
\]
Cette relation permet de relier la cause physique à l'effet observé. On vérifie les unités, puis on insère des ordres de grandeur réalistes.

La conclusion attendue est la suivante : Relier perte d'énergie et hauteur de rebond. La réponse doit toujours préciser les hypothèses utilisées. Si le phénomène dépend d'un milieu réel, d'une géométrie complexe, de frottements, de pertes ou de gradients, le résultat obtenu est un ordre de grandeur et non une valeur exacte.
\end{corrigebox}

\subsubsection*{Partie D — Approfondissement niveau Bac+3}

\begin{approfondissementbox}[Réponse Bac+3]
Au niveau Bac+3, on ne se contente plus d'une formule globale. On remplace le modèle simple par une loi locale, différentielle ou énergétique.

Une forme générale possible est :
\[
\frac{\mathrm d X}{\mathrm dt}=\text{apport}-\text{pertes},
\]
ou, pour un phénomène propagatif :
\[
\frac{\partial^2 \psi}{\partial t^2}=c^2\Delta\psi.
\]
Selon le cas, on peut aussi utiliser un principe variationnel, par exemple le principe de Fermat en optique :
\[
\delta \int n\,\mathrm ds=0.
\]

Pour ce phénomène, le modèle Bac+3 consiste à partir de la relation centrale
\[
e=\frac{v_{\mathrm{après}}}{v_{\mathrm{avant}}}
\]
puis à introduire les corrections négligées : dépendance spatiale, dépendance temporelle, non-linéarités, pertes, conditions aux limites ou fluctuations. Cela permet de passer d'une réponse d'ordre de grandeur à une prédiction plus robuste.
\end{approfondissementbox}

\subsubsection*{Partie E — Limites, critique et prolongement}

\begin{enumerate}
  \item Citer au moins deux limites du modèle Terminale.
  \item Dire quelles grandeurs seraient difficiles à mesurer expérimentalement.
  \item Proposer une expérience simple permettant de tester le modèle.
  \item Expliquer ce qui changerait si l'on doublait une grandeur clé.
  \item Rédiger une conclusion courte, accessible à un élève de Terminale.
\end{enumerate}

\begin{corrigebox}[Synthèse rédigée]
Le phénomène « Rebond d'une balle » s'explique par un petit nombre de lois physiques simples, mais sa manifestation réelle dépend souvent de nombreux paramètres. Le modèle Terminale donne l'intuition et les ordres de grandeur ; le modèle Bac+3 introduit les variations locales, les pertes, les équations différentielles ou les principes variationnels nécessaires pour décrire finement la réalité.
\end{corrigebox}

% ------------------------------------------------------------
\subsection{Problème corrigé 67 — Frottements solides}
% ------------------------------------------------------------

\begin{exercicebox}[Question-problème]
\textbf{Observation.} Un objet glisse puis s'arrête.

\medskip
\textbf{Question.} Calculer une distance d'arrêt simple.

\medskip
\textbf{Thèmes mobilisés :} Frottement de Coulomb.

\medskip
\textbf{Relation ou idée centrale :}
\[
F_f=\mu N
\]
\end{exercicebox}

\subsubsection*{Partie A — Modélisation niveau Terminale}

\begin{enumerate}
  \item Décrire précisément le phénomène observé.
  \item Identifier le système physique étudié.
  \item Lister les grandeurs importantes et leurs unités.
  \item Écrire la loi physique simple qui permet de commencer le raisonnement.
  \item Expliquer les hypothèses simplificatrices du modèle Terminale.
\end{enumerate}

\subsubsection*{Partie B — Mise en équation et calcul}

\begin{enumerate}
  \item Définir les variables utiles.
  \item Écrire la relation mathématique principale.
  \item Vérifier l'homogénéité de la formule.
  \item Choisir un ordre de grandeur réaliste pour les données manquantes.
  \item Effectuer le calcul ou établir le raisonnement qualitatif demandé.
  \item Interpréter le résultat en une phrase claire.
\end{enumerate}

\subsubsection*{Partie C — Réponse accessible niveau Terminale}

\begin{corrigebox}[Réponse Terminale]
On part de l'observation concrète : Un objet glisse puis s'arrête.

Au niveau Terminale, on cherche d'abord la grandeur qui contrôle le phénomène. Ici, l'idée centrale est donnée par :
\[
F_f=\mu N
\]
Cette relation permet de relier la cause physique à l'effet observé. On vérifie les unités, puis on insère des ordres de grandeur réalistes.

La conclusion attendue est la suivante : Calculer une distance d'arrêt simple. La réponse doit toujours préciser les hypothèses utilisées. Si le phénomène dépend d'un milieu réel, d'une géométrie complexe, de frottements, de pertes ou de gradients, le résultat obtenu est un ordre de grandeur et non une valeur exacte.
\end{corrigebox}

\subsubsection*{Partie D — Approfondissement niveau Bac+3}

\begin{approfondissementbox}[Réponse Bac+3]
Au niveau Bac+3, on ne se contente plus d'une formule globale. On remplace le modèle simple par une loi locale, différentielle ou énergétique.

Une forme générale possible est :
\[
\frac{\mathrm d X}{\mathrm dt}=\text{apport}-\text{pertes},
\]
ou, pour un phénomène propagatif :
\[
\frac{\partial^2 \psi}{\partial t^2}=c^2\Delta\psi.
\]
Selon le cas, on peut aussi utiliser un principe variationnel, par exemple le principe de Fermat en optique :
\[
\delta \int n\,\mathrm ds=0.
\]

Pour ce phénomène, le modèle Bac+3 consiste à partir de la relation centrale
\[
F_f=\mu N
\]
puis à introduire les corrections négligées : dépendance spatiale, dépendance temporelle, non-linéarités, pertes, conditions aux limites ou fluctuations. Cela permet de passer d'une réponse d'ordre de grandeur à une prédiction plus robuste.
\end{approfondissementbox}

\subsubsection*{Partie E — Limites, critique et prolongement}

\begin{enumerate}
  \item Citer au moins deux limites du modèle Terminale.
  \item Dire quelles grandeurs seraient difficiles à mesurer expérimentalement.
  \item Proposer une expérience simple permettant de tester le modèle.
  \item Expliquer ce qui changerait si l'on doublait une grandeur clé.
  \item Rédiger une conclusion courte, accessible à un élève de Terminale.
\end{enumerate}

\begin{corrigebox}[Synthèse rédigée]
Le phénomène « Frottements solides » s'explique par un petit nombre de lois physiques simples, mais sa manifestation réelle dépend souvent de nombreux paramètres. Le modèle Terminale donne l'intuition et les ordres de grandeur ; le modèle Bac+3 introduit les variations locales, les pertes, les équations différentielles ou les principes variationnels nécessaires pour décrire finement la réalité.
\end{corrigebox}

% ------------------------------------------------------------
\subsection{Problème corrigé 68 — Frottements fluides}
% ------------------------------------------------------------

\begin{exercicebox}[Question-problème]
\textbf{Observation.} La résistance de l'air augmente avec la vitesse.

\medskip
\textbf{Question.} Expliquer la consommation à grande vitesse.

\medskip
\textbf{Thèmes mobilisés :} Traînée.

\medskip
\textbf{Relation ou idée centrale :}
\[
F_D=\frac12\rho C_DSv^2
\]
\end{exercicebox}

\subsubsection*{Partie A — Modélisation niveau Terminale}

\begin{enumerate}
  \item Décrire précisément le phénomène observé.
  \item Identifier le système physique étudié.
  \item Lister les grandeurs importantes et leurs unités.
  \item Écrire la loi physique simple qui permet de commencer le raisonnement.
  \item Expliquer les hypothèses simplificatrices du modèle Terminale.
\end{enumerate}

\subsubsection*{Partie B — Mise en équation et calcul}

\begin{enumerate}
  \item Définir les variables utiles.
  \item Écrire la relation mathématique principale.
  \item Vérifier l'homogénéité de la formule.
  \item Choisir un ordre de grandeur réaliste pour les données manquantes.
  \item Effectuer le calcul ou établir le raisonnement qualitatif demandé.
  \item Interpréter le résultat en une phrase claire.
\end{enumerate}

\subsubsection*{Partie C — Réponse accessible niveau Terminale}

\begin{corrigebox}[Réponse Terminale]
On part de l'observation concrète : La résistance de l'air augmente avec la vitesse.

Au niveau Terminale, on cherche d'abord la grandeur qui contrôle le phénomène. Ici, l'idée centrale est donnée par :
\[
F_D=\frac12\rho C_DSv^2
\]
Cette relation permet de relier la cause physique à l'effet observé. On vérifie les unités, puis on insère des ordres de grandeur réalistes.

La conclusion attendue est la suivante : Expliquer la consommation à grande vitesse. La réponse doit toujours préciser les hypothèses utilisées. Si le phénomène dépend d'un milieu réel, d'une géométrie complexe, de frottements, de pertes ou de gradients, le résultat obtenu est un ordre de grandeur et non une valeur exacte.
\end{corrigebox}

\subsubsection*{Partie D — Approfondissement niveau Bac+3}

\begin{approfondissementbox}[Réponse Bac+3]
Au niveau Bac+3, on ne se contente plus d'une formule globale. On remplace le modèle simple par une loi locale, différentielle ou énergétique.

Une forme générale possible est :
\[
\frac{\mathrm d X}{\mathrm dt}=\text{apport}-\text{pertes},
\]
ou, pour un phénomène propagatif :
\[
\frac{\partial^2 \psi}{\partial t^2}=c^2\Delta\psi.
\]
Selon le cas, on peut aussi utiliser un principe variationnel, par exemple le principe de Fermat en optique :
\[
\delta \int n\,\mathrm ds=0.
\]

Pour ce phénomène, le modèle Bac+3 consiste à partir de la relation centrale
\[
F_D=\frac12\rho C_DSv^2
\]
puis à introduire les corrections négligées : dépendance spatiale, dépendance temporelle, non-linéarités, pertes, conditions aux limites ou fluctuations. Cela permet de passer d'une réponse d'ordre de grandeur à une prédiction plus robuste.
\end{approfondissementbox}

\subsubsection*{Partie E — Limites, critique et prolongement}

\begin{enumerate}
  \item Citer au moins deux limites du modèle Terminale.
  \item Dire quelles grandeurs seraient difficiles à mesurer expérimentalement.
  \item Proposer une expérience simple permettant de tester le modèle.
  \item Expliquer ce qui changerait si l'on doublait une grandeur clé.
  \item Rédiger une conclusion courte, accessible à un élève de Terminale.
\end{enumerate}

\begin{corrigebox}[Synthèse rédigée]
Le phénomène « Frottements fluides » s'explique par un petit nombre de lois physiques simples, mais sa manifestation réelle dépend souvent de nombreux paramètres. Le modèle Terminale donne l'intuition et les ordres de grandeur ; le modèle Bac+3 introduit les variations locales, les pertes, les équations différentielles ou les principes variationnels nécessaires pour décrire finement la réalité.
\end{corrigebox}

% ------------------------------------------------------------
\subsection{Problème corrigé 69 — Aquaplaning}
% ------------------------------------------------------------

\begin{exercicebox}[Question-problème]
\textbf{Observation.} Un pneu perd l'adhérence sur l'eau.

\medskip
\textbf{Question.} Relier risque et vitesse.

\medskip
\textbf{Thèmes mobilisés :} Pression dynamique, évacuation d'eau.

\medskip
\textbf{Relation ou idée centrale :}
\[
P_{\mathrm{dyn}}\sim\rho v^2
\]
\end{exercicebox}

\subsubsection*{Partie A — Modélisation niveau Terminale}

\begin{enumerate}
  \item Décrire précisément le phénomène observé.
  \item Identifier le système physique étudié.
  \item Lister les grandeurs importantes et leurs unités.
  \item Écrire la loi physique simple qui permet de commencer le raisonnement.
  \item Expliquer les hypothèses simplificatrices du modèle Terminale.
\end{enumerate}

\subsubsection*{Partie B — Mise en équation et calcul}

\begin{enumerate}
  \item Définir les variables utiles.
  \item Écrire la relation mathématique principale.
  \item Vérifier l'homogénéité de la formule.
  \item Choisir un ordre de grandeur réaliste pour les données manquantes.
  \item Effectuer le calcul ou établir le raisonnement qualitatif demandé.
  \item Interpréter le résultat en une phrase claire.
\end{enumerate}

\subsubsection*{Partie C — Réponse accessible niveau Terminale}

\begin{corrigebox}[Réponse Terminale]
On part de l'observation concrète : Un pneu perd l'adhérence sur l'eau.

Au niveau Terminale, on cherche d'abord la grandeur qui contrôle le phénomène. Ici, l'idée centrale est donnée par :
\[
P_{\mathrm{dyn}}\sim\rho v^2
\]
Cette relation permet de relier la cause physique à l'effet observé. On vérifie les unités, puis on insère des ordres de grandeur réalistes.

La conclusion attendue est la suivante : Relier risque et vitesse. La réponse doit toujours préciser les hypothèses utilisées. Si le phénomène dépend d'un milieu réel, d'une géométrie complexe, de frottements, de pertes ou de gradients, le résultat obtenu est un ordre de grandeur et non une valeur exacte.
\end{corrigebox}

\subsubsection*{Partie D — Approfondissement niveau Bac+3}

\begin{approfondissementbox}[Réponse Bac+3]
Au niveau Bac+3, on ne se contente plus d'une formule globale. On remplace le modèle simple par une loi locale, différentielle ou énergétique.

Une forme générale possible est :
\[
\frac{\mathrm d X}{\mathrm dt}=\text{apport}-\text{pertes},
\]
ou, pour un phénomène propagatif :
\[
\frac{\partial^2 \psi}{\partial t^2}=c^2\Delta\psi.
\]
Selon le cas, on peut aussi utiliser un principe variationnel, par exemple le principe de Fermat en optique :
\[
\delta \int n\,\mathrm ds=0.
\]

Pour ce phénomène, le modèle Bac+3 consiste à partir de la relation centrale
\[
P_{\mathrm{dyn}}\sim\rho v^2
\]
puis à introduire les corrections négligées : dépendance spatiale, dépendance temporelle, non-linéarités, pertes, conditions aux limites ou fluctuations. Cela permet de passer d'une réponse d'ordre de grandeur à une prédiction plus robuste.
\end{approfondissementbox}

\subsubsection*{Partie E — Limites, critique et prolongement}

\begin{enumerate}
  \item Citer au moins deux limites du modèle Terminale.
  \item Dire quelles grandeurs seraient difficiles à mesurer expérimentalement.
  \item Proposer une expérience simple permettant de tester le modèle.
  \item Expliquer ce qui changerait si l'on doublait une grandeur clé.
  \item Rédiger une conclusion courte, accessible à un élève de Terminale.
\end{enumerate}

\begin{corrigebox}[Synthèse rédigée]
Le phénomène « Aquaplaning » s'explique par un petit nombre de lois physiques simples, mais sa manifestation réelle dépend souvent de nombreux paramètres. Le modèle Terminale donne l'intuition et les ordres de grandeur ; le modèle Bac+3 introduit les variations locales, les pertes, les équations différentielles ou les principes variationnels nécessaires pour décrire finement la réalité.
\end{corrigebox}

% ------------------------------------------------------------
\subsection{Problème corrigé 70 — Portance d'un avion}
% ------------------------------------------------------------

\begin{exercicebox}[Question-problème]
\textbf{Observation.} Une aile porte un avion.

\medskip
\textbf{Question.} Relier vitesse et portance.

\medskip
\textbf{Thèmes mobilisés :} Portance aérodynamique.

\medskip
\textbf{Relation ou idée centrale :}
\[
L=\frac12\rho v^2SC_L
\]
\end{exercicebox}

\subsubsection*{Partie A — Modélisation niveau Terminale}

\begin{enumerate}
  \item Décrire précisément le phénomène observé.
  \item Identifier le système physique étudié.
  \item Lister les grandeurs importantes et leurs unités.
  \item Écrire la loi physique simple qui permet de commencer le raisonnement.
  \item Expliquer les hypothèses simplificatrices du modèle Terminale.
\end{enumerate}

\subsubsection*{Partie B — Mise en équation et calcul}

\begin{enumerate}
  \item Définir les variables utiles.
  \item Écrire la relation mathématique principale.
  \item Vérifier l'homogénéité de la formule.
  \item Choisir un ordre de grandeur réaliste pour les données manquantes.
  \item Effectuer le calcul ou établir le raisonnement qualitatif demandé.
  \item Interpréter le résultat en une phrase claire.
\end{enumerate}

\subsubsection*{Partie C — Réponse accessible niveau Terminale}

\begin{corrigebox}[Réponse Terminale]
On part de l'observation concrète : Une aile porte un avion.

Au niveau Terminale, on cherche d'abord la grandeur qui contrôle le phénomène. Ici, l'idée centrale est donnée par :
\[
L=\frac12\rho v^2SC_L
\]
Cette relation permet de relier la cause physique à l'effet observé. On vérifie les unités, puis on insère des ordres de grandeur réalistes.

La conclusion attendue est la suivante : Relier vitesse et portance. La réponse doit toujours préciser les hypothèses utilisées. Si le phénomène dépend d'un milieu réel, d'une géométrie complexe, de frottements, de pertes ou de gradients, le résultat obtenu est un ordre de grandeur et non une valeur exacte.
\end{corrigebox}

\subsubsection*{Partie D — Approfondissement niveau Bac+3}

\begin{approfondissementbox}[Réponse Bac+3]
Au niveau Bac+3, on ne se contente plus d'une formule globale. On remplace le modèle simple par une loi locale, différentielle ou énergétique.

Une forme générale possible est :
\[
\frac{\mathrm d X}{\mathrm dt}=\text{apport}-\text{pertes},
\]
ou, pour un phénomène propagatif :
\[
\frac{\partial^2 \psi}{\partial t^2}=c^2\Delta\psi.
\]
Selon le cas, on peut aussi utiliser un principe variationnel, par exemple le principe de Fermat en optique :
\[
\delta \int n\,\mathrm ds=0.
\]

Pour ce phénomène, le modèle Bac+3 consiste à partir de la relation centrale
\[
L=\frac12\rho v^2SC_L
\]
puis à introduire les corrections négligées : dépendance spatiale, dépendance temporelle, non-linéarités, pertes, conditions aux limites ou fluctuations. Cela permet de passer d'une réponse d'ordre de grandeur à une prédiction plus robuste.
\end{approfondissementbox}

\subsubsection*{Partie E — Limites, critique et prolongement}

\begin{enumerate}
  \item Citer au moins deux limites du modèle Terminale.
  \item Dire quelles grandeurs seraient difficiles à mesurer expérimentalement.
  \item Proposer une expérience simple permettant de tester le modèle.
  \item Expliquer ce qui changerait si l'on doublait une grandeur clé.
  \item Rédiger une conclusion courte, accessible à un élève de Terminale.
\end{enumerate}

\begin{corrigebox}[Synthèse rédigée]
Le phénomène « Portance d'un avion » s'explique par un petit nombre de lois physiques simples, mais sa manifestation réelle dépend souvent de nombreux paramètres. Le modèle Terminale donne l'intuition et les ordres de grandeur ; le modèle Bac+3 introduit les variations locales, les pertes, les équations différentielles ou les principes variationnels nécessaires pour décrire finement la réalité.
\end{corrigebox}

% ------------------------------------------------------------
\subsection{Problème corrigé 71 — Traînée aérodynamique}
% ------------------------------------------------------------

\begin{exercicebox}[Question-problème]
\textbf{Observation.} Un véhicule résiste à l'air.

\medskip
\textbf{Question.} Comprendre l'intérêt du profilage.

\medskip
\textbf{Thèmes mobilisés :} Transfert de quantité de mouvement.

\medskip
\textbf{Relation ou idée centrale :}
\[
D=\frac12\rho v^2SC_D
\]
\end{exercicebox}

\subsubsection*{Partie A — Modélisation niveau Terminale}

\begin{enumerate}
  \item Décrire précisément le phénomène observé.
  \item Identifier le système physique étudié.
  \item Lister les grandeurs importantes et leurs unités.
  \item Écrire la loi physique simple qui permet de commencer le raisonnement.
  \item Expliquer les hypothèses simplificatrices du modèle Terminale.
\end{enumerate}

\subsubsection*{Partie B — Mise en équation et calcul}

\begin{enumerate}
  \item Définir les variables utiles.
  \item Écrire la relation mathématique principale.
  \item Vérifier l'homogénéité de la formule.
  \item Choisir un ordre de grandeur réaliste pour les données manquantes.
  \item Effectuer le calcul ou établir le raisonnement qualitatif demandé.
  \item Interpréter le résultat en une phrase claire.
\end{enumerate}

\subsubsection*{Partie C — Réponse accessible niveau Terminale}

\begin{corrigebox}[Réponse Terminale]
On part de l'observation concrète : Un véhicule résiste à l'air.

Au niveau Terminale, on cherche d'abord la grandeur qui contrôle le phénomène. Ici, l'idée centrale est donnée par :
\[
D=\frac12\rho v^2SC_D
\]
Cette relation permet de relier la cause physique à l'effet observé. On vérifie les unités, puis on insère des ordres de grandeur réalistes.

La conclusion attendue est la suivante : Comprendre l'intérêt du profilage. La réponse doit toujours préciser les hypothèses utilisées. Si le phénomène dépend d'un milieu réel, d'une géométrie complexe, de frottements, de pertes ou de gradients, le résultat obtenu est un ordre de grandeur et non une valeur exacte.
\end{corrigebox}

\subsubsection*{Partie D — Approfondissement niveau Bac+3}

\begin{approfondissementbox}[Réponse Bac+3]
Au niveau Bac+3, on ne se contente plus d'une formule globale. On remplace le modèle simple par une loi locale, différentielle ou énergétique.

Une forme générale possible est :
\[
\frac{\mathrm d X}{\mathrm dt}=\text{apport}-\text{pertes},
\]
ou, pour un phénomène propagatif :
\[
\frac{\partial^2 \psi}{\partial t^2}=c^2\Delta\psi.
\]
Selon le cas, on peut aussi utiliser un principe variationnel, par exemple le principe de Fermat en optique :
\[
\delta \int n\,\mathrm ds=0.
\]

Pour ce phénomène, le modèle Bac+3 consiste à partir de la relation centrale
\[
D=\frac12\rho v^2SC_D
\]
puis à introduire les corrections négligées : dépendance spatiale, dépendance temporelle, non-linéarités, pertes, conditions aux limites ou fluctuations. Cela permet de passer d'une réponse d'ordre de grandeur à une prédiction plus robuste.
\end{approfondissementbox}

\subsubsection*{Partie E — Limites, critique et prolongement}

\begin{enumerate}
  \item Citer au moins deux limites du modèle Terminale.
  \item Dire quelles grandeurs seraient difficiles à mesurer expérimentalement.
  \item Proposer une expérience simple permettant de tester le modèle.
  \item Expliquer ce qui changerait si l'on doublait une grandeur clé.
  \item Rédiger une conclusion courte, accessible à un élève de Terminale.
\end{enumerate}

\begin{corrigebox}[Synthèse rédigée]
Le phénomène « Traînée aérodynamique » s'explique par un petit nombre de lois physiques simples, mais sa manifestation réelle dépend souvent de nombreux paramètres. Le modèle Terminale donne l'intuition et les ordres de grandeur ; le modèle Bac+3 introduit les variations locales, les pertes, les équations différentielles ou les principes variationnels nécessaires pour décrire finement la réalité.
\end{corrigebox}

% ------------------------------------------------------------
\subsection{Problème corrigé 72 — Parachute}
% ------------------------------------------------------------

\begin{exercicebox}[Question-problème]
\textbf{Observation.} Un parachute ralentit une chute.

\medskip
\textbf{Question.} Calculer une vitesse terminale.

\medskip
\textbf{Thèmes mobilisés :} Vitesse terminale.

\medskip
\textbf{Relation ou idée centrale :}
\[
mg=\frac12\rho C_DSv_T^2
\]
\end{exercicebox}

\subsubsection*{Partie A — Modélisation niveau Terminale}

\begin{enumerate}
  \item Décrire précisément le phénomène observé.
  \item Identifier le système physique étudié.
  \item Lister les grandeurs importantes et leurs unités.
  \item Écrire la loi physique simple qui permet de commencer le raisonnement.
  \item Expliquer les hypothèses simplificatrices du modèle Terminale.
\end{enumerate}

\subsubsection*{Partie B — Mise en équation et calcul}

\begin{enumerate}
  \item Définir les variables utiles.
  \item Écrire la relation mathématique principale.
  \item Vérifier l'homogénéité de la formule.
  \item Choisir un ordre de grandeur réaliste pour les données manquantes.
  \item Effectuer le calcul ou établir le raisonnement qualitatif demandé.
  \item Interpréter le résultat en une phrase claire.
\end{enumerate}

\subsubsection*{Partie C — Réponse accessible niveau Terminale}

\begin{corrigebox}[Réponse Terminale]
On part de l'observation concrète : Un parachute ralentit une chute.

Au niveau Terminale, on cherche d'abord la grandeur qui contrôle le phénomène. Ici, l'idée centrale est donnée par :
\[
mg=\frac12\rho C_DSv_T^2
\]
Cette relation permet de relier la cause physique à l'effet observé. On vérifie les unités, puis on insère des ordres de grandeur réalistes.

La conclusion attendue est la suivante : Calculer une vitesse terminale. La réponse doit toujours préciser les hypothèses utilisées. Si le phénomène dépend d'un milieu réel, d'une géométrie complexe, de frottements, de pertes ou de gradients, le résultat obtenu est un ordre de grandeur et non une valeur exacte.
\end{corrigebox}

\subsubsection*{Partie D — Approfondissement niveau Bac+3}

\begin{approfondissementbox}[Réponse Bac+3]
Au niveau Bac+3, on ne se contente plus d'une formule globale. On remplace le modèle simple par une loi locale, différentielle ou énergétique.

Une forme générale possible est :
\[
\frac{\mathrm d X}{\mathrm dt}=\text{apport}-\text{pertes},
\]
ou, pour un phénomène propagatif :
\[
\frac{\partial^2 \psi}{\partial t^2}=c^2\Delta\psi.
\]
Selon le cas, on peut aussi utiliser un principe variationnel, par exemple le principe de Fermat en optique :
\[
\delta \int n\,\mathrm ds=0.
\]

Pour ce phénomène, le modèle Bac+3 consiste à partir de la relation centrale
\[
mg=\frac12\rho C_DSv_T^2
\]
puis à introduire les corrections négligées : dépendance spatiale, dépendance temporelle, non-linéarités, pertes, conditions aux limites ou fluctuations. Cela permet de passer d'une réponse d'ordre de grandeur à une prédiction plus robuste.
\end{approfondissementbox}

\subsubsection*{Partie E — Limites, critique et prolongement}

\begin{enumerate}
  \item Citer au moins deux limites du modèle Terminale.
  \item Dire quelles grandeurs seraient difficiles à mesurer expérimentalement.
  \item Proposer une expérience simple permettant de tester le modèle.
  \item Expliquer ce qui changerait si l'on doublait une grandeur clé.
  \item Rédiger une conclusion courte, accessible à un élève de Terminale.
\end{enumerate}

\begin{corrigebox}[Synthèse rédigée]
Le phénomène « Parachute » s'explique par un petit nombre de lois physiques simples, mais sa manifestation réelle dépend souvent de nombreux paramètres. Le modèle Terminale donne l'intuition et les ordres de grandeur ; le modèle Bac+3 introduit les variations locales, les pertes, les équations différentielles ou les principes variationnels nécessaires pour décrire finement la réalité.
\end{corrigebox}

% ------------------------------------------------------------
\subsection{Problème corrigé 73 — Fusée}
% ------------------------------------------------------------

\begin{exercicebox}[Question-problème]
\textbf{Observation.} Une fusée avance en éjectant des gaz.

\medskip
\textbf{Question.} Expliquer la propulsion dans le vide.

\medskip
\textbf{Thèmes mobilisés :} Quantité de mouvement.

\medskip
\textbf{Relation ou idée centrale :}
\[
F\simeq\dot m v_e
\]
\end{exercicebox}

\subsubsection*{Partie A — Modélisation niveau Terminale}

\begin{enumerate}
  \item Décrire précisément le phénomène observé.
  \item Identifier le système physique étudié.
  \item Lister les grandeurs importantes et leurs unités.
  \item Écrire la loi physique simple qui permet de commencer le raisonnement.
  \item Expliquer les hypothèses simplificatrices du modèle Terminale.
\end{enumerate}

\subsubsection*{Partie B — Mise en équation et calcul}

\begin{enumerate}
  \item Définir les variables utiles.
  \item Écrire la relation mathématique principale.
  \item Vérifier l'homogénéité de la formule.
  \item Choisir un ordre de grandeur réaliste pour les données manquantes.
  \item Effectuer le calcul ou établir le raisonnement qualitatif demandé.
  \item Interpréter le résultat en une phrase claire.
\end{enumerate}

\subsubsection*{Partie C — Réponse accessible niveau Terminale}

\begin{corrigebox}[Réponse Terminale]
On part de l'observation concrète : Une fusée avance en éjectant des gaz.

Au niveau Terminale, on cherche d'abord la grandeur qui contrôle le phénomène. Ici, l'idée centrale est donnée par :
\[
F\simeq\dot m v_e
\]
Cette relation permet de relier la cause physique à l'effet observé. On vérifie les unités, puis on insère des ordres de grandeur réalistes.

La conclusion attendue est la suivante : Expliquer la propulsion dans le vide. La réponse doit toujours préciser les hypothèses utilisées. Si le phénomène dépend d'un milieu réel, d'une géométrie complexe, de frottements, de pertes ou de gradients, le résultat obtenu est un ordre de grandeur et non une valeur exacte.
\end{corrigebox}

\subsubsection*{Partie D — Approfondissement niveau Bac+3}

\begin{approfondissementbox}[Réponse Bac+3]
Au niveau Bac+3, on ne se contente plus d'une formule globale. On remplace le modèle simple par une loi locale, différentielle ou énergétique.

Une forme générale possible est :
\[
\frac{\mathrm d X}{\mathrm dt}=\text{apport}-\text{pertes},
\]
ou, pour un phénomène propagatif :
\[
\frac{\partial^2 \psi}{\partial t^2}=c^2\Delta\psi.
\]
Selon le cas, on peut aussi utiliser un principe variationnel, par exemple le principe de Fermat en optique :
\[
\delta \int n\,\mathrm ds=0.
\]

Pour ce phénomène, le modèle Bac+3 consiste à partir de la relation centrale
\[
F\simeq\dot m v_e
\]
puis à introduire les corrections négligées : dépendance spatiale, dépendance temporelle, non-linéarités, pertes, conditions aux limites ou fluctuations. Cela permet de passer d'une réponse d'ordre de grandeur à une prédiction plus robuste.
\end{approfondissementbox}

\subsubsection*{Partie E — Limites, critique et prolongement}

\begin{enumerate}
  \item Citer au moins deux limites du modèle Terminale.
  \item Dire quelles grandeurs seraient difficiles à mesurer expérimentalement.
  \item Proposer une expérience simple permettant de tester le modèle.
  \item Expliquer ce qui changerait si l'on doublait une grandeur clé.
  \item Rédiger une conclusion courte, accessible à un élève de Terminale.
\end{enumerate}

\begin{corrigebox}[Synthèse rédigée]
Le phénomène « Fusée » s'explique par un petit nombre de lois physiques simples, mais sa manifestation réelle dépend souvent de nombreux paramètres. Le modèle Terminale donne l'intuition et les ordres de grandeur ; le modèle Bac+3 introduit les variations locales, les pertes, les équations différentielles ou les principes variationnels nécessaires pour décrire finement la réalité.
\end{corrigebox}

% ------------------------------------------------------------
\subsection{Problème corrigé 74 — Satellite}
% ------------------------------------------------------------

\begin{exercicebox}[Question-problème]
\textbf{Observation.} Un satellite tombe sans toucher la Terre.

\medskip
\textbf{Question.} Calculer une vitesse orbitale.

\medskip
\textbf{Thèmes mobilisés :} Gravitation et vitesse orbitale.

\medskip
\textbf{Relation ou idée centrale :}
\[
v=\sqrt{\frac{GM}{r}}
\]
\end{exercicebox}

\subsubsection*{Partie A — Modélisation niveau Terminale}

\begin{enumerate}
  \item Décrire précisément le phénomène observé.
  \item Identifier le système physique étudié.
  \item Lister les grandeurs importantes et leurs unités.
  \item Écrire la loi physique simple qui permet de commencer le raisonnement.
  \item Expliquer les hypothèses simplificatrices du modèle Terminale.
\end{enumerate}

\subsubsection*{Partie B — Mise en équation et calcul}

\begin{enumerate}
  \item Définir les variables utiles.
  \item Écrire la relation mathématique principale.
  \item Vérifier l'homogénéité de la formule.
  \item Choisir un ordre de grandeur réaliste pour les données manquantes.
  \item Effectuer le calcul ou établir le raisonnement qualitatif demandé.
  \item Interpréter le résultat en une phrase claire.
\end{enumerate}

\subsubsection*{Partie C — Réponse accessible niveau Terminale}

\begin{corrigebox}[Réponse Terminale]
On part de l'observation concrète : Un satellite tombe sans toucher la Terre.

Au niveau Terminale, on cherche d'abord la grandeur qui contrôle le phénomène. Ici, l'idée centrale est donnée par :
\[
v=\sqrt{\frac{GM}{r}}
\]
Cette relation permet de relier la cause physique à l'effet observé. On vérifie les unités, puis on insère des ordres de grandeur réalistes.

La conclusion attendue est la suivante : Calculer une vitesse orbitale. La réponse doit toujours préciser les hypothèses utilisées. Si le phénomène dépend d'un milieu réel, d'une géométrie complexe, de frottements, de pertes ou de gradients, le résultat obtenu est un ordre de grandeur et non une valeur exacte.
\end{corrigebox}

\subsubsection*{Partie D — Approfondissement niveau Bac+3}

\begin{approfondissementbox}[Réponse Bac+3]
Au niveau Bac+3, on ne se contente plus d'une formule globale. On remplace le modèle simple par une loi locale, différentielle ou énergétique.

Une forme générale possible est :
\[
\frac{\mathrm d X}{\mathrm dt}=\text{apport}-\text{pertes},
\]
ou, pour un phénomène propagatif :
\[
\frac{\partial^2 \psi}{\partial t^2}=c^2\Delta\psi.
\]
Selon le cas, on peut aussi utiliser un principe variationnel, par exemple le principe de Fermat en optique :
\[
\delta \int n\,\mathrm ds=0.
\]

Pour ce phénomène, le modèle Bac+3 consiste à partir de la relation centrale
\[
v=\sqrt{\frac{GM}{r}}
\]
puis à introduire les corrections négligées : dépendance spatiale, dépendance temporelle, non-linéarités, pertes, conditions aux limites ou fluctuations. Cela permet de passer d'une réponse d'ordre de grandeur à une prédiction plus robuste.
\end{approfondissementbox}

\subsubsection*{Partie E — Limites, critique et prolongement}

\begin{enumerate}
  \item Citer au moins deux limites du modèle Terminale.
  \item Dire quelles grandeurs seraient difficiles à mesurer expérimentalement.
  \item Proposer une expérience simple permettant de tester le modèle.
  \item Expliquer ce qui changerait si l'on doublait une grandeur clé.
  \item Rédiger une conclusion courte, accessible à un élève de Terminale.
\end{enumerate}

\begin{corrigebox}[Synthèse rédigée]
Le phénomène « Satellite » s'explique par un petit nombre de lois physiques simples, mais sa manifestation réelle dépend souvent de nombreux paramètres. Le modèle Terminale donne l'intuition et les ordres de grandeur ; le modèle Bac+3 introduit les variations locales, les pertes, les équations différentielles ou les principes variationnels nécessaires pour décrire finement la réalité.
\end{corrigebox}

% ------------------------------------------------------------
\subsection{Problème corrigé 75 — Marées}
% ------------------------------------------------------------

\begin{exercicebox}[Question-problème]
\textbf{Observation.} La mer monte et descend.

\medskip
\textbf{Question.} Expliquer le caractère différentiel.

\medskip
\textbf{Thèmes mobilisés :} Gradient gravitationnel lunaire.

\medskip
\textbf{Relation ou idée centrale :}
\[
F\propto\frac1{r^2}
\]
\end{exercicebox}

\subsubsection*{Partie A — Modélisation niveau Terminale}

\begin{enumerate}
  \item Décrire précisément le phénomène observé.
  \item Identifier le système physique étudié.
  \item Lister les grandeurs importantes et leurs unités.
  \item Écrire la loi physique simple qui permet de commencer le raisonnement.
  \item Expliquer les hypothèses simplificatrices du modèle Terminale.
\end{enumerate}

\subsubsection*{Partie B — Mise en équation et calcul}

\begin{enumerate}
  \item Définir les variables utiles.
  \item Écrire la relation mathématique principale.
  \item Vérifier l'homogénéité de la formule.
  \item Choisir un ordre de grandeur réaliste pour les données manquantes.
  \item Effectuer le calcul ou établir le raisonnement qualitatif demandé.
  \item Interpréter le résultat en une phrase claire.
\end{enumerate}

\subsubsection*{Partie C — Réponse accessible niveau Terminale}

\begin{corrigebox}[Réponse Terminale]
On part de l'observation concrète : La mer monte et descend.

Au niveau Terminale, on cherche d'abord la grandeur qui contrôle le phénomène. Ici, l'idée centrale est donnée par :
\[
F\propto\frac1{r^2}
\]
Cette relation permet de relier la cause physique à l'effet observé. On vérifie les unités, puis on insère des ordres de grandeur réalistes.

La conclusion attendue est la suivante : Expliquer le caractère différentiel. La réponse doit toujours préciser les hypothèses utilisées. Si le phénomène dépend d'un milieu réel, d'une géométrie complexe, de frottements, de pertes ou de gradients, le résultat obtenu est un ordre de grandeur et non une valeur exacte.
\end{corrigebox}

\subsubsection*{Partie D — Approfondissement niveau Bac+3}

\begin{approfondissementbox}[Réponse Bac+3]
Au niveau Bac+3, on ne se contente plus d'une formule globale. On remplace le modèle simple par une loi locale, différentielle ou énergétique.

Une forme générale possible est :
\[
\frac{\mathrm d X}{\mathrm dt}=\text{apport}-\text{pertes},
\]
ou, pour un phénomène propagatif :
\[
\frac{\partial^2 \psi}{\partial t^2}=c^2\Delta\psi.
\]
Selon le cas, on peut aussi utiliser un principe variationnel, par exemple le principe de Fermat en optique :
\[
\delta \int n\,\mathrm ds=0.
\]

Pour ce phénomène, le modèle Bac+3 consiste à partir de la relation centrale
\[
F\propto\frac1{r^2}
\]
puis à introduire les corrections négligées : dépendance spatiale, dépendance temporelle, non-linéarités, pertes, conditions aux limites ou fluctuations. Cela permet de passer d'une réponse d'ordre de grandeur à une prédiction plus robuste.
\end{approfondissementbox}

\subsubsection*{Partie E — Limites, critique et prolongement}

\begin{enumerate}
  \item Citer au moins deux limites du modèle Terminale.
  \item Dire quelles grandeurs seraient difficiles à mesurer expérimentalement.
  \item Proposer une expérience simple permettant de tester le modèle.
  \item Expliquer ce qui changerait si l'on doublait une grandeur clé.
  \item Rédiger une conclusion courte, accessible à un élève de Terminale.
\end{enumerate}

\begin{corrigebox}[Synthèse rédigée]
Le phénomène « Marées » s'explique par un petit nombre de lois physiques simples, mais sa manifestation réelle dépend souvent de nombreux paramètres. Le modèle Terminale donne l'intuition et les ordres de grandeur ; le modèle Bac+3 introduit les variations locales, les pertes, les équations différentielles ou les principes variationnels nécessaires pour décrire finement la réalité.
\end{corrigebox}

% ------------------------------------------------------------
\subsection{Problème corrigé 76 — Pendule}
% ------------------------------------------------------------

\begin{exercicebox}[Question-problème]
\textbf{Observation.} Un pendule oscille.

\medskip
\textbf{Question.} Calculer une période.

\medskip
\textbf{Thèmes mobilisés :} Oscillateur gravitationnel.

\medskip
\textbf{Relation ou idée centrale :}
\[
T\simeq2\pi\sqrt{\frac{\ell}{g}}
\]
\end{exercicebox}

\subsubsection*{Partie A — Modélisation niveau Terminale}

\begin{enumerate}
  \item Décrire précisément le phénomène observé.
  \item Identifier le système physique étudié.
  \item Lister les grandeurs importantes et leurs unités.
  \item Écrire la loi physique simple qui permet de commencer le raisonnement.
  \item Expliquer les hypothèses simplificatrices du modèle Terminale.
\end{enumerate}

\subsubsection*{Partie B — Mise en équation et calcul}

\begin{enumerate}
  \item Définir les variables utiles.
  \item Écrire la relation mathématique principale.
  \item Vérifier l'homogénéité de la formule.
  \item Choisir un ordre de grandeur réaliste pour les données manquantes.
  \item Effectuer le calcul ou établir le raisonnement qualitatif demandé.
  \item Interpréter le résultat en une phrase claire.
\end{enumerate}

\subsubsection*{Partie C — Réponse accessible niveau Terminale}

\begin{corrigebox}[Réponse Terminale]
On part de l'observation concrète : Un pendule oscille.

Au niveau Terminale, on cherche d'abord la grandeur qui contrôle le phénomène. Ici, l'idée centrale est donnée par :
\[
T\simeq2\pi\sqrt{\frac{\ell}{g}}
\]
Cette relation permet de relier la cause physique à l'effet observé. On vérifie les unités, puis on insère des ordres de grandeur réalistes.

La conclusion attendue est la suivante : Calculer une période. La réponse doit toujours préciser les hypothèses utilisées. Si le phénomène dépend d'un milieu réel, d'une géométrie complexe, de frottements, de pertes ou de gradients, le résultat obtenu est un ordre de grandeur et non une valeur exacte.
\end{corrigebox}

\subsubsection*{Partie D — Approfondissement niveau Bac+3}

\begin{approfondissementbox}[Réponse Bac+3]
Au niveau Bac+3, on ne se contente plus d'une formule globale. On remplace le modèle simple par une loi locale, différentielle ou énergétique.

Une forme générale possible est :
\[
\frac{\mathrm d X}{\mathrm dt}=\text{apport}-\text{pertes},
\]
ou, pour un phénomène propagatif :
\[
\frac{\partial^2 \psi}{\partial t^2}=c^2\Delta\psi.
\]
Selon le cas, on peut aussi utiliser un principe variationnel, par exemple le principe de Fermat en optique :
\[
\delta \int n\,\mathrm ds=0.
\]

Pour ce phénomène, le modèle Bac+3 consiste à partir de la relation centrale
\[
T\simeq2\pi\sqrt{\frac{\ell}{g}}
\]
puis à introduire les corrections négligées : dépendance spatiale, dépendance temporelle, non-linéarités, pertes, conditions aux limites ou fluctuations. Cela permet de passer d'une réponse d'ordre de grandeur à une prédiction plus robuste.
\end{approfondissementbox}

\subsubsection*{Partie E — Limites, critique et prolongement}

\begin{enumerate}
  \item Citer au moins deux limites du modèle Terminale.
  \item Dire quelles grandeurs seraient difficiles à mesurer expérimentalement.
  \item Proposer une expérience simple permettant de tester le modèle.
  \item Expliquer ce qui changerait si l'on doublait une grandeur clé.
  \item Rédiger une conclusion courte, accessible à un élève de Terminale.
\end{enumerate}

\begin{corrigebox}[Synthèse rédigée]
Le phénomène « Pendule » s'explique par un petit nombre de lois physiques simples, mais sa manifestation réelle dépend souvent de nombreux paramètres. Le modèle Terminale donne l'intuition et les ordres de grandeur ; le modèle Bac+3 introduit les variations locales, les pertes, les équations différentielles ou les principes variationnels nécessaires pour décrire finement la réalité.
\end{corrigebox}

% ------------------------------------------------------------
\subsection{Problème corrigé 77 — Gyroscope}
% ------------------------------------------------------------

\begin{exercicebox}[Question-problème]
\textbf{Observation.} Un axe en rotation reste stable.

\medskip
\textbf{Question.} Expliquer la stabilité gyroscopique.

\medskip
\textbf{Thèmes mobilisés :} Moment cinétique.

\medskip
\textbf{Relation ou idée centrale :}
\[
\vec L=I\vec\omega
\]
\end{exercicebox}

\subsubsection*{Partie A — Modélisation niveau Terminale}

\begin{enumerate}
  \item Décrire précisément le phénomène observé.
  \item Identifier le système physique étudié.
  \item Lister les grandeurs importantes et leurs unités.
  \item Écrire la loi physique simple qui permet de commencer le raisonnement.
  \item Expliquer les hypothèses simplificatrices du modèle Terminale.
\end{enumerate}

\subsubsection*{Partie B — Mise en équation et calcul}

\begin{enumerate}
  \item Définir les variables utiles.
  \item Écrire la relation mathématique principale.
  \item Vérifier l'homogénéité de la formule.
  \item Choisir un ordre de grandeur réaliste pour les données manquantes.
  \item Effectuer le calcul ou établir le raisonnement qualitatif demandé.
  \item Interpréter le résultat en une phrase claire.
\end{enumerate}

\subsubsection*{Partie C — Réponse accessible niveau Terminale}

\begin{corrigebox}[Réponse Terminale]
On part de l'observation concrète : Un axe en rotation reste stable.

Au niveau Terminale, on cherche d'abord la grandeur qui contrôle le phénomène. Ici, l'idée centrale est donnée par :
\[
\vec L=I\vec\omega
\]
Cette relation permet de relier la cause physique à l'effet observé. On vérifie les unités, puis on insère des ordres de grandeur réalistes.

La conclusion attendue est la suivante : Expliquer la stabilité gyroscopique. La réponse doit toujours préciser les hypothèses utilisées. Si le phénomène dépend d'un milieu réel, d'une géométrie complexe, de frottements, de pertes ou de gradients, le résultat obtenu est un ordre de grandeur et non une valeur exacte.
\end{corrigebox}

\subsubsection*{Partie D — Approfondissement niveau Bac+3}

\begin{approfondissementbox}[Réponse Bac+3]
Au niveau Bac+3, on ne se contente plus d'une formule globale. On remplace le modèle simple par une loi locale, différentielle ou énergétique.

Une forme générale possible est :
\[
\frac{\mathrm d X}{\mathrm dt}=\text{apport}-\text{pertes},
\]
ou, pour un phénomène propagatif :
\[
\frac{\partial^2 \psi}{\partial t^2}=c^2\Delta\psi.
\]
Selon le cas, on peut aussi utiliser un principe variationnel, par exemple le principe de Fermat en optique :
\[
\delta \int n\,\mathrm ds=0.
\]

Pour ce phénomène, le modèle Bac+3 consiste à partir de la relation centrale
\[
\vec L=I\vec\omega
\]
puis à introduire les corrections négligées : dépendance spatiale, dépendance temporelle, non-linéarités, pertes, conditions aux limites ou fluctuations. Cela permet de passer d'une réponse d'ordre de grandeur à une prédiction plus robuste.
\end{approfondissementbox}

\subsubsection*{Partie E — Limites, critique et prolongement}

\begin{enumerate}
  \item Citer au moins deux limites du modèle Terminale.
  \item Dire quelles grandeurs seraient difficiles à mesurer expérimentalement.
  \item Proposer une expérience simple permettant de tester le modèle.
  \item Expliquer ce qui changerait si l'on doublait une grandeur clé.
  \item Rédiger une conclusion courte, accessible à un élève de Terminale.
\end{enumerate}

\begin{corrigebox}[Synthèse rédigée]
Le phénomène « Gyroscope » s'explique par un petit nombre de lois physiques simples, mais sa manifestation réelle dépend souvent de nombreux paramètres. Le modèle Terminale donne l'intuition et les ordres de grandeur ; le modèle Bac+3 introduit les variations locales, les pertes, les équations différentielles ou les principes variationnels nécessaires pour décrire finement la réalité.
\end{corrigebox}

% ------------------------------------------------------------
\subsection{Problème corrigé 78 — Toupie}
% ------------------------------------------------------------

\begin{exercicebox}[Question-problème]
\textbf{Observation.} Une toupie tient debout en tournant.

\medskip
\textbf{Question.} Relier rotation et stabilité.

\medskip
\textbf{Thèmes mobilisés :} Précession.

\medskip
\textbf{Relation ou idée centrale :}
\[
\Omega\simeq\frac{mgr}{L}
\]
\end{exercicebox}

\subsubsection*{Partie A — Modélisation niveau Terminale}

\begin{enumerate}
  \item Décrire précisément le phénomène observé.
  \item Identifier le système physique étudié.
  \item Lister les grandeurs importantes et leurs unités.
  \item Écrire la loi physique simple qui permet de commencer le raisonnement.
  \item Expliquer les hypothèses simplificatrices du modèle Terminale.
\end{enumerate}

\subsubsection*{Partie B — Mise en équation et calcul}

\begin{enumerate}
  \item Définir les variables utiles.
  \item Écrire la relation mathématique principale.
  \item Vérifier l'homogénéité de la formule.
  \item Choisir un ordre de grandeur réaliste pour les données manquantes.
  \item Effectuer le calcul ou établir le raisonnement qualitatif demandé.
  \item Interpréter le résultat en une phrase claire.
\end{enumerate}

\subsubsection*{Partie C — Réponse accessible niveau Terminale}

\begin{corrigebox}[Réponse Terminale]
On part de l'observation concrète : Une toupie tient debout en tournant.

Au niveau Terminale, on cherche d'abord la grandeur qui contrôle le phénomène. Ici, l'idée centrale est donnée par :
\[
\Omega\simeq\frac{mgr}{L}
\]
Cette relation permet de relier la cause physique à l'effet observé. On vérifie les unités, puis on insère des ordres de grandeur réalistes.

La conclusion attendue est la suivante : Relier rotation et stabilité. La réponse doit toujours préciser les hypothèses utilisées. Si le phénomène dépend d'un milieu réel, d'une géométrie complexe, de frottements, de pertes ou de gradients, le résultat obtenu est un ordre de grandeur et non une valeur exacte.
\end{corrigebox}

\subsubsection*{Partie D — Approfondissement niveau Bac+3}

\begin{approfondissementbox}[Réponse Bac+3]
Au niveau Bac+3, on ne se contente plus d'une formule globale. On remplace le modèle simple par une loi locale, différentielle ou énergétique.

Une forme générale possible est :
\[
\frac{\mathrm d X}{\mathrm dt}=\text{apport}-\text{pertes},
\]
ou, pour un phénomène propagatif :
\[
\frac{\partial^2 \psi}{\partial t^2}=c^2\Delta\psi.
\]
Selon le cas, on peut aussi utiliser un principe variationnel, par exemple le principe de Fermat en optique :
\[
\delta \int n\,\mathrm ds=0.
\]

Pour ce phénomène, le modèle Bac+3 consiste à partir de la relation centrale
\[
\Omega\simeq\frac{mgr}{L}
\]
puis à introduire les corrections négligées : dépendance spatiale, dépendance temporelle, non-linéarités, pertes, conditions aux limites ou fluctuations. Cela permet de passer d'une réponse d'ordre de grandeur à une prédiction plus robuste.
\end{approfondissementbox}

\subsubsection*{Partie E — Limites, critique et prolongement}

\begin{enumerate}
  \item Citer au moins deux limites du modèle Terminale.
  \item Dire quelles grandeurs seraient difficiles à mesurer expérimentalement.
  \item Proposer une expérience simple permettant de tester le modèle.
  \item Expliquer ce qui changerait si l'on doublait une grandeur clé.
  \item Rédiger une conclusion courte, accessible à un élève de Terminale.
\end{enumerate}

\begin{corrigebox}[Synthèse rédigée]
Le phénomène « Toupie » s'explique par un petit nombre de lois physiques simples, mais sa manifestation réelle dépend souvent de nombreux paramètres. Le modèle Terminale donne l'intuition et les ordres de grandeur ; le modèle Bac+3 introduit les variations locales, les pertes, les équations différentielles ou les principes variationnels nécessaires pour décrire finement la réalité.
\end{corrigebox}

% ------------------------------------------------------------
\subsection{Problème corrigé 79 — Stabilité d'un vélo}
% ------------------------------------------------------------

\begin{exercicebox}[Question-problème]
\textbf{Observation.} Un vélo en mouvement est plus stable.

\medskip
\textbf{Question.} Expliquer le contrôle de l'équilibre.

\medskip
\textbf{Thèmes mobilisés :} Roulis, direction, gyroscopie.

\medskip
\textbf{Relation ou idée centrale :}
\[
\vec\tau=\frac{\mathrm d\vec L}{\mathrm dt}
\]
\end{exercicebox}

\subsubsection*{Partie A — Modélisation niveau Terminale}

\begin{enumerate}
  \item Décrire précisément le phénomène observé.
  \item Identifier le système physique étudié.
  \item Lister les grandeurs importantes et leurs unités.
  \item Écrire la loi physique simple qui permet de commencer le raisonnement.
  \item Expliquer les hypothèses simplificatrices du modèle Terminale.
\end{enumerate}

\subsubsection*{Partie B — Mise en équation et calcul}

\begin{enumerate}
  \item Définir les variables utiles.
  \item Écrire la relation mathématique principale.
  \item Vérifier l'homogénéité de la formule.
  \item Choisir un ordre de grandeur réaliste pour les données manquantes.
  \item Effectuer le calcul ou établir le raisonnement qualitatif demandé.
  \item Interpréter le résultat en une phrase claire.
\end{enumerate}

\subsubsection*{Partie C — Réponse accessible niveau Terminale}

\begin{corrigebox}[Réponse Terminale]
On part de l'observation concrète : Un vélo en mouvement est plus stable.

Au niveau Terminale, on cherche d'abord la grandeur qui contrôle le phénomène. Ici, l'idée centrale est donnée par :
\[
\vec\tau=\frac{\mathrm d\vec L}{\mathrm dt}
\]
Cette relation permet de relier la cause physique à l'effet observé. On vérifie les unités, puis on insère des ordres de grandeur réalistes.

La conclusion attendue est la suivante : Expliquer le contrôle de l'équilibre. La réponse doit toujours préciser les hypothèses utilisées. Si le phénomène dépend d'un milieu réel, d'une géométrie complexe, de frottements, de pertes ou de gradients, le résultat obtenu est un ordre de grandeur et non une valeur exacte.
\end{corrigebox}

\subsubsection*{Partie D — Approfondissement niveau Bac+3}

\begin{approfondissementbox}[Réponse Bac+3]
Au niveau Bac+3, on ne se contente plus d'une formule globale. On remplace le modèle simple par une loi locale, différentielle ou énergétique.

Une forme générale possible est :
\[
\frac{\mathrm d X}{\mathrm dt}=\text{apport}-\text{pertes},
\]
ou, pour un phénomène propagatif :
\[
\frac{\partial^2 \psi}{\partial t^2}=c^2\Delta\psi.
\]
Selon le cas, on peut aussi utiliser un principe variationnel, par exemple le principe de Fermat en optique :
\[
\delta \int n\,\mathrm ds=0.
\]

Pour ce phénomène, le modèle Bac+3 consiste à partir de la relation centrale
\[
\vec\tau=\frac{\mathrm d\vec L}{\mathrm dt}
\]
puis à introduire les corrections négligées : dépendance spatiale, dépendance temporelle, non-linéarités, pertes, conditions aux limites ou fluctuations. Cela permet de passer d'une réponse d'ordre de grandeur à une prédiction plus robuste.
\end{approfondissementbox}

\subsubsection*{Partie E — Limites, critique et prolongement}

\begin{enumerate}
  \item Citer au moins deux limites du modèle Terminale.
  \item Dire quelles grandeurs seraient difficiles à mesurer expérimentalement.
  \item Proposer une expérience simple permettant de tester le modèle.
  \item Expliquer ce qui changerait si l'on doublait une grandeur clé.
  \item Rédiger une conclusion courte, accessible à un élève de Terminale.
\end{enumerate}

\begin{corrigebox}[Synthèse rédigée]
Le phénomène « Stabilité d'un vélo » s'explique par un petit nombre de lois physiques simples, mais sa manifestation réelle dépend souvent de nombreux paramètres. Le modèle Terminale donne l'intuition et les ordres de grandeur ; le modèle Bac+3 introduit les variations locales, les pertes, les équations différentielles ou les principes variationnels nécessaires pour décrire finement la réalité.
\end{corrigebox}



\section{Partie XVI — Problèmes entièrement rédigés : une question, un phénomène, deux niveaux de correction}
% ============================================================

\begin{attentionbox}[But de cette nouvelle partie]
Dans cette partie, chaque phénomène est traité comme un vrai problème long. La question n'est pas seulement « appliquer une formule », mais répondre à un \textbf{pourquoi physique}. 

Chaque problème contient :
\begin{itemize}
  \item une question centrale ;
  \item un énoncé guidé ;
  \item une loi physique centrale ;
  \item une correction rédigée niveau Terminale ;
  \item une correction approfondie niveau Bac+3 ;
  \item un bilan pédagogique.
\end{itemize}
\end{attentionbox}

% ------------------------------------------------------------
\subsection{Problème détaillé 1 — Se frotter les mains peut-il réchauffer une pièce de 1 degré ?}
% ------------------------------------------------------------

\begin{exercicebox}[Question centrale]
\textbf{Phénomène étudié :} échauffement par frottement des mains.

\medskip
\textbf{Question à résoudre :} Pourquoi se frotter les mains chauffe localement très vite, alors qu'il faudrait énormément de temps pour réchauffer une pièce entière ?

\medskip
\textbf{Contexte.} Une personne dans une pièce de dimensions 4,0 m par 3,0 m par 2,5 m se frotte les mains fortement. On suppose que la puissance mécanique dissipée par frottement vaut 20 W et que toute cette puissance finit, dans un premier modèle très optimiste, par chauffer l'air de la pièce.

\medskip
\textbf{Données.} $\rho_{\rm air}=1{,}2\,\mathrm{kg.m^{-3}}$, $c_{\rm air}=1000\,\mathrm{J.kg^{-1}.K^{-1}}$, $\Delta T=1\,\mathrm K$, $P=20\,\mathrm W$.
\end{exercicebox}

\subsubsection*{Énoncé guidé}

\begin{enumerate}
  \item Calculer le volume de la pièce.
  \item Calculer la masse d'air contenue dans la pièce.
  \item Calculer l'énergie nécessaire pour augmenter de 1 K la température de l'air seul.
  \item En déduire le temps minimal si la puissance fournie vaut 20 W.
  \item Expliquer pourquoi ce temps est en réalité sous-estimé.
\end{enumerate}

\subsubsection*{Modélisation}

\begin{theorembox}[Loi physique centrale]
On utilise comme point de départ :
\[
Q=mc\Delta T,\qquad m=\rho V,\qquad t=\frac{Q}{P}
\]
Cette relation doit être lue physiquement : elle indique quelles grandeurs contrôlent le phénomène et comment une modification d'un paramètre change le résultat.
\end{theorembox}

\subsubsection*{Correction niveau Terminale}

\begin{corrigebox}[Réponse accessible Terminale]

Le volume de la pièce vaut
\[
V=4{,}0\times3{,}0\times2{,}5=30\,\mathrm{m^3}.
\]
La masse d'air est donc
\[
m=\rho V=1{,}2\times30=36\,\mathrm{kg}.
\]
Pour réchauffer cette masse d'air de \(1\,\mathrm K\), il faut
\[
Q=mc\Delta T=36\times1000\times1=36000\,\mathrm J.
\]
Si les mains dissipent \(20\,\mathrm W=20\,\mathrm{J.s^{-1}}\), le temps minimal est
\[
t=\frac{36000}{20}=1800\,\mathrm s=30\,\mathrm{min}.
\]
La réponse est donc : dans un modèle extrêmement favorable, il faudrait déjà environ trente minutes pour réchauffer uniquement l'air de la pièce de \(1^\circ\mathrm C\). En réalité, c'est beaucoup plus long car l'énergie chauffe aussi les mains, le corps, les murs, les meubles, le sol, et une partie s'échappe vers l'extérieur.

\end{corrigebox}

\subsubsection*{Correction approfondie niveau Bac+3}

\begin{approfondissementbox}[Réponse Bac+3]

Un modèle plus réaliste remplace l'air seul par une capacité thermique effective \(C_{\rm eff}\), comprenant l'air, les murs, le sol, les meubles et les échanges avec l'extérieur. On peut écrire
\[
C_{\rm eff}\frac{\mathrm dT}{\mathrm dt}=P-G(T-T_{\rm ext}),
\]
où \(G\) est la conductance thermique globale vers l'extérieur. La solution tend vers
\[
T_\infty=T_{\rm ext}+\frac{P}{G}.
\]
Ainsi, même en se frottant les mains indéfiniment, la température ne monte pas sans limite : elle se stabilise quand les pertes compensent l'apport de puissance. Comme \(P=20\,\mathrm W\) est faible devant la puissance d'un radiateur domestique, la hausse de température réelle serait très faible.

\end{approfondissementbox}

\subsubsection*{Bilan pédagogique}

\begin{methodebox}[Ce qu'il faut retenir]
\begin{itemize}
  \item La bonne démarche consiste à partir de l'observation, puis à identifier le système.
  \item La loi physique n'est pas une formule isolée : elle explique le sens du phénomène.
  \item Le calcul Terminale donne un ordre de grandeur.
  \item Le modèle Bac+3 explique pourquoi le modèle simple marche, et où il cesse de marcher.
  \item Une réponse complète doit toujours contenir une phrase de conclusion en français.
\end{itemize}
\end{methodebox}

% ------------------------------------------------------------
\subsection{Problème détaillé 2 — Pourquoi voit-on une flaque imaginaire sur une route chaude ?}
% ------------------------------------------------------------

\begin{exercicebox}[Question centrale]
\textbf{Phénomène étudié :} mirage inférieur.

\medskip
\textbf{Question à résoudre :} Pourquoi une route très chaude peut-elle donner l'impression qu'il y a de l'eau au loin ?

\medskip
\textbf{Contexte.} Sur une route en été, l'air au contact du bitume est beaucoup plus chaud que l'air situé un peu plus haut. Le conducteur voit parfois une zone brillante qui ressemble à une flaque d'eau.

\medskip
\textbf{Données.} Indice de l'air proche de 1 ; variation typique d'indice : de l'ordre de \(10^{-4}\).
\end{exercicebox}

\subsubsection*{Énoncé guidé}

\begin{enumerate}
  \item Expliquer pourquoi l'indice de l'air dépend de la température.
  \item Modéliser l'air par couches horizontales d'indices différents.
  \item Utiliser Snell-Descartes pour déterminer le sens de déviation des rayons.
  \item Expliquer pourquoi le rayon semble provenir du sol.
  \item Conclure sur la nature réelle ou virtuelle de l'image.
\end{enumerate}

\subsubsection*{Modélisation}

\begin{theorembox}[Loi physique centrale]
On utilise comme point de départ :
\[
n_1\sin i_1=n_2\sin i_2,\qquad n=n(z)
\]
Cette relation doit être lue physiquement : elle indique quelles grandeurs contrôlent le phénomène et comment une modification d'un paramètre change le résultat.
\end{theorembox}

\subsubsection*{Correction niveau Terminale}

\begin{corrigebox}[Réponse accessible Terminale]

L'air chaud est moins dense que l'air froid. Or l'indice optique de l'air dépend légèrement de sa densité : l'air chaud est donc un peu moins réfringent. On peut modéliser l'atmosphère proche du sol par plusieurs couches superposées, l'indice diminuant quand on se rapproche du bitume brûlant.

À chaque passage entre deux couches, la loi de Snell-Descartes s'applique :
\[
n_1\sin i_1=n_2\sin i_2.
\]
En descendant vers des couches de plus faible indice, le rayon s'éloigne progressivement de la normale. Comme l'indice varie continûment, le rayon n'est pas cassé brutalement mais courbé. Un rayon venant du ciel peut être courbé vers l'œil de l'observateur. Le cerveau prolonge ce rayon en ligne droite et croit qu'il vient du sol : on voit donc une image virtuelle du ciel sur la route, qui ressemble à une flaque d'eau.

\end{corrigebox}

\subsubsection*{Correction approfondie niveau Bac+3}

\begin{approfondissementbox}[Réponse Bac+3]

Dans un milieu stratifié où \(n=n(z)\), l'optique géométrique conduit à une équation de rayon. Une quantité analogue à \(n(z)\sin i(z)\) se conserve localement pour une stratification horizontale. La trajectoire suit le principe de Fermat :
\[
\delta\int n\,\mathrm ds=0.
\]
Le rayon se courbe vers les régions d'indice plus élevé. Ici, l'indice augmente avec l'altitude, donc le rayon peut se courber vers le haut après avoir frôlé les couches chaudes. Le mirage est donc un phénomène de géodésique optique dans un milieu inhomogène, pas une réflexion réelle sur le bitume.

\end{approfondissementbox}

\subsubsection*{Bilan pédagogique}

\begin{methodebox}[Ce qu'il faut retenir]
\begin{itemize}
  \item La bonne démarche consiste à partir de l'observation, puis à identifier le système.
  \item La loi physique n'est pas une formule isolée : elle explique le sens du phénomène.
  \item Le calcul Terminale donne un ordre de grandeur.
  \item Le modèle Bac+3 explique pourquoi le modèle simple marche, et où il cesse de marcher.
  \item Une réponse complète doit toujours contenir une phrase de conclusion en français.
\end{itemize}
\end{methodebox}

% ------------------------------------------------------------
\subsection{Problème détaillé 3 — Pourquoi l'arc-en-ciel apparaît-il autour de 42 degrés ?}
% ------------------------------------------------------------

\begin{exercicebox}[Question centrale]
\textbf{Phénomène étudié :} arc-en-ciel primaire.

\medskip
\textbf{Question à résoudre :} Pourquoi l'arc-en-ciel n'apparaît-il pas n'importe où dans le ciel, mais autour d'un angle bien déterminé ?

\medskip
\textbf{Contexte.} Après la pluie, des gouttes d'eau sphériques restent en suspension. La lumière blanche du Soleil entre dans les gouttes, s'y réfléchit et ressort vers l'observateur.

\medskip
\textbf{Données.} Indice de l'eau : \(n\simeq1{,}33\). Angle du rouge primaire : environ \(42^\circ\).
\end{exercicebox}

\subsubsection*{Énoncé guidé}

\begin{enumerate}
  \item Décrire le trajet d'un rayon dans une goutte pour l'arc primaire.
  \item Expliquer le rôle de la réfraction à l'entrée.
  \item Expliquer le rôle de la réflexion interne.
  \item Écrire l'expression de la déviation totale.
  \item Expliquer pourquoi un extremum de déviation concentre la lumière.
\end{enumerate}

\subsubsection*{Modélisation}

\begin{theorembox}[Loi physique centrale]
On utilise comme point de départ :
\[
D(i)=\pi+2i-4r,\qquad \sin i=n\sin r
\]
Cette relation doit être lue physiquement : elle indique quelles grandeurs contrôlent le phénomène et comment une modification d'un paramètre change le résultat.
\end{theorembox}

\subsubsection*{Correction niveau Terminale}

\begin{corrigebox}[Réponse accessible Terminale]

Pour l'arc primaire, un rayon solaire entre dans une goutte d'eau : il est réfracté. Il subit ensuite une réflexion interne sur la face opposée de la goutte, puis ressort en étant de nouveau réfracté. La loi de Snell-Descartes donne
\[
\sin i=n\sin r.
\]
La déviation totale d'un rayon primaire est
\[
D(i)=\pi+2i-4r.
\]
Tous les rayons ne ressortent pas avec la même déviation. Mais autour d'une certaine incidence, la déviation varie très peu quand l'incidence varie : on parle d'un extremum de déviation. Beaucoup de rayons sortent alors dans des directions proches, ce qui concentre la lumière. Pour le rouge, cela donne un angle d'observation d'environ \(42^\circ\). L'arc-en-ciel est donc visible sur un cône de direction autour de l'axe antisolaire.

\end{corrigebox}

\subsubsection*{Correction approfondie niveau Bac+3}

\begin{approfondissementbox}[Réponse Bac+3]

Le calcul fin consiste à dériver \(D(i)\) en tenant compte de la relation implicite
\[
\sin i=n\sin r.
\]
On cherche
\[
\frac{\mathrm dD}{\mathrm di}=0.
\]
Cette condition de stationnarité produit une caustique : la densité de rayons émergents devient grande près de la direction stationnaire. L'arc-en-ciel est donc une caustique optique. La dispersion \(n(\lambda)\) sépare les extremums selon la longueur d'onde, d'où la séparation des couleurs.

\end{approfondissementbox}

\subsubsection*{Bilan pédagogique}

\begin{methodebox}[Ce qu'il faut retenir]
\begin{itemize}
  \item La bonne démarche consiste à partir de l'observation, puis à identifier le système.
  \item La loi physique n'est pas une formule isolée : elle explique le sens du phénomène.
  \item Le calcul Terminale donne un ordre de grandeur.
  \item Le modèle Bac+3 explique pourquoi le modèle simple marche, et où il cesse de marcher.
  \item Une réponse complète doit toujours contenir une phrase de conclusion en français.
\end{itemize}
\end{methodebox}

% ------------------------------------------------------------
\subsection{Problème détaillé 4 — Pourquoi une fibre optique guide-t-elle la lumière ?}
% ------------------------------------------------------------

\begin{exercicebox}[Question centrale]
\textbf{Phénomène étudié :} fibre optique.

\medskip
\textbf{Question à résoudre :} Pourquoi un rayon lumineux peut-il rester piégé dans un fil de verre sur des kilomètres ?

\medskip
\textbf{Contexte.} Une fibre optique possède un cœur d'indice \(n_c\) entouré d'une gaine d'indice \(n_g<n_c\).

\medskip
\textbf{Données.} Exemple : \(n_c=1{,}48\), \(n_g=1{,}46\), \(L=10\,\mathrm{km}\).
\end{exercicebox}

\subsubsection*{Énoncé guidé}

\begin{enumerate}
  \item Expliquer la condition \(n_c>n_g\).
  \item Définir l'angle limite de réflexion totale.
  \item Calculer l'angle limite pour les valeurs données.
  \item Calculer le temps de propagation dans 10 km.
  \item Expliquer les limites du modèle rayon.
\end{enumerate}

\subsubsection*{Modélisation}

\begin{theorembox}[Loi physique centrale]
On utilise comme point de départ :
\[
\sin i_c=\frac{n_g}{n_c},\qquad t=\frac{nL}{c}
\]
Cette relation doit être lue physiquement : elle indique quelles grandeurs contrôlent le phénomène et comment une modification d'un paramètre change le résultat.
\end{theorembox}

\subsubsection*{Correction niveau Terminale}

\begin{corrigebox}[Réponse accessible Terminale]

Quand un rayon passe d'un milieu d'indice élevé vers un milieu d'indice plus faible, il peut exister un angle limite au-delà duquel il n'y a plus de rayon réfracté : toute la lumière est réfléchie. Dans une fibre, le cœur a un indice légèrement supérieur à celui de la gaine. Si l'incidence sur l'interface cœur-gaine est supérieure à l'angle limite \(i_c\), alors
\[
\sin i_c=\frac{n_g}{n_c}.
\]
Avec \(n_c=1{,}48\) et \(n_g=1{,}46\),
\[
\sin i_c=\frac{1{,}46}{1{,}48}\simeq0{,}986.
\]
L'angle limite est donc grand, proche de \(80^\circ\). Les rayons suffisamment rasants restent guidés par réflexions totales successives.

Le temps de propagation dans une fibre de longueur \(L\) vaut environ
\[
t=\frac{nL}{c}.
\]
Pour \(L=10\,\mathrm{km}\) et \(n\simeq1{,}5\),
\[
t\simeq\frac{1{,}5\times10^4}{3{,}0\times10^8}
=5{,}0\times10^{-5}\,\mathrm s.
\]
Soit environ \(50\,\mu\mathrm s\).

\end{corrigebox}

\subsubsection*{Correction approfondie niveau Bac+3}

\begin{approfondissementbox}[Réponse Bac+3]

Au niveau Bac+3, le modèle du rayon est remplacé par un modèle modal. Le champ électromagnétique dans la fibre vérifie une équation de Helmholtz avec conditions aux limites à l'interface cœur-gaine. Les modes guidés correspondent à des solutions confinées transversalement et propagatives longitudinalement. La fibre monomode impose un diamètre de cœur assez petit pour ne laisser qu'un seul mode se propager. La dispersion modale et la dispersion chromatique expliquent l'élargissement des impulsions et limitent le débit sur de longues distances.

\end{approfondissementbox}

\subsubsection*{Bilan pédagogique}

\begin{methodebox}[Ce qu'il faut retenir]
\begin{itemize}
  \item La bonne démarche consiste à partir de l'observation, puis à identifier le système.
  \item La loi physique n'est pas une formule isolée : elle explique le sens du phénomène.
  \item Le calcul Terminale donne un ordre de grandeur.
  \item Le modèle Bac+3 explique pourquoi le modèle simple marche, et où il cesse de marcher.
  \item Une réponse complète doit toujours contenir une phrase de conclusion en français.
\end{itemize}
\end{methodebox}

% ------------------------------------------------------------
\subsection{Problème détaillé 5 — Pourquoi le ciel est-il bleu ?}
% ------------------------------------------------------------

\begin{exercicebox}[Question centrale]
\textbf{Phénomène étudié :} diffusion Rayleigh.

\medskip
\textbf{Question à résoudre :} Pourquoi la lumière bleue est-elle plus visible dans le ciel que la lumière rouge ?

\medskip
\textbf{Contexte.} La lumière blanche du Soleil traverse l'atmosphère. Les molécules de l'air diffusent une partie de cette lumière dans toutes les directions.

\medskip
\textbf{Données.} \(\lambda_{\rm bleu}=450\,\mathrm{nm}\), \(\lambda_{\rm rouge}=650\,\mathrm{nm}\).
\end{exercicebox}

\subsubsection*{Énoncé guidé}

\begin{enumerate}
  \item Écrire la loi de diffusion Rayleigh.
  \item Comparer quantitativement bleu et rouge.
  \item Calculer le rapport des intensités diffusées.
  \item Expliquer pourquoi le ciel n'est pas violet.
  \item Relier cette explication au coucher du Soleil.
\end{enumerate}

\subsubsection*{Modélisation}

\begin{theorembox}[Loi physique centrale]
On utilise comme point de départ :
\[
I_{\rm diffusée}\propto \frac{1}{\lambda^4}
\]
Cette relation doit être lue physiquement : elle indique quelles grandeurs contrôlent le phénomène et comment une modification d'un paramètre change le résultat.
\end{theorembox}

\subsubsection*{Correction niveau Terminale}

\begin{corrigebox}[Réponse accessible Terminale]

La diffusion Rayleigh indique que l'intensité diffusée est d'autant plus grande que la longueur d'onde est petite :
\[
I\propto \frac{1}{\lambda^4}.
\]
Donc
\[
\frac{I_{\rm bleu}}{I_{\rm rouge}}
=
\left(\frac{\lambda_{\rm rouge}}{\lambda_{\rm bleu}}\right)^4
=
\left(\frac{650}{450}\right)^4
\simeq4{,}4.
\]
La lumière bleue est donc diffusée plusieurs fois plus efficacement que la lumière rouge. Comme cette lumière diffusée arrive à notre œil depuis toutes les directions du ciel, le ciel apparaît bleu. Il n'est pas violet car le Soleil émet moins dans le violet, l'œil humain y est moins sensible, et une partie du proche ultraviolet est absorbée par l'atmosphère.

\end{corrigebox}

\subsubsection*{Correction approfondie niveau Bac+3}

\begin{approfondissementbox}[Réponse Bac+3]

La diffusion Rayleigh s'explique par l'oscillation de dipôles induits dans les molécules d'air par le champ électrique de l'onde lumineuse. La puissance rayonnée par un dipôle accéléré varie fortement avec la fréquence, ce qui conduit à une section efficace
\[
\sigma\propto \frac{1}{\lambda^4}.
\]
Un traitement plus complet nécessite l'électromagnétisme classique, la polarisabilité moléculaire et l'équation de transfert radiatif dans l'atmosphère.

\end{approfondissementbox}

\subsubsection*{Bilan pédagogique}

\begin{methodebox}[Ce qu'il faut retenir]
\begin{itemize}
  \item La bonne démarche consiste à partir de l'observation, puis à identifier le système.
  \item La loi physique n'est pas une formule isolée : elle explique le sens du phénomène.
  \item Le calcul Terminale donne un ordre de grandeur.
  \item Le modèle Bac+3 explique pourquoi le modèle simple marche, et où il cesse de marcher.
  \item Une réponse complète doit toujours contenir une phrase de conclusion en français.
\end{itemize}
\end{methodebox}

% ------------------------------------------------------------
\subsection{Problème détaillé 6 — Pourquoi le Soleil devient-il rouge au coucher ?}
% ------------------------------------------------------------

\begin{exercicebox}[Question centrale]
\textbf{Phénomène étudié :} coucher de soleil rouge.

\medskip
\textbf{Question à résoudre :} Pourquoi la lumière directe du Soleil est-elle enrichie en rouge lorsqu'il est proche de l'horizon ?

\medskip
\textbf{Contexte.} Au coucher, les rayons solaires traversent une épaisseur d'atmosphère beaucoup plus grande qu'en plein midi.

\medskip
\textbf{Données.} Atténuation simplifiée : \(I=I_0e^{-\alpha(\lambda)L}\), avec \(\alpha\) plus grand pour le bleu.
\end{exercicebox}

\subsubsection*{Énoncé guidé}

\begin{enumerate}
  \item Comparer le trajet atmosphérique à midi et au coucher.
  \item Expliquer pourquoi le bleu est davantage retiré du faisceau direct.
  \item Interpréter la loi exponentielle d'atténuation.
  \item Expliquer le rôle des aérosols.
  \item Relier ce phénomène au ciel bleu.
\end{enumerate}

\subsubsection*{Modélisation}

\begin{theorembox}[Loi physique centrale]
On utilise comme point de départ :
\[
I=I_0e^{-\alpha(\lambda)L}
\]
Cette relation doit être lue physiquement : elle indique quelles grandeurs contrôlent le phénomène et comment une modification d'un paramètre change le résultat.
\end{theorembox}

\subsubsection*{Correction niveau Terminale}

\begin{corrigebox}[Réponse accessible Terminale]

Au coucher du Soleil, la lumière traverse l'atmosphère presque tangentiellement. Le trajet \(L\) dans l'air est donc beaucoup plus long. Or les courtes longueurs d'onde, comme le bleu, sont plus diffusées que les grandes longueurs d'onde, comme le rouge. La lumière bleue est donc en grande partie retirée du faisceau direct avant d'arriver à l'œil.

La loi
\[
I=I_0e^{-\alpha(\lambda)L}
\]
traduit cette atténuation. Comme \(\alpha\) est plus grand pour le bleu, \(I_{\rm bleu}\) diminue plus vite que \(I_{\rm rouge}\). La lumière transmise est donc enrichie en rouge.

\end{corrigebox}

\subsubsection*{Correction approfondie niveau Bac+3}

\begin{approfondissementbox}[Réponse Bac+3]

Un modèle avancé relève du transfert radiatif. L'intensité spectrale \(I_\lambda\) vérifie une équation intégrant absorption, diffusion hors du faisceau et diffusion multiple. La profondeur optique
\[
\tau_\lambda=\int \alpha_\lambda(s)\,\mathrm ds
\]
contrôle l'atténuation :
\[
I_\lambda=I_{\lambda,0}e^{-\tau_\lambda}.
\]
Au coucher, la profondeur optique augmente fortement, et sa dépendance en \(\lambda\) colore le Soleil.

\end{approfondissementbox}

\subsubsection*{Bilan pédagogique}

\begin{methodebox}[Ce qu'il faut retenir]
\begin{itemize}
  \item La bonne démarche consiste à partir de l'observation, puis à identifier le système.
  \item La loi physique n'est pas une formule isolée : elle explique le sens du phénomène.
  \item Le calcul Terminale donne un ordre de grandeur.
  \item Le modèle Bac+3 explique pourquoi le modèle simple marche, et où il cesse de marcher.
  \item Une réponse complète doit toujours contenir une phrase de conclusion en français.
\end{itemize}
\end{methodebox}

% ------------------------------------------------------------
\subsection{Problème détaillé 7 — Pourquoi une voiture consomme-t-elle beaucoup plus vite sur autoroute ?}
% ------------------------------------------------------------

\begin{exercicebox}[Question centrale]
\textbf{Phénomène étudié :} traînée aérodynamique.

\medskip
\textbf{Question à résoudre :} Pourquoi passer de 90 km/h à 130 km/h augmente-t-il fortement la puissance nécessaire ?

\medskip
\textbf{Contexte.} À grande vitesse, une voiture doit pousser l'air devant elle et créer un sillage turbulent.

\medskip
\textbf{Données.} Force de traînée : \(F_D=\frac12\rho C_DS v^2\).
\end{exercicebox}

\subsubsection*{Énoncé guidé}

\begin{enumerate}
  \item Montrer que la traînée varie comme \(v^2\).
  \item Montrer que la puissance varie comme \(v^3\).
  \item Comparer 90 km/h et 130 km/h.
  \item Interpréter le résultat en termes de consommation.
  \item Donner les limites du modèle.
\end{enumerate}

\subsubsection*{Modélisation}

\begin{theorembox}[Loi physique centrale]
On utilise comme point de départ :
\[
F_D=\frac12\rho C_DSv^2,\qquad P=F_Dv
\]
Cette relation doit être lue physiquement : elle indique quelles grandeurs contrôlent le phénomène et comment une modification d'un paramètre change le résultat.
\end{theorembox}

\subsubsection*{Correction niveau Terminale}

\begin{corrigebox}[Réponse accessible Terminale]

La force de traînée aérodynamique est
\[
F_D=\frac12\rho C_DSv^2.
\]
Elle varie donc comme le carré de la vitesse. Mais la puissance nécessaire pour vaincre cette force est
\[
P=F_Dv.
\]
Ainsi
\[
P\propto v^3.
\]
Le rapport entre les puissances à \(130\) et \(90\,\mathrm{km.h^{-1}}\) vaut environ
\[
\left(\frac{130}{90}\right)^3\simeq3.
\]
La puissance aérodynamique nécessaire est donc environ trois fois plus grande à 130 km/h qu'à 90 km/h. Cela explique pourquoi la consommation augmente fortement sur autoroute.

\end{corrigebox}

\subsubsection*{Correction approfondie niveau Bac+3}

\begin{approfondissementbox}[Réponse Bac+3]

La loi quadratique vient d'une analyse dimensionnelle du transfert de quantité de mouvement à grand nombre de Reynolds :
\[
F_D\sim \rho S v^2.
\]
Le coefficient \(C_D\) dépend de la géométrie, du régime de couche limite, du décollement et du sillage turbulent. La puissance \(P\sim v^3\) montre que les gains aérodynamiques sont cruciaux à haute vitesse. Un modèle complet couple traînée aérodynamique, résistance au roulement, rendement moteur et stratégie de conduite.

\end{approfondissementbox}

\subsubsection*{Bilan pédagogique}

\begin{methodebox}[Ce qu'il faut retenir]
\begin{itemize}
  \item La bonne démarche consiste à partir de l'observation, puis à identifier le système.
  \item La loi physique n'est pas une formule isolée : elle explique le sens du phénomène.
  \item Le calcul Terminale donne un ordre de grandeur.
  \item Le modèle Bac+3 explique pourquoi le modèle simple marche, et où il cesse de marcher.
  \item Une réponse complète doit toujours contenir une phrase de conclusion en français.
\end{itemize}
\end{methodebox}

% ------------------------------------------------------------
\subsection{Problème détaillé 8 — Pourquoi une balle liftée retombe-t-elle plus vite ?}
% ------------------------------------------------------------

\begin{exercicebox}[Question centrale]
\textbf{Phénomène étudié :} effet Magnus au tennis.

\medskip
\textbf{Question à résoudre :} Pourquoi un joueur peut-il frapper fort avec du lift sans sortir la balle du terrain ?

\medskip
\textbf{Contexte.} Une balle de tennis en topspin tourne rapidement sur elle-même pendant son vol.

\medskip
\textbf{Données.} Force de Magnus : \(\vec F_M\propto \vec\omega\times\vec v\).
\end{exercicebox}

\subsubsection*{Énoncé guidé}

\begin{enumerate}
  \item Identifier les forces sur la balle.
  \item Déterminer le sens qualitatif de la force de Magnus.
  \item Comparer une balle liftée et une balle sans rotation.
  \item Expliquer l'intérêt tactique du topspin.
  \item Écrire un modèle Bac+3 avec traînée et Magnus.
\end{enumerate}

\subsubsection*{Modélisation}

\begin{theorembox}[Loi physique centrale]
On utilise comme point de départ :
\[
\vec F_M\propto \vec\omega\times\vec v
\]
Cette relation doit être lue physiquement : elle indique quelles grandeurs contrôlent le phénomène et comment une modification d'un paramètre change le résultat.
\end{theorembox}

\subsubsection*{Correction niveau Terminale}

\begin{corrigebox}[Réponse accessible Terminale]

Une balle sans rotation est essentiellement soumise à son poids et aux frottements de l'air. Une balle liftée possède en plus une rotation. Cette rotation modifie l'écoulement de l'air autour de la balle et crée une force de Magnus :
\[
\vec F_M\propto \vec\omega\times\vec v.
\]
Pour un topspin, cette force est dirigée vers le bas. Elle courbe donc davantage la trajectoire. Le joueur peut frapper plus fort : la balle passe au-dessus du filet puis retombe plus vite dans le court.

\end{corrigebox}

\subsubsection*{Correction approfondie niveau Bac+3}

\begin{approfondissementbox}[Réponse Bac+3]

Un modèle dynamique complet s'écrit
\[
m\frac{\mathrm d\vec v}{\mathrm dt}
=
m\vec g
-\frac12\rho C_DSv\vec v
+
C_M\rho R^3\vec\omega\times\vec v.
\]
Cette équation différentielle vectorielle couple poids, traînée et force de Magnus. La rotation peut elle-même décroître par couple de frottement fluide. La trajectoire n'est donc pas une parabole mais une courbe obtenue numériquement.

\end{approfondissementbox}

\subsubsection*{Bilan pédagogique}

\begin{methodebox}[Ce qu'il faut retenir]
\begin{itemize}
  \item La bonne démarche consiste à partir de l'observation, puis à identifier le système.
  \item La loi physique n'est pas une formule isolée : elle explique le sens du phénomène.
  \item Le calcul Terminale donne un ordre de grandeur.
  \item Le modèle Bac+3 explique pourquoi le modèle simple marche, et où il cesse de marcher.
  \item Une réponse complète doit toujours contenir une phrase de conclusion en français.
\end{itemize}
\end{methodebox}

% ------------------------------------------------------------
\subsection{Problème détaillé 9 — Pourquoi un avion peut-il voler ?}
% ------------------------------------------------------------

\begin{exercicebox}[Question centrale]
\textbf{Phénomène étudié :} portance.

\medskip
\textbf{Question à résoudre :} Pourquoi une aile peut-elle produire une force verticale suffisante pour porter un avion ?

\medskip
\textbf{Contexte.} Une aile en mouvement dans l'air dévie l'écoulement et crée une différence de pression entre intrados et extrados.

\medskip
\textbf{Données.} Portance : \(L=\frac12\rho v^2SC_L\).
\end{exercicebox}

\subsubsection*{Énoncé guidé}

\begin{enumerate}
  \item Identifier les grandeurs de la formule de portance.
  \item Expliquer pourquoi la portance augmente avec la vitesse.
  \item Déterminer la condition de décollage.
  \item Expliquer le rôle de l'angle d'attaque.
  \item Discuter le décrochage.
\end{enumerate}

\subsubsection*{Modélisation}

\begin{theorembox}[Loi physique centrale]
On utilise comme point de départ :
\[
L=\frac12\rho v^2SC_L
\]
Cette relation doit être lue physiquement : elle indique quelles grandeurs contrôlent le phénomène et comment une modification d'un paramètre change le résultat.
\end{theorembox}

\subsubsection*{Correction niveau Terminale}

\begin{corrigebox}[Réponse accessible Terminale]

La portance est modélisée par
\[
L=\frac12\rho v^2SC_L.
\]
Elle augmente avec la masse volumique de l'air, la surface de l'aile, le coefficient de portance et surtout le carré de la vitesse. Pour que l'avion vole horizontalement, il faut que la portance compense le poids :
\[
L\simeq mg.
\]
Au décollage, l'avion doit donc atteindre une vitesse suffisante. L'angle d'attaque augmente \(C_L\), mais s'il devient trop grand, l'écoulement se décolle : c'est le décrochage.

\end{corrigebox}

\subsubsection*{Correction approfondie niveau Bac+3}

\begin{approfondissementbox}[Réponse Bac+3]

Au niveau avancé, la portance est liée à la circulation de l'écoulement autour de l'aile. Le théorème de Kutta-Joukowski donne, par unité de longueur,
\[
L'=\rho v\Gamma.
\]
La condition de Kutta fixe la circulation pour un profil donné. La portance réelle dépend aussi de la couche limite, de la turbulence, de la compressibilité et des tourbillons marginaux.

\end{approfondissementbox}

\subsubsection*{Bilan pédagogique}

\begin{methodebox}[Ce qu'il faut retenir]
\begin{itemize}
  \item La bonne démarche consiste à partir de l'observation, puis à identifier le système.
  \item La loi physique n'est pas une formule isolée : elle explique le sens du phénomène.
  \item Le calcul Terminale donne un ordre de grandeur.
  \item Le modèle Bac+3 explique pourquoi le modèle simple marche, et où il cesse de marcher.
  \item Une réponse complète doit toujours contenir une phrase de conclusion en français.
\end{itemize}
\end{methodebox}

% ------------------------------------------------------------
\subsection{Problème détaillé 10 — Pourquoi un satellite ne tombe-t-il pas sur Terre ?}
% ------------------------------------------------------------

\begin{exercicebox}[Question centrale]
\textbf{Phénomène étudié :} orbite circulaire.

\medskip
\textbf{Question à résoudre :} Pourquoi dit-on qu'un satellite est en chute libre permanente ?

\medskip
\textbf{Contexte.} Un satellite en orbite basse subit essentiellement la gravitation terrestre mais possède une vitesse horizontale très élevée.

\medskip
\textbf{Données.} Vitesse orbitale : \(v=\sqrt{GM/r}\).
\end{exercicebox}

\subsubsection*{Énoncé guidé}

\begin{enumerate}
  \item Identifier la force qui agit sur le satellite.
  \item Écrire la deuxième loi de Newton pour un mouvement circulaire.
  \item Retrouver la vitesse orbitale.
  \item Expliquer pourquoi le satellite tombe sans toucher le sol.
  \item Donner une limite du modèle.
\end{enumerate}

\subsubsection*{Modélisation}

\begin{theorembox}[Loi physique centrale]
On utilise comme point de départ :
\[
v=\sqrt{\frac{GM}{r}}
\]
Cette relation doit être lue physiquement : elle indique quelles grandeurs contrôlent le phénomène et comment une modification d'un paramètre change le résultat.
\end{theorembox}

\subsubsection*{Correction niveau Terminale}

\begin{corrigebox}[Réponse accessible Terminale]

La seule force principale sur un satellite est la gravitation :
\[
F=\frac{GMm}{r^2}.
\]
Pour une orbite circulaire, l'accélération centripète vaut
\[
a=\frac{v^2}{r}.
\]
La deuxième loi de Newton donne
\[
\frac{GMm}{r^2}=m\frac{v^2}{r}.
\]
On simplifie par \(m\) et on obtient
\[
v=\sqrt{\frac{GM}{r}}.
\]
Le satellite tombe bien vers la Terre, mais sa vitesse horizontale est telle que la surface terrestre se dérobe sous lui : il reste en orbite.

\end{corrigebox}

\subsubsection*{Correction approfondie niveau Bac+3}

\begin{approfondissementbox}[Réponse Bac+3]

Le mouvement orbital est un problème à force centrale dans un potentiel
\[
U(r)=-\frac{GMm}{r}.
\]
La conservation de l'énergie mécanique et du moment cinétique conduit aux trajectoires coniques de Kepler. Une orbite circulaire correspond à un minimum effectif du potentiel
\[
U_{\rm eff}(r)=\frac{L^2}{2mr^2}-\frac{GMm}{r}.
\]
Les perturbations réelles incluent frottement atmosphérique, aplatissement terrestre et influences lunaires/solaires.

\end{approfondissementbox}

\subsubsection*{Bilan pédagogique}

\begin{methodebox}[Ce qu'il faut retenir]
\begin{itemize}
  \item La bonne démarche consiste à partir de l'observation, puis à identifier le système.
  \item La loi physique n'est pas une formule isolée : elle explique le sens du phénomène.
  \item Le calcul Terminale donne un ordre de grandeur.
  \item Le modèle Bac+3 explique pourquoi le modèle simple marche, et où il cesse de marcher.
  \item Une réponse complète doit toujours contenir une phrase de conclusion en français.
\end{itemize}
\end{methodebox}

% ------------------------------------------------------------
\subsection{Problème détaillé 11 — Pourquoi le métal paraît-il plus froid que le bois ?}
% ------------------------------------------------------------

\begin{exercicebox}[Question centrale]
\textbf{Phénomène étudié :} sensation thermique.

\medskip
\textbf{Question à résoudre :} Pourquoi le métal paraît-il plus froid que le bois ?

\medskip
\textbf{Contexte.} On observe le phénomène suivant : sensation thermique. On cherche à comprendre non seulement ce que l'on voit, mais pourquoi la loi physique impose ce comportement.

\medskip
\textbf{Données.} On utilisera des ordres de grandeur réalistes et on précisera les hypothèses.
\end{exercicebox}

\subsubsection*{Énoncé guidé}

\begin{enumerate}
  \item Décrire l'observation.
  \item Identifier le système et les grandeurs.
  \item Écrire la loi physique principale.
  \item Faire un raisonnement d'ordre de grandeur.
  \item Discuter les limites du modèle.
\end{enumerate}

\subsubsection*{Modélisation}

\begin{theorembox}[Loi physique centrale]
On utilise comme point de départ :
\[
\Phi\propto \lambda_{\rm th}\Delta T
\]
Cette relation doit être lue physiquement : elle indique quelles grandeurs contrôlent le phénomène et comment une modification d'un paramètre change le résultat.
\end{theorembox}

\subsubsection*{Correction niveau Terminale}

\begin{corrigebox}[Réponse accessible Terminale]

Au niveau Terminale, on explique ce phénomène en isolant la grandeur dominante. La relation utile est
\[
\Phi\propto \lambda_{\rm th}\Delta T
\]
Elle permet de comprendre le sens de variation : si la grandeur qui apparaît au numérateur augmente, l'effet augmente ; si elle apparaît au dénominateur, l'effet diminue.

La réponse qualitative est donc que le phénomène « sensation thermique » n'est pas magique : il résulte d'un transfert d'énergie, d'une propagation d'onde, d'une force ou d'une contrainte géométrique. Le modèle Terminale donne l'explication principale et permet de calculer un ordre de grandeur.

\end{corrigebox}

\subsubsection*{Correction approfondie niveau Bac+3}

\begin{approfondissementbox}[Réponse Bac+3]

Au niveau Bac+3, on raffine le modèle en introduisant une loi locale ou différentielle. La relation
\[
\Phi\propto \lambda_{\rm th}\Delta T
\]
devient alors une conséquence d'un bilan plus général : bilan d'énergie, équation d'onde, équation de diffusion, dynamique du solide ou conservation de la quantité de mouvement. Les corrections importantes sont les pertes, les conditions aux limites, la géométrie réelle et les non-linéarités.

\end{approfondissementbox}

\subsubsection*{Bilan pédagogique}

\begin{methodebox}[Ce qu'il faut retenir]
\begin{itemize}
  \item La bonne démarche consiste à partir de l'observation, puis à identifier le système.
  \item La loi physique n'est pas une formule isolée : elle explique le sens du phénomène.
  \item Le calcul Terminale donne un ordre de grandeur.
  \item Le modèle Bac+3 explique pourquoi le modèle simple marche, et où il cesse de marcher.
  \item Une réponse complète doit toujours contenir une phrase de conclusion en français.
\end{itemize}
\end{methodebox}

% ------------------------------------------------------------
\subsection{Problème détaillé 12 — Pourquoi la transpiration refroidit-elle le corps ?}
% ------------------------------------------------------------

\begin{exercicebox}[Question centrale]
\textbf{Phénomène étudié :} transpiration.

\medskip
\textbf{Question à résoudre :} Pourquoi la transpiration refroidit-elle le corps ?

\medskip
\textbf{Contexte.} On observe le phénomène suivant : transpiration. On cherche à comprendre non seulement ce que l'on voit, mais pourquoi la loi physique impose ce comportement.

\medskip
\textbf{Données.} On utilisera des ordres de grandeur réalistes et on précisera les hypothèses.
\end{exercicebox}

\subsubsection*{Énoncé guidé}

\begin{enumerate}
  \item Décrire l'observation.
  \item Identifier le système et les grandeurs.
  \item Écrire la loi physique principale.
  \item Faire un raisonnement d'ordre de grandeur.
  \item Discuter les limites du modèle.
\end{enumerate}

\subsubsection*{Modélisation}

\begin{theorembox}[Loi physique centrale]
On utilise comme point de départ :
\[
P=\dot m L_v
\]
Cette relation doit être lue physiquement : elle indique quelles grandeurs contrôlent le phénomène et comment une modification d'un paramètre change le résultat.
\end{theorembox}

\subsubsection*{Correction niveau Terminale}

\begin{corrigebox}[Réponse accessible Terminale]

Au niveau Terminale, on explique ce phénomène en isolant la grandeur dominante. La relation utile est
\[
P=\dot m L_v
\]
Elle permet de comprendre le sens de variation : si la grandeur qui apparaît au numérateur augmente, l'effet augmente ; si elle apparaît au dénominateur, l'effet diminue.

La réponse qualitative est donc que le phénomène « transpiration » n'est pas magique : il résulte d'un transfert d'énergie, d'une propagation d'onde, d'une force ou d'une contrainte géométrique. Le modèle Terminale donne l'explication principale et permet de calculer un ordre de grandeur.

\end{corrigebox}

\subsubsection*{Correction approfondie niveau Bac+3}

\begin{approfondissementbox}[Réponse Bac+3]

Au niveau Bac+3, on raffine le modèle en introduisant une loi locale ou différentielle. La relation
\[
P=\dot m L_v
\]
devient alors une conséquence d'un bilan plus général : bilan d'énergie, équation d'onde, équation de diffusion, dynamique du solide ou conservation de la quantité de mouvement. Les corrections importantes sont les pertes, les conditions aux limites, la géométrie réelle et les non-linéarités.

\end{approfondissementbox}

\subsubsection*{Bilan pédagogique}

\begin{methodebox}[Ce qu'il faut retenir]
\begin{itemize}
  \item La bonne démarche consiste à partir de l'observation, puis à identifier le système.
  \item La loi physique n'est pas une formule isolée : elle explique le sens du phénomène.
  \item Le calcul Terminale donne un ordre de grandeur.
  \item Le modèle Bac+3 explique pourquoi le modèle simple marche, et où il cesse de marcher.
  \item Une réponse complète doit toujours contenir une phrase de conclusion en français.
\end{itemize}
\end{methodebox}

% ------------------------------------------------------------
\subsection{Problème détaillé 13 — Pourquoi un glaçon met-il du temps à fondre ?}
% ------------------------------------------------------------

\begin{exercicebox}[Question centrale]
\textbf{Phénomène étudié :} fusion.

\medskip
\textbf{Question à résoudre :} Pourquoi un glaçon met-il du temps à fondre ?

\medskip
\textbf{Contexte.} On observe le phénomène suivant : fusion. On cherche à comprendre non seulement ce que l'on voit, mais pourquoi la loi physique impose ce comportement.

\medskip
\textbf{Données.} On utilisera des ordres de grandeur réalistes et on précisera les hypothèses.
\end{exercicebox}

\subsubsection*{Énoncé guidé}

\begin{enumerate}
  \item Décrire l'observation.
  \item Identifier le système et les grandeurs.
  \item Écrire la loi physique principale.
  \item Faire un raisonnement d'ordre de grandeur.
  \item Discuter les limites du modèle.
\end{enumerate}

\subsubsection*{Modélisation}

\begin{theorembox}[Loi physique centrale]
On utilise comme point de départ :
\[
Q=mL_f
\]
Cette relation doit être lue physiquement : elle indique quelles grandeurs contrôlent le phénomène et comment une modification d'un paramètre change le résultat.
\end{theorembox}

\subsubsection*{Correction niveau Terminale}

\begin{corrigebox}[Réponse accessible Terminale]

Au niveau Terminale, on explique ce phénomène en isolant la grandeur dominante. La relation utile est
\[
Q=mL_f
\]
Elle permet de comprendre le sens de variation : si la grandeur qui apparaît au numérateur augmente, l'effet augmente ; si elle apparaît au dénominateur, l'effet diminue.

La réponse qualitative est donc que le phénomène « fusion » n'est pas magique : il résulte d'un transfert d'énergie, d'une propagation d'onde, d'une force ou d'une contrainte géométrique. Le modèle Terminale donne l'explication principale et permet de calculer un ordre de grandeur.

\end{corrigebox}

\subsubsection*{Correction approfondie niveau Bac+3}

\begin{approfondissementbox}[Réponse Bac+3]

Au niveau Bac+3, on raffine le modèle en introduisant une loi locale ou différentielle. La relation
\[
Q=mL_f
\]
devient alors une conséquence d'un bilan plus général : bilan d'énergie, équation d'onde, équation de diffusion, dynamique du solide ou conservation de la quantité de mouvement. Les corrections importantes sont les pertes, les conditions aux limites, la géométrie réelle et les non-linéarités.

\end{approfondissementbox}

\subsubsection*{Bilan pédagogique}

\begin{methodebox}[Ce qu'il faut retenir]
\begin{itemize}
  \item La bonne démarche consiste à partir de l'observation, puis à identifier le système.
  \item La loi physique n'est pas une formule isolée : elle explique le sens du phénomène.
  \item Le calcul Terminale donne un ordre de grandeur.
  \item Le modèle Bac+3 explique pourquoi le modèle simple marche, et où il cesse de marcher.
  \item Une réponse complète doit toujours contenir une phrase de conclusion en français.
\end{itemize}
\end{methodebox}

% ------------------------------------------------------------
\subsection{Problème détaillé 14 — Pourquoi une cocotte-minute cuit-elle plus vite ?}
% ------------------------------------------------------------

\begin{exercicebox}[Question centrale]
\textbf{Phénomène étudié :} cocotte-minute.

\medskip
\textbf{Question à résoudre :} Pourquoi une cocotte-minute cuit-elle plus vite ?

\medskip
\textbf{Contexte.} On observe le phénomène suivant : cocotte-minute. On cherche à comprendre non seulement ce que l'on voit, mais pourquoi la loi physique impose ce comportement.

\medskip
\textbf{Données.} On utilisera des ordres de grandeur réalistes et on précisera les hypothèses.
\end{exercicebox}

\subsubsection*{Énoncé guidé}

\begin{enumerate}
  \item Décrire l'observation.
  \item Identifier le système et les grandeurs.
  \item Écrire la loi physique principale.
  \item Faire un raisonnement d'ordre de grandeur.
  \item Discuter les limites du modèle.
\end{enumerate}

\subsubsection*{Modélisation}

\begin{theorembox}[Loi physique centrale]
On utilise comme point de départ :
\[
P_{\rm sat}(T)=P_{\rm ext}
\]
Cette relation doit être lue physiquement : elle indique quelles grandeurs contrôlent le phénomène et comment une modification d'un paramètre change le résultat.
\end{theorembox}

\subsubsection*{Correction niveau Terminale}

\begin{corrigebox}[Réponse accessible Terminale]

Au niveau Terminale, on explique ce phénomène en isolant la grandeur dominante. La relation utile est
\[
P_{\rm sat}(T)=P_{\rm ext}
\]
Elle permet de comprendre le sens de variation : si la grandeur qui apparaît au numérateur augmente, l'effet augmente ; si elle apparaît au dénominateur, l'effet diminue.

La réponse qualitative est donc que le phénomène « cocotte-minute » n'est pas magique : il résulte d'un transfert d'énergie, d'une propagation d'onde, d'une force ou d'une contrainte géométrique. Le modèle Terminale donne l'explication principale et permet de calculer un ordre de grandeur.

\end{corrigebox}

\subsubsection*{Correction approfondie niveau Bac+3}

\begin{approfondissementbox}[Réponse Bac+3]

Au niveau Bac+3, on raffine le modèle en introduisant une loi locale ou différentielle. La relation
\[
P_{\rm sat}(T)=P_{\rm ext}
\]
devient alors une conséquence d'un bilan plus général : bilan d'énergie, équation d'onde, équation de diffusion, dynamique du solide ou conservation de la quantité de mouvement. Les corrections importantes sont les pertes, les conditions aux limites, la géométrie réelle et les non-linéarités.

\end{approfondissementbox}

\subsubsection*{Bilan pédagogique}

\begin{methodebox}[Ce qu'il faut retenir]
\begin{itemize}
  \item La bonne démarche consiste à partir de l'observation, puis à identifier le système.
  \item La loi physique n'est pas une formule isolée : elle explique le sens du phénomène.
  \item Le calcul Terminale donne un ordre de grandeur.
  \item Le modèle Bac+3 explique pourquoi le modèle simple marche, et où il cesse de marcher.
  \item Une réponse complète doit toujours contenir une phrase de conclusion en français.
\end{itemize}
\end{methodebox}

% ------------------------------------------------------------
\subsection{Problème détaillé 15 — Pourquoi les bulles de savon sont-elles colorées ?}
% ------------------------------------------------------------

\begin{exercicebox}[Question centrale]
\textbf{Phénomène étudié :} bulles de savon.

\medskip
\textbf{Question à résoudre :} Pourquoi les bulles de savon sont-elles colorées ?

\medskip
\textbf{Contexte.} On observe le phénomène suivant : bulles de savon. On cherche à comprendre non seulement ce que l'on voit, mais pourquoi la loi physique impose ce comportement.

\medskip
\textbf{Données.} On utilisera des ordres de grandeur réalistes et on précisera les hypothèses.
\end{exercicebox}

\subsubsection*{Énoncé guidé}

\begin{enumerate}
  \item Décrire l'observation.
  \item Identifier le système et les grandeurs.
  \item Écrire la loi physique principale.
  \item Faire un raisonnement d'ordre de grandeur.
  \item Discuter les limites du modèle.
\end{enumerate}

\subsubsection*{Modélisation}

\begin{theorembox}[Loi physique centrale]
On utilise comme point de départ :
\[
\delta=2ne\cos r
\]
Cette relation doit être lue physiquement : elle indique quelles grandeurs contrôlent le phénomène et comment une modification d'un paramètre change le résultat.
\end{theorembox}

\subsubsection*{Correction niveau Terminale}

\begin{corrigebox}[Réponse accessible Terminale]

Au niveau Terminale, on explique ce phénomène en isolant la grandeur dominante. La relation utile est
\[
\delta=2ne\cos r
\]
Elle permet de comprendre le sens de variation : si la grandeur qui apparaît au numérateur augmente, l'effet augmente ; si elle apparaît au dénominateur, l'effet diminue.

La réponse qualitative est donc que le phénomène « bulles de savon » n'est pas magique : il résulte d'un transfert d'énergie, d'une propagation d'onde, d'une force ou d'une contrainte géométrique. Le modèle Terminale donne l'explication principale et permet de calculer un ordre de grandeur.

\end{corrigebox}

\subsubsection*{Correction approfondie niveau Bac+3}

\begin{approfondissementbox}[Réponse Bac+3]

Au niveau Bac+3, on raffine le modèle en introduisant une loi locale ou différentielle. La relation
\[
\delta=2ne\cos r
\]
devient alors une conséquence d'un bilan plus général : bilan d'énergie, équation d'onde, équation de diffusion, dynamique du solide ou conservation de la quantité de mouvement. Les corrections importantes sont les pertes, les conditions aux limites, la géométrie réelle et les non-linéarités.

\end{approfondissementbox}

\subsubsection*{Bilan pédagogique}

\begin{methodebox}[Ce qu'il faut retenir]
\begin{itemize}
  \item La bonne démarche consiste à partir de l'observation, puis à identifier le système.
  \item La loi physique n'est pas une formule isolée : elle explique le sens du phénomène.
  \item Le calcul Terminale donne un ordre de grandeur.
  \item Le modèle Bac+3 explique pourquoi le modèle simple marche, et où il cesse de marcher.
  \item Une réponse complète doit toujours contenir une phrase de conclusion en français.
\end{itemize}
\end{methodebox}

% ------------------------------------------------------------
\subsection{Problème détaillé 16 — Pourquoi les lunettes polarisantes suppriment-elles des reflets ?}
% ------------------------------------------------------------

\begin{exercicebox}[Question centrale]
\textbf{Phénomène étudié :} polarisation.

\medskip
\textbf{Question à résoudre :} Pourquoi les lunettes polarisantes suppriment-elles des reflets ?

\medskip
\textbf{Contexte.} On observe le phénomène suivant : polarisation. On cherche à comprendre non seulement ce que l'on voit, mais pourquoi la loi physique impose ce comportement.

\medskip
\textbf{Données.} On utilisera des ordres de grandeur réalistes et on précisera les hypothèses.
\end{exercicebox}

\subsubsection*{Énoncé guidé}

\begin{enumerate}
  \item Décrire l'observation.
  \item Identifier le système et les grandeurs.
  \item Écrire la loi physique principale.
  \item Faire un raisonnement d'ordre de grandeur.
  \item Discuter les limites du modèle.
\end{enumerate}

\subsubsection*{Modélisation}

\begin{theorembox}[Loi physique centrale]
On utilise comme point de départ :
\[
I=I_0\cos^2\theta
\]
Cette relation doit être lue physiquement : elle indique quelles grandeurs contrôlent le phénomène et comment une modification d'un paramètre change le résultat.
\end{theorembox}

\subsubsection*{Correction niveau Terminale}

\begin{corrigebox}[Réponse accessible Terminale]

Au niveau Terminale, on explique ce phénomène en isolant la grandeur dominante. La relation utile est
\[
I=I_0\cos^2\theta
\]
Elle permet de comprendre le sens de variation : si la grandeur qui apparaît au numérateur augmente, l'effet augmente ; si elle apparaît au dénominateur, l'effet diminue.

La réponse qualitative est donc que le phénomène « polarisation » n'est pas magique : il résulte d'un transfert d'énergie, d'une propagation d'onde, d'une force ou d'une contrainte géométrique. Le modèle Terminale donne l'explication principale et permet de calculer un ordre de grandeur.

\end{corrigebox}

\subsubsection*{Correction approfondie niveau Bac+3}

\begin{approfondissementbox}[Réponse Bac+3]

Au niveau Bac+3, on raffine le modèle en introduisant une loi locale ou différentielle. La relation
\[
I=I_0\cos^2\theta
\]
devient alors une conséquence d'un bilan plus général : bilan d'énergie, équation d'onde, équation de diffusion, dynamique du solide ou conservation de la quantité de mouvement. Les corrections importantes sont les pertes, les conditions aux limites, la géométrie réelle et les non-linéarités.

\end{approfondissementbox}

\subsubsection*{Bilan pédagogique}

\begin{methodebox}[Ce qu'il faut retenir]
\begin{itemize}
  \item La bonne démarche consiste à partir de l'observation, puis à identifier le système.
  \item La loi physique n'est pas une formule isolée : elle explique le sens du phénomène.
  \item Le calcul Terminale donne un ordre de grandeur.
  \item Le modèle Bac+3 explique pourquoi le modèle simple marche, et où il cesse de marcher.
  \item Une réponse complète doit toujours contenir une phrase de conclusion en français.
\end{itemize}
\end{methodebox}

% ------------------------------------------------------------
\subsection{Problème détaillé 17 — Pourquoi une flûte change-t-elle de note quand on bouche un trou ?}
% ------------------------------------------------------------

\begin{exercicebox}[Question centrale]
\textbf{Phénomène étudié :} instrument à vent.

\medskip
\textbf{Question à résoudre :} Pourquoi une flûte change-t-elle de note quand on bouche un trou ?

\medskip
\textbf{Contexte.} On observe le phénomène suivant : instrument à vent. On cherche à comprendre non seulement ce que l'on voit, mais pourquoi la loi physique impose ce comportement.

\medskip
\textbf{Données.} On utilisera des ordres de grandeur réalistes et on précisera les hypothèses.
\end{exercicebox}

\subsubsection*{Énoncé guidé}

\begin{enumerate}
  \item Décrire l'observation.
  \item Identifier le système et les grandeurs.
  \item Écrire la loi physique principale.
  \item Faire un raisonnement d'ordre de grandeur.
  \item Discuter les limites du modèle.
\end{enumerate}

\subsubsection*{Modélisation}

\begin{theorembox}[Loi physique centrale]
On utilise comme point de départ :
\[
f_n=n\frac{c}{2L}
\]
Cette relation doit être lue physiquement : elle indique quelles grandeurs contrôlent le phénomène et comment une modification d'un paramètre change le résultat.
\end{theorembox}

\subsubsection*{Correction niveau Terminale}

\begin{corrigebox}[Réponse accessible Terminale]

Au niveau Terminale, on explique ce phénomène en isolant la grandeur dominante. La relation utile est
\[
f_n=n\frac{c}{2L}
\]
Elle permet de comprendre le sens de variation : si la grandeur qui apparaît au numérateur augmente, l'effet augmente ; si elle apparaît au dénominateur, l'effet diminue.

La réponse qualitative est donc que le phénomène « instrument à vent » n'est pas magique : il résulte d'un transfert d'énergie, d'une propagation d'onde, d'une force ou d'une contrainte géométrique. Le modèle Terminale donne l'explication principale et permet de calculer un ordre de grandeur.

\end{corrigebox}

\subsubsection*{Correction approfondie niveau Bac+3}

\begin{approfondissementbox}[Réponse Bac+3]

Au niveau Bac+3, on raffine le modèle en introduisant une loi locale ou différentielle. La relation
\[
f_n=n\frac{c}{2L}
\]
devient alors une conséquence d'un bilan plus général : bilan d'énergie, équation d'onde, équation de diffusion, dynamique du solide ou conservation de la quantité de mouvement. Les corrections importantes sont les pertes, les conditions aux limites, la géométrie réelle et les non-linéarités.

\end{approfondissementbox}

\subsubsection*{Bilan pédagogique}

\begin{methodebox}[Ce qu'il faut retenir]
\begin{itemize}
  \item La bonne démarche consiste à partir de l'observation, puis à identifier le système.
  \item La loi physique n'est pas une formule isolée : elle explique le sens du phénomène.
  \item Le calcul Terminale donne un ordre de grandeur.
  \item Le modèle Bac+3 explique pourquoi le modèle simple marche, et où il cesse de marcher.
  \item Une réponse complète doit toujours contenir une phrase de conclusion en français.
\end{itemize}
\end{methodebox}

% ------------------------------------------------------------
\subsection{Problème détaillé 18 — Pourquoi un pont peut-il entrer en résonance ?}
% ------------------------------------------------------------

\begin{exercicebox}[Question centrale]
\textbf{Phénomène étudié :} résonance.

\medskip
\textbf{Question à résoudre :} Pourquoi un pont peut-il entrer en résonance ?

\medskip
\textbf{Contexte.} On observe le phénomène suivant : résonance. On cherche à comprendre non seulement ce que l'on voit, mais pourquoi la loi physique impose ce comportement.

\medskip
\textbf{Données.} On utilisera des ordres de grandeur réalistes et on précisera les hypothèses.
\end{exercicebox}

\subsubsection*{Énoncé guidé}

\begin{enumerate}
  \item Décrire l'observation.
  \item Identifier le système et les grandeurs.
  \item Écrire la loi physique principale.
  \item Faire un raisonnement d'ordre de grandeur.
  \item Discuter les limites du modèle.
\end{enumerate}

\subsubsection*{Modélisation}

\begin{theorembox}[Loi physique centrale]
On utilise comme point de départ :
\[
\omega\simeq\omega_0
\]
Cette relation doit être lue physiquement : elle indique quelles grandeurs contrôlent le phénomène et comment une modification d'un paramètre change le résultat.
\end{theorembox}

\subsubsection*{Correction niveau Terminale}

\begin{corrigebox}[Réponse accessible Terminale]

Au niveau Terminale, on explique ce phénomène en isolant la grandeur dominante. La relation utile est
\[
\omega\simeq\omega_0
\]
Elle permet de comprendre le sens de variation : si la grandeur qui apparaît au numérateur augmente, l'effet augmente ; si elle apparaît au dénominateur, l'effet diminue.

La réponse qualitative est donc que le phénomène « résonance » n'est pas magique : il résulte d'un transfert d'énergie, d'une propagation d'onde, d'une force ou d'une contrainte géométrique. Le modèle Terminale donne l'explication principale et permet de calculer un ordre de grandeur.

\end{corrigebox}

\subsubsection*{Correction approfondie niveau Bac+3}

\begin{approfondissementbox}[Réponse Bac+3]

Au niveau Bac+3, on raffine le modèle en introduisant une loi locale ou différentielle. La relation
\[
\omega\simeq\omega_0
\]
devient alors une conséquence d'un bilan plus général : bilan d'énergie, équation d'onde, équation de diffusion, dynamique du solide ou conservation de la quantité de mouvement. Les corrections importantes sont les pertes, les conditions aux limites, la géométrie réelle et les non-linéarités.

\end{approfondissementbox}

\subsubsection*{Bilan pédagogique}

\begin{methodebox}[Ce qu'il faut retenir]
\begin{itemize}
  \item La bonne démarche consiste à partir de l'observation, puis à identifier le système.
  \item La loi physique n'est pas une formule isolée : elle explique le sens du phénomène.
  \item Le calcul Terminale donne un ordre de grandeur.
  \item Le modèle Bac+3 explique pourquoi le modèle simple marche, et où il cesse de marcher.
  \item Une réponse complète doit toujours contenir une phrase de conclusion en français.
\end{itemize}
\end{methodebox}

% ------------------------------------------------------------
\subsection{Problème détaillé 19 — Pourquoi un parachute ralentit-il la chute ?}
% ------------------------------------------------------------

\begin{exercicebox}[Question centrale]
\textbf{Phénomène étudié :} parachute.

\medskip
\textbf{Question à résoudre :} Pourquoi un parachute ralentit-il la chute ?

\medskip
\textbf{Contexte.} On observe le phénomène suivant : parachute. On cherche à comprendre non seulement ce que l'on voit, mais pourquoi la loi physique impose ce comportement.

\medskip
\textbf{Données.} On utilisera des ordres de grandeur réalistes et on précisera les hypothèses.
\end{exercicebox}

\subsubsection*{Énoncé guidé}

\begin{enumerate}
  \item Décrire l'observation.
  \item Identifier le système et les grandeurs.
  \item Écrire la loi physique principale.
  \item Faire un raisonnement d'ordre de grandeur.
  \item Discuter les limites du modèle.
\end{enumerate}

\subsubsection*{Modélisation}

\begin{theorembox}[Loi physique centrale]
On utilise comme point de départ :
\[
mg=\frac12\rho C_DSv_T^2
\]
Cette relation doit être lue physiquement : elle indique quelles grandeurs contrôlent le phénomène et comment une modification d'un paramètre change le résultat.
\end{theorembox}

\subsubsection*{Correction niveau Terminale}

\begin{corrigebox}[Réponse accessible Terminale]

Au niveau Terminale, on explique ce phénomène en isolant la grandeur dominante. La relation utile est
\[
mg=\frac12\rho C_DSv_T^2
\]
Elle permet de comprendre le sens de variation : si la grandeur qui apparaît au numérateur augmente, l'effet augmente ; si elle apparaît au dénominateur, l'effet diminue.

La réponse qualitative est donc que le phénomène « parachute » n'est pas magique : il résulte d'un transfert d'énergie, d'une propagation d'onde, d'une force ou d'une contrainte géométrique. Le modèle Terminale donne l'explication principale et permet de calculer un ordre de grandeur.

\end{corrigebox}

\subsubsection*{Correction approfondie niveau Bac+3}

\begin{approfondissementbox}[Réponse Bac+3]

Au niveau Bac+3, on raffine le modèle en introduisant une loi locale ou différentielle. La relation
\[
mg=\frac12\rho C_DSv_T^2
\]
devient alors une conséquence d'un bilan plus général : bilan d'énergie, équation d'onde, équation de diffusion, dynamique du solide ou conservation de la quantité de mouvement. Les corrections importantes sont les pertes, les conditions aux limites, la géométrie réelle et les non-linéarités.

\end{approfondissementbox}

\subsubsection*{Bilan pédagogique}

\begin{methodebox}[Ce qu'il faut retenir]
\begin{itemize}
  \item La bonne démarche consiste à partir de l'observation, puis à identifier le système.
  \item La loi physique n'est pas une formule isolée : elle explique le sens du phénomène.
  \item Le calcul Terminale donne un ordre de grandeur.
  \item Le modèle Bac+3 explique pourquoi le modèle simple marche, et où il cesse de marcher.
  \item Une réponse complète doit toujours contenir une phrase de conclusion en français.
\end{itemize}
\end{methodebox}

% ------------------------------------------------------------
\subsection{Problème détaillé 20 — Pourquoi un vélo est-il plus stable quand il roule ?}
% ------------------------------------------------------------

\begin{exercicebox}[Question centrale]
\textbf{Phénomène étudié :} vélo.

\medskip
\textbf{Question à résoudre :} Pourquoi un vélo est-il plus stable quand il roule ?

\medskip
\textbf{Contexte.} On observe le phénomène suivant : vélo. On cherche à comprendre non seulement ce que l'on voit, mais pourquoi la loi physique impose ce comportement.

\medskip
\textbf{Données.} On utilisera des ordres de grandeur réalistes et on précisera les hypothèses.
\end{exercicebox}

\subsubsection*{Énoncé guidé}

\begin{enumerate}
  \item Décrire l'observation.
  \item Identifier le système et les grandeurs.
  \item Écrire la loi physique principale.
  \item Faire un raisonnement d'ordre de grandeur.
  \item Discuter les limites du modèle.
\end{enumerate}

\subsubsection*{Modélisation}

\begin{theorembox}[Loi physique centrale]
On utilise comme point de départ :
\[
\vec\tau=\frac{\mathrm d\vec L}{\mathrm dt}
\]
Cette relation doit être lue physiquement : elle indique quelles grandeurs contrôlent le phénomène et comment une modification d'un paramètre change le résultat.
\end{theorembox}

\subsubsection*{Correction niveau Terminale}

\begin{corrigebox}[Réponse accessible Terminale]

Au niveau Terminale, on explique ce phénomène en isolant la grandeur dominante. La relation utile est
\[
\vec\tau=\frac{\mathrm d\vec L}{\mathrm dt}
\]
Elle permet de comprendre le sens de variation : si la grandeur qui apparaît au numérateur augmente, l'effet augmente ; si elle apparaît au dénominateur, l'effet diminue.

La réponse qualitative est donc que le phénomène « vélo » n'est pas magique : il résulte d'un transfert d'énergie, d'une propagation d'onde, d'une force ou d'une contrainte géométrique. Le modèle Terminale donne l'explication principale et permet de calculer un ordre de grandeur.

\end{corrigebox}

\subsubsection*{Correction approfondie niveau Bac+3}

\begin{approfondissementbox}[Réponse Bac+3]

Au niveau Bac+3, on raffine le modèle en introduisant une loi locale ou différentielle. La relation
\[
\vec\tau=\frac{\mathrm d\vec L}{\mathrm dt}
\]
devient alors une conséquence d'un bilan plus général : bilan d'énergie, équation d'onde, équation de diffusion, dynamique du solide ou conservation de la quantité de mouvement. Les corrections importantes sont les pertes, les conditions aux limites, la géométrie réelle et les non-linéarités.

\end{approfondissementbox}

\subsubsection*{Bilan pédagogique}

\begin{methodebox}[Ce qu'il faut retenir]
\begin{itemize}
  \item La bonne démarche consiste à partir de l'observation, puis à identifier le système.
  \item La loi physique n'est pas une formule isolée : elle explique le sens du phénomène.
  \item Le calcul Terminale donne un ordre de grandeur.
  \item Le modèle Bac+3 explique pourquoi le modèle simple marche, et où il cesse de marcher.
  \item Une réponse complète doit toujours contenir une phrase de conclusion en français.
\end{itemize}
\end{methodebox}

\end{document}

